\documentclass{article}
\usepackage[paperwidth=16cm,paperheight=14cm,margin=0cm]{geometry}
\usepackage[utf8]{inputenc}\usepackage[T1]{fontenc}
\usepackage{mathpazo}\usepackage{amsmath,amssymb}
\usepackage[most]{tcolorbox}
\pagestyle{empty}\setlength{\parindent}{0pt}

\definecolor{csui}{HTML}{1F6FA8}
\definecolor{cfon}{HTML}{6A4C93}
\definecolor{cexl}{HTML}{2E7D6B}
\definecolor{cint}{HTML}{C77F0A}
\definecolor{cesp}{HTML}{C0503A}
\definecolor{cprb}{HTML}{8A5A2B}
\begin{document}
\begin{tcolorbox}[enhanced,width=16cm,height=14cm,colback=white,colframe=csui,boxrule=2pt,arc=0mm,
left=6mm,right=6mm,top=0mm,bottom=4mm,sharp corners,
overlay={\fill[csui] (frame.north west) rectangle ([yshift=-1.25cm]frame.north east);
\node[anchor=west,text=white,font=\sffamily\bfseries\LARGE] at ([xshift=5mm,yshift=-0.62cm]frame.north west) {Question 1};
\node[anchor=east,text=white,font=\sffamily\itshape\large] at ([xshift=-5mm,yshift=-0.62cm]frame.north east) {Suites};}]
\vspace*{1.75cm}
\begin{minipage}[t][6.9cm][c]{\linewidth}\LARGE\centering
\begin{minipage}{0.95\linewidth}\raggedright\linespread{1.25}\selectfont $(u_n)$ est la suite arithmétique de premier terme $u_0=3$ et de raison $2$. Que vaut $u_{10}$~?\end{minipage}
\end{minipage}

\vfill
\begin{tcolorbox}[enhanced,nobeforeafter,width=7.1cm,height=2.1cm,valign=center,colback=csui!6,colframe=csui,arc=2.5mm,boxrule=0.9pt,left=3mm]{\Large\bfseries\color{csui}a)}\quad{\Large $23$}\end{tcolorbox}\hfill\begin{tcolorbox}[enhanced,nobeforeafter,width=7.1cm,height=2.1cm,valign=center,colback=csui!6,colframe=csui,arc=2.5mm,boxrule=0.9pt,left=3mm]{\Large\bfseries\color{csui}b)}\quad{\Large $20$}\end{tcolorbox}\par\vspace{2.5mm}
\begin{tcolorbox}[enhanced,nobeforeafter,width=7.1cm,height=2.1cm,valign=center,colback=csui!6,colframe=csui,arc=2.5mm,boxrule=0.9pt,left=3mm]{\Large\bfseries\color{csui}c)}\quad{\Large $21$}\end{tcolorbox}\hfill\begin{tcolorbox}[enhanced,nobeforeafter,width=7.1cm,height=2.1cm,valign=center,colback=csui!6,colframe=csui,arc=2.5mm,boxrule=0.9pt,left=3mm]{\Large\bfseries\color{csui}d)}\quad{\Large $25$}\end{tcolorbox}
\end{tcolorbox}
\newpage
\begin{tcolorbox}[enhanced,width=16cm,height=14cm,colback=white,colframe=csui,boxrule=2pt,arc=0mm,
left=6mm,right=6mm,top=0mm,bottom=4mm,sharp corners,
overlay={\fill[csui] (frame.north west) rectangle ([yshift=-1.25cm]frame.north east);
\node[anchor=west,text=white,font=\sffamily\bfseries\LARGE] at ([xshift=5mm,yshift=-0.62cm]frame.north west) {Question 2};
\node[anchor=east,text=white,font=\sffamily\itshape\large] at ([xshift=-5mm,yshift=-0.62cm]frame.north east) {Suites};}]
\vspace*{1.75cm}
\begin{minipage}[t][6.9cm][c]{\linewidth}\LARGE\centering
\begin{minipage}{0.95\linewidth}\raggedright\linespread{1.25}\selectfont $(u_n)$ est la suite arithmétique de premier terme $u_0=5$ et de raison $-3$. Que vaut $u_{8}$~?\end{minipage}
\end{minipage}

\vfill
\begin{tcolorbox}[enhanced,nobeforeafter,width=7.1cm,height=2.1cm,valign=center,colback=csui!6,colframe=csui,arc=2.5mm,boxrule=0.9pt,left=3mm]{\Large\bfseries\color{csui}a)}\quad{\Large $-19$}\end{tcolorbox}\hfill\begin{tcolorbox}[enhanced,nobeforeafter,width=7.1cm,height=2.1cm,valign=center,colback=csui!6,colframe=csui,arc=2.5mm,boxrule=0.9pt,left=3mm]{\Large\bfseries\color{csui}b)}\quad{\Large $-16$}\end{tcolorbox}\par\vspace{2.5mm}
\begin{tcolorbox}[enhanced,nobeforeafter,width=7.1cm,height=2.1cm,valign=center,colback=csui!6,colframe=csui,arc=2.5mm,boxrule=0.9pt,left=3mm]{\Large\bfseries\color{csui}c)}\quad{\Large $-22$}\end{tcolorbox}\hfill\begin{tcolorbox}[enhanced,nobeforeafter,width=7.1cm,height=2.1cm,valign=center,colback=csui!6,colframe=csui,arc=2.5mm,boxrule=0.9pt,left=3mm]{\Large\bfseries\color{csui}d)}\quad{\Large $-24$}\end{tcolorbox}
\end{tcolorbox}
\newpage
\begin{tcolorbox}[enhanced,width=16cm,height=14cm,colback=white,colframe=csui,boxrule=2pt,arc=0mm,
left=6mm,right=6mm,top=0mm,bottom=4mm,sharp corners,
overlay={\fill[csui] (frame.north west) rectangle ([yshift=-1.25cm]frame.north east);
\node[anchor=west,text=white,font=\sffamily\bfseries\LARGE] at ([xshift=5mm,yshift=-0.62cm]frame.north west) {Question 3};
\node[anchor=east,text=white,font=\sffamily\itshape\large] at ([xshift=-5mm,yshift=-0.62cm]frame.north east) {Suites};}]
\vspace*{1.75cm}
\begin{minipage}[t][6.9cm][c]{\linewidth}\LARGE\centering
\begin{minipage}{0.95\linewidth}\raggedright\linespread{1.25}\selectfont Que vaut la somme $1+2+2^2+\dots+2^{5}$~?\end{minipage}
\end{minipage}

\vfill
\begin{tcolorbox}[enhanced,nobeforeafter,width=7.1cm,height=2.1cm,valign=center,colback=csui!6,colframe=csui,arc=2.5mm,boxrule=0.9pt,left=3mm]{\Large\bfseries\color{csui}a)}\quad{\Large $31$}\end{tcolorbox}\hfill\begin{tcolorbox}[enhanced,nobeforeafter,width=7.1cm,height=2.1cm,valign=center,colback=csui!6,colframe=csui,arc=2.5mm,boxrule=0.9pt,left=3mm]{\Large\bfseries\color{csui}b)}\quad{\Large $63$}\end{tcolorbox}\par\vspace{2.5mm}
\begin{tcolorbox}[enhanced,nobeforeafter,width=7.1cm,height=2.1cm,valign=center,colback=csui!6,colframe=csui,arc=2.5mm,boxrule=0.9pt,left=3mm]{\Large\bfseries\color{csui}c)}\quad{\Large $64$}\end{tcolorbox}\hfill\begin{tcolorbox}[enhanced,nobeforeafter,width=7.1cm,height=2.1cm,valign=center,colback=csui!6,colframe=csui,arc=2.5mm,boxrule=0.9pt,left=3mm]{\Large\bfseries\color{csui}d)}\quad{\Large $32$}\end{tcolorbox}
\end{tcolorbox}
\newpage
\begin{tcolorbox}[enhanced,width=16cm,height=14cm,colback=white,colframe=csui,boxrule=2pt,arc=0mm,
left=6mm,right=6mm,top=0mm,bottom=4mm,sharp corners,
overlay={\fill[csui] (frame.north west) rectangle ([yshift=-1.25cm]frame.north east);
\node[anchor=west,text=white,font=\sffamily\bfseries\LARGE] at ([xshift=5mm,yshift=-0.62cm]frame.north west) {Question 4};
\node[anchor=east,text=white,font=\sffamily\itshape\large] at ([xshift=-5mm,yshift=-0.62cm]frame.north east) {Suites};}]
\vspace*{1.75cm}
\begin{minipage}[t][6.9cm][c]{\linewidth}\LARGE\centering
\begin{minipage}{0.95\linewidth}\raggedright\linespread{1.25}\selectfont Que vaut la somme $1+3+3^2+\dots+3^{4}$~?\end{minipage}
\end{minipage}

\vfill
\begin{tcolorbox}[enhanced,nobeforeafter,width=7.1cm,height=2.1cm,valign=center,colback=csui!6,colframe=csui,arc=2.5mm,boxrule=0.9pt,left=3mm]{\Large\bfseries\color{csui}a)}\quad{\Large $242$}\end{tcolorbox}\hfill\begin{tcolorbox}[enhanced,nobeforeafter,width=7.1cm,height=2.1cm,valign=center,colback=csui!6,colframe=csui,arc=2.5mm,boxrule=0.9pt,left=3mm]{\Large\bfseries\color{csui}b)}\quad{\Large $40$}\end{tcolorbox}\par\vspace{2.5mm}
\begin{tcolorbox}[enhanced,nobeforeafter,width=7.1cm,height=2.1cm,valign=center,colback=csui!6,colframe=csui,arc=2.5mm,boxrule=0.9pt,left=3mm]{\Large\bfseries\color{csui}c)}\quad{\Large $243$}\end{tcolorbox}\hfill\begin{tcolorbox}[enhanced,nobeforeafter,width=7.1cm,height=2.1cm,valign=center,colback=csui!6,colframe=csui,arc=2.5mm,boxrule=0.9pt,left=3mm]{\Large\bfseries\color{csui}d)}\quad{\Large $121$}\end{tcolorbox}
\end{tcolorbox}
\newpage
\begin{tcolorbox}[enhanced,width=16cm,height=14cm,colback=white,colframe=csui,boxrule=2pt,arc=0mm,
left=6mm,right=6mm,top=0mm,bottom=4mm,sharp corners,
overlay={\fill[csui] (frame.north west) rectangle ([yshift=-1.25cm]frame.north east);
\node[anchor=west,text=white,font=\sffamily\bfseries\LARGE] at ([xshift=5mm,yshift=-0.62cm]frame.north west) {Question 5};
\node[anchor=east,text=white,font=\sffamily\itshape\large] at ([xshift=-5mm,yshift=-0.62cm]frame.north east) {Suites};}]
\vspace*{1.75cm}
\begin{minipage}[t][6.9cm][c]{\linewidth}\LARGE\centering
\begin{minipage}{0.95\linewidth}\raggedright\linespread{1.25}\selectfont Que vaut $\displaystyle\lim_{n\to+\infty}\frac{3n^2+1}{2n^2+5n}$~?\end{minipage}
\end{minipage}

\vfill
\begin{tcolorbox}[enhanced,nobeforeafter,width=7.1cm,height=2.1cm,valign=center,colback=csui!6,colframe=csui,arc=2.5mm,boxrule=0.9pt,left=3mm]{\Large\bfseries\color{csui}a)}\quad{\Large $\dfrac{2}{3}$}\end{tcolorbox}\hfill\begin{tcolorbox}[enhanced,nobeforeafter,width=7.1cm,height=2.1cm,valign=center,colback=csui!6,colframe=csui,arc=2.5mm,boxrule=0.9pt,left=3mm]{\Large\bfseries\color{csui}b)}\quad{\Large $\dfrac{3}{2}$}\end{tcolorbox}\par\vspace{2.5mm}
\begin{tcolorbox}[enhanced,nobeforeafter,width=7.1cm,height=2.1cm,valign=center,colback=csui!6,colframe=csui,arc=2.5mm,boxrule=0.9pt,left=3mm]{\Large\bfseries\color{csui}c)}\quad{\Large $+\infty$}\end{tcolorbox}\hfill\begin{tcolorbox}[enhanced,nobeforeafter,width=7.1cm,height=2.1cm,valign=center,colback=csui!6,colframe=csui,arc=2.5mm,boxrule=0.9pt,left=3mm]{\Large\bfseries\color{csui}d)}\quad{\Large $0$}\end{tcolorbox}
\end{tcolorbox}
\newpage
\begin{tcolorbox}[enhanced,width=16cm,height=14cm,colback=white,colframe=csui,boxrule=2pt,arc=0mm,
left=6mm,right=6mm,top=0mm,bottom=4mm,sharp corners,
overlay={\fill[csui] (frame.north west) rectangle ([yshift=-1.25cm]frame.north east);
\node[anchor=west,text=white,font=\sffamily\bfseries\LARGE] at ([xshift=5mm,yshift=-0.62cm]frame.north west) {Question 6};
\node[anchor=east,text=white,font=\sffamily\itshape\large] at ([xshift=-5mm,yshift=-0.62cm]frame.north east) {Suites};}]
\vspace*{1.75cm}
\begin{minipage}[t][6.9cm][c]{\linewidth}\LARGE\centering
\begin{minipage}{0.95\linewidth}\raggedright\linespread{1.25}\selectfont Que vaut $\displaystyle\lim_{n\to+\infty}\frac{4n^2-2}{7n^2+n}$~?\end{minipage}
\end{minipage}

\vfill
\begin{tcolorbox}[enhanced,nobeforeafter,width=7.1cm,height=2.1cm,valign=center,colback=csui!6,colframe=csui,arc=2.5mm,boxrule=0.9pt,left=3mm]{\Large\bfseries\color{csui}a)}\quad{\Large $+\infty$}\end{tcolorbox}\hfill\begin{tcolorbox}[enhanced,nobeforeafter,width=7.1cm,height=2.1cm,valign=center,colback=csui!6,colframe=csui,arc=2.5mm,boxrule=0.9pt,left=3mm]{\Large\bfseries\color{csui}b)}\quad{\Large $\dfrac{4}{7}$}\end{tcolorbox}\par\vspace{2.5mm}
\begin{tcolorbox}[enhanced,nobeforeafter,width=7.1cm,height=2.1cm,valign=center,colback=csui!6,colframe=csui,arc=2.5mm,boxrule=0.9pt,left=3mm]{\Large\bfseries\color{csui}c)}\quad{\Large $\dfrac{7}{4}$}\end{tcolorbox}\hfill\begin{tcolorbox}[enhanced,nobeforeafter,width=7.1cm,height=2.1cm,valign=center,colback=csui!6,colframe=csui,arc=2.5mm,boxrule=0.9pt,left=3mm]{\Large\bfseries\color{csui}d)}\quad{\Large $0$}\end{tcolorbox}
\end{tcolorbox}
\newpage
\begin{tcolorbox}[enhanced,width=16cm,height=14cm,colback=white,colframe=csui,boxrule=2pt,arc=0mm,
left=6mm,right=6mm,top=0mm,bottom=4mm,sharp corners,
overlay={\fill[csui] (frame.north west) rectangle ([yshift=-1.25cm]frame.north east);
\node[anchor=west,text=white,font=\sffamily\bfseries\LARGE] at ([xshift=5mm,yshift=-0.62cm]frame.north west) {Question 7};
\node[anchor=east,text=white,font=\sffamily\itshape\large] at ([xshift=-5mm,yshift=-0.62cm]frame.north east) {Suites};}]
\vspace*{1.75cm}
\begin{minipage}[t][6.9cm][c]{\linewidth}\LARGE\centering
\begin{minipage}{0.95\linewidth}\raggedright\linespread{1.25}\selectfont La suite définie par $u_0=0$ et $u_{n+1}=\frac{1}{2}\,u_n+3$ converge. Quelle est sa limite~?\end{minipage}
\end{minipage}

\vfill
\begin{tcolorbox}[enhanced,nobeforeafter,width=7.1cm,height=2.1cm,valign=center,colback=csui!6,colframe=csui,arc=2.5mm,boxrule=0.9pt,left=3mm]{\Large\bfseries\color{csui}a)}\quad{\Large $6$}\end{tcolorbox}\hfill\begin{tcolorbox}[enhanced,nobeforeafter,width=7.1cm,height=2.1cm,valign=center,colback=csui!6,colframe=csui,arc=2.5mm,boxrule=0.9pt,left=3mm]{\Large\bfseries\color{csui}b)}\quad{\Large $3$}\end{tcolorbox}\par\vspace{2.5mm}
\begin{tcolorbox}[enhanced,nobeforeafter,width=7.1cm,height=2.1cm,valign=center,colback=csui!6,colframe=csui,arc=2.5mm,boxrule=0.9pt,left=3mm]{\Large\bfseries\color{csui}c)}\quad{\Large $+\infty$}\end{tcolorbox}\hfill\begin{tcolorbox}[enhanced,nobeforeafter,width=7.1cm,height=2.1cm,valign=center,colback=csui!6,colframe=csui,arc=2.5mm,boxrule=0.9pt,left=3mm]{\Large\bfseries\color{csui}d)}\quad{\Large $2$}\end{tcolorbox}
\end{tcolorbox}
\newpage
\begin{tcolorbox}[enhanced,width=16cm,height=14cm,colback=white,colframe=csui,boxrule=2pt,arc=0mm,
left=6mm,right=6mm,top=0mm,bottom=4mm,sharp corners,
overlay={\fill[csui] (frame.north west) rectangle ([yshift=-1.25cm]frame.north east);
\node[anchor=west,text=white,font=\sffamily\bfseries\LARGE] at ([xshift=5mm,yshift=-0.62cm]frame.north west) {Question 8};
\node[anchor=east,text=white,font=\sffamily\itshape\large] at ([xshift=-5mm,yshift=-0.62cm]frame.north east) {Suites};}]
\vspace*{1.75cm}
\begin{minipage}[t][6.9cm][c]{\linewidth}\LARGE\centering
\begin{minipage}{0.95\linewidth}\raggedright\linespread{1.25}\selectfont La suite définie par $u_0=0$ et $u_{n+1}=\frac{1}{3}\,u_n+4$ converge. Quelle est sa limite~?\end{minipage}
\end{minipage}

\vfill
\begin{tcolorbox}[enhanced,nobeforeafter,width=7.1cm,height=2.1cm,valign=center,colback=csui!6,colframe=csui,arc=2.5mm,boxrule=0.9pt,left=3mm]{\Large\bfseries\color{csui}a)}\quad{\Large $6$}\end{tcolorbox}\hfill\begin{tcolorbox}[enhanced,nobeforeafter,width=7.1cm,height=2.1cm,valign=center,colback=csui!6,colframe=csui,arc=2.5mm,boxrule=0.9pt,left=3mm]{\Large\bfseries\color{csui}b)}\quad{\Large $3$}\end{tcolorbox}\par\vspace{2.5mm}
\begin{tcolorbox}[enhanced,nobeforeafter,width=7.1cm,height=2.1cm,valign=center,colback=csui!6,colframe=csui,arc=2.5mm,boxrule=0.9pt,left=3mm]{\Large\bfseries\color{csui}c)}\quad{\Large $4$}\end{tcolorbox}\hfill\begin{tcolorbox}[enhanced,nobeforeafter,width=7.1cm,height=2.1cm,valign=center,colback=csui!6,colframe=csui,arc=2.5mm,boxrule=0.9pt,left=3mm]{\Large\bfseries\color{csui}d)}\quad{\Large $+\infty$}\end{tcolorbox}
\end{tcolorbox}
\newpage
\begin{tcolorbox}[enhanced,width=16cm,height=14cm,colback=white,colframe=csui,boxrule=2pt,arc=0mm,
left=6mm,right=6mm,top=0mm,bottom=4mm,sharp corners,
overlay={\fill[csui] (frame.north west) rectangle ([yshift=-1.25cm]frame.north east);
\node[anchor=west,text=white,font=\sffamily\bfseries\LARGE] at ([xshift=5mm,yshift=-0.62cm]frame.north west) {Question 9};
\node[anchor=east,text=white,font=\sffamily\itshape\large] at ([xshift=-5mm,yshift=-0.62cm]frame.north east) {Suites};}]
\vspace*{1.75cm}
\begin{minipage}[t][6.9cm][c]{\linewidth}\LARGE\centering
\begin{minipage}{0.95\linewidth}\raggedright\linespread{1.25}\selectfont On pose $u_0=1$ et, pour tout entier $n$, $u_{n+1}=u_n+2n+1$. Que vaut $u_{5}$~?\end{minipage}
\end{minipage}

\vfill
\begin{tcolorbox}[enhanced,nobeforeafter,width=7.1cm,height=2.1cm,valign=center,colback=csui!6,colframe=csui,arc=2.5mm,boxrule=0.9pt,left=3mm]{\Large\bfseries\color{csui}a)}\quad{\Large $37$}\end{tcolorbox}\hfill\begin{tcolorbox}[enhanced,nobeforeafter,width=7.1cm,height=2.1cm,valign=center,colback=csui!6,colframe=csui,arc=2.5mm,boxrule=0.9pt,left=3mm]{\Large\bfseries\color{csui}b)}\quad{\Large $12$}\end{tcolorbox}\par\vspace{2.5mm}
\begin{tcolorbox}[enhanced,nobeforeafter,width=7.1cm,height=2.1cm,valign=center,colback=csui!6,colframe=csui,arc=2.5mm,boxrule=0.9pt,left=3mm]{\Large\bfseries\color{csui}c)}\quad{\Large $25$}\end{tcolorbox}\hfill\begin{tcolorbox}[enhanced,nobeforeafter,width=7.1cm,height=2.1cm,valign=center,colback=csui!6,colframe=csui,arc=2.5mm,boxrule=0.9pt,left=3mm]{\Large\bfseries\color{csui}d)}\quad{\Large $26$}\end{tcolorbox}
\end{tcolorbox}
\newpage
\begin{tcolorbox}[enhanced,width=16cm,height=14cm,colback=white,colframe=csui,boxrule=2pt,arc=0mm,
left=6mm,right=6mm,top=0mm,bottom=4mm,sharp corners,
overlay={\fill[csui] (frame.north west) rectangle ([yshift=-1.25cm]frame.north east);
\node[anchor=west,text=white,font=\sffamily\bfseries\LARGE] at ([xshift=5mm,yshift=-0.62cm]frame.north west) {Question 10};
\node[anchor=east,text=white,font=\sffamily\itshape\large] at ([xshift=-5mm,yshift=-0.62cm]frame.north east) {Suites};}]
\vspace*{1.75cm}
\begin{minipage}[t][6.9cm][c]{\linewidth}\LARGE\centering
\begin{minipage}{0.95\linewidth}\raggedright\linespread{1.25}\selectfont On pose $u_0=3$ et, pour tout entier $n$, $u_{n+1}=u_n+2n+1$. Que vaut $u_{6}$~?\end{minipage}
\end{minipage}

\vfill
\begin{tcolorbox}[enhanced,nobeforeafter,width=7.1cm,height=2.1cm,valign=center,colback=csui!6,colframe=csui,arc=2.5mm,boxrule=0.9pt,left=3mm]{\Large\bfseries\color{csui}a)}\quad{\Large $52$}\end{tcolorbox}\hfill\begin{tcolorbox}[enhanced,nobeforeafter,width=7.1cm,height=2.1cm,valign=center,colback=csui!6,colframe=csui,arc=2.5mm,boxrule=0.9pt,left=3mm]{\Large\bfseries\color{csui}b)}\quad{\Large $38$}\end{tcolorbox}\par\vspace{2.5mm}
\begin{tcolorbox}[enhanced,nobeforeafter,width=7.1cm,height=2.1cm,valign=center,colback=csui!6,colframe=csui,arc=2.5mm,boxrule=0.9pt,left=3mm]{\Large\bfseries\color{csui}c)}\quad{\Large $39$}\end{tcolorbox}\hfill\begin{tcolorbox}[enhanced,nobeforeafter,width=7.1cm,height=2.1cm,valign=center,colback=csui!6,colframe=csui,arc=2.5mm,boxrule=0.9pt,left=3mm]{\Large\bfseries\color{csui}d)}\quad{\Large $16$}\end{tcolorbox}
\end{tcolorbox}
\newpage
\begin{tcolorbox}[enhanced,width=16cm,height=14cm,colback=white,colframe=csui,boxrule=2pt,arc=0mm,
left=6mm,right=6mm,top=0mm,bottom=4mm,sharp corners,
overlay={\fill[csui] (frame.north west) rectangle ([yshift=-1.25cm]frame.north east);
\node[anchor=west,text=white,font=\sffamily\bfseries\LARGE] at ([xshift=5mm,yshift=-0.62cm]frame.north west) {Question 11};
\node[anchor=east,text=white,font=\sffamily\itshape\large] at ([xshift=-5mm,yshift=-0.62cm]frame.north east) {Suites};}]
\vspace*{1.75cm}
\begin{minipage}[t][6.9cm][c]{\linewidth}\LARGE\centering
\begin{minipage}{0.95\linewidth}\raggedright\linespread{1.25}\selectfont Que vaut $\displaystyle\lim_{n\to+\infty}\frac{2n+(-1)^n}{3n+1}$~?\end{minipage}
\end{minipage}

\vfill
\begin{tcolorbox}[enhanced,nobeforeafter,width=7.1cm,height=2.1cm,valign=center,colback=csui!6,colframe=csui,arc=2.5mm,boxrule=0.9pt,left=3mm]{\Large\bfseries\color{csui}a)}\quad{\Large $+\infty$}\end{tcolorbox}\hfill\begin{tcolorbox}[enhanced,nobeforeafter,width=7.1cm,height=2.1cm,valign=center,colback=csui!6,colframe=csui,arc=2.5mm,boxrule=0.9pt,left=3mm]{\Large\bfseries\color{csui}b)}\quad{\Large $\dfrac{3}{2}$}\end{tcolorbox}\par\vspace{2.5mm}
\begin{tcolorbox}[enhanced,nobeforeafter,width=7.1cm,height=2.1cm,valign=center,colback=csui!6,colframe=csui,arc=2.5mm,boxrule=0.9pt,left=3mm]{\Large\bfseries\color{csui}c)}\quad{\Large $0$}\end{tcolorbox}\hfill\begin{tcolorbox}[enhanced,nobeforeafter,width=7.1cm,height=2.1cm,valign=center,colback=csui!6,colframe=csui,arc=2.5mm,boxrule=0.9pt,left=3mm]{\Large\bfseries\color{csui}d)}\quad{\Large $\dfrac{2}{3}$}\end{tcolorbox}
\end{tcolorbox}
\newpage
\begin{tcolorbox}[enhanced,width=16cm,height=14cm,colback=white,colframe=csui,boxrule=2pt,arc=0mm,
left=6mm,right=6mm,top=0mm,bottom=4mm,sharp corners,
overlay={\fill[csui] (frame.north west) rectangle ([yshift=-1.25cm]frame.north east);
\node[anchor=west,text=white,font=\sffamily\bfseries\LARGE] at ([xshift=5mm,yshift=-0.62cm]frame.north west) {Question 12};
\node[anchor=east,text=white,font=\sffamily\itshape\large] at ([xshift=-5mm,yshift=-0.62cm]frame.north east) {Suites};}]
\vspace*{1.75cm}
\begin{minipage}[t][6.9cm][c]{\linewidth}\LARGE\centering
\begin{minipage}{0.95\linewidth}\raggedright\linespread{1.25}\selectfont Que vaut $\displaystyle\lim_{n\to+\infty}\frac{5n+(-1)^n}{2n+1}$~?\end{minipage}
\end{minipage}

\vfill
\begin{tcolorbox}[enhanced,nobeforeafter,width=7.1cm,height=2.1cm,valign=center,colback=csui!6,colframe=csui,arc=2.5mm,boxrule=0.9pt,left=3mm]{\Large\bfseries\color{csui}a)}\quad{\Large $0$}\end{tcolorbox}\hfill\begin{tcolorbox}[enhanced,nobeforeafter,width=7.1cm,height=2.1cm,valign=center,colback=csui!6,colframe=csui,arc=2.5mm,boxrule=0.9pt,left=3mm]{\Large\bfseries\color{csui}b)}\quad{\Large $\dfrac{2}{5}$}\end{tcolorbox}\par\vspace{2.5mm}
\begin{tcolorbox}[enhanced,nobeforeafter,width=7.1cm,height=2.1cm,valign=center,colback=csui!6,colframe=csui,arc=2.5mm,boxrule=0.9pt,left=3mm]{\Large\bfseries\color{csui}c)}\quad{\Large $+\infty$}\end{tcolorbox}\hfill\begin{tcolorbox}[enhanced,nobeforeafter,width=7.1cm,height=2.1cm,valign=center,colback=csui!6,colframe=csui,arc=2.5mm,boxrule=0.9pt,left=3mm]{\Large\bfseries\color{csui}d)}\quad{\Large $\dfrac{5}{2}$}\end{tcolorbox}
\end{tcolorbox}
\newpage
\begin{tcolorbox}[enhanced,width=16cm,height=14cm,colback=white,colframe=csui,boxrule=2pt,arc=0mm,
left=6mm,right=6mm,top=0mm,bottom=4mm,sharp corners,
overlay={\fill[csui] (frame.north west) rectangle ([yshift=-1.25cm]frame.north east);
\node[anchor=west,text=white,font=\sffamily\bfseries\LARGE] at ([xshift=5mm,yshift=-0.62cm]frame.north west) {Question 13};
\node[anchor=east,text=white,font=\sffamily\itshape\large] at ([xshift=-5mm,yshift=-0.62cm]frame.north east) {Suites};}]
\vspace*{1.75cm}
\begin{minipage}[t][6.9cm][c]{\linewidth}\LARGE\centering
\begin{minipage}{0.95\linewidth}\raggedright\linespread{1.25}\selectfont $(u_n)$ est une suite géométrique de raison $2$ telle que $u_3=5$. Que vaut $u_5$~?\end{minipage}
\end{minipage}

\vfill
\begin{tcolorbox}[enhanced,nobeforeafter,width=7.1cm,height=2.1cm,valign=center,colback=csui!6,colframe=csui,arc=2.5mm,boxrule=0.9pt,left=3mm]{\Large\bfseries\color{csui}a)}\quad{\Large $10$}\end{tcolorbox}\hfill\begin{tcolorbox}[enhanced,nobeforeafter,width=7.1cm,height=2.1cm,valign=center,colback=csui!6,colframe=csui,arc=2.5mm,boxrule=0.9pt,left=3mm]{\Large\bfseries\color{csui}b)}\quad{\Large $20$}\end{tcolorbox}\par\vspace{2.5mm}
\begin{tcolorbox}[enhanced,nobeforeafter,width=7.1cm,height=2.1cm,valign=center,colback=csui!6,colframe=csui,arc=2.5mm,boxrule=0.9pt,left=3mm]{\Large\bfseries\color{csui}c)}\quad{\Large $40$}\end{tcolorbox}\hfill\begin{tcolorbox}[enhanced,nobeforeafter,width=7.1cm,height=2.1cm,valign=center,colback=csui!6,colframe=csui,arc=2.5mm,boxrule=0.9pt,left=3mm]{\Large\bfseries\color{csui}d)}\quad{\Large $9$}\end{tcolorbox}
\end{tcolorbox}
\newpage
\begin{tcolorbox}[enhanced,width=16cm,height=14cm,colback=white,colframe=csui,boxrule=2pt,arc=0mm,
left=6mm,right=6mm,top=0mm,bottom=4mm,sharp corners,
overlay={\fill[csui] (frame.north west) rectangle ([yshift=-1.25cm]frame.north east);
\node[anchor=west,text=white,font=\sffamily\bfseries\LARGE] at ([xshift=5mm,yshift=-0.62cm]frame.north west) {Question 14};
\node[anchor=east,text=white,font=\sffamily\itshape\large] at ([xshift=-5mm,yshift=-0.62cm]frame.north east) {Suites};}]
\vspace*{1.75cm}
\begin{minipage}[t][6.9cm][c]{\linewidth}\LARGE\centering
\begin{minipage}{0.95\linewidth}\raggedright\linespread{1.25}\selectfont $(u_n)$ est une suite géométrique de raison $3$ telle que $u_3=2$. Que vaut $u_5$~?\end{minipage}
\end{minipage}

\vfill
\begin{tcolorbox}[enhanced,nobeforeafter,width=7.1cm,height=2.1cm,valign=center,colback=csui!6,colframe=csui,arc=2.5mm,boxrule=0.9pt,left=3mm]{\Large\bfseries\color{csui}a)}\quad{\Large $54$}\end{tcolorbox}\hfill\begin{tcolorbox}[enhanced,nobeforeafter,width=7.1cm,height=2.1cm,valign=center,colback=csui!6,colframe=csui,arc=2.5mm,boxrule=0.9pt,left=3mm]{\Large\bfseries\color{csui}b)}\quad{\Large $6$}\end{tcolorbox}\par\vspace{2.5mm}
\begin{tcolorbox}[enhanced,nobeforeafter,width=7.1cm,height=2.1cm,valign=center,colback=csui!6,colframe=csui,arc=2.5mm,boxrule=0.9pt,left=3mm]{\Large\bfseries\color{csui}c)}\quad{\Large $18$}\end{tcolorbox}\hfill\begin{tcolorbox}[enhanced,nobeforeafter,width=7.1cm,height=2.1cm,valign=center,colback=csui!6,colframe=csui,arc=2.5mm,boxrule=0.9pt,left=3mm]{\Large\bfseries\color{csui}d)}\quad{\Large $8$}\end{tcolorbox}
\end{tcolorbox}
\newpage
\begin{tcolorbox}[enhanced,width=16cm,height=14cm,colback=white,colframe=csui,boxrule=2pt,arc=0mm,
left=6mm,right=6mm,top=0mm,bottom=4mm,sharp corners,
overlay={\fill[csui] (frame.north west) rectangle ([yshift=-1.25cm]frame.north east);
\node[anchor=west,text=white,font=\sffamily\bfseries\LARGE] at ([xshift=5mm,yshift=-0.62cm]frame.north west) {Question 15};
\node[anchor=east,text=white,font=\sffamily\itshape\large] at ([xshift=-5mm,yshift=-0.62cm]frame.north east) {Suites};}]
\vspace*{1.75cm}
\begin{minipage}[t][6.9cm][c]{\linewidth}\LARGE\centering
\begin{minipage}{0.95\linewidth}\raggedright\linespread{1.25}\selectfont Quel est le plus petit entier naturel $n$ tel que $2^n>1000$~?\end{minipage}
\end{minipage}

\vfill
\begin{tcolorbox}[enhanced,nobeforeafter,width=7.1cm,height=2.1cm,valign=center,colback=csui!6,colframe=csui,arc=2.5mm,boxrule=0.9pt,left=3mm]{\Large\bfseries\color{csui}a)}\quad{\Large $10$}\end{tcolorbox}\hfill\begin{tcolorbox}[enhanced,nobeforeafter,width=7.1cm,height=2.1cm,valign=center,colback=csui!6,colframe=csui,arc=2.5mm,boxrule=0.9pt,left=3mm]{\Large\bfseries\color{csui}b)}\quad{\Large $9$}\end{tcolorbox}\par\vspace{2.5mm}
\begin{tcolorbox}[enhanced,nobeforeafter,width=7.1cm,height=2.1cm,valign=center,colback=csui!6,colframe=csui,arc=2.5mm,boxrule=0.9pt,left=3mm]{\Large\bfseries\color{csui}c)}\quad{\Large $12$}\end{tcolorbox}\hfill\begin{tcolorbox}[enhanced,nobeforeafter,width=7.1cm,height=2.1cm,valign=center,colback=csui!6,colframe=csui,arc=2.5mm,boxrule=0.9pt,left=3mm]{\Large\bfseries\color{csui}d)}\quad{\Large $11$}\end{tcolorbox}
\end{tcolorbox}
\newpage
\begin{tcolorbox}[enhanced,width=16cm,height=14cm,colback=white,colframe=csui,boxrule=2pt,arc=0mm,
left=6mm,right=6mm,top=0mm,bottom=4mm,sharp corners,
overlay={\fill[csui] (frame.north west) rectangle ([yshift=-1.25cm]frame.north east);
\node[anchor=west,text=white,font=\sffamily\bfseries\LARGE] at ([xshift=5mm,yshift=-0.62cm]frame.north west) {Question 16};
\node[anchor=east,text=white,font=\sffamily\itshape\large] at ([xshift=-5mm,yshift=-0.62cm]frame.north east) {Suites};}]
\vspace*{1.75cm}
\begin{minipage}[t][6.9cm][c]{\linewidth}\LARGE\centering
\begin{minipage}{0.95\linewidth}\raggedright\linespread{1.25}\selectfont Quel est le plus petit entier naturel $n$ tel que $3^n>500$~?\end{minipage}
\end{minipage}

\vfill
\begin{tcolorbox}[enhanced,nobeforeafter,width=7.1cm,height=2.1cm,valign=center,colback=csui!6,colframe=csui,arc=2.5mm,boxrule=0.9pt,left=3mm]{\Large\bfseries\color{csui}a)}\quad{\Large $7$}\end{tcolorbox}\hfill\begin{tcolorbox}[enhanced,nobeforeafter,width=7.1cm,height=2.1cm,valign=center,colback=csui!6,colframe=csui,arc=2.5mm,boxrule=0.9pt,left=3mm]{\Large\bfseries\color{csui}b)}\quad{\Large $5$}\end{tcolorbox}\par\vspace{2.5mm}
\begin{tcolorbox}[enhanced,nobeforeafter,width=7.1cm,height=2.1cm,valign=center,colback=csui!6,colframe=csui,arc=2.5mm,boxrule=0.9pt,left=3mm]{\Large\bfseries\color{csui}c)}\quad{\Large $8$}\end{tcolorbox}\hfill\begin{tcolorbox}[enhanced,nobeforeafter,width=7.1cm,height=2.1cm,valign=center,colback=csui!6,colframe=csui,arc=2.5mm,boxrule=0.9pt,left=3mm]{\Large\bfseries\color{csui}d)}\quad{\Large $6$}\end{tcolorbox}
\end{tcolorbox}
\newpage
\begin{tcolorbox}[enhanced,width=16cm,height=14cm,colback=white,colframe=cfon,boxrule=2pt,arc=0mm,
left=6mm,right=6mm,top=0mm,bottom=4mm,sharp corners,
overlay={\fill[cfon] (frame.north west) rectangle ([yshift=-1.25cm]frame.north east);
\node[anchor=west,text=white,font=\sffamily\bfseries\LARGE] at ([xshift=5mm,yshift=-0.62cm]frame.north west) {Question 17};
\node[anchor=east,text=white,font=\sffamily\itshape\large] at ([xshift=-5mm,yshift=-0.62cm]frame.north east) {Fonctions};}]
\vspace*{1.75cm}
\begin{minipage}[t][6.9cm][c]{\linewidth}\LARGE\centering
\begin{minipage}{0.95\linewidth}\raggedright\linespread{1.25}\selectfont Soit $f(x)=(2x+1)^3$. Que vaut $f'(0)$~?\end{minipage}
\end{minipage}

\vfill
\begin{tcolorbox}[enhanced,nobeforeafter,width=7.1cm,height=2.1cm,valign=center,colback=cfon!6,colframe=cfon,arc=2.5mm,boxrule=0.9pt,left=3mm]{\Large\bfseries\color{cfon}a)}\quad{\Large $1$}\end{tcolorbox}\hfill\begin{tcolorbox}[enhanced,nobeforeafter,width=7.1cm,height=2.1cm,valign=center,colback=cfon!6,colframe=cfon,arc=2.5mm,boxrule=0.9pt,left=3mm]{\Large\bfseries\color{cfon}b)}\quad{\Large $6$}\end{tcolorbox}\par\vspace{2.5mm}
\begin{tcolorbox}[enhanced,nobeforeafter,width=7.1cm,height=2.1cm,valign=center,colback=cfon!6,colframe=cfon,arc=2.5mm,boxrule=0.9pt,left=3mm]{\Large\bfseries\color{cfon}c)}\quad{\Large $3$}\end{tcolorbox}\hfill\begin{tcolorbox}[enhanced,nobeforeafter,width=7.1cm,height=2.1cm,valign=center,colback=cfon!6,colframe=cfon,arc=2.5mm,boxrule=0.9pt,left=3mm]{\Large\bfseries\color{cfon}d)}\quad{\Large $2$}\end{tcolorbox}
\end{tcolorbox}
\newpage
\begin{tcolorbox}[enhanced,width=16cm,height=14cm,colback=white,colframe=cfon,boxrule=2pt,arc=0mm,
left=6mm,right=6mm,top=0mm,bottom=4mm,sharp corners,
overlay={\fill[cfon] (frame.north west) rectangle ([yshift=-1.25cm]frame.north east);
\node[anchor=west,text=white,font=\sffamily\bfseries\LARGE] at ([xshift=5mm,yshift=-0.62cm]frame.north west) {Question 18};
\node[anchor=east,text=white,font=\sffamily\itshape\large] at ([xshift=-5mm,yshift=-0.62cm]frame.north east) {Fonctions};}]
\vspace*{1.75cm}
\begin{minipage}[t][6.9cm][c]{\linewidth}\LARGE\centering
\begin{minipage}{0.95\linewidth}\raggedright\linespread{1.25}\selectfont Soit $f(x)=(3x-2)^3$. Que vaut $f'(0)$~?\end{minipage}
\end{minipage}

\vfill
\begin{tcolorbox}[enhanced,nobeforeafter,width=7.1cm,height=2.1cm,valign=center,colback=cfon!6,colframe=cfon,arc=2.5mm,boxrule=0.9pt,left=3mm]{\Large\bfseries\color{cfon}a)}\quad{\Large $-24$}\end{tcolorbox}\hfill\begin{tcolorbox}[enhanced,nobeforeafter,width=7.1cm,height=2.1cm,valign=center,colback=cfon!6,colframe=cfon,arc=2.5mm,boxrule=0.9pt,left=3mm]{\Large\bfseries\color{cfon}b)}\quad{\Large $12$}\end{tcolorbox}\par\vspace{2.5mm}
\begin{tcolorbox}[enhanced,nobeforeafter,width=7.1cm,height=2.1cm,valign=center,colback=cfon!6,colframe=cfon,arc=2.5mm,boxrule=0.9pt,left=3mm]{\Large\bfseries\color{cfon}c)}\quad{\Large $-18$}\end{tcolorbox}\hfill\begin{tcolorbox}[enhanced,nobeforeafter,width=7.1cm,height=2.1cm,valign=center,colback=cfon!6,colframe=cfon,arc=2.5mm,boxrule=0.9pt,left=3mm]{\Large\bfseries\color{cfon}d)}\quad{\Large $36$}\end{tcolorbox}
\end{tcolorbox}
\newpage
\begin{tcolorbox}[enhanced,width=16cm,height=14cm,colback=white,colframe=cfon,boxrule=2pt,arc=0mm,
left=6mm,right=6mm,top=0mm,bottom=4mm,sharp corners,
overlay={\fill[cfon] (frame.north west) rectangle ([yshift=-1.25cm]frame.north east);
\node[anchor=west,text=white,font=\sffamily\bfseries\LARGE] at ([xshift=5mm,yshift=-0.62cm]frame.north west) {Question 19};
\node[anchor=east,text=white,font=\sffamily\itshape\large] at ([xshift=-5mm,yshift=-0.62cm]frame.north east) {Fonctions};}]
\vspace*{1.75cm}
\begin{minipage}[t][6.9cm][c]{\linewidth}\LARGE\centering
\begin{minipage}{0.95\linewidth}\raggedright\linespread{1.25}\selectfont Soit $f(x)=\dfrac{x+2}{x+1}$. Que vaut $f'(0)$~?\end{minipage}
\end{minipage}

\vfill
\begin{tcolorbox}[enhanced,nobeforeafter,width=7.1cm,height=2.1cm,valign=center,colback=cfon!6,colframe=cfon,arc=2.5mm,boxrule=0.9pt,left=3mm]{\Large\bfseries\color{cfon}a)}\quad{\Large $-1$}\end{tcolorbox}\hfill\begin{tcolorbox}[enhanced,nobeforeafter,width=7.1cm,height=2.1cm,valign=center,colback=cfon!6,colframe=cfon,arc=2.5mm,boxrule=0.9pt,left=3mm]{\Large\bfseries\color{cfon}b)}\quad{\Large $0$}\end{tcolorbox}\par\vspace{2.5mm}
\begin{tcolorbox}[enhanced,nobeforeafter,width=7.1cm,height=2.1cm,valign=center,colback=cfon!6,colframe=cfon,arc=2.5mm,boxrule=0.9pt,left=3mm]{\Large\bfseries\color{cfon}c)}\quad{\Large $1$}\end{tcolorbox}\hfill\begin{tcolorbox}[enhanced,nobeforeafter,width=7.1cm,height=2.1cm,valign=center,colback=cfon!6,colframe=cfon,arc=2.5mm,boxrule=0.9pt,left=3mm]{\Large\bfseries\color{cfon}d)}\quad{\Large $3$}\end{tcolorbox}
\end{tcolorbox}
\newpage
\begin{tcolorbox}[enhanced,width=16cm,height=14cm,colback=white,colframe=cfon,boxrule=2pt,arc=0mm,
left=6mm,right=6mm,top=0mm,bottom=4mm,sharp corners,
overlay={\fill[cfon] (frame.north west) rectangle ([yshift=-1.25cm]frame.north east);
\node[anchor=west,text=white,font=\sffamily\bfseries\LARGE] at ([xshift=5mm,yshift=-0.62cm]frame.north west) {Question 20};
\node[anchor=east,text=white,font=\sffamily\itshape\large] at ([xshift=-5mm,yshift=-0.62cm]frame.north east) {Fonctions};}]
\vspace*{1.75cm}
\begin{minipage}[t][6.9cm][c]{\linewidth}\LARGE\centering
\begin{minipage}{0.95\linewidth}\raggedright\linespread{1.25}\selectfont Soit $f(x)=\dfrac{2x+1}{x+3}$. Que vaut $f'(0)$~?\end{minipage}
\end{minipage}

\vfill
\begin{tcolorbox}[enhanced,nobeforeafter,width=7.1cm,height=2.1cm,valign=center,colback=cfon!6,colframe=cfon,arc=2.5mm,boxrule=0.9pt,left=3mm]{\Large\bfseries\color{cfon}a)}\quad{\Large $\dfrac{7}{9}$}\end{tcolorbox}\hfill\begin{tcolorbox}[enhanced,nobeforeafter,width=7.1cm,height=2.1cm,valign=center,colback=cfon!6,colframe=cfon,arc=2.5mm,boxrule=0.9pt,left=3mm]{\Large\bfseries\color{cfon}b)}\quad{\Large $2$}\end{tcolorbox}\par\vspace{2.5mm}
\begin{tcolorbox}[enhanced,nobeforeafter,width=7.1cm,height=2.1cm,valign=center,colback=cfon!6,colframe=cfon,arc=2.5mm,boxrule=0.9pt,left=3mm]{\Large\bfseries\color{cfon}c)}\quad{\Large $\dfrac{5}{9}$}\end{tcolorbox}\hfill\begin{tcolorbox}[enhanced,nobeforeafter,width=7.1cm,height=2.1cm,valign=center,colback=cfon!6,colframe=cfon,arc=2.5mm,boxrule=0.9pt,left=3mm]{\Large\bfseries\color{cfon}d)}\quad{\Large $-\dfrac{5}{9}$}\end{tcolorbox}
\end{tcolorbox}
\newpage
\begin{tcolorbox}[enhanced,width=16cm,height=14cm,colback=white,colframe=cfon,boxrule=2pt,arc=0mm,
left=6mm,right=6mm,top=0mm,bottom=4mm,sharp corners,
overlay={\fill[cfon] (frame.north west) rectangle ([yshift=-1.25cm]frame.north east);
\node[anchor=west,text=white,font=\sffamily\bfseries\LARGE] at ([xshift=5mm,yshift=-0.62cm]frame.north west) {Question 21};
\node[anchor=east,text=white,font=\sffamily\itshape\large] at ([xshift=-5mm,yshift=-0.62cm]frame.north east) {Fonctions};}]
\vspace*{1.75cm}
\begin{minipage}[t][6.9cm][c]{\linewidth}\LARGE\centering
\begin{minipage}{0.95\linewidth}\raggedright\linespread{1.25}\selectfont Quelle est l'équation de la tangente à la courbe de $f(x)=x^2+x+2$ au point d'abscisse $1$~?\end{minipage}
\end{minipage}

\vfill
\begin{tcolorbox}[enhanced,nobeforeafter,width=7.1cm,height=2.1cm,valign=center,colback=cfon!6,colframe=cfon,arc=2.5mm,boxrule=0.9pt,left=3mm]{\Large\bfseries\color{cfon}a)}\quad{\Large $y=3x+1$}\end{tcolorbox}\hfill\begin{tcolorbox}[enhanced,nobeforeafter,width=7.1cm,height=2.1cm,valign=center,colback=cfon!6,colframe=cfon,arc=2.5mm,boxrule=0.9pt,left=3mm]{\Large\bfseries\color{cfon}b)}\quad{\Large $y=3x-1$}\end{tcolorbox}\par\vspace{2.5mm}
\begin{tcolorbox}[enhanced,nobeforeafter,width=7.1cm,height=2.1cm,valign=center,colback=cfon!6,colframe=cfon,arc=2.5mm,boxrule=0.9pt,left=3mm]{\Large\bfseries\color{cfon}c)}\quad{\Large $y=3x+2$}\end{tcolorbox}\hfill\begin{tcolorbox}[enhanced,nobeforeafter,width=7.1cm,height=2.1cm,valign=center,colback=cfon!6,colframe=cfon,arc=2.5mm,boxrule=0.9pt,left=3mm]{\Large\bfseries\color{cfon}d)}\quad{\Large $y=2x+1$}\end{tcolorbox}
\end{tcolorbox}
\newpage
\begin{tcolorbox}[enhanced,width=16cm,height=14cm,colback=white,colframe=cfon,boxrule=2pt,arc=0mm,
left=6mm,right=6mm,top=0mm,bottom=4mm,sharp corners,
overlay={\fill[cfon] (frame.north west) rectangle ([yshift=-1.25cm]frame.north east);
\node[anchor=west,text=white,font=\sffamily\bfseries\LARGE] at ([xshift=5mm,yshift=-0.62cm]frame.north west) {Question 22};
\node[anchor=east,text=white,font=\sffamily\itshape\large] at ([xshift=-5mm,yshift=-0.62cm]frame.north east) {Fonctions};}]
\vspace*{1.75cm}
\begin{minipage}[t][6.9cm][c]{\linewidth}\LARGE\centering
\begin{minipage}{0.95\linewidth}\raggedright\linespread{1.25}\selectfont Quelle est l'équation de la tangente à la courbe de $f(x)=x^2-2x+3$ au point d'abscisse $2$~?\end{minipage}
\end{minipage}

\vfill
\begin{tcolorbox}[enhanced,nobeforeafter,width=7.1cm,height=2.1cm,valign=center,colback=cfon!6,colframe=cfon,arc=2.5mm,boxrule=0.9pt,left=3mm]{\Large\bfseries\color{cfon}a)}\quad{\Large $y=-1$}\end{tcolorbox}\hfill\begin{tcolorbox}[enhanced,nobeforeafter,width=7.1cm,height=2.1cm,valign=center,colback=cfon!6,colframe=cfon,arc=2.5mm,boxrule=0.9pt,left=3mm]{\Large\bfseries\color{cfon}b)}\quad{\Large $y=2x-1$}\end{tcolorbox}\par\vspace{2.5mm}
\begin{tcolorbox}[enhanced,nobeforeafter,width=7.1cm,height=2.1cm,valign=center,colback=cfon!6,colframe=cfon,arc=2.5mm,boxrule=0.9pt,left=3mm]{\Large\bfseries\color{cfon}c)}\quad{\Large $y=2x+3$}\end{tcolorbox}\hfill\begin{tcolorbox}[enhanced,nobeforeafter,width=7.1cm,height=2.1cm,valign=center,colback=cfon!6,colframe=cfon,arc=2.5mm,boxrule=0.9pt,left=3mm]{\Large\bfseries\color{cfon}d)}\quad{\Large $y=2x+1$}\end{tcolorbox}
\end{tcolorbox}
\newpage
\begin{tcolorbox}[enhanced,width=16cm,height=14cm,colback=white,colframe=cfon,boxrule=2pt,arc=0mm,
left=6mm,right=6mm,top=0mm,bottom=4mm,sharp corners,
overlay={\fill[cfon] (frame.north west) rectangle ([yshift=-1.25cm]frame.north east);
\node[anchor=west,text=white,font=\sffamily\bfseries\LARGE] at ([xshift=5mm,yshift=-0.62cm]frame.north west) {Question 23};
\node[anchor=east,text=white,font=\sffamily\itshape\large] at ([xshift=-5mm,yshift=-0.62cm]frame.north east) {Fonctions};}]
\vspace*{1.75cm}
\begin{minipage}[t][6.9cm][c]{\linewidth}\LARGE\centering
\begin{minipage}{0.95\linewidth}\raggedright\linespread{1.25}\selectfont La courbe de $f(x)=\dfrac{2x+1}{x-3}$ admet une asymptote verticale. Laquelle~?\end{minipage}
\end{minipage}

\vfill
\begin{tcolorbox}[enhanced,nobeforeafter,width=7.1cm,height=2.1cm,valign=center,colback=cfon!6,colframe=cfon,arc=2.5mm,boxrule=0.9pt,left=3mm]{\Large\bfseries\color{cfon}a)}\quad{\Large $y=2$}\end{tcolorbox}\hfill\begin{tcolorbox}[enhanced,nobeforeafter,width=7.1cm,height=2.1cm,valign=center,colback=cfon!6,colframe=cfon,arc=2.5mm,boxrule=0.9pt,left=3mm]{\Large\bfseries\color{cfon}b)}\quad{\Large $x=-3$}\end{tcolorbox}\par\vspace{2.5mm}
\begin{tcolorbox}[enhanced,nobeforeafter,width=7.1cm,height=2.1cm,valign=center,colback=cfon!6,colframe=cfon,arc=2.5mm,boxrule=0.9pt,left=3mm]{\Large\bfseries\color{cfon}c)}\quad{\Large $x=-\frac{1}{2}$}\end{tcolorbox}\hfill\begin{tcolorbox}[enhanced,nobeforeafter,width=7.1cm,height=2.1cm,valign=center,colback=cfon!6,colframe=cfon,arc=2.5mm,boxrule=0.9pt,left=3mm]{\Large\bfseries\color{cfon}d)}\quad{\Large $x=3$}\end{tcolorbox}
\end{tcolorbox}
\newpage
\begin{tcolorbox}[enhanced,width=16cm,height=14cm,colback=white,colframe=cfon,boxrule=2pt,arc=0mm,
left=6mm,right=6mm,top=0mm,bottom=4mm,sharp corners,
overlay={\fill[cfon] (frame.north west) rectangle ([yshift=-1.25cm]frame.north east);
\node[anchor=west,text=white,font=\sffamily\bfseries\LARGE] at ([xshift=5mm,yshift=-0.62cm]frame.north west) {Question 24};
\node[anchor=east,text=white,font=\sffamily\itshape\large] at ([xshift=-5mm,yshift=-0.62cm]frame.north east) {Fonctions};}]
\vspace*{1.75cm}
\begin{minipage}[t][6.9cm][c]{\linewidth}\LARGE\centering
\begin{minipage}{0.95\linewidth}\raggedright\linespread{1.25}\selectfont La courbe de $f(x)=\dfrac{3x-1}{x+2}$ admet une asymptote verticale. Laquelle~?\end{minipage}
\end{minipage}

\vfill
\begin{tcolorbox}[enhanced,nobeforeafter,width=7.1cm,height=2.1cm,valign=center,colback=cfon!6,colframe=cfon,arc=2.5mm,boxrule=0.9pt,left=3mm]{\Large\bfseries\color{cfon}a)}\quad{\Large $x=2$}\end{tcolorbox}\hfill\begin{tcolorbox}[enhanced,nobeforeafter,width=7.1cm,height=2.1cm,valign=center,colback=cfon!6,colframe=cfon,arc=2.5mm,boxrule=0.9pt,left=3mm]{\Large\bfseries\color{cfon}b)}\quad{\Large $y=3$}\end{tcolorbox}\par\vspace{2.5mm}
\begin{tcolorbox}[enhanced,nobeforeafter,width=7.1cm,height=2.1cm,valign=center,colback=cfon!6,colframe=cfon,arc=2.5mm,boxrule=0.9pt,left=3mm]{\Large\bfseries\color{cfon}c)}\quad{\Large $x=\frac{1}{3}$}\end{tcolorbox}\hfill\begin{tcolorbox}[enhanced,nobeforeafter,width=7.1cm,height=2.1cm,valign=center,colback=cfon!6,colframe=cfon,arc=2.5mm,boxrule=0.9pt,left=3mm]{\Large\bfseries\color{cfon}d)}\quad{\Large $x=-2$}\end{tcolorbox}
\end{tcolorbox}
\newpage
\begin{tcolorbox}[enhanced,width=16cm,height=14cm,colback=white,colframe=cfon,boxrule=2pt,arc=0mm,
left=6mm,right=6mm,top=0mm,bottom=4mm,sharp corners,
overlay={\fill[cfon] (frame.north west) rectangle ([yshift=-1.25cm]frame.north east);
\node[anchor=west,text=white,font=\sffamily\bfseries\LARGE] at ([xshift=5mm,yshift=-0.62cm]frame.north west) {Question 25};
\node[anchor=east,text=white,font=\sffamily\itshape\large] at ([xshift=-5mm,yshift=-0.62cm]frame.north east) {Fonctions};}]
\vspace*{1.75cm}
\begin{minipage}[t][6.9cm][c]{\linewidth}\LARGE\centering
\begin{minipage}{0.95\linewidth}\raggedright\linespread{1.25}\selectfont La courbe de $f(x)=x^3+3x^2+x$ admet un point d'inflexion. Quelle est son abscisse~?\end{minipage}
\end{minipage}

\vfill
\begin{tcolorbox}[enhanced,nobeforeafter,width=7.1cm,height=2.1cm,valign=center,colback=cfon!6,colframe=cfon,arc=2.5mm,boxrule=0.9pt,left=3mm]{\Large\bfseries\color{cfon}a)}\quad{\Large $-\dfrac{1}{2}$}\end{tcolorbox}\hfill\begin{tcolorbox}[enhanced,nobeforeafter,width=7.1cm,height=2.1cm,valign=center,colback=cfon!6,colframe=cfon,arc=2.5mm,boxrule=0.9pt,left=3mm]{\Large\bfseries\color{cfon}b)}\quad{\Large $-1$}\end{tcolorbox}\par\vspace{2.5mm}
\begin{tcolorbox}[enhanced,nobeforeafter,width=7.1cm,height=2.1cm,valign=center,colback=cfon!6,colframe=cfon,arc=2.5mm,boxrule=0.9pt,left=3mm]{\Large\bfseries\color{cfon}c)}\quad{\Large $1$}\end{tcolorbox}\hfill\begin{tcolorbox}[enhanced,nobeforeafter,width=7.1cm,height=2.1cm,valign=center,colback=cfon!6,colframe=cfon,arc=2.5mm,boxrule=0.9pt,left=3mm]{\Large\bfseries\color{cfon}d)}\quad{\Large $-\dfrac{3}{2}$}\end{tcolorbox}
\end{tcolorbox}
\newpage
\begin{tcolorbox}[enhanced,width=16cm,height=14cm,colback=white,colframe=cfon,boxrule=2pt,arc=0mm,
left=6mm,right=6mm,top=0mm,bottom=4mm,sharp corners,
overlay={\fill[cfon] (frame.north west) rectangle ([yshift=-1.25cm]frame.north east);
\node[anchor=west,text=white,font=\sffamily\bfseries\LARGE] at ([xshift=5mm,yshift=-0.62cm]frame.north west) {Question 26};
\node[anchor=east,text=white,font=\sffamily\itshape\large] at ([xshift=-5mm,yshift=-0.62cm]frame.north east) {Fonctions};}]
\vspace*{1.75cm}
\begin{minipage}[t][6.9cm][c]{\linewidth}\LARGE\centering
\begin{minipage}{0.95\linewidth}\raggedright\linespread{1.25}\selectfont La courbe de $f(x)=x^3-6x^2+2x$ admet un point d'inflexion. Quelle est son abscisse~?\end{minipage}
\end{minipage}

\vfill
\begin{tcolorbox}[enhanced,nobeforeafter,width=7.1cm,height=2.1cm,valign=center,colback=cfon!6,colframe=cfon,arc=2.5mm,boxrule=0.9pt,left=3mm]{\Large\bfseries\color{cfon}a)}\quad{\Large $-2$}\end{tcolorbox}\hfill\begin{tcolorbox}[enhanced,nobeforeafter,width=7.1cm,height=2.1cm,valign=center,colback=cfon!6,colframe=cfon,arc=2.5mm,boxrule=0.9pt,left=3mm]{\Large\bfseries\color{cfon}b)}\quad{\Large $2$}\end{tcolorbox}\par\vspace{2.5mm}
\begin{tcolorbox}[enhanced,nobeforeafter,width=7.1cm,height=2.1cm,valign=center,colback=cfon!6,colframe=cfon,arc=2.5mm,boxrule=0.9pt,left=3mm]{\Large\bfseries\color{cfon}c)}\quad{\Large $3$}\end{tcolorbox}\hfill\begin{tcolorbox}[enhanced,nobeforeafter,width=7.1cm,height=2.1cm,valign=center,colback=cfon!6,colframe=cfon,arc=2.5mm,boxrule=0.9pt,left=3mm]{\Large\bfseries\color{cfon}d)}\quad{\Large $1$}\end{tcolorbox}
\end{tcolorbox}
\newpage
\begin{tcolorbox}[enhanced,width=16cm,height=14cm,colback=white,colframe=cfon,boxrule=2pt,arc=0mm,
left=6mm,right=6mm,top=0mm,bottom=4mm,sharp corners,
overlay={\fill[cfon] (frame.north west) rectangle ([yshift=-1.25cm]frame.north east);
\node[anchor=west,text=white,font=\sffamily\bfseries\LARGE] at ([xshift=5mm,yshift=-0.62cm]frame.north west) {Question 27};
\node[anchor=east,text=white,font=\sffamily\itshape\large] at ([xshift=-5mm,yshift=-0.62cm]frame.north east) {Fonctions};}]
\vspace*{1.75cm}
\begin{minipage}[t][6.9cm][c]{\linewidth}\LARGE\centering
\begin{minipage}{0.95\linewidth}\raggedright\linespread{1.25}\selectfont Combien de solutions réelles l'équation $x^3-3x=1$ admet-elle~?\end{minipage}
\end{minipage}

\vfill
\begin{tcolorbox}[enhanced,nobeforeafter,width=7.1cm,height=2.1cm,valign=center,colback=cfon!6,colframe=cfon,arc=2.5mm,boxrule=0.9pt,left=3mm]{\Large\bfseries\color{cfon}a)}\quad{\Large $1$}\end{tcolorbox}\hfill\begin{tcolorbox}[enhanced,nobeforeafter,width=7.1cm,height=2.1cm,valign=center,colback=cfon!6,colframe=cfon,arc=2.5mm,boxrule=0.9pt,left=3mm]{\Large\bfseries\color{cfon}b)}\quad{\Large $3$}\end{tcolorbox}\par\vspace{2.5mm}
\begin{tcolorbox}[enhanced,nobeforeafter,width=7.1cm,height=2.1cm,valign=center,colback=cfon!6,colframe=cfon,arc=2.5mm,boxrule=0.9pt,left=3mm]{\Large\bfseries\color{cfon}c)}\quad{\Large $0$}\end{tcolorbox}\hfill\begin{tcolorbox}[enhanced,nobeforeafter,width=7.1cm,height=2.1cm,valign=center,colback=cfon!6,colframe=cfon,arc=2.5mm,boxrule=0.9pt,left=3mm]{\Large\bfseries\color{cfon}d)}\quad{\Large $2$}\end{tcolorbox}
\end{tcolorbox}
\newpage
\begin{tcolorbox}[enhanced,width=16cm,height=14cm,colback=white,colframe=cfon,boxrule=2pt,arc=0mm,
left=6mm,right=6mm,top=0mm,bottom=4mm,sharp corners,
overlay={\fill[cfon] (frame.north west) rectangle ([yshift=-1.25cm]frame.north east);
\node[anchor=west,text=white,font=\sffamily\bfseries\LARGE] at ([xshift=5mm,yshift=-0.62cm]frame.north west) {Question 28};
\node[anchor=east,text=white,font=\sffamily\itshape\large] at ([xshift=-5mm,yshift=-0.62cm]frame.north east) {Fonctions};}]
\vspace*{1.75cm}
\begin{minipage}[t][6.9cm][c]{\linewidth}\LARGE\centering
\begin{minipage}{0.95\linewidth}\raggedright\linespread{1.25}\selectfont Combien de solutions réelles l'équation $x^3-3x=3$ admet-elle~?\end{minipage}
\end{minipage}

\vfill
\begin{tcolorbox}[enhanced,nobeforeafter,width=7.1cm,height=2.1cm,valign=center,colback=cfon!6,colframe=cfon,arc=2.5mm,boxrule=0.9pt,left=3mm]{\Large\bfseries\color{cfon}a)}\quad{\Large $1$}\end{tcolorbox}\hfill\begin{tcolorbox}[enhanced,nobeforeafter,width=7.1cm,height=2.1cm,valign=center,colback=cfon!6,colframe=cfon,arc=2.5mm,boxrule=0.9pt,left=3mm]{\Large\bfseries\color{cfon}b)}\quad{\Large $3$}\end{tcolorbox}\par\vspace{2.5mm}
\begin{tcolorbox}[enhanced,nobeforeafter,width=7.1cm,height=2.1cm,valign=center,colback=cfon!6,colframe=cfon,arc=2.5mm,boxrule=0.9pt,left=3mm]{\Large\bfseries\color{cfon}c)}\quad{\Large $2$}\end{tcolorbox}\hfill\begin{tcolorbox}[enhanced,nobeforeafter,width=7.1cm,height=2.1cm,valign=center,colback=cfon!6,colframe=cfon,arc=2.5mm,boxrule=0.9pt,left=3mm]{\Large\bfseries\color{cfon}d)}\quad{\Large $0$}\end{tcolorbox}
\end{tcolorbox}
\newpage
\begin{tcolorbox}[enhanced,width=16cm,height=14cm,colback=white,colframe=cfon,boxrule=2pt,arc=0mm,
left=6mm,right=6mm,top=0mm,bottom=4mm,sharp corners,
overlay={\fill[cfon] (frame.north west) rectangle ([yshift=-1.25cm]frame.north east);
\node[anchor=west,text=white,font=\sffamily\bfseries\LARGE] at ([xshift=5mm,yshift=-0.62cm]frame.north west) {Question 29};
\node[anchor=east,text=white,font=\sffamily\itshape\large] at ([xshift=-5mm,yshift=-0.62cm]frame.north east) {Fonctions};}]
\vspace*{1.75cm}
\begin{minipage}[t][6.9cm][c]{\linewidth}\LARGE\centering
\begin{minipage}{0.95\linewidth}\raggedright\linespread{1.25}\selectfont Quel est le minimum de la fonction $f(x)=x+\dfrac{4}{x}$ sur $]0;+\infty[$~?\end{minipage}
\end{minipage}

\vfill
\begin{tcolorbox}[enhanced,nobeforeafter,width=7.1cm,height=2.1cm,valign=center,colback=cfon!6,colframe=cfon,arc=2.5mm,boxrule=0.9pt,left=3mm]{\Large\bfseries\color{cfon}a)}\quad{\Large $4$}\end{tcolorbox}\hfill\begin{tcolorbox}[enhanced,nobeforeafter,width=7.1cm,height=2.1cm,valign=center,colback=cfon!6,colframe=cfon,arc=2.5mm,boxrule=0.9pt,left=3mm]{\Large\bfseries\color{cfon}b)}\quad{\Large $8$}\end{tcolorbox}\par\vspace{2.5mm}
\begin{tcolorbox}[enhanced,nobeforeafter,width=7.1cm,height=2.1cm,valign=center,colback=cfon!6,colframe=cfon,arc=2.5mm,boxrule=0.9pt,left=3mm]{\Large\bfseries\color{cfon}c)}\quad{\Large $5$}\end{tcolorbox}\hfill\begin{tcolorbox}[enhanced,nobeforeafter,width=7.1cm,height=2.1cm,valign=center,colback=cfon!6,colframe=cfon,arc=2.5mm,boxrule=0.9pt,left=3mm]{\Large\bfseries\color{cfon}d)}\quad{\Large $2$}\end{tcolorbox}
\end{tcolorbox}
\newpage
\begin{tcolorbox}[enhanced,width=16cm,height=14cm,colback=white,colframe=cfon,boxrule=2pt,arc=0mm,
left=6mm,right=6mm,top=0mm,bottom=4mm,sharp corners,
overlay={\fill[cfon] (frame.north west) rectangle ([yshift=-1.25cm]frame.north east);
\node[anchor=west,text=white,font=\sffamily\bfseries\LARGE] at ([xshift=5mm,yshift=-0.62cm]frame.north west) {Question 30};
\node[anchor=east,text=white,font=\sffamily\itshape\large] at ([xshift=-5mm,yshift=-0.62cm]frame.north east) {Fonctions};}]
\vspace*{1.75cm}
\begin{minipage}[t][6.9cm][c]{\linewidth}\LARGE\centering
\begin{minipage}{0.95\linewidth}\raggedright\linespread{1.25}\selectfont Quel est le minimum de la fonction $f(x)=x+\dfrac{9}{x}$ sur $]0;+\infty[$~?\end{minipage}
\end{minipage}

\vfill
\begin{tcolorbox}[enhanced,nobeforeafter,width=7.1cm,height=2.1cm,valign=center,colback=cfon!6,colframe=cfon,arc=2.5mm,boxrule=0.9pt,left=3mm]{\Large\bfseries\color{cfon}a)}\quad{\Large $9$}\end{tcolorbox}\hfill\begin{tcolorbox}[enhanced,nobeforeafter,width=7.1cm,height=2.1cm,valign=center,colback=cfon!6,colframe=cfon,arc=2.5mm,boxrule=0.9pt,left=3mm]{\Large\bfseries\color{cfon}b)}\quad{\Large $10$}\end{tcolorbox}\par\vspace{2.5mm}
\begin{tcolorbox}[enhanced,nobeforeafter,width=7.1cm,height=2.1cm,valign=center,colback=cfon!6,colframe=cfon,arc=2.5mm,boxrule=0.9pt,left=3mm]{\Large\bfseries\color{cfon}c)}\quad{\Large $6$}\end{tcolorbox}\hfill\begin{tcolorbox}[enhanced,nobeforeafter,width=7.1cm,height=2.1cm,valign=center,colback=cfon!6,colframe=cfon,arc=2.5mm,boxrule=0.9pt,left=3mm]{\Large\bfseries\color{cfon}d)}\quad{\Large $3$}\end{tcolorbox}
\end{tcolorbox}
\newpage
\begin{tcolorbox}[enhanced,width=16cm,height=14cm,colback=white,colframe=cfon,boxrule=2pt,arc=0mm,
left=6mm,right=6mm,top=0mm,bottom=4mm,sharp corners,
overlay={\fill[cfon] (frame.north west) rectangle ([yshift=-1.25cm]frame.north east);
\node[anchor=west,text=white,font=\sffamily\bfseries\LARGE] at ([xshift=5mm,yshift=-0.62cm]frame.north west) {Question 31};
\node[anchor=east,text=white,font=\sffamily\itshape\large] at ([xshift=-5mm,yshift=-0.62cm]frame.north east) {Fonctions};}]
\vspace*{1.75cm}
\begin{minipage}[t][6.9cm][c]{\linewidth}\LARGE\centering
\begin{minipage}{0.95\linewidth}\raggedright\linespread{1.25}\selectfont Soit $f(x)=\sqrt{2x+1}$. Que vaut $f'(4)$~?\end{minipage}
\end{minipage}

\vfill
\begin{tcolorbox}[enhanced,nobeforeafter,width=7.1cm,height=2.1cm,valign=center,colback=cfon!6,colframe=cfon,arc=2.5mm,boxrule=0.9pt,left=3mm]{\Large\bfseries\color{cfon}a)}\quad{\Large $1$}\end{tcolorbox}\hfill\begin{tcolorbox}[enhanced,nobeforeafter,width=7.1cm,height=2.1cm,valign=center,colback=cfon!6,colframe=cfon,arc=2.5mm,boxrule=0.9pt,left=3mm]{\Large\bfseries\color{cfon}b)}\quad{\Large $\dfrac{1}{3}$}\end{tcolorbox}\par\vspace{2.5mm}
\begin{tcolorbox}[enhanced,nobeforeafter,width=7.1cm,height=2.1cm,valign=center,colback=cfon!6,colframe=cfon,arc=2.5mm,boxrule=0.9pt,left=3mm]{\Large\bfseries\color{cfon}c)}\quad{\Large $\dfrac{2}{3}$}\end{tcolorbox}\hfill\begin{tcolorbox}[enhanced,nobeforeafter,width=7.1cm,height=2.1cm,valign=center,colback=cfon!6,colframe=cfon,arc=2.5mm,boxrule=0.9pt,left=3mm]{\Large\bfseries\color{cfon}d)}\quad{\Large $\dfrac{1}{6}$}\end{tcolorbox}
\end{tcolorbox}
\newpage
\begin{tcolorbox}[enhanced,width=16cm,height=14cm,colback=white,colframe=cfon,boxrule=2pt,arc=0mm,
left=6mm,right=6mm,top=0mm,bottom=4mm,sharp corners,
overlay={\fill[cfon] (frame.north west) rectangle ([yshift=-1.25cm]frame.north east);
\node[anchor=west,text=white,font=\sffamily\bfseries\LARGE] at ([xshift=5mm,yshift=-0.62cm]frame.north west) {Question 32};
\node[anchor=east,text=white,font=\sffamily\itshape\large] at ([xshift=-5mm,yshift=-0.62cm]frame.north east) {Fonctions};}]
\vspace*{1.75cm}
\begin{minipage}[t][6.9cm][c]{\linewidth}\LARGE\centering
\begin{minipage}{0.95\linewidth}\raggedright\linespread{1.25}\selectfont Soit $f(x)=\sqrt{3x+1}$. Que vaut $f'(5)$~?\end{minipage}
\end{minipage}

\vfill
\begin{tcolorbox}[enhanced,nobeforeafter,width=7.1cm,height=2.1cm,valign=center,colback=cfon!6,colframe=cfon,arc=2.5mm,boxrule=0.9pt,left=3mm]{\Large\bfseries\color{cfon}a)}\quad{\Large $\dfrac{3}{4}$}\end{tcolorbox}\hfill\begin{tcolorbox}[enhanced,nobeforeafter,width=7.1cm,height=2.1cm,valign=center,colback=cfon!6,colframe=cfon,arc=2.5mm,boxrule=0.9pt,left=3mm]{\Large\bfseries\color{cfon}b)}\quad{\Large $\dfrac{3}{8}$}\end{tcolorbox}\par\vspace{2.5mm}
\begin{tcolorbox}[enhanced,nobeforeafter,width=7.1cm,height=2.1cm,valign=center,colback=cfon!6,colframe=cfon,arc=2.5mm,boxrule=0.9pt,left=3mm]{\Large\bfseries\color{cfon}c)}\quad{\Large $\dfrac{3}{2}$}\end{tcolorbox}\hfill\begin{tcolorbox}[enhanced,nobeforeafter,width=7.1cm,height=2.1cm,valign=center,colback=cfon!6,colframe=cfon,arc=2.5mm,boxrule=0.9pt,left=3mm]{\Large\bfseries\color{cfon}d)}\quad{\Large $\dfrac{1}{8}$}\end{tcolorbox}
\end{tcolorbox}
\newpage
\begin{tcolorbox}[enhanced,width=16cm,height=14cm,colback=white,colframe=cexl,boxrule=2pt,arc=0mm,
left=6mm,right=6mm,top=0mm,bottom=4mm,sharp corners,
overlay={\fill[cexl] (frame.north west) rectangle ([yshift=-1.25cm]frame.north east);
\node[anchor=west,text=white,font=\sffamily\bfseries\LARGE] at ([xshift=5mm,yshift=-0.62cm]frame.north west) {Question 33};
\node[anchor=east,text=white,font=\sffamily\itshape\large] at ([xshift=-5mm,yshift=-0.62cm]frame.north east) {Exponentielle et logarithme};}]
\vspace*{1.75cm}
\begin{minipage}[t][6.9cm][c]{\linewidth}\LARGE\centering
\begin{minipage}{0.95\linewidth}\raggedright\linespread{1.25}\selectfont Simplifier $\dfrac{e^{2}\times e^{3}}{e}$.\end{minipage}
\end{minipage}

\vfill
\begin{tcolorbox}[enhanced,nobeforeafter,width=7.1cm,height=2.1cm,valign=center,colback=cexl!6,colframe=cexl,arc=2.5mm,boxrule=0.9pt,left=3mm]{\Large\bfseries\color{cexl}a)}\quad{\Large $e^{5}$}\end{tcolorbox}\hfill\begin{tcolorbox}[enhanced,nobeforeafter,width=7.1cm,height=2.1cm,valign=center,colback=cexl!6,colframe=cexl,arc=2.5mm,boxrule=0.9pt,left=3mm]{\Large\bfseries\color{cexl}b)}\quad{\Large $e^{4}$}\end{tcolorbox}\par\vspace{2.5mm}
\begin{tcolorbox}[enhanced,nobeforeafter,width=7.1cm,height=2.1cm,valign=center,colback=cexl!6,colframe=cexl,arc=2.5mm,boxrule=0.9pt,left=3mm]{\Large\bfseries\color{cexl}c)}\quad{\Large $e^{6}$}\end{tcolorbox}\hfill\begin{tcolorbox}[enhanced,nobeforeafter,width=7.1cm,height=2.1cm,valign=center,colback=cexl!6,colframe=cexl,arc=2.5mm,boxrule=0.9pt,left=3mm]{\Large\bfseries\color{cexl}d)}\quad{\Large $1$}\end{tcolorbox}
\end{tcolorbox}
\newpage
\begin{tcolorbox}[enhanced,width=16cm,height=14cm,colback=white,colframe=cexl,boxrule=2pt,arc=0mm,
left=6mm,right=6mm,top=0mm,bottom=4mm,sharp corners,
overlay={\fill[cexl] (frame.north west) rectangle ([yshift=-1.25cm]frame.north east);
\node[anchor=west,text=white,font=\sffamily\bfseries\LARGE] at ([xshift=5mm,yshift=-0.62cm]frame.north west) {Question 34};
\node[anchor=east,text=white,font=\sffamily\itshape\large] at ([xshift=-5mm,yshift=-0.62cm]frame.north east) {Exponentielle et logarithme};}]
\vspace*{1.75cm}
\begin{minipage}[t][6.9cm][c]{\linewidth}\LARGE\centering
\begin{minipage}{0.95\linewidth}\raggedright\linespread{1.25}\selectfont Simplifier $\dfrac{e^{5}\times e^{-2}}{e}$.\end{minipage}
\end{minipage}

\vfill
\begin{tcolorbox}[enhanced,nobeforeafter,width=7.1cm,height=2.1cm,valign=center,colback=cexl!6,colframe=cexl,arc=2.5mm,boxrule=0.9pt,left=3mm]{\Large\bfseries\color{cexl}a)}\quad{\Large $e^{4}$}\end{tcolorbox}\hfill\begin{tcolorbox}[enhanced,nobeforeafter,width=7.1cm,height=2.1cm,valign=center,colback=cexl!6,colframe=cexl,arc=2.5mm,boxrule=0.9pt,left=3mm]{\Large\bfseries\color{cexl}b)}\quad{\Large $e^{-11}$}\end{tcolorbox}\par\vspace{2.5mm}
\begin{tcolorbox}[enhanced,nobeforeafter,width=7.1cm,height=2.1cm,valign=center,colback=cexl!6,colframe=cexl,arc=2.5mm,boxrule=0.9pt,left=3mm]{\Large\bfseries\color{cexl}c)}\quad{\Large $e^{8}$}\end{tcolorbox}\hfill\begin{tcolorbox}[enhanced,nobeforeafter,width=7.1cm,height=2.1cm,valign=center,colback=cexl!6,colframe=cexl,arc=2.5mm,boxrule=0.9pt,left=3mm]{\Large\bfseries\color{cexl}d)}\quad{\Large $e^{2}$}\end{tcolorbox}
\end{tcolorbox}
\newpage
\begin{tcolorbox}[enhanced,width=16cm,height=14cm,colback=white,colframe=cexl,boxrule=2pt,arc=0mm,
left=6mm,right=6mm,top=0mm,bottom=4mm,sharp corners,
overlay={\fill[cexl] (frame.north west) rectangle ([yshift=-1.25cm]frame.north east);
\node[anchor=west,text=white,font=\sffamily\bfseries\LARGE] at ([xshift=5mm,yshift=-0.62cm]frame.north west) {Question 35};
\node[anchor=east,text=white,font=\sffamily\itshape\large] at ([xshift=-5mm,yshift=-0.62cm]frame.north east) {Exponentielle et logarithme};}]
\vspace*{1.75cm}
\begin{minipage}[t][6.9cm][c]{\linewidth}\LARGE\centering
\begin{minipage}{0.95\linewidth}\raggedright\linespread{1.25}\selectfont Quelle est la solution de l'équation $\ln(2x+3)=0$~?\end{minipage}
\end{minipage}

\vfill
\begin{tcolorbox}[enhanced,nobeforeafter,width=7.1cm,height=2.1cm,valign=center,colback=cexl!6,colframe=cexl,arc=2.5mm,boxrule=0.9pt,left=3mm]{\Large\bfseries\color{cexl}a)}\quad{\Large $\dfrac{e-3}{2}$}\end{tcolorbox}\hfill\begin{tcolorbox}[enhanced,nobeforeafter,width=7.1cm,height=2.1cm,valign=center,colback=cexl!6,colframe=cexl,arc=2.5mm,boxrule=0.9pt,left=3mm]{\Large\bfseries\color{cexl}b)}\quad{\Large $-\dfrac{3}{2}$}\end{tcolorbox}\par\vspace{2.5mm}
\begin{tcolorbox}[enhanced,nobeforeafter,width=7.1cm,height=2.1cm,valign=center,colback=cexl!6,colframe=cexl,arc=2.5mm,boxrule=0.9pt,left=3mm]{\Large\bfseries\color{cexl}c)}\quad{\Large $2$}\end{tcolorbox}\hfill\begin{tcolorbox}[enhanced,nobeforeafter,width=7.1cm,height=2.1cm,valign=center,colback=cexl!6,colframe=cexl,arc=2.5mm,boxrule=0.9pt,left=3mm]{\Large\bfseries\color{cexl}d)}\quad{\Large $-1$}\end{tcolorbox}
\end{tcolorbox}
\newpage
\begin{tcolorbox}[enhanced,width=16cm,height=14cm,colback=white,colframe=cexl,boxrule=2pt,arc=0mm,
left=6mm,right=6mm,top=0mm,bottom=4mm,sharp corners,
overlay={\fill[cexl] (frame.north west) rectangle ([yshift=-1.25cm]frame.north east);
\node[anchor=west,text=white,font=\sffamily\bfseries\LARGE] at ([xshift=5mm,yshift=-0.62cm]frame.north west) {Question 36};
\node[anchor=east,text=white,font=\sffamily\itshape\large] at ([xshift=-5mm,yshift=-0.62cm]frame.north east) {Exponentielle et logarithme};}]
\vspace*{1.75cm}
\begin{minipage}[t][6.9cm][c]{\linewidth}\LARGE\centering
\begin{minipage}{0.95\linewidth}\raggedright\linespread{1.25}\selectfont Quelle est la solution de l'équation $\ln(3x-2)=0$~?\end{minipage}
\end{minipage}

\vfill
\begin{tcolorbox}[enhanced,nobeforeafter,width=7.1cm,height=2.1cm,valign=center,colback=cexl!6,colframe=cexl,arc=2.5mm,boxrule=0.9pt,left=3mm]{\Large\bfseries\color{cexl}a)}\quad{\Large $-\dfrac{1}{3}$}\end{tcolorbox}\hfill\begin{tcolorbox}[enhanced,nobeforeafter,width=7.1cm,height=2.1cm,valign=center,colback=cexl!6,colframe=cexl,arc=2.5mm,boxrule=0.9pt,left=3mm]{\Large\bfseries\color{cexl}b)}\quad{\Large $\dfrac{2}{3}$}\end{tcolorbox}\par\vspace{2.5mm}
\begin{tcolorbox}[enhanced,nobeforeafter,width=7.1cm,height=2.1cm,valign=center,colback=cexl!6,colframe=cexl,arc=2.5mm,boxrule=0.9pt,left=3mm]{\Large\bfseries\color{cexl}c)}\quad{\Large $1$}\end{tcolorbox}\hfill\begin{tcolorbox}[enhanced,nobeforeafter,width=7.1cm,height=2.1cm,valign=center,colback=cexl!6,colframe=cexl,arc=2.5mm,boxrule=0.9pt,left=3mm]{\Large\bfseries\color{cexl}d)}\quad{\Large $\dfrac{e+2}{3}$}\end{tcolorbox}
\end{tcolorbox}
\newpage
\begin{tcolorbox}[enhanced,width=16cm,height=14cm,colback=white,colframe=cexl,boxrule=2pt,arc=0mm,
left=6mm,right=6mm,top=0mm,bottom=4mm,sharp corners,
overlay={\fill[cexl] (frame.north west) rectangle ([yshift=-1.25cm]frame.north east);
\node[anchor=west,text=white,font=\sffamily\bfseries\LARGE] at ([xshift=5mm,yshift=-0.62cm]frame.north west) {Question 37};
\node[anchor=east,text=white,font=\sffamily\itshape\large] at ([xshift=-5mm,yshift=-0.62cm]frame.north east) {Exponentielle et logarithme};}]
\vspace*{1.75cm}
\begin{minipage}[t][6.9cm][c]{\linewidth}\LARGE\centering
\begin{minipage}{0.95\linewidth}\raggedright\linespread{1.25}\selectfont Simplifier $\ln 8-2\ln 2$.\end{minipage}
\end{minipage}

\vfill
\begin{tcolorbox}[enhanced,nobeforeafter,width=7.1cm,height=2.1cm,valign=center,colback=cexl!6,colframe=cexl,arc=2.5mm,boxrule=0.9pt,left=3mm]{\Large\bfseries\color{cexl}a)}\quad{\Large $5\ln 2$}\end{tcolorbox}\hfill\begin{tcolorbox}[enhanced,nobeforeafter,width=7.1cm,height=2.1cm,valign=center,colback=cexl!6,colframe=cexl,arc=2.5mm,boxrule=0.9pt,left=3mm]{\Large\bfseries\color{cexl}b)}\quad{\Large $\ln 2$}\end{tcolorbox}\par\vspace{2.5mm}
\begin{tcolorbox}[enhanced,nobeforeafter,width=7.1cm,height=2.1cm,valign=center,colback=cexl!6,colframe=cexl,arc=2.5mm,boxrule=0.9pt,left=3mm]{\Large\bfseries\color{cexl}c)}\quad{\Large $2\ln 2$}\end{tcolorbox}\hfill\begin{tcolorbox}[enhanced,nobeforeafter,width=7.1cm,height=2.1cm,valign=center,colback=cexl!6,colframe=cexl,arc=2.5mm,boxrule=0.9pt,left=3mm]{\Large\bfseries\color{cexl}d)}\quad{\Large $6\ln 2$}\end{tcolorbox}
\end{tcolorbox}
\newpage
\begin{tcolorbox}[enhanced,width=16cm,height=14cm,colback=white,colframe=cexl,boxrule=2pt,arc=0mm,
left=6mm,right=6mm,top=0mm,bottom=4mm,sharp corners,
overlay={\fill[cexl] (frame.north west) rectangle ([yshift=-1.25cm]frame.north east);
\node[anchor=west,text=white,font=\sffamily\bfseries\LARGE] at ([xshift=5mm,yshift=-0.62cm]frame.north west) {Question 38};
\node[anchor=east,text=white,font=\sffamily\itshape\large] at ([xshift=-5mm,yshift=-0.62cm]frame.north east) {Exponentielle et logarithme};}]
\vspace*{1.75cm}
\begin{minipage}[t][6.9cm][c]{\linewidth}\LARGE\centering
\begin{minipage}{0.95\linewidth}\raggedright\linespread{1.25}\selectfont Simplifier $\ln 81-2\ln 3$.\end{minipage}
\end{minipage}

\vfill
\begin{tcolorbox}[enhanced,nobeforeafter,width=7.1cm,height=2.1cm,valign=center,colback=cexl!6,colframe=cexl,arc=2.5mm,boxrule=0.9pt,left=3mm]{\Large\bfseries\color{cexl}a)}\quad{\Large $3\ln 3$}\end{tcolorbox}\hfill\begin{tcolorbox}[enhanced,nobeforeafter,width=7.1cm,height=2.1cm,valign=center,colback=cexl!6,colframe=cexl,arc=2.5mm,boxrule=0.9pt,left=3mm]{\Large\bfseries\color{cexl}b)}\quad{\Large $6\ln 3$}\end{tcolorbox}\par\vspace{2.5mm}
\begin{tcolorbox}[enhanced,nobeforeafter,width=7.1cm,height=2.1cm,valign=center,colback=cexl!6,colframe=cexl,arc=2.5mm,boxrule=0.9pt,left=3mm]{\Large\bfseries\color{cexl}c)}\quad{\Large $2\ln 3$}\end{tcolorbox}\hfill\begin{tcolorbox}[enhanced,nobeforeafter,width=7.1cm,height=2.1cm,valign=center,colback=cexl!6,colframe=cexl,arc=2.5mm,boxrule=0.9pt,left=3mm]{\Large\bfseries\color{cexl}d)}\quad{\Large $8\ln 3$}\end{tcolorbox}
\end{tcolorbox}
\newpage
\begin{tcolorbox}[enhanced,width=16cm,height=14cm,colback=white,colframe=cexl,boxrule=2pt,arc=0mm,
left=6mm,right=6mm,top=0mm,bottom=4mm,sharp corners,
overlay={\fill[cexl] (frame.north west) rectangle ([yshift=-1.25cm]frame.north east);
\node[anchor=west,text=white,font=\sffamily\bfseries\LARGE] at ([xshift=5mm,yshift=-0.62cm]frame.north west) {Question 39};
\node[anchor=east,text=white,font=\sffamily\itshape\large] at ([xshift=-5mm,yshift=-0.62cm]frame.north east) {Exponentielle et logarithme};}]
\vspace*{1.75cm}
\begin{minipage}[t][6.9cm][c]{\linewidth}\LARGE\centering
\begin{minipage}{0.95\linewidth}\raggedright\linespread{1.25}\selectfont Soit $f(x)=x\,e^{x}$. Que vaut $f'(1)$~?\end{minipage}
\end{minipage}

\vfill
\begin{tcolorbox}[enhanced,nobeforeafter,width=7.1cm,height=2.1cm,valign=center,colback=cexl!6,colframe=cexl,arc=2.5mm,boxrule=0.9pt,left=3mm]{\Large\bfseries\color{cexl}a)}\quad{\Large $e$}\end{tcolorbox}\hfill\begin{tcolorbox}[enhanced,nobeforeafter,width=7.1cm,height=2.1cm,valign=center,colback=cexl!6,colframe=cexl,arc=2.5mm,boxrule=0.9pt,left=3mm]{\Large\bfseries\color{cexl}b)}\quad{\Large $2e^{2}$}\end{tcolorbox}\par\vspace{2.5mm}
\begin{tcolorbox}[enhanced,nobeforeafter,width=7.1cm,height=2.1cm,valign=center,colback=cexl!6,colframe=cexl,arc=2.5mm,boxrule=0.9pt,left=3mm]{\Large\bfseries\color{cexl}c)}\quad{\Large $2e$}\end{tcolorbox}\hfill\begin{tcolorbox}[enhanced,nobeforeafter,width=7.1cm,height=2.1cm,valign=center,colback=cexl!6,colframe=cexl,arc=2.5mm,boxrule=0.9pt,left=3mm]{\Large\bfseries\color{cexl}d)}\quad{\Large $0$}\end{tcolorbox}
\end{tcolorbox}
\newpage
\begin{tcolorbox}[enhanced,width=16cm,height=14cm,colback=white,colframe=cexl,boxrule=2pt,arc=0mm,
left=6mm,right=6mm,top=0mm,bottom=4mm,sharp corners,
overlay={\fill[cexl] (frame.north west) rectangle ([yshift=-1.25cm]frame.north east);
\node[anchor=west,text=white,font=\sffamily\bfseries\LARGE] at ([xshift=5mm,yshift=-0.62cm]frame.north west) {Question 40};
\node[anchor=east,text=white,font=\sffamily\itshape\large] at ([xshift=-5mm,yshift=-0.62cm]frame.north east) {Exponentielle et logarithme};}]
\vspace*{1.75cm}
\begin{minipage}[t][6.9cm][c]{\linewidth}\LARGE\centering
\begin{minipage}{0.95\linewidth}\raggedright\linespread{1.25}\selectfont Soit $f(x)=x\,e^{2x}$. Que vaut $f'(1)$~?\end{minipage}
\end{minipage}

\vfill
\begin{tcolorbox}[enhanced,nobeforeafter,width=7.1cm,height=2.1cm,valign=center,colback=cexl!6,colframe=cexl,arc=2.5mm,boxrule=0.9pt,left=3mm]{\Large\bfseries\color{cexl}a)}\quad{\Large $e^{2}$}\end{tcolorbox}\hfill\begin{tcolorbox}[enhanced,nobeforeafter,width=7.1cm,height=2.1cm,valign=center,colback=cexl!6,colframe=cexl,arc=2.5mm,boxrule=0.9pt,left=3mm]{\Large\bfseries\color{cexl}b)}\quad{\Large $2e^{2}$}\end{tcolorbox}\par\vspace{2.5mm}
\begin{tcolorbox}[enhanced,nobeforeafter,width=7.1cm,height=2.1cm,valign=center,colback=cexl!6,colframe=cexl,arc=2.5mm,boxrule=0.9pt,left=3mm]{\Large\bfseries\color{cexl}c)}\quad{\Large $3e^{2}$}\end{tcolorbox}\hfill\begin{tcolorbox}[enhanced,nobeforeafter,width=7.1cm,height=2.1cm,valign=center,colback=cexl!6,colframe=cexl,arc=2.5mm,boxrule=0.9pt,left=3mm]{\Large\bfseries\color{cexl}d)}\quad{\Large $-e^{2}$}\end{tcolorbox}
\end{tcolorbox}
\newpage
\begin{tcolorbox}[enhanced,width=16cm,height=14cm,colback=white,colframe=cexl,boxrule=2pt,arc=0mm,
left=6mm,right=6mm,top=0mm,bottom=4mm,sharp corners,
overlay={\fill[cexl] (frame.north west) rectangle ([yshift=-1.25cm]frame.north east);
\node[anchor=west,text=white,font=\sffamily\bfseries\LARGE] at ([xshift=5mm,yshift=-0.62cm]frame.north west) {Question 41};
\node[anchor=east,text=white,font=\sffamily\itshape\large] at ([xshift=-5mm,yshift=-0.62cm]frame.north east) {Exponentielle et logarithme};}]
\vspace*{1.75cm}
\begin{minipage}[t][6.9cm][c]{\linewidth}\LARGE\centering
\begin{minipage}{0.95\linewidth}\raggedright\linespread{1.25}\selectfont Que vaut $\displaystyle\lim_{x\to+\infty}\frac{2e^x+x^2}{3e^x+1}$~?\end{minipage}
\end{minipage}

\vfill
\begin{tcolorbox}[enhanced,nobeforeafter,width=7.1cm,height=2.1cm,valign=center,colback=cexl!6,colframe=cexl,arc=2.5mm,boxrule=0.9pt,left=3mm]{\Large\bfseries\color{cexl}a)}\quad{\Large $\dfrac{2}{3}$}\end{tcolorbox}\hfill\begin{tcolorbox}[enhanced,nobeforeafter,width=7.1cm,height=2.1cm,valign=center,colback=cexl!6,colframe=cexl,arc=2.5mm,boxrule=0.9pt,left=3mm]{\Large\bfseries\color{cexl}b)}\quad{\Large $+\infty$}\end{tcolorbox}\par\vspace{2.5mm}
\begin{tcolorbox}[enhanced,nobeforeafter,width=7.1cm,height=2.1cm,valign=center,colback=cexl!6,colframe=cexl,arc=2.5mm,boxrule=0.9pt,left=3mm]{\Large\bfseries\color{cexl}c)}\quad{\Large $0$}\end{tcolorbox}\hfill\begin{tcolorbox}[enhanced,nobeforeafter,width=7.1cm,height=2.1cm,valign=center,colback=cexl!6,colframe=cexl,arc=2.5mm,boxrule=0.9pt,left=3mm]{\Large\bfseries\color{cexl}d)}\quad{\Large $\dfrac{3}{2}$}\end{tcolorbox}
\end{tcolorbox}
\newpage
\begin{tcolorbox}[enhanced,width=16cm,height=14cm,colback=white,colframe=cexl,boxrule=2pt,arc=0mm,
left=6mm,right=6mm,top=0mm,bottom=4mm,sharp corners,
overlay={\fill[cexl] (frame.north west) rectangle ([yshift=-1.25cm]frame.north east);
\node[anchor=west,text=white,font=\sffamily\bfseries\LARGE] at ([xshift=5mm,yshift=-0.62cm]frame.north west) {Question 42};
\node[anchor=east,text=white,font=\sffamily\itshape\large] at ([xshift=-5mm,yshift=-0.62cm]frame.north east) {Exponentielle et logarithme};}]
\vspace*{1.75cm}
\begin{minipage}[t][6.9cm][c]{\linewidth}\LARGE\centering
\begin{minipage}{0.95\linewidth}\raggedright\linespread{1.25}\selectfont Que vaut $\displaystyle\lim_{x\to+\infty}\frac{5e^x+x^2}{2e^x+1}$~?\end{minipage}
\end{minipage}

\vfill
\begin{tcolorbox}[enhanced,nobeforeafter,width=7.1cm,height=2.1cm,valign=center,colback=cexl!6,colframe=cexl,arc=2.5mm,boxrule=0.9pt,left=3mm]{\Large\bfseries\color{cexl}a)}\quad{\Large $\dfrac{5}{2}$}\end{tcolorbox}\hfill\begin{tcolorbox}[enhanced,nobeforeafter,width=7.1cm,height=2.1cm,valign=center,colback=cexl!6,colframe=cexl,arc=2.5mm,boxrule=0.9pt,left=3mm]{\Large\bfseries\color{cexl}b)}\quad{\Large $\dfrac{2}{5}$}\end{tcolorbox}\par\vspace{2.5mm}
\begin{tcolorbox}[enhanced,nobeforeafter,width=7.1cm,height=2.1cm,valign=center,colback=cexl!6,colframe=cexl,arc=2.5mm,boxrule=0.9pt,left=3mm]{\Large\bfseries\color{cexl}c)}\quad{\Large $0$}\end{tcolorbox}\hfill\begin{tcolorbox}[enhanced,nobeforeafter,width=7.1cm,height=2.1cm,valign=center,colback=cexl!6,colframe=cexl,arc=2.5mm,boxrule=0.9pt,left=3mm]{\Large\bfseries\color{cexl}d)}\quad{\Large $+\infty$}\end{tcolorbox}
\end{tcolorbox}
\newpage
\begin{tcolorbox}[enhanced,width=16cm,height=14cm,colback=white,colframe=cexl,boxrule=2pt,arc=0mm,
left=6mm,right=6mm,top=0mm,bottom=4mm,sharp corners,
overlay={\fill[cexl] (frame.north west) rectangle ([yshift=-1.25cm]frame.north east);
\node[anchor=west,text=white,font=\sffamily\bfseries\LARGE] at ([xshift=5mm,yshift=-0.62cm]frame.north west) {Question 43};
\node[anchor=east,text=white,font=\sffamily\itshape\large] at ([xshift=-5mm,yshift=-0.62cm]frame.north east) {Exponentielle et logarithme};}]
\vspace*{1.75cm}
\begin{minipage}[t][6.9cm][c]{\linewidth}\LARGE\centering
\begin{minipage}{0.95\linewidth}\raggedright\linespread{1.25}\selectfont Soit $f(x)=\ln(x^2+1)$. Que vaut $f'(1)$~?\end{minipage}
\end{minipage}

\vfill
\begin{tcolorbox}[enhanced,nobeforeafter,width=7.1cm,height=2.1cm,valign=center,colback=cexl!6,colframe=cexl,arc=2.5mm,boxrule=0.9pt,left=3mm]{\Large\bfseries\color{cexl}a)}\quad{\Large $\dfrac{1}{2}$}\end{tcolorbox}\hfill\begin{tcolorbox}[enhanced,nobeforeafter,width=7.1cm,height=2.1cm,valign=center,colback=cexl!6,colframe=cexl,arc=2.5mm,boxrule=0.9pt,left=3mm]{\Large\bfseries\color{cexl}b)}\quad{\Large $\dfrac{3}{2}$}\end{tcolorbox}\par\vspace{2.5mm}
\begin{tcolorbox}[enhanced,nobeforeafter,width=7.1cm,height=2.1cm,valign=center,colback=cexl!6,colframe=cexl,arc=2.5mm,boxrule=0.9pt,left=3mm]{\Large\bfseries\color{cexl}c)}\quad{\Large $1$}\end{tcolorbox}\hfill\begin{tcolorbox}[enhanced,nobeforeafter,width=7.1cm,height=2.1cm,valign=center,colback=cexl!6,colframe=cexl,arc=2.5mm,boxrule=0.9pt,left=3mm]{\Large\bfseries\color{cexl}d)}\quad{\Large $2$}\end{tcolorbox}
\end{tcolorbox}
\newpage
\begin{tcolorbox}[enhanced,width=16cm,height=14cm,colback=white,colframe=cexl,boxrule=2pt,arc=0mm,
left=6mm,right=6mm,top=0mm,bottom=4mm,sharp corners,
overlay={\fill[cexl] (frame.north west) rectangle ([yshift=-1.25cm]frame.north east);
\node[anchor=west,text=white,font=\sffamily\bfseries\LARGE] at ([xshift=5mm,yshift=-0.62cm]frame.north west) {Question 44};
\node[anchor=east,text=white,font=\sffamily\itshape\large] at ([xshift=-5mm,yshift=-0.62cm]frame.north east) {Exponentielle et logarithme};}]
\vspace*{1.75cm}
\begin{minipage}[t][6.9cm][c]{\linewidth}\LARGE\centering
\begin{minipage}{0.95\linewidth}\raggedright\linespread{1.25}\selectfont Soit $f(x)=\ln(2x^2+3)$. Que vaut $f'(1)$~?\end{minipage}
\end{minipage}

\vfill
\begin{tcolorbox}[enhanced,nobeforeafter,width=7.1cm,height=2.1cm,valign=center,colback=cexl!6,colframe=cexl,arc=2.5mm,boxrule=0.9pt,left=3mm]{\Large\bfseries\color{cexl}a)}\quad{\Large $\dfrac{4}{5}$}\end{tcolorbox}\hfill\begin{tcolorbox}[enhanced,nobeforeafter,width=7.1cm,height=2.1cm,valign=center,colback=cexl!6,colframe=cexl,arc=2.5mm,boxrule=0.9pt,left=3mm]{\Large\bfseries\color{cexl}b)}\quad{\Large $\dfrac{2}{5}$}\end{tcolorbox}\par\vspace{2.5mm}
\begin{tcolorbox}[enhanced,nobeforeafter,width=7.1cm,height=2.1cm,valign=center,colback=cexl!6,colframe=cexl,arc=2.5mm,boxrule=0.9pt,left=3mm]{\Large\bfseries\color{cexl}c)}\quad{\Large $4$}\end{tcolorbox}\hfill\begin{tcolorbox}[enhanced,nobeforeafter,width=7.1cm,height=2.1cm,valign=center,colback=cexl!6,colframe=cexl,arc=2.5mm,boxrule=0.9pt,left=3mm]{\Large\bfseries\color{cexl}d)}\quad{\Large $\dfrac{1}{5}$}\end{tcolorbox}
\end{tcolorbox}
\newpage
\begin{tcolorbox}[enhanced,width=16cm,height=14cm,colback=white,colframe=cexl,boxrule=2pt,arc=0mm,
left=6mm,right=6mm,top=0mm,bottom=4mm,sharp corners,
overlay={\fill[cexl] (frame.north west) rectangle ([yshift=-1.25cm]frame.north east);
\node[anchor=west,text=white,font=\sffamily\bfseries\LARGE] at ([xshift=5mm,yshift=-0.62cm]frame.north west) {Question 45};
\node[anchor=east,text=white,font=\sffamily\itshape\large] at ([xshift=-5mm,yshift=-0.62cm]frame.north east) {Exponentielle et logarithme};}]
\vspace*{1.75cm}
\begin{minipage}[t][6.9cm][c]{\linewidth}\LARGE\centering
\begin{minipage}{0.95\linewidth}\raggedright\linespread{1.25}\selectfont Quel est le plus petit entier naturel $n$ tel que $0{,}9^n<0{,}5$~?\end{minipage}
\end{minipage}

\vfill
\begin{tcolorbox}[enhanced,nobeforeafter,width=7.1cm,height=2.1cm,valign=center,colback=cexl!6,colframe=cexl,arc=2.5mm,boxrule=0.9pt,left=3mm]{\Large\bfseries\color{cexl}a)}\quad{\Large $7$}\end{tcolorbox}\hfill\begin{tcolorbox}[enhanced,nobeforeafter,width=7.1cm,height=2.1cm,valign=center,colback=cexl!6,colframe=cexl,arc=2.5mm,boxrule=0.9pt,left=3mm]{\Large\bfseries\color{cexl}b)}\quad{\Large $6$}\end{tcolorbox}\par\vspace{2.5mm}
\begin{tcolorbox}[enhanced,nobeforeafter,width=7.1cm,height=2.1cm,valign=center,colback=cexl!6,colframe=cexl,arc=2.5mm,boxrule=0.9pt,left=3mm]{\Large\bfseries\color{cexl}c)}\quad{\Large $8$}\end{tcolorbox}\hfill\begin{tcolorbox}[enhanced,nobeforeafter,width=7.1cm,height=2.1cm,valign=center,colback=cexl!6,colframe=cexl,arc=2.5mm,boxrule=0.9pt,left=3mm]{\Large\bfseries\color{cexl}d)}\quad{\Large $9$}\end{tcolorbox}
\end{tcolorbox}
\newpage
\begin{tcolorbox}[enhanced,width=16cm,height=14cm,colback=white,colframe=cexl,boxrule=2pt,arc=0mm,
left=6mm,right=6mm,top=0mm,bottom=4mm,sharp corners,
overlay={\fill[cexl] (frame.north west) rectangle ([yshift=-1.25cm]frame.north east);
\node[anchor=west,text=white,font=\sffamily\bfseries\LARGE] at ([xshift=5mm,yshift=-0.62cm]frame.north west) {Question 46};
\node[anchor=east,text=white,font=\sffamily\itshape\large] at ([xshift=-5mm,yshift=-0.62cm]frame.north east) {Exponentielle et logarithme};}]
\vspace*{1.75cm}
\begin{minipage}[t][6.9cm][c]{\linewidth}\LARGE\centering
\begin{minipage}{0.95\linewidth}\raggedright\linespread{1.25}\selectfont Quel est le plus petit entier naturel $n$ tel que $0{,}8^n<0{,}1$~?\end{minipage}
\end{minipage}

\vfill
\begin{tcolorbox}[enhanced,nobeforeafter,width=7.1cm,height=2.1cm,valign=center,colback=cexl!6,colframe=cexl,arc=2.5mm,boxrule=0.9pt,left=3mm]{\Large\bfseries\color{cexl}a)}\quad{\Large $12$}\end{tcolorbox}\hfill\begin{tcolorbox}[enhanced,nobeforeafter,width=7.1cm,height=2.1cm,valign=center,colback=cexl!6,colframe=cexl,arc=2.5mm,boxrule=0.9pt,left=3mm]{\Large\bfseries\color{cexl}b)}\quad{\Large $13$}\end{tcolorbox}\par\vspace{2.5mm}
\begin{tcolorbox}[enhanced,nobeforeafter,width=7.1cm,height=2.1cm,valign=center,colback=cexl!6,colframe=cexl,arc=2.5mm,boxrule=0.9pt,left=3mm]{\Large\bfseries\color{cexl}c)}\quad{\Large $10$}\end{tcolorbox}\hfill\begin{tcolorbox}[enhanced,nobeforeafter,width=7.1cm,height=2.1cm,valign=center,colback=cexl!6,colframe=cexl,arc=2.5mm,boxrule=0.9pt,left=3mm]{\Large\bfseries\color{cexl}d)}\quad{\Large $11$}\end{tcolorbox}
\end{tcolorbox}
\newpage
\begin{tcolorbox}[enhanced,width=16cm,height=14cm,colback=white,colframe=cexl,boxrule=2pt,arc=0mm,
left=6mm,right=6mm,top=0mm,bottom=4mm,sharp corners,
overlay={\fill[cexl] (frame.north west) rectangle ([yshift=-1.25cm]frame.north east);
\node[anchor=west,text=white,font=\sffamily\bfseries\LARGE] at ([xshift=5mm,yshift=-0.62cm]frame.north west) {Question 47};
\node[anchor=east,text=white,font=\sffamily\itshape\large] at ([xshift=-5mm,yshift=-0.62cm]frame.north east) {Exponentielle et logarithme};}]
\vspace*{1.75cm}
\begin{minipage}[t][6.9cm][c]{\linewidth}\LARGE\centering
\begin{minipage}{0.95\linewidth}\raggedright\linespread{1.25}\selectfont Quelle est la solution de l'équation $e^{2x+1}=3$~?\end{minipage}
\end{minipage}

\vfill
\begin{tcolorbox}[enhanced,nobeforeafter,width=7.1cm,height=2.1cm,valign=center,colback=cexl!6,colframe=cexl,arc=2.5mm,boxrule=0.9pt,left=3mm]{\Large\bfseries\color{cexl}a)}\quad{\Large $\dfrac{\ln 3}{2}-1$}\end{tcolorbox}\hfill\begin{tcolorbox}[enhanced,nobeforeafter,width=7.1cm,height=2.1cm,valign=center,colback=cexl!6,colframe=cexl,arc=2.5mm,boxrule=0.9pt,left=3mm]{\Large\bfseries\color{cexl}b)}\quad{\Large $\dfrac{\ln 3+1}{2}$}\end{tcolorbox}\par\vspace{2.5mm}
\begin{tcolorbox}[enhanced,nobeforeafter,width=7.1cm,height=2.1cm,valign=center,colback=cexl!6,colframe=cexl,arc=2.5mm,boxrule=0.9pt,left=3mm]{\Large\bfseries\color{cexl}c)}\quad{\Large $\dfrac{e^{3}}{2}$}\end{tcolorbox}\hfill\begin{tcolorbox}[enhanced,nobeforeafter,width=7.1cm,height=2.1cm,valign=center,colback=cexl!6,colframe=cexl,arc=2.5mm,boxrule=0.9pt,left=3mm]{\Large\bfseries\color{cexl}d)}\quad{\Large $\dfrac{\ln 3-1}{2}$}\end{tcolorbox}
\end{tcolorbox}
\newpage
\begin{tcolorbox}[enhanced,width=16cm,height=14cm,colback=white,colframe=cexl,boxrule=2pt,arc=0mm,
left=6mm,right=6mm,top=0mm,bottom=4mm,sharp corners,
overlay={\fill[cexl] (frame.north west) rectangle ([yshift=-1.25cm]frame.north east);
\node[anchor=west,text=white,font=\sffamily\bfseries\LARGE] at ([xshift=5mm,yshift=-0.62cm]frame.north west) {Question 48};
\node[anchor=east,text=white,font=\sffamily\itshape\large] at ([xshift=-5mm,yshift=-0.62cm]frame.north east) {Exponentielle et logarithme};}]
\vspace*{1.75cm}
\begin{minipage}[t][6.9cm][c]{\linewidth}\LARGE\centering
\begin{minipage}{0.95\linewidth}\raggedright\linespread{1.25}\selectfont Quelle est la solution de l'équation $e^{3x-2}=5$~?\end{minipage}
\end{minipage}

\vfill
\begin{tcolorbox}[enhanced,nobeforeafter,width=7.1cm,height=2.1cm,valign=center,colback=cexl!6,colframe=cexl,arc=2.5mm,boxrule=0.9pt,left=3mm]{\Large\bfseries\color{cexl}a)}\quad{\Large $\dfrac{\ln 5}{3}+2$}\end{tcolorbox}\hfill\begin{tcolorbox}[enhanced,nobeforeafter,width=7.1cm,height=2.1cm,valign=center,colback=cexl!6,colframe=cexl,arc=2.5mm,boxrule=0.9pt,left=3mm]{\Large\bfseries\color{cexl}b)}\quad{\Large $\dfrac{\ln 5-2}{3}$}\end{tcolorbox}\par\vspace{2.5mm}
\begin{tcolorbox}[enhanced,nobeforeafter,width=7.1cm,height=2.1cm,valign=center,colback=cexl!6,colframe=cexl,arc=2.5mm,boxrule=0.9pt,left=3mm]{\Large\bfseries\color{cexl}c)}\quad{\Large $\dfrac{\ln 5+2}{3}$}\end{tcolorbox}\hfill\begin{tcolorbox}[enhanced,nobeforeafter,width=7.1cm,height=2.1cm,valign=center,colback=cexl!6,colframe=cexl,arc=2.5mm,boxrule=0.9pt,left=3mm]{\Large\bfseries\color{cexl}d)}\quad{\Large $\dfrac{e^{5}}{3}$}\end{tcolorbox}
\end{tcolorbox}
\newpage
\begin{tcolorbox}[enhanced,width=16cm,height=14cm,colback=white,colframe=cint,boxrule=2pt,arc=0mm,
left=6mm,right=6mm,top=0mm,bottom=4mm,sharp corners,
overlay={\fill[cint] (frame.north west) rectangle ([yshift=-1.25cm]frame.north east);
\node[anchor=west,text=white,font=\sffamily\bfseries\LARGE] at ([xshift=5mm,yshift=-0.62cm]frame.north west) {Question 49};
\node[anchor=east,text=white,font=\sffamily\itshape\large] at ([xshift=-5mm,yshift=-0.62cm]frame.north east) {Intégrales et équations différentielles};}]
\vspace*{1.75cm}
\begin{minipage}[t][6.9cm][c]{\linewidth}\LARGE\centering
\begin{minipage}{0.95\linewidth}\raggedright\linespread{1.25}\selectfont Que vaut $\displaystyle\int_0^{2} x^{2}\,\mathrm{d}x$~?\end{minipage}
\end{minipage}

\vfill
\begin{tcolorbox}[enhanced,nobeforeafter,width=7.1cm,height=2.1cm,valign=center,colback=cint!6,colframe=cint,arc=2.5mm,boxrule=0.9pt,left=3mm]{\Large\bfseries\color{cint}a)}\quad{\Large $4$}\end{tcolorbox}\hfill\begin{tcolorbox}[enhanced,nobeforeafter,width=7.1cm,height=2.1cm,valign=center,colback=cint!6,colframe=cint,arc=2.5mm,boxrule=0.9pt,left=3mm]{\Large\bfseries\color{cint}b)}\quad{\Large $8$}\end{tcolorbox}\par\vspace{2.5mm}
\begin{tcolorbox}[enhanced,nobeforeafter,width=7.1cm,height=2.1cm,valign=center,colback=cint!6,colframe=cint,arc=2.5mm,boxrule=0.9pt,left=3mm]{\Large\bfseries\color{cint}c)}\quad{\Large $\dfrac{8}{3}$}\end{tcolorbox}\hfill\begin{tcolorbox}[enhanced,nobeforeafter,width=7.1cm,height=2.1cm,valign=center,colback=cint!6,colframe=cint,arc=2.5mm,boxrule=0.9pt,left=3mm]{\Large\bfseries\color{cint}d)}\quad{\Large $2$}\end{tcolorbox}
\end{tcolorbox}
\newpage
\begin{tcolorbox}[enhanced,width=16cm,height=14cm,colback=white,colframe=cint,boxrule=2pt,arc=0mm,
left=6mm,right=6mm,top=0mm,bottom=4mm,sharp corners,
overlay={\fill[cint] (frame.north west) rectangle ([yshift=-1.25cm]frame.north east);
\node[anchor=west,text=white,font=\sffamily\bfseries\LARGE] at ([xshift=5mm,yshift=-0.62cm]frame.north west) {Question 50};
\node[anchor=east,text=white,font=\sffamily\itshape\large] at ([xshift=-5mm,yshift=-0.62cm]frame.north east) {Intégrales et équations différentielles};}]
\vspace*{1.75cm}
\begin{minipage}[t][6.9cm][c]{\linewidth}\LARGE\centering
\begin{minipage}{0.95\linewidth}\raggedright\linespread{1.25}\selectfont Que vaut $\displaystyle\int_0^{3} x\,\mathrm{d}x$~?\end{minipage}
\end{minipage}

\vfill
\begin{tcolorbox}[enhanced,nobeforeafter,width=7.1cm,height=2.1cm,valign=center,colback=cint!6,colframe=cint,arc=2.5mm,boxrule=0.9pt,left=3mm]{\Large\bfseries\color{cint}a)}\quad{\Large $1$}\end{tcolorbox}\hfill\begin{tcolorbox}[enhanced,nobeforeafter,width=7.1cm,height=2.1cm,valign=center,colback=cint!6,colframe=cint,arc=2.5mm,boxrule=0.9pt,left=3mm]{\Large\bfseries\color{cint}b)}\quad{\Large $\dfrac{9}{2}$}\end{tcolorbox}\par\vspace{2.5mm}
\begin{tcolorbox}[enhanced,nobeforeafter,width=7.1cm,height=2.1cm,valign=center,colback=cint!6,colframe=cint,arc=2.5mm,boxrule=0.9pt,left=3mm]{\Large\bfseries\color{cint}c)}\quad{\Large $3$}\end{tcolorbox}\hfill\begin{tcolorbox}[enhanced,nobeforeafter,width=7.1cm,height=2.1cm,valign=center,colback=cint!6,colframe=cint,arc=2.5mm,boxrule=0.9pt,left=3mm]{\Large\bfseries\color{cint}d)}\quad{\Large $9$}\end{tcolorbox}
\end{tcolorbox}
\newpage
\begin{tcolorbox}[enhanced,width=16cm,height=14cm,colback=white,colframe=cint,boxrule=2pt,arc=0mm,
left=6mm,right=6mm,top=0mm,bottom=4mm,sharp corners,
overlay={\fill[cint] (frame.north west) rectangle ([yshift=-1.25cm]frame.north east);
\node[anchor=west,text=white,font=\sffamily\bfseries\LARGE] at ([xshift=5mm,yshift=-0.62cm]frame.north west) {Question 51};
\node[anchor=east,text=white,font=\sffamily\itshape\large] at ([xshift=-5mm,yshift=-0.62cm]frame.north east) {Intégrales et équations différentielles};}]
\vspace*{1.75cm}
\begin{minipage}[t][6.9cm][c]{\linewidth}\LARGE\centering
\begin{minipage}{0.95\linewidth}\raggedright\linespread{1.25}\selectfont Que vaut $\displaystyle\int_0^{1} e^{2x}\,\mathrm{d}x$~?\end{minipage}
\end{minipage}

\vfill
\begin{tcolorbox}[enhanced,nobeforeafter,width=7.1cm,height=2.1cm,valign=center,colback=cint!6,colframe=cint,arc=2.5mm,boxrule=0.9pt,left=3mm]{\Large\bfseries\color{cint}a)}\quad{\Large $\dfrac{e^{2}-1}{2}$}\end{tcolorbox}\hfill\begin{tcolorbox}[enhanced,nobeforeafter,width=7.1cm,height=2.1cm,valign=center,colback=cint!6,colframe=cint,arc=2.5mm,boxrule=0.9pt,left=3mm]{\Large\bfseries\color{cint}b)}\quad{\Large $e^{2}-1$}\end{tcolorbox}\par\vspace{2.5mm}
\begin{tcolorbox}[enhanced,nobeforeafter,width=7.1cm,height=2.1cm,valign=center,colback=cint!6,colframe=cint,arc=2.5mm,boxrule=0.9pt,left=3mm]{\Large\bfseries\color{cint}c)}\quad{\Large $2(e^{2}-1)$}\end{tcolorbox}\hfill\begin{tcolorbox}[enhanced,nobeforeafter,width=7.1cm,height=2.1cm,valign=center,colback=cint!6,colframe=cint,arc=2.5mm,boxrule=0.9pt,left=3mm]{\Large\bfseries\color{cint}d)}\quad{\Large $\dfrac{e^{2}}{2}$}\end{tcolorbox}
\end{tcolorbox}
\newpage
\begin{tcolorbox}[enhanced,width=16cm,height=14cm,colback=white,colframe=cint,boxrule=2pt,arc=0mm,
left=6mm,right=6mm,top=0mm,bottom=4mm,sharp corners,
overlay={\fill[cint] (frame.north west) rectangle ([yshift=-1.25cm]frame.north east);
\node[anchor=west,text=white,font=\sffamily\bfseries\LARGE] at ([xshift=5mm,yshift=-0.62cm]frame.north west) {Question 52};
\node[anchor=east,text=white,font=\sffamily\itshape\large] at ([xshift=-5mm,yshift=-0.62cm]frame.north east) {Intégrales et équations différentielles};}]
\vspace*{1.75cm}
\begin{minipage}[t][6.9cm][c]{\linewidth}\LARGE\centering
\begin{minipage}{0.95\linewidth}\raggedright\linespread{1.25}\selectfont Que vaut $\displaystyle\int_0^{1} e^{3x}\,\mathrm{d}x$~?\end{minipage}
\end{minipage}

\vfill
\begin{tcolorbox}[enhanced,nobeforeafter,width=7.1cm,height=2.1cm,valign=center,colback=cint!6,colframe=cint,arc=2.5mm,boxrule=0.9pt,left=3mm]{\Large\bfseries\color{cint}a)}\quad{\Large $3(e^{3}-1)$}\end{tcolorbox}\hfill\begin{tcolorbox}[enhanced,nobeforeafter,width=7.1cm,height=2.1cm,valign=center,colback=cint!6,colframe=cint,arc=2.5mm,boxrule=0.9pt,left=3mm]{\Large\bfseries\color{cint}b)}\quad{\Large $\dfrac{e^{3}-1}{3}$}\end{tcolorbox}\par\vspace{2.5mm}
\begin{tcolorbox}[enhanced,nobeforeafter,width=7.1cm,height=2.1cm,valign=center,colback=cint!6,colframe=cint,arc=2.5mm,boxrule=0.9pt,left=3mm]{\Large\bfseries\color{cint}c)}\quad{\Large $e^{3}-1$}\end{tcolorbox}\hfill\begin{tcolorbox}[enhanced,nobeforeafter,width=7.1cm,height=2.1cm,valign=center,colback=cint!6,colframe=cint,arc=2.5mm,boxrule=0.9pt,left=3mm]{\Large\bfseries\color{cint}d)}\quad{\Large $\dfrac{e^{3}}{3}$}\end{tcolorbox}
\end{tcolorbox}
\newpage
\begin{tcolorbox}[enhanced,width=16cm,height=14cm,colback=white,colframe=cint,boxrule=2pt,arc=0mm,
left=6mm,right=6mm,top=0mm,bottom=4mm,sharp corners,
overlay={\fill[cint] (frame.north west) rectangle ([yshift=-1.25cm]frame.north east);
\node[anchor=west,text=white,font=\sffamily\bfseries\LARGE] at ([xshift=5mm,yshift=-0.62cm]frame.north west) {Question 53};
\node[anchor=east,text=white,font=\sffamily\itshape\large] at ([xshift=-5mm,yshift=-0.62cm]frame.north east) {Intégrales et équations différentielles};}]
\vspace*{1.75cm}
\begin{minipage}[t][6.9cm][c]{\linewidth}\LARGE\centering
\begin{minipage}{0.95\linewidth}\raggedright\linespread{1.25}\selectfont Que vaut $\displaystyle\int_1^{e} \left(\frac{2}{x}+1\right)\mathrm{d}x$~?\end{minipage}
\end{minipage}

\vfill
\begin{tcolorbox}[enhanced,nobeforeafter,width=7.1cm,height=2.1cm,valign=center,colback=cint!6,colframe=cint,arc=2.5mm,boxrule=0.9pt,left=3mm]{\Large\bfseries\color{cint}a)}\quad{\Large $2e+1$}\end{tcolorbox}\hfill\begin{tcolorbox}[enhanced,nobeforeafter,width=7.1cm,height=2.1cm,valign=center,colback=cint!6,colframe=cint,arc=2.5mm,boxrule=0.9pt,left=3mm]{\Large\bfseries\color{cint}b)}\quad{\Large $e+1$}\end{tcolorbox}\par\vspace{2.5mm}
\begin{tcolorbox}[enhanced,nobeforeafter,width=7.1cm,height=2.1cm,valign=center,colback=cint!6,colframe=cint,arc=2.5mm,boxrule=0.9pt,left=3mm]{\Large\bfseries\color{cint}c)}\quad{\Large $e+2$}\end{tcolorbox}\hfill\begin{tcolorbox}[enhanced,nobeforeafter,width=7.1cm,height=2.1cm,valign=center,colback=cint!6,colframe=cint,arc=2.5mm,boxrule=0.9pt,left=3mm]{\Large\bfseries\color{cint}d)}\quad{\Large $3e$}\end{tcolorbox}
\end{tcolorbox}
\newpage
\begin{tcolorbox}[enhanced,width=16cm,height=14cm,colback=white,colframe=cint,boxrule=2pt,arc=0mm,
left=6mm,right=6mm,top=0mm,bottom=4mm,sharp corners,
overlay={\fill[cint] (frame.north west) rectangle ([yshift=-1.25cm]frame.north east);
\node[anchor=west,text=white,font=\sffamily\bfseries\LARGE] at ([xshift=5mm,yshift=-0.62cm]frame.north west) {Question 54};
\node[anchor=east,text=white,font=\sffamily\itshape\large] at ([xshift=-5mm,yshift=-0.62cm]frame.north east) {Intégrales et équations différentielles};}]
\vspace*{1.75cm}
\begin{minipage}[t][6.9cm][c]{\linewidth}\LARGE\centering
\begin{minipage}{0.95\linewidth}\raggedright\linespread{1.25}\selectfont Que vaut $\displaystyle\int_1^{e} \left(\frac{3}{x}+2\right)\mathrm{d}x$~?\end{minipage}
\end{minipage}

\vfill
\begin{tcolorbox}[enhanced,nobeforeafter,width=7.1cm,height=2.1cm,valign=center,colback=cint!6,colframe=cint,arc=2.5mm,boxrule=0.9pt,left=3mm]{\Large\bfseries\color{cint}a)}\quad{\Large $3e+2$}\end{tcolorbox}\hfill\begin{tcolorbox}[enhanced,nobeforeafter,width=7.1cm,height=2.1cm,valign=center,colback=cint!6,colframe=cint,arc=2.5mm,boxrule=0.9pt,left=3mm]{\Large\bfseries\color{cint}b)}\quad{\Large $5e$}\end{tcolorbox}\par\vspace{2.5mm}
\begin{tcolorbox}[enhanced,nobeforeafter,width=7.1cm,height=2.1cm,valign=center,colback=cint!6,colframe=cint,arc=2.5mm,boxrule=0.9pt,left=3mm]{\Large\bfseries\color{cint}c)}\quad{\Large $2e+3$}\end{tcolorbox}\hfill\begin{tcolorbox}[enhanced,nobeforeafter,width=7.1cm,height=2.1cm,valign=center,colback=cint!6,colframe=cint,arc=2.5mm,boxrule=0.9pt,left=3mm]{\Large\bfseries\color{cint}d)}\quad{\Large $2e+1$}\end{tcolorbox}
\end{tcolorbox}
\newpage
\begin{tcolorbox}[enhanced,width=16cm,height=14cm,colback=white,colframe=cint,boxrule=2pt,arc=0mm,
left=6mm,right=6mm,top=0mm,bottom=4mm,sharp corners,
overlay={\fill[cint] (frame.north west) rectangle ([yshift=-1.25cm]frame.north east);
\node[anchor=west,text=white,font=\sffamily\bfseries\LARGE] at ([xshift=5mm,yshift=-0.62cm]frame.north west) {Question 55};
\node[anchor=east,text=white,font=\sffamily\itshape\large] at ([xshift=-5mm,yshift=-0.62cm]frame.north east) {Intégrales et équations différentielles};}]
\vspace*{1.75cm}
\begin{minipage}[t][6.9cm][c]{\linewidth}\LARGE\centering
\begin{minipage}{0.95\linewidth}\raggedright\linespread{1.25}\selectfont Quelle est la primitive $F$ de $f(x)=e^{2x}$ sur $\mathbb{R}$ telle que $F(0)=1$~?\end{minipage}
\end{minipage}

\vfill
\begin{tcolorbox}[enhanced,nobeforeafter,width=7.1cm,height=2.1cm,valign=center,colback=cint!6,colframe=cint,arc=2.5mm,boxrule=0.9pt,left=3mm]{\Large\bfseries\color{cint}a)}\quad{\Large $F(x)=\frac{1}{2}\,e^{2x}+\frac{1}{2}$}\end{tcolorbox}\hfill\begin{tcolorbox}[enhanced,nobeforeafter,width=7.1cm,height=2.1cm,valign=center,colback=cint!6,colframe=cint,arc=2.5mm,boxrule=0.9pt,left=3mm]{\Large\bfseries\color{cint}b)}\quad{\Large $F(x)=2\,e^{2x}-1$}\end{tcolorbox}\par\vspace{2.5mm}
\begin{tcolorbox}[enhanced,nobeforeafter,width=7.1cm,height=2.1cm,valign=center,colback=cint!6,colframe=cint,arc=2.5mm,boxrule=0.9pt,left=3mm]{\Large\bfseries\color{cint}c)}\quad{\Large $F(x)=\frac{1}{2}\,e^{2x}+1$}\end{tcolorbox}\hfill\begin{tcolorbox}[enhanced,nobeforeafter,width=7.1cm,height=2.1cm,valign=center,colback=cint!6,colframe=cint,arc=2.5mm,boxrule=0.9pt,left=3mm]{\Large\bfseries\color{cint}d)}\quad{\Large $F(x)=\frac{1}{2}\,e^{2x}+\frac{3}{2}$}\end{tcolorbox}
\end{tcolorbox}
\newpage
\begin{tcolorbox}[enhanced,width=16cm,height=14cm,colback=white,colframe=cint,boxrule=2pt,arc=0mm,
left=6mm,right=6mm,top=0mm,bottom=4mm,sharp corners,
overlay={\fill[cint] (frame.north west) rectangle ([yshift=-1.25cm]frame.north east);
\node[anchor=west,text=white,font=\sffamily\bfseries\LARGE] at ([xshift=5mm,yshift=-0.62cm]frame.north west) {Question 56};
\node[anchor=east,text=white,font=\sffamily\itshape\large] at ([xshift=-5mm,yshift=-0.62cm]frame.north east) {Intégrales et équations différentielles};}]
\vspace*{1.75cm}
\begin{minipage}[t][6.9cm][c]{\linewidth}\LARGE\centering
\begin{minipage}{0.95\linewidth}\raggedright\linespread{1.25}\selectfont Quelle est la primitive $F$ de $f(x)=e^{3x}$ sur $\mathbb{R}$ telle que $F(0)=0$~?\end{minipage}
\end{minipage}

\vfill
\begin{tcolorbox}[enhanced,nobeforeafter,width=7.1cm,height=2.1cm,valign=center,colback=cint!6,colframe=cint,arc=2.5mm,boxrule=0.9pt,left=3mm]{\Large\bfseries\color{cint}a)}\quad{\Large $F(x)=\frac{1}{3}\,e^{3x}-\frac{1}{3}$}\end{tcolorbox}\hfill\begin{tcolorbox}[enhanced,nobeforeafter,width=7.1cm,height=2.1cm,valign=center,colback=cint!6,colframe=cint,arc=2.5mm,boxrule=0.9pt,left=3mm]{\Large\bfseries\color{cint}b)}\quad{\Large $F(x)=\frac{1}{3}\,e^{3x}$}\end{tcolorbox}\par\vspace{2.5mm}
\begin{tcolorbox}[enhanced,nobeforeafter,width=7.1cm,height=2.1cm,valign=center,colback=cint!6,colframe=cint,arc=2.5mm,boxrule=0.9pt,left=3mm]{\Large\bfseries\color{cint}c)}\quad{\Large $F(x)=\frac{1}{3}\,e^{3x}+\frac{1}{3}$}\end{tcolorbox}\hfill\begin{tcolorbox}[enhanced,nobeforeafter,width=7.1cm,height=2.1cm,valign=center,colback=cint!6,colframe=cint,arc=2.5mm,boxrule=0.9pt,left=3mm]{\Large\bfseries\color{cint}d)}\quad{\Large $F(x)=3\,e^{3x}-3$}\end{tcolorbox}
\end{tcolorbox}
\newpage
\begin{tcolorbox}[enhanced,width=16cm,height=14cm,colback=white,colframe=cint,boxrule=2pt,arc=0mm,
left=6mm,right=6mm,top=0mm,bottom=4mm,sharp corners,
overlay={\fill[cint] (frame.north west) rectangle ([yshift=-1.25cm]frame.north east);
\node[anchor=west,text=white,font=\sffamily\bfseries\LARGE] at ([xshift=5mm,yshift=-0.62cm]frame.north west) {Question 57};
\node[anchor=east,text=white,font=\sffamily\itshape\large] at ([xshift=-5mm,yshift=-0.62cm]frame.north east) {Intégrales et équations différentielles};}]
\vspace*{1.75cm}
\begin{minipage}[t][6.9cm][c]{\linewidth}\LARGE\centering
\begin{minipage}{0.95\linewidth}\raggedright\linespread{1.25}\selectfont Quelle est la valeur moyenne de la fonction $x\mapsto x^2$ sur l'intervalle $[0;3]$~?\end{minipage}
\end{minipage}

\vfill
\begin{tcolorbox}[enhanced,nobeforeafter,width=7.1cm,height=2.1cm,valign=center,colback=cint!6,colframe=cint,arc=2.5mm,boxrule=0.9pt,left=3mm]{\Large\bfseries\color{cint}a)}\quad{\Large $\dfrac{9}{2}$}\end{tcolorbox}\hfill\begin{tcolorbox}[enhanced,nobeforeafter,width=7.1cm,height=2.1cm,valign=center,colback=cint!6,colframe=cint,arc=2.5mm,boxrule=0.9pt,left=3mm]{\Large\bfseries\color{cint}b)}\quad{\Large $1$}\end{tcolorbox}\par\vspace{2.5mm}
\begin{tcolorbox}[enhanced,nobeforeafter,width=7.1cm,height=2.1cm,valign=center,colback=cint!6,colframe=cint,arc=2.5mm,boxrule=0.9pt,left=3mm]{\Large\bfseries\color{cint}c)}\quad{\Large $3$}\end{tcolorbox}\hfill\begin{tcolorbox}[enhanced,nobeforeafter,width=7.1cm,height=2.1cm,valign=center,colback=cint!6,colframe=cint,arc=2.5mm,boxrule=0.9pt,left=3mm]{\Large\bfseries\color{cint}d)}\quad{\Large $9$}\end{tcolorbox}
\end{tcolorbox}
\newpage
\begin{tcolorbox}[enhanced,width=16cm,height=14cm,colback=white,colframe=cint,boxrule=2pt,arc=0mm,
left=6mm,right=6mm,top=0mm,bottom=4mm,sharp corners,
overlay={\fill[cint] (frame.north west) rectangle ([yshift=-1.25cm]frame.north east);
\node[anchor=west,text=white,font=\sffamily\bfseries\LARGE] at ([xshift=5mm,yshift=-0.62cm]frame.north west) {Question 58};
\node[anchor=east,text=white,font=\sffamily\itshape\large] at ([xshift=-5mm,yshift=-0.62cm]frame.north east) {Intégrales et équations différentielles};}]
\vspace*{1.75cm}
\begin{minipage}[t][6.9cm][c]{\linewidth}\LARGE\centering
\begin{minipage}{0.95\linewidth}\raggedright\linespread{1.25}\selectfont Quelle est la valeur moyenne de la fonction $x\mapsto x^2$ sur l'intervalle $[0;6]$~?\end{minipage}
\end{minipage}

\vfill
\begin{tcolorbox}[enhanced,nobeforeafter,width=7.1cm,height=2.1cm,valign=center,colback=cint!6,colframe=cint,arc=2.5mm,boxrule=0.9pt,left=3mm]{\Large\bfseries\color{cint}a)}\quad{\Large $12$}\end{tcolorbox}\hfill\begin{tcolorbox}[enhanced,nobeforeafter,width=7.1cm,height=2.1cm,valign=center,colback=cint!6,colframe=cint,arc=2.5mm,boxrule=0.9pt,left=3mm]{\Large\bfseries\color{cint}b)}\quad{\Large $18$}\end{tcolorbox}\par\vspace{2.5mm}
\begin{tcolorbox}[enhanced,nobeforeafter,width=7.1cm,height=2.1cm,valign=center,colback=cint!6,colframe=cint,arc=2.5mm,boxrule=0.9pt,left=3mm]{\Large\bfseries\color{cint}c)}\quad{\Large $72$}\end{tcolorbox}\hfill\begin{tcolorbox}[enhanced,nobeforeafter,width=7.1cm,height=2.1cm,valign=center,colback=cint!6,colframe=cint,arc=2.5mm,boxrule=0.9pt,left=3mm]{\Large\bfseries\color{cint}d)}\quad{\Large $36$}\end{tcolorbox}
\end{tcolorbox}
\newpage
\begin{tcolorbox}[enhanced,width=16cm,height=14cm,colback=white,colframe=cint,boxrule=2pt,arc=0mm,
left=6mm,right=6mm,top=0mm,bottom=4mm,sharp corners,
overlay={\fill[cint] (frame.north west) rectangle ([yshift=-1.25cm]frame.north east);
\node[anchor=west,text=white,font=\sffamily\bfseries\LARGE] at ([xshift=5mm,yshift=-0.62cm]frame.north west) {Question 59};
\node[anchor=east,text=white,font=\sffamily\itshape\large] at ([xshift=-5mm,yshift=-0.62cm]frame.north east) {Intégrales et équations différentielles};}]
\vspace*{1.75cm}
\begin{minipage}[t][6.9cm][c]{\linewidth}\LARGE\centering
\begin{minipage}{0.95\linewidth}\raggedright\linespread{1.25}\selectfont Quelle est la solution constante de l'équation différentielle $y'=2y+6$~?\end{minipage}
\end{minipage}

\vfill
\begin{tcolorbox}[enhanced,nobeforeafter,width=7.1cm,height=2.1cm,valign=center,colback=cint!6,colframe=cint,arc=2.5mm,boxrule=0.9pt,left=3mm]{\Large\bfseries\color{cint}a)}\quad{\Large $y=-3$}\end{tcolorbox}\hfill\begin{tcolorbox}[enhanced,nobeforeafter,width=7.1cm,height=2.1cm,valign=center,colback=cint!6,colframe=cint,arc=2.5mm,boxrule=0.9pt,left=3mm]{\Large\bfseries\color{cint}b)}\quad{\Large $y=3$}\end{tcolorbox}\par\vspace{2.5mm}
\begin{tcolorbox}[enhanced,nobeforeafter,width=7.1cm,height=2.1cm,valign=center,colback=cint!6,colframe=cint,arc=2.5mm,boxrule=0.9pt,left=3mm]{\Large\bfseries\color{cint}c)}\quad{\Large $y=-\dfrac{1}{3}$}\end{tcolorbox}\hfill\begin{tcolorbox}[enhanced,nobeforeafter,width=7.1cm,height=2.1cm,valign=center,colback=cint!6,colframe=cint,arc=2.5mm,boxrule=0.9pt,left=3mm]{\Large\bfseries\color{cint}d)}\quad{\Large $y=6$}\end{tcolorbox}
\end{tcolorbox}
\newpage
\begin{tcolorbox}[enhanced,width=16cm,height=14cm,colback=white,colframe=cint,boxrule=2pt,arc=0mm,
left=6mm,right=6mm,top=0mm,bottom=4mm,sharp corners,
overlay={\fill[cint] (frame.north west) rectangle ([yshift=-1.25cm]frame.north east);
\node[anchor=west,text=white,font=\sffamily\bfseries\LARGE] at ([xshift=5mm,yshift=-0.62cm]frame.north west) {Question 60};
\node[anchor=east,text=white,font=\sffamily\itshape\large] at ([xshift=-5mm,yshift=-0.62cm]frame.north east) {Intégrales et équations différentielles};}]
\vspace*{1.75cm}
\begin{minipage}[t][6.9cm][c]{\linewidth}\LARGE\centering
\begin{minipage}{0.95\linewidth}\raggedright\linespread{1.25}\selectfont Quelle est la solution constante de l'équation différentielle $y'=3y-1$~?\end{minipage}
\end{minipage}

\vfill
\begin{tcolorbox}[enhanced,nobeforeafter,width=7.1cm,height=2.1cm,valign=center,colback=cint!6,colframe=cint,arc=2.5mm,boxrule=0.9pt,left=3mm]{\Large\bfseries\color{cint}a)}\quad{\Large $y=\dfrac{1}{3}$}\end{tcolorbox}\hfill\begin{tcolorbox}[enhanced,nobeforeafter,width=7.1cm,height=2.1cm,valign=center,colback=cint!6,colframe=cint,arc=2.5mm,boxrule=0.9pt,left=3mm]{\Large\bfseries\color{cint}b)}\quad{\Large $y=-1$}\end{tcolorbox}\par\vspace{2.5mm}
\begin{tcolorbox}[enhanced,nobeforeafter,width=7.1cm,height=2.1cm,valign=center,colback=cint!6,colframe=cint,arc=2.5mm,boxrule=0.9pt,left=3mm]{\Large\bfseries\color{cint}c)}\quad{\Large $y=3$}\end{tcolorbox}\hfill\begin{tcolorbox}[enhanced,nobeforeafter,width=7.1cm,height=2.1cm,valign=center,colback=cint!6,colframe=cint,arc=2.5mm,boxrule=0.9pt,left=3mm]{\Large\bfseries\color{cint}d)}\quad{\Large $y=-\dfrac{1}{3}$}\end{tcolorbox}
\end{tcolorbox}
\newpage
\begin{tcolorbox}[enhanced,width=16cm,height=14cm,colback=white,colframe=cint,boxrule=2pt,arc=0mm,
left=6mm,right=6mm,top=0mm,bottom=4mm,sharp corners,
overlay={\fill[cint] (frame.north west) rectangle ([yshift=-1.25cm]frame.north east);
\node[anchor=west,text=white,font=\sffamily\bfseries\LARGE] at ([xshift=5mm,yshift=-0.62cm]frame.north west) {Question 61};
\node[anchor=east,text=white,font=\sffamily\itshape\large] at ([xshift=-5mm,yshift=-0.62cm]frame.north east) {Intégrales et équations différentielles};}]
\vspace*{1.75cm}
\begin{minipage}[t][6.9cm][c]{\linewidth}\LARGE\centering
\begin{minipage}{0.95\linewidth}\raggedright\linespread{1.25}\selectfont $f$ est la solution de $y'=2y$ telle que $f(0)=3$. Que vaut $f(1)$~?\end{minipage}
\end{minipage}

\vfill
\begin{tcolorbox}[enhanced,nobeforeafter,width=7.1cm,height=2.1cm,valign=center,colback=cint!6,colframe=cint,arc=2.5mm,boxrule=0.9pt,left=3mm]{\Large\bfseries\color{cint}a)}\quad{\Large $3+e^{2}$}\end{tcolorbox}\hfill\begin{tcolorbox}[enhanced,nobeforeafter,width=7.1cm,height=2.1cm,valign=center,colback=cint!6,colframe=cint,arc=2.5mm,boxrule=0.9pt,left=3mm]{\Large\bfseries\color{cint}b)}\quad{\Large $3e^{-2}$}\end{tcolorbox}\par\vspace{2.5mm}
\begin{tcolorbox}[enhanced,nobeforeafter,width=7.1cm,height=2.1cm,valign=center,colback=cint!6,colframe=cint,arc=2.5mm,boxrule=0.9pt,left=3mm]{\Large\bfseries\color{cint}c)}\quad{\Large $2e^{3}$}\end{tcolorbox}\hfill\begin{tcolorbox}[enhanced,nobeforeafter,width=7.1cm,height=2.1cm,valign=center,colback=cint!6,colframe=cint,arc=2.5mm,boxrule=0.9pt,left=3mm]{\Large\bfseries\color{cint}d)}\quad{\Large $3e^{2}$}\end{tcolorbox}
\end{tcolorbox}
\newpage
\begin{tcolorbox}[enhanced,width=16cm,height=14cm,colback=white,colframe=cint,boxrule=2pt,arc=0mm,
left=6mm,right=6mm,top=0mm,bottom=4mm,sharp corners,
overlay={\fill[cint] (frame.north west) rectangle ([yshift=-1.25cm]frame.north east);
\node[anchor=west,text=white,font=\sffamily\bfseries\LARGE] at ([xshift=5mm,yshift=-0.62cm]frame.north west) {Question 62};
\node[anchor=east,text=white,font=\sffamily\itshape\large] at ([xshift=-5mm,yshift=-0.62cm]frame.north east) {Intégrales et équations différentielles};}]
\vspace*{1.75cm}
\begin{minipage}[t][6.9cm][c]{\linewidth}\LARGE\centering
\begin{minipage}{0.95\linewidth}\raggedright\linespread{1.25}\selectfont $f$ est la solution de $y'=-y$ telle que $f(0)=5$. Que vaut $f(1)$~?\end{minipage}
\end{minipage}

\vfill
\begin{tcolorbox}[enhanced,nobeforeafter,width=7.1cm,height=2.1cm,valign=center,colback=cint!6,colframe=cint,arc=2.5mm,boxrule=0.9pt,left=3mm]{\Large\bfseries\color{cint}a)}\quad{\Large $-e^{5}$}\end{tcolorbox}\hfill\begin{tcolorbox}[enhanced,nobeforeafter,width=7.1cm,height=2.1cm,valign=center,colback=cint!6,colframe=cint,arc=2.5mm,boxrule=0.9pt,left=3mm]{\Large\bfseries\color{cint}b)}\quad{\Large $5e$}\end{tcolorbox}\par\vspace{2.5mm}
\begin{tcolorbox}[enhanced,nobeforeafter,width=7.1cm,height=2.1cm,valign=center,colback=cint!6,colframe=cint,arc=2.5mm,boxrule=0.9pt,left=3mm]{\Large\bfseries\color{cint}c)}\quad{\Large $5e^{-1}$}\end{tcolorbox}\hfill\begin{tcolorbox}[enhanced,nobeforeafter,width=7.1cm,height=2.1cm,valign=center,colback=cint!6,colframe=cint,arc=2.5mm,boxrule=0.9pt,left=3mm]{\Large\bfseries\color{cint}d)}\quad{\Large $5+e^{-1}$}\end{tcolorbox}
\end{tcolorbox}
\newpage
\begin{tcolorbox}[enhanced,width=16cm,height=14cm,colback=white,colframe=cint,boxrule=2pt,arc=0mm,
left=6mm,right=6mm,top=0mm,bottom=4mm,sharp corners,
overlay={\fill[cint] (frame.north west) rectangle ([yshift=-1.25cm]frame.north east);
\node[anchor=west,text=white,font=\sffamily\bfseries\LARGE] at ([xshift=5mm,yshift=-0.62cm]frame.north west) {Question 63};
\node[anchor=east,text=white,font=\sffamily\itshape\large] at ([xshift=-5mm,yshift=-0.62cm]frame.north east) {Intégrales et équations différentielles};}]
\vspace*{1.75cm}
\begin{minipage}[t][6.9cm][c]{\linewidth}\LARGE\centering
\begin{minipage}{0.95\linewidth}\raggedright\linespread{1.25}\selectfont Que vaut $\displaystyle\int_0^{1} x\,e^{x}\,\mathrm{d}x$~?\end{minipage}
\end{minipage}

\vfill
\begin{tcolorbox}[enhanced,nobeforeafter,width=7.1cm,height=2.1cm,valign=center,colback=cint!6,colframe=cint,arc=2.5mm,boxrule=0.9pt,left=3mm]{\Large\bfseries\color{cint}a)}\quad{\Large $1$}\end{tcolorbox}\hfill\begin{tcolorbox}[enhanced,nobeforeafter,width=7.1cm,height=2.1cm,valign=center,colback=cint!6,colframe=cint,arc=2.5mm,boxrule=0.9pt,left=3mm]{\Large\bfseries\color{cint}b)}\quad{\Large $2e-1$}\end{tcolorbox}\par\vspace{2.5mm}
\begin{tcolorbox}[enhanced,nobeforeafter,width=7.1cm,height=2.1cm,valign=center,colback=cint!6,colframe=cint,arc=2.5mm,boxrule=0.9pt,left=3mm]{\Large\bfseries\color{cint}c)}\quad{\Large $e$}\end{tcolorbox}\hfill\begin{tcolorbox}[enhanced,nobeforeafter,width=7.1cm,height=2.1cm,valign=center,colback=cint!6,colframe=cint,arc=2.5mm,boxrule=0.9pt,left=3mm]{\Large\bfseries\color{cint}d)}\quad{\Large $0$}\end{tcolorbox}
\end{tcolorbox}
\newpage
\begin{tcolorbox}[enhanced,width=16cm,height=14cm,colback=white,colframe=cint,boxrule=2pt,arc=0mm,
left=6mm,right=6mm,top=0mm,bottom=4mm,sharp corners,
overlay={\fill[cint] (frame.north west) rectangle ([yshift=-1.25cm]frame.north east);
\node[anchor=west,text=white,font=\sffamily\bfseries\LARGE] at ([xshift=5mm,yshift=-0.62cm]frame.north west) {Question 64};
\node[anchor=east,text=white,font=\sffamily\itshape\large] at ([xshift=-5mm,yshift=-0.62cm]frame.north east) {Intégrales et équations différentielles};}]
\vspace*{1.75cm}
\begin{minipage}[t][6.9cm][c]{\linewidth}\LARGE\centering
\begin{minipage}{0.95\linewidth}\raggedright\linespread{1.25}\selectfont Que vaut $\displaystyle\int_0^{2} x\,e^{x}\,\mathrm{d}x$~?\end{minipage}
\end{minipage}

\vfill
\begin{tcolorbox}[enhanced,nobeforeafter,width=7.1cm,height=2.1cm,valign=center,colback=cint!6,colframe=cint,arc=2.5mm,boxrule=0.9pt,left=3mm]{\Large\bfseries\color{cint}a)}\quad{\Large $2e^{2}$}\end{tcolorbox}\hfill\begin{tcolorbox}[enhanced,nobeforeafter,width=7.1cm,height=2.1cm,valign=center,colback=cint!6,colframe=cint,arc=2.5mm,boxrule=0.9pt,left=3mm]{\Large\bfseries\color{cint}b)}\quad{\Large $3e^{2}-1$}\end{tcolorbox}\par\vspace{2.5mm}
\begin{tcolorbox}[enhanced,nobeforeafter,width=7.1cm,height=2.1cm,valign=center,colback=cint!6,colframe=cint,arc=2.5mm,boxrule=0.9pt,left=3mm]{\Large\bfseries\color{cint}c)}\quad{\Large $e^{2}+1$}\end{tcolorbox}\hfill\begin{tcolorbox}[enhanced,nobeforeafter,width=7.1cm,height=2.1cm,valign=center,colback=cint!6,colframe=cint,arc=2.5mm,boxrule=0.9pt,left=3mm]{\Large\bfseries\color{cint}d)}\quad{\Large $e^{2}$}\end{tcolorbox}
\end{tcolorbox}
\newpage
\begin{tcolorbox}[enhanced,width=16cm,height=14cm,colback=white,colframe=cesp,boxrule=2pt,arc=0mm,
left=6mm,right=6mm,top=0mm,bottom=4mm,sharp corners,
overlay={\fill[cesp] (frame.north west) rectangle ([yshift=-1.25cm]frame.north east);
\node[anchor=west,text=white,font=\sffamily\bfseries\LARGE] at ([xshift=5mm,yshift=-0.62cm]frame.north west) {Question 65};
\node[anchor=east,text=white,font=\sffamily\itshape\large] at ([xshift=-5mm,yshift=-0.62cm]frame.north east) {Géométrie dans l'espace};}]
\vspace*{1.75cm}
\begin{minipage}[t][6.9cm][c]{\linewidth}\Large\centering
\begin{minipage}{0.95\linewidth}\raggedright\linespread{1.25}\selectfont Dans un repère orthonormé, que vaut le produit scalaire $\vec u\cdot\vec v$ pour $\vec u\begin{pmatrix}1\\2\\-1\end{pmatrix}$ et $\vec v\begin{pmatrix}3\\0\\2\end{pmatrix}$~?\end{minipage}
\end{minipage}

\vfill
\begin{tcolorbox}[enhanced,nobeforeafter,width=7.1cm,height=2.1cm,valign=center,colback=cesp!6,colframe=cesp,arc=2.5mm,boxrule=0.9pt,left=3mm]{\Large\bfseries\color{cesp}a)}\quad{\Large $10$}\end{tcolorbox}\hfill\begin{tcolorbox}[enhanced,nobeforeafter,width=7.1cm,height=2.1cm,valign=center,colback=cesp!6,colframe=cesp,arc=2.5mm,boxrule=0.9pt,left=3mm]{\Large\bfseries\color{cesp}b)}\quad{\Large $-1$}\end{tcolorbox}\par\vspace{2.5mm}
\begin{tcolorbox}[enhanced,nobeforeafter,width=7.1cm,height=2.1cm,valign=center,colback=cesp!6,colframe=cesp,arc=2.5mm,boxrule=0.9pt,left=3mm]{\Large\bfseries\color{cesp}c)}\quad{\Large $1$}\end{tcolorbox}\hfill\begin{tcolorbox}[enhanced,nobeforeafter,width=7.1cm,height=2.1cm,valign=center,colback=cesp!6,colframe=cesp,arc=2.5mm,boxrule=0.9pt,left=3mm]{\Large\bfseries\color{cesp}d)}\quad{\Large $3$}\end{tcolorbox}
\end{tcolorbox}
\newpage
\begin{tcolorbox}[enhanced,width=16cm,height=14cm,colback=white,colframe=cesp,boxrule=2pt,arc=0mm,
left=6mm,right=6mm,top=0mm,bottom=4mm,sharp corners,
overlay={\fill[cesp] (frame.north west) rectangle ([yshift=-1.25cm]frame.north east);
\node[anchor=west,text=white,font=\sffamily\bfseries\LARGE] at ([xshift=5mm,yshift=-0.62cm]frame.north west) {Question 66};
\node[anchor=east,text=white,font=\sffamily\itshape\large] at ([xshift=-5mm,yshift=-0.62cm]frame.north east) {Géométrie dans l'espace};}]
\vspace*{1.75cm}
\begin{minipage}[t][6.9cm][c]{\linewidth}\Large\centering
\begin{minipage}{0.95\linewidth}\raggedright\linespread{1.25}\selectfont Dans un repère orthonormé, que vaut le produit scalaire $\vec u\cdot\vec v$ pour $\vec u\begin{pmatrix}2\\-1\\4\end{pmatrix}$ et $\vec v\begin{pmatrix}1\\3\\1\end{pmatrix}$~?\end{minipage}
\end{minipage}

\vfill
\begin{tcolorbox}[enhanced,nobeforeafter,width=7.1cm,height=2.1cm,valign=center,colback=cesp!6,colframe=cesp,arc=2.5mm,boxrule=0.9pt,left=3mm]{\Large\bfseries\color{cesp}a)}\quad{\Large $5$}\end{tcolorbox}\hfill\begin{tcolorbox}[enhanced,nobeforeafter,width=7.1cm,height=2.1cm,valign=center,colback=cesp!6,colframe=cesp,arc=2.5mm,boxrule=0.9pt,left=3mm]{\Large\bfseries\color{cesp}b)}\quad{\Large $25$}\end{tcolorbox}\par\vspace{2.5mm}
\begin{tcolorbox}[enhanced,nobeforeafter,width=7.1cm,height=2.1cm,valign=center,colback=cesp!6,colframe=cesp,arc=2.5mm,boxrule=0.9pt,left=3mm]{\Large\bfseries\color{cesp}c)}\quad{\Large $-3$}\end{tcolorbox}\hfill\begin{tcolorbox}[enhanced,nobeforeafter,width=7.1cm,height=2.1cm,valign=center,colback=cesp!6,colframe=cesp,arc=2.5mm,boxrule=0.9pt,left=3mm]{\Large\bfseries\color{cesp}d)}\quad{\Large $3$}\end{tcolorbox}
\end{tcolorbox}
\newpage
\begin{tcolorbox}[enhanced,width=16cm,height=14cm,colback=white,colframe=cesp,boxrule=2pt,arc=0mm,
left=6mm,right=6mm,top=0mm,bottom=4mm,sharp corners,
overlay={\fill[cesp] (frame.north west) rectangle ([yshift=-1.25cm]frame.north east);
\node[anchor=west,text=white,font=\sffamily\bfseries\LARGE] at ([xshift=5mm,yshift=-0.62cm]frame.north west) {Question 67};
\node[anchor=east,text=white,font=\sffamily\itshape\large] at ([xshift=-5mm,yshift=-0.62cm]frame.north east) {Géométrie dans l'espace};}]
\vspace*{1.75cm}
\begin{minipage}[t][6.9cm][c]{\linewidth}\LARGE\centering
\begin{minipage}{0.95\linewidth}\raggedright\linespread{1.25}\selectfont Dans un repère orthonormé, quelle est la norme du vecteur $\vec u\begin{pmatrix}1\\2\\2\end{pmatrix}$~?\end{minipage}
\end{minipage}

\vfill
\begin{tcolorbox}[enhanced,nobeforeafter,width=7.1cm,height=2.1cm,valign=center,colback=cesp!6,colframe=cesp,arc=2.5mm,boxrule=0.9pt,left=3mm]{\Large\bfseries\color{cesp}a)}\quad{\Large $5$}\end{tcolorbox}\hfill\begin{tcolorbox}[enhanced,nobeforeafter,width=7.1cm,height=2.1cm,valign=center,colback=cesp!6,colframe=cesp,arc=2.5mm,boxrule=0.9pt,left=3mm]{\Large\bfseries\color{cesp}b)}\quad{\Large $3$}\end{tcolorbox}\par\vspace{2.5mm}
\begin{tcolorbox}[enhanced,nobeforeafter,width=7.1cm,height=2.1cm,valign=center,colback=cesp!6,colframe=cesp,arc=2.5mm,boxrule=0.9pt,left=3mm]{\Large\bfseries\color{cesp}c)}\quad{\Large $9$}\end{tcolorbox}\hfill\begin{tcolorbox}[enhanced,nobeforeafter,width=7.1cm,height=2.1cm,valign=center,colback=cesp!6,colframe=cesp,arc=2.5mm,boxrule=0.9pt,left=3mm]{\Large\bfseries\color{cesp}d)}\quad{\Large $\sqrt{10}$}\end{tcolorbox}
\end{tcolorbox}
\newpage
\begin{tcolorbox}[enhanced,width=16cm,height=14cm,colback=white,colframe=cesp,boxrule=2pt,arc=0mm,
left=6mm,right=6mm,top=0mm,bottom=4mm,sharp corners,
overlay={\fill[cesp] (frame.north west) rectangle ([yshift=-1.25cm]frame.north east);
\node[anchor=west,text=white,font=\sffamily\bfseries\LARGE] at ([xshift=5mm,yshift=-0.62cm]frame.north west) {Question 68};
\node[anchor=east,text=white,font=\sffamily\itshape\large] at ([xshift=-5mm,yshift=-0.62cm]frame.north east) {Géométrie dans l'espace};}]
\vspace*{1.75cm}
\begin{minipage}[t][6.9cm][c]{\linewidth}\LARGE\centering
\begin{minipage}{0.95\linewidth}\raggedright\linespread{1.25}\selectfont Dans un repère orthonormé, quelle est la norme du vecteur $\vec u\begin{pmatrix}2\\3\\6\end{pmatrix}$~?\end{minipage}
\end{minipage}

\vfill
\begin{tcolorbox}[enhanced,nobeforeafter,width=7.1cm,height=2.1cm,valign=center,colback=cesp!6,colframe=cesp,arc=2.5mm,boxrule=0.9pt,left=3mm]{\Large\bfseries\color{cesp}a)}\quad{\Large $5\sqrt{2}$}\end{tcolorbox}\hfill\begin{tcolorbox}[enhanced,nobeforeafter,width=7.1cm,height=2.1cm,valign=center,colback=cesp!6,colframe=cesp,arc=2.5mm,boxrule=0.9pt,left=3mm]{\Large\bfseries\color{cesp}b)}\quad{\Large $49$}\end{tcolorbox}\par\vspace{2.5mm}
\begin{tcolorbox}[enhanced,nobeforeafter,width=7.1cm,height=2.1cm,valign=center,colback=cesp!6,colframe=cesp,arc=2.5mm,boxrule=0.9pt,left=3mm]{\Large\bfseries\color{cesp}c)}\quad{\Large $7$}\end{tcolorbox}\hfill\begin{tcolorbox}[enhanced,nobeforeafter,width=7.1cm,height=2.1cm,valign=center,colback=cesp!6,colframe=cesp,arc=2.5mm,boxrule=0.9pt,left=3mm]{\Large\bfseries\color{cesp}d)}\quad{\Large $11$}\end{tcolorbox}
\end{tcolorbox}
\newpage
\begin{tcolorbox}[enhanced,width=16cm,height=14cm,colback=white,colframe=cesp,boxrule=2pt,arc=0mm,
left=6mm,right=6mm,top=0mm,bottom=4mm,sharp corners,
overlay={\fill[cesp] (frame.north west) rectangle ([yshift=-1.25cm]frame.north east);
\node[anchor=west,text=white,font=\sffamily\bfseries\LARGE] at ([xshift=5mm,yshift=-0.62cm]frame.north west) {Question 69};
\node[anchor=east,text=white,font=\sffamily\itshape\large] at ([xshift=-5mm,yshift=-0.62cm]frame.north east) {Géométrie dans l'espace};}]
\vspace*{1.75cm}
\begin{minipage}[t][6.9cm][c]{\linewidth}\Large\centering
\begin{minipage}{0.95\linewidth}\raggedright\linespread{1.25}\selectfont Pour quelle valeur de $m$ les vecteurs $\vec u\begin{pmatrix}1\\m\\2\end{pmatrix}$ et $\vec v\begin{pmatrix}2\\1\\-3\end{pmatrix}$ sont-ils orthogonaux~?\end{minipage}
\end{minipage}

\vfill
\begin{tcolorbox}[enhanced,nobeforeafter,width=7.1cm,height=2.1cm,valign=center,colback=cesp!6,colframe=cesp,arc=2.5mm,boxrule=0.9pt,left=3mm]{\Large\bfseries\color{cesp}a)}\quad{\Large $m=1$}\end{tcolorbox}\hfill\begin{tcolorbox}[enhanced,nobeforeafter,width=7.1cm,height=2.1cm,valign=center,colback=cesp!6,colframe=cesp,arc=2.5mm,boxrule=0.9pt,left=3mm]{\Large\bfseries\color{cesp}b)}\quad{\Large $m=5$}\end{tcolorbox}\par\vspace{2.5mm}
\begin{tcolorbox}[enhanced,nobeforeafter,width=7.1cm,height=2.1cm,valign=center,colback=cesp!6,colframe=cesp,arc=2.5mm,boxrule=0.9pt,left=3mm]{\Large\bfseries\color{cesp}c)}\quad{\Large $m=4$}\end{tcolorbox}\hfill\begin{tcolorbox}[enhanced,nobeforeafter,width=7.1cm,height=2.1cm,valign=center,colback=cesp!6,colframe=cesp,arc=2.5mm,boxrule=0.9pt,left=3mm]{\Large\bfseries\color{cesp}d)}\quad{\Large $m=-4$}\end{tcolorbox}
\end{tcolorbox}
\newpage
\begin{tcolorbox}[enhanced,width=16cm,height=14cm,colback=white,colframe=cesp,boxrule=2pt,arc=0mm,
left=6mm,right=6mm,top=0mm,bottom=4mm,sharp corners,
overlay={\fill[cesp] (frame.north west) rectangle ([yshift=-1.25cm]frame.north east);
\node[anchor=west,text=white,font=\sffamily\bfseries\LARGE] at ([xshift=5mm,yshift=-0.62cm]frame.north west) {Question 70};
\node[anchor=east,text=white,font=\sffamily\itshape\large] at ([xshift=-5mm,yshift=-0.62cm]frame.north east) {Géométrie dans l'espace};}]
\vspace*{1.75cm}
\begin{minipage}[t][6.9cm][c]{\linewidth}\Large\centering
\begin{minipage}{0.95\linewidth}\raggedright\linespread{1.25}\selectfont Pour quelle valeur de $m$ les vecteurs $\vec u\begin{pmatrix}1\\m\\2\end{pmatrix}$ et $\vec v\begin{pmatrix}1\\2\\3\end{pmatrix}$ sont-ils orthogonaux~?\end{minipage}
\end{minipage}

\vfill
\begin{tcolorbox}[enhanced,nobeforeafter,width=7.1cm,height=2.1cm,valign=center,colback=cesp!6,colframe=cesp,arc=2.5mm,boxrule=0.9pt,left=3mm]{\Large\bfseries\color{cesp}a)}\quad{\Large $m=\dfrac{7}{2}$}\end{tcolorbox}\hfill\begin{tcolorbox}[enhanced,nobeforeafter,width=7.1cm,height=2.1cm,valign=center,colback=cesp!6,colframe=cesp,arc=2.5mm,boxrule=0.9pt,left=3mm]{\Large\bfseries\color{cesp}b)}\quad{\Large $m=-2$}\end{tcolorbox}\par\vspace{2.5mm}
\begin{tcolorbox}[enhanced,nobeforeafter,width=7.1cm,height=2.1cm,valign=center,colback=cesp!6,colframe=cesp,arc=2.5mm,boxrule=0.9pt,left=3mm]{\Large\bfseries\color{cesp}c)}\quad{\Large $m=-\dfrac{7}{2}$}\end{tcolorbox}\hfill\begin{tcolorbox}[enhanced,nobeforeafter,width=7.1cm,height=2.1cm,valign=center,colback=cesp!6,colframe=cesp,arc=2.5mm,boxrule=0.9pt,left=3mm]{\Large\bfseries\color{cesp}d)}\quad{\Large $m=-\dfrac{5}{2}$}\end{tcolorbox}
\end{tcolorbox}
\newpage
\begin{tcolorbox}[enhanced,width=16cm,height=14cm,colback=white,colframe=cesp,boxrule=2pt,arc=0mm,
left=6mm,right=6mm,top=0mm,bottom=4mm,sharp corners,
overlay={\fill[cesp] (frame.north west) rectangle ([yshift=-1.25cm]frame.north east);
\node[anchor=west,text=white,font=\sffamily\bfseries\LARGE] at ([xshift=5mm,yshift=-0.62cm]frame.north west) {Question 71};
\node[anchor=east,text=white,font=\sffamily\itshape\large] at ([xshift=-5mm,yshift=-0.62cm]frame.north east) {Géométrie dans l'espace};}]
\vspace*{1.75cm}
\begin{minipage}[t][6.9cm][c]{\linewidth}\Large\centering
\begin{minipage}{0.95\linewidth}\raggedright\linespread{1.25}\selectfont Quelle est une équation du plan passant par $A(1;0;2)$ et de vecteur normal $\vec n\begin{pmatrix}1\\2\\-1\end{pmatrix}$~?\end{minipage}
\end{minipage}

\vfill
\begin{tcolorbox}[enhanced,nobeforeafter,width=7.1cm,height=2.1cm,valign=center,colback=cesp!6,colframe=cesp,arc=2.5mm,boxrule=0.9pt,left=3mm]{\Large\bfseries\color{cesp}a)}\quad{\Large $x+2z+1=0$}\end{tcolorbox}\hfill\begin{tcolorbox}[enhanced,nobeforeafter,width=7.1cm,height=2.1cm,valign=center,colback=cesp!6,colframe=cesp,arc=2.5mm,boxrule=0.9pt,left=3mm]{\Large\bfseries\color{cesp}b)}\quad{\Large $x+2y-z-1=0$}\end{tcolorbox}\par\vspace{2.5mm}
\begin{tcolorbox}[enhanced,nobeforeafter,width=7.1cm,height=2.1cm,valign=center,colback=cesp!6,colframe=cesp,arc=2.5mm,boxrule=0.9pt,left=3mm]{\Large\bfseries\color{cesp}c)}\quad{\Large $x+2y-z+1=0$}\end{tcolorbox}\hfill\begin{tcolorbox}[enhanced,nobeforeafter,width=7.1cm,height=2.1cm,valign=center,colback=cesp!6,colframe=cesp,arc=2.5mm,boxrule=0.9pt,left=3mm]{\Large\bfseries\color{cesp}d)}\quad{\Large $x+2y-z+2=0$}\end{tcolorbox}
\end{tcolorbox}
\newpage
\begin{tcolorbox}[enhanced,width=16cm,height=14cm,colback=white,colframe=cesp,boxrule=2pt,arc=0mm,
left=6mm,right=6mm,top=0mm,bottom=4mm,sharp corners,
overlay={\fill[cesp] (frame.north west) rectangle ([yshift=-1.25cm]frame.north east);
\node[anchor=west,text=white,font=\sffamily\bfseries\LARGE] at ([xshift=5mm,yshift=-0.62cm]frame.north west) {Question 72};
\node[anchor=east,text=white,font=\sffamily\itshape\large] at ([xshift=-5mm,yshift=-0.62cm]frame.north east) {Géométrie dans l'espace};}]
\vspace*{1.75cm}
\begin{minipage}[t][6.9cm][c]{\linewidth}\Large\centering
\begin{minipage}{0.95\linewidth}\raggedright\linespread{1.25}\selectfont Quelle est une équation du plan passant par $A(0;1;1)$ et de vecteur normal $\vec n\begin{pmatrix}2\\-1\\3\end{pmatrix}$~?\end{minipage}
\end{minipage}

\vfill
\begin{tcolorbox}[enhanced,nobeforeafter,width=7.1cm,height=2.1cm,valign=center,colback=cesp!6,colframe=cesp,arc=2.5mm,boxrule=0.9pt,left=3mm]{\Large\bfseries\color{cesp}a)}\quad{\Large $2x-y+3z+2=0$}\end{tcolorbox}\hfill\begin{tcolorbox}[enhanced,nobeforeafter,width=7.1cm,height=2.1cm,valign=center,colback=cesp!6,colframe=cesp,arc=2.5mm,boxrule=0.9pt,left=3mm]{\Large\bfseries\color{cesp}b)}\quad{\Large $2x-y+3z-1=0$}\end{tcolorbox}\par\vspace{2.5mm}
\begin{tcolorbox}[enhanced,nobeforeafter,width=7.1cm,height=2.1cm,valign=center,colback=cesp!6,colframe=cesp,arc=2.5mm,boxrule=0.9pt,left=3mm]{\Large\bfseries\color{cesp}c)}\quad{\Large $2x-y+3z-2=0$}\end{tcolorbox}\hfill\begin{tcolorbox}[enhanced,nobeforeafter,width=7.1cm,height=2.1cm,valign=center,colback=cesp!6,colframe=cesp,arc=2.5mm,boxrule=0.9pt,left=3mm]{\Large\bfseries\color{cesp}d)}\quad{\Large $y+z-2=0$}\end{tcolorbox}
\end{tcolorbox}
\newpage
\begin{tcolorbox}[enhanced,width=16cm,height=14cm,colback=white,colframe=cesp,boxrule=2pt,arc=0mm,
left=6mm,right=6mm,top=0mm,bottom=4mm,sharp corners,
overlay={\fill[cesp] (frame.north west) rectangle ([yshift=-1.25cm]frame.north east);
\node[anchor=west,text=white,font=\sffamily\bfseries\LARGE] at ([xshift=5mm,yshift=-0.62cm]frame.north west) {Question 73};
\node[anchor=east,text=white,font=\sffamily\itshape\large] at ([xshift=-5mm,yshift=-0.62cm]frame.north east) {Géométrie dans l'espace};}]
\vspace*{1.75cm}
\begin{minipage}[t][6.9cm][c]{\linewidth}\Large\centering
\begin{minipage}{0.95\linewidth}\raggedright\linespread{1.25}\selectfont Lequel de ces points appartient à la droite de représentation paramétrique $\left\{\begin{array}{l}x=1+2t\\y=-1+t\\z=3-t\end{array}\right.$ ($t\in\mathbb{R}$)~?\end{minipage}
\end{minipage}

\vfill
\begin{tcolorbox}[enhanced,nobeforeafter,width=7.1cm,height=2.1cm,valign=center,colback=cesp!6,colframe=cesp,arc=2.5mm,boxrule=0.9pt,left=3mm]{\Large\bfseries\color{cesp}a)}\quad{\Large $(5;1;1)$}\end{tcolorbox}\hfill\begin{tcolorbox}[enhanced,nobeforeafter,width=7.1cm,height=2.1cm,valign=center,colback=cesp!6,colframe=cesp,arc=2.5mm,boxrule=0.9pt,left=3mm]{\Large\bfseries\color{cesp}b)}\quad{\Large $(5;1;2)$}\end{tcolorbox}\par\vspace{2.5mm}
\begin{tcolorbox}[enhanced,nobeforeafter,width=7.1cm,height=2.1cm,valign=center,colback=cesp!6,colframe=cesp,arc=2.5mm,boxrule=0.9pt,left=3mm]{\Large\bfseries\color{cesp}c)}\quad{\Large $(6;1;1)$}\end{tcolorbox}\hfill\begin{tcolorbox}[enhanced,nobeforeafter,width=7.1cm,height=2.1cm,valign=center,colback=cesp!6,colframe=cesp,arc=2.5mm,boxrule=0.9pt,left=3mm]{\Large\bfseries\color{cesp}d)}\quad{\Large $(5;2;1)$}\end{tcolorbox}
\end{tcolorbox}
\newpage
\begin{tcolorbox}[enhanced,width=16cm,height=14cm,colback=white,colframe=cesp,boxrule=2pt,arc=0mm,
left=6mm,right=6mm,top=0mm,bottom=4mm,sharp corners,
overlay={\fill[cesp] (frame.north west) rectangle ([yshift=-1.25cm]frame.north east);
\node[anchor=west,text=white,font=\sffamily\bfseries\LARGE] at ([xshift=5mm,yshift=-0.62cm]frame.north west) {Question 74};
\node[anchor=east,text=white,font=\sffamily\itshape\large] at ([xshift=-5mm,yshift=-0.62cm]frame.north east) {Géométrie dans l'espace};}]
\vspace*{1.75cm}
\begin{minipage}[t][6.9cm][c]{\linewidth}\Large\centering
\begin{minipage}{0.95\linewidth}\raggedright\linespread{1.25}\selectfont Lequel de ces points appartient à la droite de représentation paramétrique $\left\{\begin{array}{l}x=t\\y=2-t\\z=1+3t\end{array}\right.$ ($t\in\mathbb{R}$)~?\end{minipage}
\end{minipage}

\vfill
\begin{tcolorbox}[enhanced,nobeforeafter,width=7.1cm,height=2.1cm,valign=center,colback=cesp!6,colframe=cesp,arc=2.5mm,boxrule=0.9pt,left=3mm]{\Large\bfseries\color{cesp}a)}\quad{\Large $(0;3;-2)$}\end{tcolorbox}\hfill\begin{tcolorbox}[enhanced,nobeforeafter,width=7.1cm,height=2.1cm,valign=center,colback=cesp!6,colframe=cesp,arc=2.5mm,boxrule=0.9pt,left=3mm]{\Large\bfseries\color{cesp}b)}\quad{\Large $(-1;3;-1)$}\end{tcolorbox}\par\vspace{2.5mm}
\begin{tcolorbox}[enhanced,nobeforeafter,width=7.1cm,height=2.1cm,valign=center,colback=cesp!6,colframe=cesp,arc=2.5mm,boxrule=0.9pt,left=3mm]{\Large\bfseries\color{cesp}c)}\quad{\Large $(-1;3;-2)$}\end{tcolorbox}\hfill\begin{tcolorbox}[enhanced,nobeforeafter,width=7.1cm,height=2.1cm,valign=center,colback=cesp!6,colframe=cesp,arc=2.5mm,boxrule=0.9pt,left=3mm]{\Large\bfseries\color{cesp}d)}\quad{\Large $(-1;4;-2)$}\end{tcolorbox}
\end{tcolorbox}
\newpage
\begin{tcolorbox}[enhanced,width=16cm,height=14cm,colback=white,colframe=cesp,boxrule=2pt,arc=0mm,
left=6mm,right=6mm,top=0mm,bottom=4mm,sharp corners,
overlay={\fill[cesp] (frame.north west) rectangle ([yshift=-1.25cm]frame.north east);
\node[anchor=west,text=white,font=\sffamily\bfseries\LARGE] at ([xshift=5mm,yshift=-0.62cm]frame.north west) {Question 75};
\node[anchor=east,text=white,font=\sffamily\itshape\large] at ([xshift=-5mm,yshift=-0.62cm]frame.north east) {Géométrie dans l'espace};}]
\vspace*{1.75cm}
\begin{minipage}[t][6.9cm][c]{\linewidth}\Large\centering
\begin{minipage}{0.95\linewidth}\raggedright\linespread{1.25}\selectfont Pour quelle valeur de $m$ les vecteurs $\vec u\begin{pmatrix}2\\m\\6\end{pmatrix}$ et $\vec v\begin{pmatrix}1\\3\\3\end{pmatrix}$ sont-ils colinéaires~?\end{minipage}
\end{minipage}

\vfill
\begin{tcolorbox}[enhanced,nobeforeafter,width=7.1cm,height=2.1cm,valign=center,colback=cesp!6,colframe=cesp,arc=2.5mm,boxrule=0.9pt,left=3mm]{\Large\bfseries\color{cesp}a)}\quad{\Large $m=\dfrac{3}{2}$}\end{tcolorbox}\hfill\begin{tcolorbox}[enhanced,nobeforeafter,width=7.1cm,height=2.1cm,valign=center,colback=cesp!6,colframe=cesp,arc=2.5mm,boxrule=0.9pt,left=3mm]{\Large\bfseries\color{cesp}b)}\quad{\Large $m=-6$}\end{tcolorbox}\par\vspace{2.5mm}
\begin{tcolorbox}[enhanced,nobeforeafter,width=7.1cm,height=2.1cm,valign=center,colback=cesp!6,colframe=cesp,arc=2.5mm,boxrule=0.9pt,left=3mm]{\Large\bfseries\color{cesp}c)}\quad{\Large $m=5$}\end{tcolorbox}\hfill\begin{tcolorbox}[enhanced,nobeforeafter,width=7.1cm,height=2.1cm,valign=center,colback=cesp!6,colframe=cesp,arc=2.5mm,boxrule=0.9pt,left=3mm]{\Large\bfseries\color{cesp}d)}\quad{\Large $m=6$}\end{tcolorbox}
\end{tcolorbox}
\newpage
\begin{tcolorbox}[enhanced,width=16cm,height=14cm,colback=white,colframe=cesp,boxrule=2pt,arc=0mm,
left=6mm,right=6mm,top=0mm,bottom=4mm,sharp corners,
overlay={\fill[cesp] (frame.north west) rectangle ([yshift=-1.25cm]frame.north east);
\node[anchor=west,text=white,font=\sffamily\bfseries\LARGE] at ([xshift=5mm,yshift=-0.62cm]frame.north west) {Question 76};
\node[anchor=east,text=white,font=\sffamily\itshape\large] at ([xshift=-5mm,yshift=-0.62cm]frame.north east) {Géométrie dans l'espace};}]
\vspace*{1.75cm}
\begin{minipage}[t][6.9cm][c]{\linewidth}\Large\centering
\begin{minipage}{0.95\linewidth}\raggedright\linespread{1.25}\selectfont Pour quelle valeur de $m$ les vecteurs $\vec u\begin{pmatrix}6\\m\\3\end{pmatrix}$ et $\vec v\begin{pmatrix}2\\-1\\1\end{pmatrix}$ sont-ils colinéaires~?\end{minipage}
\end{minipage}

\vfill
\begin{tcolorbox}[enhanced,nobeforeafter,width=7.1cm,height=2.1cm,valign=center,colback=cesp!6,colframe=cesp,arc=2.5mm,boxrule=0.9pt,left=3mm]{\Large\bfseries\color{cesp}a)}\quad{\Large $m=-\dfrac{1}{3}$}\end{tcolorbox}\hfill\begin{tcolorbox}[enhanced,nobeforeafter,width=7.1cm,height=2.1cm,valign=center,colback=cesp!6,colframe=cesp,arc=2.5mm,boxrule=0.9pt,left=3mm]{\Large\bfseries\color{cesp}b)}\quad{\Large $m=2$}\end{tcolorbox}\par\vspace{2.5mm}
\begin{tcolorbox}[enhanced,nobeforeafter,width=7.1cm,height=2.1cm,valign=center,colback=cesp!6,colframe=cesp,arc=2.5mm,boxrule=0.9pt,left=3mm]{\Large\bfseries\color{cesp}c)}\quad{\Large $m=-3$}\end{tcolorbox}\hfill\begin{tcolorbox}[enhanced,nobeforeafter,width=7.1cm,height=2.1cm,valign=center,colback=cesp!6,colframe=cesp,arc=2.5mm,boxrule=0.9pt,left=3mm]{\Large\bfseries\color{cesp}d)}\quad{\Large $m=3$}\end{tcolorbox}
\end{tcolorbox}
\newpage
\begin{tcolorbox}[enhanced,width=16cm,height=14cm,colback=white,colframe=cesp,boxrule=2pt,arc=0mm,
left=6mm,right=6mm,top=0mm,bottom=4mm,sharp corners,
overlay={\fill[cesp] (frame.north west) rectangle ([yshift=-1.25cm]frame.north east);
\node[anchor=west,text=white,font=\sffamily\bfseries\LARGE] at ([xshift=5mm,yshift=-0.62cm]frame.north west) {Question 77};
\node[anchor=east,text=white,font=\sffamily\itshape\large] at ([xshift=-5mm,yshift=-0.62cm]frame.north east) {Géométrie dans l'espace};}]
\vspace*{1.75cm}
\begin{minipage}[t][6.9cm][c]{\linewidth}\LARGE\centering
\begin{minipage}{0.95\linewidth}\raggedright\linespread{1.25}\selectfont Dans un repère orthonormé, quelle est la distance entre $A(1;2;3)$ et $B(3;3;5)$~?\end{minipage}
\end{minipage}

\vfill
\begin{tcolorbox}[enhanced,nobeforeafter,width=7.1cm,height=2.1cm,valign=center,colback=cesp!6,colframe=cesp,arc=2.5mm,boxrule=0.9pt,left=3mm]{\Large\bfseries\color{cesp}a)}\quad{\Large $\sqrt{11}$}\end{tcolorbox}\hfill\begin{tcolorbox}[enhanced,nobeforeafter,width=7.1cm,height=2.1cm,valign=center,colback=cesp!6,colframe=cesp,arc=2.5mm,boxrule=0.9pt,left=3mm]{\Large\bfseries\color{cesp}b)}\quad{\Large $5$}\end{tcolorbox}\par\vspace{2.5mm}
\begin{tcolorbox}[enhanced,nobeforeafter,width=7.1cm,height=2.1cm,valign=center,colback=cesp!6,colframe=cesp,arc=2.5mm,boxrule=0.9pt,left=3mm]{\Large\bfseries\color{cesp}c)}\quad{\Large $9$}\end{tcolorbox}\hfill\begin{tcolorbox}[enhanced,nobeforeafter,width=7.1cm,height=2.1cm,valign=center,colback=cesp!6,colframe=cesp,arc=2.5mm,boxrule=0.9pt,left=3mm]{\Large\bfseries\color{cesp}d)}\quad{\Large $3$}\end{tcolorbox}
\end{tcolorbox}
\newpage
\begin{tcolorbox}[enhanced,width=16cm,height=14cm,colback=white,colframe=cesp,boxrule=2pt,arc=0mm,
left=6mm,right=6mm,top=0mm,bottom=4mm,sharp corners,
overlay={\fill[cesp] (frame.north west) rectangle ([yshift=-1.25cm]frame.north east);
\node[anchor=west,text=white,font=\sffamily\bfseries\LARGE] at ([xshift=5mm,yshift=-0.62cm]frame.north west) {Question 78};
\node[anchor=east,text=white,font=\sffamily\itshape\large] at ([xshift=-5mm,yshift=-0.62cm]frame.north east) {Géométrie dans l'espace};}]
\vspace*{1.75cm}
\begin{minipage}[t][6.9cm][c]{\linewidth}\LARGE\centering
\begin{minipage}{0.95\linewidth}\raggedright\linespread{1.25}\selectfont Dans un repère orthonormé, quelle est la distance entre $A(0;-1;2)$ et $B(2;1;3)$~?\end{minipage}
\end{minipage}

\vfill
\begin{tcolorbox}[enhanced,nobeforeafter,width=7.1cm,height=2.1cm,valign=center,colback=cesp!6,colframe=cesp,arc=2.5mm,boxrule=0.9pt,left=3mm]{\Large\bfseries\color{cesp}a)}\quad{\Large $3$}\end{tcolorbox}\hfill\begin{tcolorbox}[enhanced,nobeforeafter,width=7.1cm,height=2.1cm,valign=center,colback=cesp!6,colframe=cesp,arc=2.5mm,boxrule=0.9pt,left=3mm]{\Large\bfseries\color{cesp}b)}\quad{\Large $9$}\end{tcolorbox}\par\vspace{2.5mm}
\begin{tcolorbox}[enhanced,nobeforeafter,width=7.1cm,height=2.1cm,valign=center,colback=cesp!6,colframe=cesp,arc=2.5mm,boxrule=0.9pt,left=3mm]{\Large\bfseries\color{cesp}c)}\quad{\Large $\sqrt{11}$}\end{tcolorbox}\hfill\begin{tcolorbox}[enhanced,nobeforeafter,width=7.1cm,height=2.1cm,valign=center,colback=cesp!6,colframe=cesp,arc=2.5mm,boxrule=0.9pt,left=3mm]{\Large\bfseries\color{cesp}d)}\quad{\Large $5$}\end{tcolorbox}
\end{tcolorbox}
\newpage
\begin{tcolorbox}[enhanced,width=16cm,height=14cm,colback=white,colframe=cesp,boxrule=2pt,arc=0mm,
left=6mm,right=6mm,top=0mm,bottom=4mm,sharp corners,
overlay={\fill[cesp] (frame.north west) rectangle ([yshift=-1.25cm]frame.north east);
\node[anchor=west,text=white,font=\sffamily\bfseries\LARGE] at ([xshift=5mm,yshift=-0.62cm]frame.north west) {Question 79};
\node[anchor=east,text=white,font=\sffamily\itshape\large] at ([xshift=-5mm,yshift=-0.62cm]frame.north east) {Géométrie dans l'espace};}]
\vspace*{1.75cm}
\begin{minipage}[t][6.9cm][c]{\linewidth}\Large\centering
\begin{minipage}{0.95\linewidth}\raggedright\linespread{1.25}\selectfont Quel est le point d'intersection de la droite $\left\{\begin{array}{l}x=1+t\\y=1+2t\\z=0\end{array}\right.$ ($t\in\mathbb{R}$) et du plan $x+y+z=8$~?\end{minipage}
\end{minipage}

\vfill
\begin{tcolorbox}[enhanced,nobeforeafter,width=7.1cm,height=2.1cm,valign=center,colback=cesp!6,colframe=cesp,arc=2.5mm,boxrule=0.9pt,left=3mm]{\Large\bfseries\color{cesp}a)}\quad{\Large $(3;5;0)$}\end{tcolorbox}\hfill\begin{tcolorbox}[enhanced,nobeforeafter,width=7.1cm,height=2.1cm,valign=center,colback=cesp!6,colframe=cesp,arc=2.5mm,boxrule=0.9pt,left=3mm]{\Large\bfseries\color{cesp}b)}\quad{\Large $(2;3;0)$}\end{tcolorbox}\par\vspace{2.5mm}
\begin{tcolorbox}[enhanced,nobeforeafter,width=7.1cm,height=2.1cm,valign=center,colback=cesp!6,colframe=cesp,arc=2.5mm,boxrule=0.9pt,left=3mm]{\Large\bfseries\color{cesp}c)}\quad{\Large $(1;1;0)$}\end{tcolorbox}\hfill\begin{tcolorbox}[enhanced,nobeforeafter,width=7.1cm,height=2.1cm,valign=center,colback=cesp!6,colframe=cesp,arc=2.5mm,boxrule=0.9pt,left=3mm]{\Large\bfseries\color{cesp}d)}\quad{\Large $(4;7;0)$}\end{tcolorbox}
\end{tcolorbox}
\newpage
\begin{tcolorbox}[enhanced,width=16cm,height=14cm,colback=white,colframe=cesp,boxrule=2pt,arc=0mm,
left=6mm,right=6mm,top=0mm,bottom=4mm,sharp corners,
overlay={\fill[cesp] (frame.north west) rectangle ([yshift=-1.25cm]frame.north east);
\node[anchor=west,text=white,font=\sffamily\bfseries\LARGE] at ([xshift=5mm,yshift=-0.62cm]frame.north west) {Question 80};
\node[anchor=east,text=white,font=\sffamily\itshape\large] at ([xshift=-5mm,yshift=-0.62cm]frame.north east) {Géométrie dans l'espace};}]
\vspace*{1.75cm}
\begin{minipage}[t][6.9cm][c]{\linewidth}\Large\centering
\begin{minipage}{0.95\linewidth}\raggedright\linespread{1.25}\selectfont Quel est le point d'intersection de la droite $\left\{\begin{array}{l}x=t\\y=1\\z=2+t\end{array}\right.$ ($t\in\mathbb{R}$) et du plan $x+y+z=5$~?\end{minipage}
\end{minipage}

\vfill
\begin{tcolorbox}[enhanced,nobeforeafter,width=7.1cm,height=2.1cm,valign=center,colback=cesp!6,colframe=cesp,arc=2.5mm,boxrule=0.9pt,left=3mm]{\Large\bfseries\color{cesp}a)}\quad{\Large $(-1;1;1)$}\end{tcolorbox}\hfill\begin{tcolorbox}[enhanced,nobeforeafter,width=7.1cm,height=2.1cm,valign=center,colback=cesp!6,colframe=cesp,arc=2.5mm,boxrule=0.9pt,left=3mm]{\Large\bfseries\color{cesp}b)}\quad{\Large $(1;1;3)$}\end{tcolorbox}\par\vspace{2.5mm}
\begin{tcolorbox}[enhanced,nobeforeafter,width=7.1cm,height=2.1cm,valign=center,colback=cesp!6,colframe=cesp,arc=2.5mm,boxrule=0.9pt,left=3mm]{\Large\bfseries\color{cesp}c)}\quad{\Large $(2;1;4)$}\end{tcolorbox}\hfill\begin{tcolorbox}[enhanced,nobeforeafter,width=7.1cm,height=2.1cm,valign=center,colback=cesp!6,colframe=cesp,arc=2.5mm,boxrule=0.9pt,left=3mm]{\Large\bfseries\color{cesp}d)}\quad{\Large $(0;1;2)$}\end{tcolorbox}
\end{tcolorbox}
\newpage
\begin{tcolorbox}[enhanced,width=16cm,height=14cm,colback=white,colframe=cprb,boxrule=2pt,arc=0mm,
left=6mm,right=6mm,top=0mm,bottom=4mm,sharp corners,
overlay={\fill[cprb] (frame.north west) rectangle ([yshift=-1.25cm]frame.north east);
\node[anchor=west,text=white,font=\sffamily\bfseries\LARGE] at ([xshift=5mm,yshift=-0.62cm]frame.north west) {Question 81};
\node[anchor=east,text=white,font=\sffamily\itshape\large] at ([xshift=-5mm,yshift=-0.62cm]frame.north east) {Probabilités et dénombrement};}]
\vspace*{1.75cm}
\begin{minipage}[t][6.9cm][c]{\linewidth}\LARGE\centering
\begin{minipage}{0.95\linewidth}\raggedright\linespread{1.25}\selectfont $X$ suit la loi binomiale $\mathcal{B}\left(4;\frac12\right)$. Que vaut $P(X=2)$~?\end{minipage}
\end{minipage}

\vfill
\begin{tcolorbox}[enhanced,nobeforeafter,width=7.1cm,height=2.1cm,valign=center,colback=cprb!6,colframe=cprb,arc=2.5mm,boxrule=0.9pt,left=3mm]{\Large\bfseries\color{cprb}a)}\quad{\Large $\dfrac{3}{2}$}\end{tcolorbox}\hfill\begin{tcolorbox}[enhanced,nobeforeafter,width=7.1cm,height=2.1cm,valign=center,colback=cprb!6,colframe=cprb,arc=2.5mm,boxrule=0.9pt,left=3mm]{\Large\bfseries\color{cprb}b)}\quad{\Large $\dfrac{3}{8}$}\end{tcolorbox}\par\vspace{2.5mm}
\begin{tcolorbox}[enhanced,nobeforeafter,width=7.1cm,height=2.1cm,valign=center,colback=cprb!6,colframe=cprb,arc=2.5mm,boxrule=0.9pt,left=3mm]{\Large\bfseries\color{cprb}c)}\quad{\Large $\dfrac{1}{2}$}\end{tcolorbox}\hfill\begin{tcolorbox}[enhanced,nobeforeafter,width=7.1cm,height=2.1cm,valign=center,colback=cprb!6,colframe=cprb,arc=2.5mm,boxrule=0.9pt,left=3mm]{\Large\bfseries\color{cprb}d)}\quad{\Large $\dfrac{1}{16}$}\end{tcolorbox}
\end{tcolorbox}
\newpage
\begin{tcolorbox}[enhanced,width=16cm,height=14cm,colback=white,colframe=cprb,boxrule=2pt,arc=0mm,
left=6mm,right=6mm,top=0mm,bottom=4mm,sharp corners,
overlay={\fill[cprb] (frame.north west) rectangle ([yshift=-1.25cm]frame.north east);
\node[anchor=west,text=white,font=\sffamily\bfseries\LARGE] at ([xshift=5mm,yshift=-0.62cm]frame.north west) {Question 82};
\node[anchor=east,text=white,font=\sffamily\itshape\large] at ([xshift=-5mm,yshift=-0.62cm]frame.north east) {Probabilités et dénombrement};}]
\vspace*{1.75cm}
\begin{minipage}[t][6.9cm][c]{\linewidth}\LARGE\centering
\begin{minipage}{0.95\linewidth}\raggedright\linespread{1.25}\selectfont $X$ suit la loi binomiale $\mathcal{B}\left(5;\frac12\right)$. Que vaut $P(X=1)$~?\end{minipage}
\end{minipage}

\vfill
\begin{tcolorbox}[enhanced,nobeforeafter,width=7.1cm,height=2.1cm,valign=center,colback=cprb!6,colframe=cprb,arc=2.5mm,boxrule=0.9pt,left=3mm]{\Large\bfseries\color{cprb}a)}\quad{\Large $\dfrac{1}{32}$}\end{tcolorbox}\hfill\begin{tcolorbox}[enhanced,nobeforeafter,width=7.1cm,height=2.1cm,valign=center,colback=cprb!6,colframe=cprb,arc=2.5mm,boxrule=0.9pt,left=3mm]{\Large\bfseries\color{cprb}b)}\quad{\Large $\dfrac{1}{5}$}\end{tcolorbox}\par\vspace{2.5mm}
\begin{tcolorbox}[enhanced,nobeforeafter,width=7.1cm,height=2.1cm,valign=center,colback=cprb!6,colframe=cprb,arc=2.5mm,boxrule=0.9pt,left=3mm]{\Large\bfseries\color{cprb}c)}\quad{\Large $\dfrac{5}{32}$}\end{tcolorbox}\hfill\begin{tcolorbox}[enhanced,nobeforeafter,width=7.1cm,height=2.1cm,valign=center,colback=cprb!6,colframe=cprb,arc=2.5mm,boxrule=0.9pt,left=3mm]{\Large\bfseries\color{cprb}d)}\quad{\Large $\dfrac{5}{2}$}\end{tcolorbox}
\end{tcolorbox}
\newpage
\begin{tcolorbox}[enhanced,width=16cm,height=14cm,colback=white,colframe=cprb,boxrule=2pt,arc=0mm,
left=6mm,right=6mm,top=0mm,bottom=4mm,sharp corners,
overlay={\fill[cprb] (frame.north west) rectangle ([yshift=-1.25cm]frame.north east);
\node[anchor=west,text=white,font=\sffamily\bfseries\LARGE] at ([xshift=5mm,yshift=-0.62cm]frame.north west) {Question 83};
\node[anchor=east,text=white,font=\sffamily\itshape\large] at ([xshift=-5mm,yshift=-0.62cm]frame.north east) {Probabilités et dénombrement};}]
\vspace*{1.75cm}
\begin{minipage}[t][6.9cm][c]{\linewidth}\LARGE\centering
\begin{minipage}{0.95\linewidth}\raggedright\linespread{1.25}\selectfont $X$ suit la loi binomiale $\mathcal{B}\left(20;\frac{1}{4}\right)$. Quelle est son espérance $E(X)$~?\end{minipage}
\end{minipage}

\vfill
\begin{tcolorbox}[enhanced,nobeforeafter,width=7.1cm,height=2.1cm,valign=center,colback=cprb!6,colframe=cprb,arc=2.5mm,boxrule=0.9pt,left=3mm]{\Large\bfseries\color{cprb}a)}\quad{\Large $\dfrac{15}{4}$}\end{tcolorbox}\hfill\begin{tcolorbox}[enhanced,nobeforeafter,width=7.1cm,height=2.1cm,valign=center,colback=cprb!6,colframe=cprb,arc=2.5mm,boxrule=0.9pt,left=3mm]{\Large\bfseries\color{cprb}b)}\quad{\Large $5$}\end{tcolorbox}\par\vspace{2.5mm}
\begin{tcolorbox}[enhanced,nobeforeafter,width=7.1cm,height=2.1cm,valign=center,colback=cprb!6,colframe=cprb,arc=2.5mm,boxrule=0.9pt,left=3mm]{\Large\bfseries\color{cprb}c)}\quad{\Large $15$}\end{tcolorbox}\hfill\begin{tcolorbox}[enhanced,nobeforeafter,width=7.1cm,height=2.1cm,valign=center,colback=cprb!6,colframe=cprb,arc=2.5mm,boxrule=0.9pt,left=3mm]{\Large\bfseries\color{cprb}d)}\quad{\Large $\dfrac{1}{4}$}\end{tcolorbox}
\end{tcolorbox}
\newpage
\begin{tcolorbox}[enhanced,width=16cm,height=14cm,colback=white,colframe=cprb,boxrule=2pt,arc=0mm,
left=6mm,right=6mm,top=0mm,bottom=4mm,sharp corners,
overlay={\fill[cprb] (frame.north west) rectangle ([yshift=-1.25cm]frame.north east);
\node[anchor=west,text=white,font=\sffamily\bfseries\LARGE] at ([xshift=5mm,yshift=-0.62cm]frame.north west) {Question 84};
\node[anchor=east,text=white,font=\sffamily\itshape\large] at ([xshift=-5mm,yshift=-0.62cm]frame.north east) {Probabilités et dénombrement};}]
\vspace*{1.75cm}
\begin{minipage}[t][6.9cm][c]{\linewidth}\LARGE\centering
\begin{minipage}{0.95\linewidth}\raggedright\linespread{1.25}\selectfont $X$ suit la loi binomiale $\mathcal{B}\left(30;\frac{1}{3}\right)$. Quelle est son espérance $E(X)$~?\end{minipage}
\end{minipage}

\vfill
\begin{tcolorbox}[enhanced,nobeforeafter,width=7.1cm,height=2.1cm,valign=center,colback=cprb!6,colframe=cprb,arc=2.5mm,boxrule=0.9pt,left=3mm]{\Large\bfseries\color{cprb}a)}\quad{\Large $\dfrac{1}{3}$}\end{tcolorbox}\hfill\begin{tcolorbox}[enhanced,nobeforeafter,width=7.1cm,height=2.1cm,valign=center,colback=cprb!6,colframe=cprb,arc=2.5mm,boxrule=0.9pt,left=3mm]{\Large\bfseries\color{cprb}b)}\quad{\Large $10$}\end{tcolorbox}\par\vspace{2.5mm}
\begin{tcolorbox}[enhanced,nobeforeafter,width=7.1cm,height=2.1cm,valign=center,colback=cprb!6,colframe=cprb,arc=2.5mm,boxrule=0.9pt,left=3mm]{\Large\bfseries\color{cprb}c)}\quad{\Large $20$}\end{tcolorbox}\hfill\begin{tcolorbox}[enhanced,nobeforeafter,width=7.1cm,height=2.1cm,valign=center,colback=cprb!6,colframe=cprb,arc=2.5mm,boxrule=0.9pt,left=3mm]{\Large\bfseries\color{cprb}d)}\quad{\Large $\dfrac{20}{3}$}\end{tcolorbox}
\end{tcolorbox}
\newpage
\begin{tcolorbox}[enhanced,width=16cm,height=14cm,colback=white,colframe=cprb,boxrule=2pt,arc=0mm,
left=6mm,right=6mm,top=0mm,bottom=4mm,sharp corners,
overlay={\fill[cprb] (frame.north west) rectangle ([yshift=-1.25cm]frame.north east);
\node[anchor=west,text=white,font=\sffamily\bfseries\LARGE] at ([xshift=5mm,yshift=-0.62cm]frame.north west) {Question 85};
\node[anchor=east,text=white,font=\sffamily\itshape\large] at ([xshift=-5mm,yshift=-0.62cm]frame.north east) {Probabilités et dénombrement};}]
\vspace*{1.75cm}
\begin{minipage}[t][6.9cm][c]{\linewidth}\LARGE\centering
\begin{minipage}{0.95\linewidth}\raggedright\linespread{1.25}\selectfont $X$ suit la loi binomiale $\mathcal{B}\left(10;\frac{1}{2}\right)$. Quelle est sa variance $V(X)$~?\end{minipage}
\end{minipage}

\vfill
\begin{tcolorbox}[enhanced,nobeforeafter,width=7.1cm,height=2.1cm,valign=center,colback=cprb!6,colframe=cprb,arc=2.5mm,boxrule=0.9pt,left=3mm]{\Large\bfseries\color{cprb}a)}\quad{\Large $\dfrac{5}{4}$}\end{tcolorbox}\hfill\begin{tcolorbox}[enhanced,nobeforeafter,width=7.1cm,height=2.1cm,valign=center,colback=cprb!6,colframe=cprb,arc=2.5mm,boxrule=0.9pt,left=3mm]{\Large\bfseries\color{cprb}b)}\quad{\Large $\dfrac{5}{2}$}\end{tcolorbox}\par\vspace{2.5mm}
\begin{tcolorbox}[enhanced,nobeforeafter,width=7.1cm,height=2.1cm,valign=center,colback=cprb!6,colframe=cprb,arc=2.5mm,boxrule=0.9pt,left=3mm]{\Large\bfseries\color{cprb}c)}\quad{\Large $5$}\end{tcolorbox}\hfill\begin{tcolorbox}[enhanced,nobeforeafter,width=7.1cm,height=2.1cm,valign=center,colback=cprb!6,colframe=cprb,arc=2.5mm,boxrule=0.9pt,left=3mm]{\Large\bfseries\color{cprb}d)}\quad{\Large $\dfrac{7}{2}$}\end{tcolorbox}
\end{tcolorbox}
\newpage
\begin{tcolorbox}[enhanced,width=16cm,height=14cm,colback=white,colframe=cprb,boxrule=2pt,arc=0mm,
left=6mm,right=6mm,top=0mm,bottom=4mm,sharp corners,
overlay={\fill[cprb] (frame.north west) rectangle ([yshift=-1.25cm]frame.north east);
\node[anchor=west,text=white,font=\sffamily\bfseries\LARGE] at ([xshift=5mm,yshift=-0.62cm]frame.north west) {Question 86};
\node[anchor=east,text=white,font=\sffamily\itshape\large] at ([xshift=-5mm,yshift=-0.62cm]frame.north east) {Probabilités et dénombrement};}]
\vspace*{1.75cm}
\begin{minipage}[t][6.9cm][c]{\linewidth}\LARGE\centering
\begin{minipage}{0.95\linewidth}\raggedright\linespread{1.25}\selectfont $X$ suit la loi binomiale $\mathcal{B}\left(20;\frac{1}{4}\right)$. Quelle est sa variance $V(X)$~?\end{minipage}
\end{minipage}

\vfill
\begin{tcolorbox}[enhanced,nobeforeafter,width=7.1cm,height=2.1cm,valign=center,colback=cprb!6,colframe=cprb,arc=2.5mm,boxrule=0.9pt,left=3mm]{\Large\bfseries\color{cprb}a)}\quad{\Large $\dfrac{15}{4}$}\end{tcolorbox}\hfill\begin{tcolorbox}[enhanced,nobeforeafter,width=7.1cm,height=2.1cm,valign=center,colback=cprb!6,colframe=cprb,arc=2.5mm,boxrule=0.9pt,left=3mm]{\Large\bfseries\color{cprb}b)}\quad{\Large $\dfrac{5}{4}$}\end{tcolorbox}\par\vspace{2.5mm}
\begin{tcolorbox}[enhanced,nobeforeafter,width=7.1cm,height=2.1cm,valign=center,colback=cprb!6,colframe=cprb,arc=2.5mm,boxrule=0.9pt,left=3mm]{\Large\bfseries\color{cprb}c)}\quad{\Large $\dfrac{45}{4}$}\end{tcolorbox}\hfill\begin{tcolorbox}[enhanced,nobeforeafter,width=7.1cm,height=2.1cm,valign=center,colback=cprb!6,colframe=cprb,arc=2.5mm,boxrule=0.9pt,left=3mm]{\Large\bfseries\color{cprb}d)}\quad{\Large $5$}\end{tcolorbox}
\end{tcolorbox}
\newpage
\begin{tcolorbox}[enhanced,width=16cm,height=14cm,colback=white,colframe=cprb,boxrule=2pt,arc=0mm,
left=6mm,right=6mm,top=0mm,bottom=4mm,sharp corners,
overlay={\fill[cprb] (frame.north west) rectangle ([yshift=-1.25cm]frame.north east);
\node[anchor=west,text=white,font=\sffamily\bfseries\LARGE] at ([xshift=5mm,yshift=-0.62cm]frame.north west) {Question 87};
\node[anchor=east,text=white,font=\sffamily\itshape\large] at ([xshift=-5mm,yshift=-0.62cm]frame.north east) {Probabilités et dénombrement};}]
\vspace*{1.75cm}
\begin{minipage}[t][6.9cm][c]{\linewidth}\LARGE\centering
\begin{minipage}{0.95\linewidth}\raggedright\linespread{1.25}\selectfont On a $P(A)=0{,}4$ et $P_A(B)=0{,}3$. Que vaut $P(A\cap B)$~?\end{minipage}
\end{minipage}

\vfill
\begin{tcolorbox}[enhanced,nobeforeafter,width=7.1cm,height=2.1cm,valign=center,colback=cprb!6,colframe=cprb,arc=2.5mm,boxrule=0.9pt,left=3mm]{\Large\bfseries\color{cprb}a)}\quad{\Large $0{,}12$}\end{tcolorbox}\hfill\begin{tcolorbox}[enhanced,nobeforeafter,width=7.1cm,height=2.1cm,valign=center,colback=cprb!6,colframe=cprb,arc=2.5mm,boxrule=0.9pt,left=3mm]{\Large\bfseries\color{cprb}b)}\quad{\Large $0{,}7$}\end{tcolorbox}\par\vspace{2.5mm}
\begin{tcolorbox}[enhanced,nobeforeafter,width=7.1cm,height=2.1cm,valign=center,colback=cprb!6,colframe=cprb,arc=2.5mm,boxrule=0.9pt,left=3mm]{\Large\bfseries\color{cprb}c)}\quad{\Large $0{,}1$}\end{tcolorbox}\hfill\begin{tcolorbox}[enhanced,nobeforeafter,width=7.1cm,height=2.1cm,valign=center,colback=cprb!6,colframe=cprb,arc=2.5mm,boxrule=0.9pt,left=3mm]{\Large\bfseries\color{cprb}d)}\quad{\Large $0{,}28$}\end{tcolorbox}
\end{tcolorbox}
\newpage
\begin{tcolorbox}[enhanced,width=16cm,height=14cm,colback=white,colframe=cprb,boxrule=2pt,arc=0mm,
left=6mm,right=6mm,top=0mm,bottom=4mm,sharp corners,
overlay={\fill[cprb] (frame.north west) rectangle ([yshift=-1.25cm]frame.north east);
\node[anchor=west,text=white,font=\sffamily\bfseries\LARGE] at ([xshift=5mm,yshift=-0.62cm]frame.north west) {Question 88};
\node[anchor=east,text=white,font=\sffamily\itshape\large] at ([xshift=-5mm,yshift=-0.62cm]frame.north east) {Probabilités et dénombrement};}]
\vspace*{1.75cm}
\begin{minipage}[t][6.9cm][c]{\linewidth}\LARGE\centering
\begin{minipage}{0.95\linewidth}\raggedright\linespread{1.25}\selectfont On a $P(A)=0{,}6$ et $P_A(B)=0{,}5$. Que vaut $P(A\cap B)$~?\end{minipage}
\end{minipage}

\vfill
\begin{tcolorbox}[enhanced,nobeforeafter,width=7.1cm,height=2.1cm,valign=center,colback=cprb!6,colframe=cprb,arc=2.5mm,boxrule=0.9pt,left=3mm]{\Large\bfseries\color{cprb}a)}\quad{\Large $0{,}6$}\end{tcolorbox}\hfill\begin{tcolorbox}[enhanced,nobeforeafter,width=7.1cm,height=2.1cm,valign=center,colback=cprb!6,colframe=cprb,arc=2.5mm,boxrule=0.9pt,left=3mm]{\Large\bfseries\color{cprb}b)}\quad{\Large $0{,}3$}\end{tcolorbox}\par\vspace{2.5mm}
\begin{tcolorbox}[enhanced,nobeforeafter,width=7.1cm,height=2.1cm,valign=center,colback=cprb!6,colframe=cprb,arc=2.5mm,boxrule=0.9pt,left=3mm]{\Large\bfseries\color{cprb}c)}\quad{\Large $0{,}5$}\end{tcolorbox}\hfill\begin{tcolorbox}[enhanced,nobeforeafter,width=7.1cm,height=2.1cm,valign=center,colback=cprb!6,colframe=cprb,arc=2.5mm,boxrule=0.9pt,left=3mm]{\Large\bfseries\color{cprb}d)}\quad{\Large $0{,}1$}\end{tcolorbox}
\end{tcolorbox}
\newpage
\begin{tcolorbox}[enhanced,width=16cm,height=14cm,colback=white,colframe=cprb,boxrule=2pt,arc=0mm,
left=6mm,right=6mm,top=0mm,bottom=4mm,sharp corners,
overlay={\fill[cprb] (frame.north west) rectangle ([yshift=-1.25cm]frame.north east);
\node[anchor=west,text=white,font=\sffamily\bfseries\LARGE] at ([xshift=5mm,yshift=-0.62cm]frame.north west) {Question 89};
\node[anchor=east,text=white,font=\sffamily\itshape\large] at ([xshift=-5mm,yshift=-0.62cm]frame.north east) {Probabilités et dénombrement};}]
\vspace*{1.75cm}
\begin{minipage}[t][6.9cm][c]{\linewidth}\LARGE\centering
\begin{minipage}{0.95\linewidth}\raggedright\linespread{1.25}\selectfont On a $P(A)=0{,}3$, $P_A(B)=0{,}5$ et $P_{\overline{A}}(B)=0{,}2$. Que vaut $P(B)$~?\end{minipage}
\end{minipage}

\vfill
\begin{tcolorbox}[enhanced,nobeforeafter,width=7.1cm,height=2.1cm,valign=center,colback=cprb!6,colframe=cprb,arc=2.5mm,boxrule=0.9pt,left=3mm]{\Large\bfseries\color{cprb}a)}\quad{\Large $0{,}7$}\end{tcolorbox}\hfill\begin{tcolorbox}[enhanced,nobeforeafter,width=7.1cm,height=2.1cm,valign=center,colback=cprb!6,colframe=cprb,arc=2.5mm,boxrule=0.9pt,left=3mm]{\Large\bfseries\color{cprb}b)}\quad{\Large $0{,}21$}\end{tcolorbox}\par\vspace{2.5mm}
\begin{tcolorbox}[enhanced,nobeforeafter,width=7.1cm,height=2.1cm,valign=center,colback=cprb!6,colframe=cprb,arc=2.5mm,boxrule=0.9pt,left=3mm]{\Large\bfseries\color{cprb}c)}\quad{\Large $0{,}15$}\end{tcolorbox}\hfill\begin{tcolorbox}[enhanced,nobeforeafter,width=7.1cm,height=2.1cm,valign=center,colback=cprb!6,colframe=cprb,arc=2.5mm,boxrule=0.9pt,left=3mm]{\Large\bfseries\color{cprb}d)}\quad{\Large $0{,}29$}\end{tcolorbox}
\end{tcolorbox}
\newpage
\begin{tcolorbox}[enhanced,width=16cm,height=14cm,colback=white,colframe=cprb,boxrule=2pt,arc=0mm,
left=6mm,right=6mm,top=0mm,bottom=4mm,sharp corners,
overlay={\fill[cprb] (frame.north west) rectangle ([yshift=-1.25cm]frame.north east);
\node[anchor=west,text=white,font=\sffamily\bfseries\LARGE] at ([xshift=5mm,yshift=-0.62cm]frame.north west) {Question 90};
\node[anchor=east,text=white,font=\sffamily\itshape\large] at ([xshift=-5mm,yshift=-0.62cm]frame.north east) {Probabilités et dénombrement};}]
\vspace*{1.75cm}
\begin{minipage}[t][6.9cm][c]{\linewidth}\LARGE\centering
\begin{minipage}{0.95\linewidth}\raggedright\linespread{1.25}\selectfont On a $P(A)=0{,}6$, $P_A(B)=0{,}4$ et $P_{\overline{A}}(B)=0{,}1$. Que vaut $P(B)$~?\end{minipage}
\end{minipage}

\vfill
\begin{tcolorbox}[enhanced,nobeforeafter,width=7.1cm,height=2.1cm,valign=center,colback=cprb!6,colframe=cprb,arc=2.5mm,boxrule=0.9pt,left=3mm]{\Large\bfseries\color{cprb}a)}\quad{\Large $0{,}5$}\end{tcolorbox}\hfill\begin{tcolorbox}[enhanced,nobeforeafter,width=7.1cm,height=2.1cm,valign=center,colback=cprb!6,colframe=cprb,arc=2.5mm,boxrule=0.9pt,left=3mm]{\Large\bfseries\color{cprb}b)}\quad{\Large $0{,}28$}\end{tcolorbox}\par\vspace{2.5mm}
\begin{tcolorbox}[enhanced,nobeforeafter,width=7.1cm,height=2.1cm,valign=center,colback=cprb!6,colframe=cprb,arc=2.5mm,boxrule=0.9pt,left=3mm]{\Large\bfseries\color{cprb}c)}\quad{\Large $0{,}3$}\end{tcolorbox}\hfill\begin{tcolorbox}[enhanced,nobeforeafter,width=7.1cm,height=2.1cm,valign=center,colback=cprb!6,colframe=cprb,arc=2.5mm,boxrule=0.9pt,left=3mm]{\Large\bfseries\color{cprb}d)}\quad{\Large $0{,}24$}\end{tcolorbox}
\end{tcolorbox}
\newpage
\begin{tcolorbox}[enhanced,width=16cm,height=14cm,colback=white,colframe=cprb,boxrule=2pt,arc=0mm,
left=6mm,right=6mm,top=0mm,bottom=4mm,sharp corners,
overlay={\fill[cprb] (frame.north west) rectangle ([yshift=-1.25cm]frame.north east);
\node[anchor=west,text=white,font=\sffamily\bfseries\LARGE] at ([xshift=5mm,yshift=-0.62cm]frame.north west) {Question 91};
\node[anchor=east,text=white,font=\sffamily\itshape\large] at ([xshift=-5mm,yshift=-0.62cm]frame.north east) {Probabilités et dénombrement};}]
\vspace*{1.75cm}
\begin{minipage}[t][6.9cm][c]{\linewidth}\LARGE\centering
\begin{minipage}{0.95\linewidth}\raggedright\linespread{1.25}\selectfont Combien y a-t-il de façons de choisir 2 élèves parmi 5 (sans tenir compte de l'ordre)~?\end{minipage}
\end{minipage}

\vfill
\begin{tcolorbox}[enhanced,nobeforeafter,width=7.1cm,height=2.1cm,valign=center,colback=cprb!6,colframe=cprb,arc=2.5mm,boxrule=0.9pt,left=3mm]{\Large\bfseries\color{cprb}a)}\quad{\Large $20$}\end{tcolorbox}\hfill\begin{tcolorbox}[enhanced,nobeforeafter,width=7.1cm,height=2.1cm,valign=center,colback=cprb!6,colframe=cprb,arc=2.5mm,boxrule=0.9pt,left=3mm]{\Large\bfseries\color{cprb}b)}\quad{\Large $15$}\end{tcolorbox}\par\vspace{2.5mm}
\begin{tcolorbox}[enhanced,nobeforeafter,width=7.1cm,height=2.1cm,valign=center,colback=cprb!6,colframe=cprb,arc=2.5mm,boxrule=0.9pt,left=3mm]{\Large\bfseries\color{cprb}c)}\quad{\Large $10$}\end{tcolorbox}\hfill\begin{tcolorbox}[enhanced,nobeforeafter,width=7.1cm,height=2.1cm,valign=center,colback=cprb!6,colframe=cprb,arc=2.5mm,boxrule=0.9pt,left=3mm]{\Large\bfseries\color{cprb}d)}\quad{\Large $25$}\end{tcolorbox}
\end{tcolorbox}
\newpage
\begin{tcolorbox}[enhanced,width=16cm,height=14cm,colback=white,colframe=cprb,boxrule=2pt,arc=0mm,
left=6mm,right=6mm,top=0mm,bottom=4mm,sharp corners,
overlay={\fill[cprb] (frame.north west) rectangle ([yshift=-1.25cm]frame.north east);
\node[anchor=west,text=white,font=\sffamily\bfseries\LARGE] at ([xshift=5mm,yshift=-0.62cm]frame.north west) {Question 92};
\node[anchor=east,text=white,font=\sffamily\itshape\large] at ([xshift=-5mm,yshift=-0.62cm]frame.north east) {Probabilités et dénombrement};}]
\vspace*{1.75cm}
\begin{minipage}[t][6.9cm][c]{\linewidth}\LARGE\centering
\begin{minipage}{0.95\linewidth}\raggedright\linespread{1.25}\selectfont Combien y a-t-il de façons de choisir 3 élèves parmi 6 (sans tenir compte de l'ordre)~?\end{minipage}
\end{minipage}

\vfill
\begin{tcolorbox}[enhanced,nobeforeafter,width=7.1cm,height=2.1cm,valign=center,colback=cprb!6,colframe=cprb,arc=2.5mm,boxrule=0.9pt,left=3mm]{\Large\bfseries\color{cprb}a)}\quad{\Large $20$}\end{tcolorbox}\hfill\begin{tcolorbox}[enhanced,nobeforeafter,width=7.1cm,height=2.1cm,valign=center,colback=cprb!6,colframe=cprb,arc=2.5mm,boxrule=0.9pt,left=3mm]{\Large\bfseries\color{cprb}b)}\quad{\Large $120$}\end{tcolorbox}\par\vspace{2.5mm}
\begin{tcolorbox}[enhanced,nobeforeafter,width=7.1cm,height=2.1cm,valign=center,colback=cprb!6,colframe=cprb,arc=2.5mm,boxrule=0.9pt,left=3mm]{\Large\bfseries\color{cprb}c)}\quad{\Large $216$}\end{tcolorbox}\hfill\begin{tcolorbox}[enhanced,nobeforeafter,width=7.1cm,height=2.1cm,valign=center,colback=cprb!6,colframe=cprb,arc=2.5mm,boxrule=0.9pt,left=3mm]{\Large\bfseries\color{cprb}d)}\quad{\Large $18$}\end{tcolorbox}
\end{tcolorbox}
\newpage
\begin{tcolorbox}[enhanced,width=16cm,height=14cm,colback=white,colframe=cprb,boxrule=2pt,arc=0mm,
left=6mm,right=6mm,top=0mm,bottom=4mm,sharp corners,
overlay={\fill[cprb] (frame.north west) rectangle ([yshift=-1.25cm]frame.north east);
\node[anchor=west,text=white,font=\sffamily\bfseries\LARGE] at ([xshift=5mm,yshift=-0.62cm]frame.north west) {Question 93};
\node[anchor=east,text=white,font=\sffamily\itshape\large] at ([xshift=-5mm,yshift=-0.62cm]frame.north east) {Probabilités et dénombrement};}]
\vspace*{1.75cm}
\begin{minipage}[t][6.9cm][c]{\linewidth}\LARGE\centering
\begin{minipage}{0.95\linewidth}\raggedright\linespread{1.25}\selectfont Combien de codes de 4 caractères peut-on former avec un alphabet de 3 symboles (répétitions autorisées)~?\end{minipage}
\end{minipage}

\vfill
\begin{tcolorbox}[enhanced,nobeforeafter,width=7.1cm,height=2.1cm,valign=center,colback=cprb!6,colframe=cprb,arc=2.5mm,boxrule=0.9pt,left=3mm]{\Large\bfseries\color{cprb}a)}\quad{\Large $81$}\end{tcolorbox}\hfill\begin{tcolorbox}[enhanced,nobeforeafter,width=7.1cm,height=2.1cm,valign=center,colback=cprb!6,colframe=cprb,arc=2.5mm,boxrule=0.9pt,left=3mm]{\Large\bfseries\color{cprb}b)}\quad{\Large $64$}\end{tcolorbox}\par\vspace{2.5mm}
\begin{tcolorbox}[enhanced,nobeforeafter,width=7.1cm,height=2.1cm,valign=center,colback=cprb!6,colframe=cprb,arc=2.5mm,boxrule=0.9pt,left=3mm]{\Large\bfseries\color{cprb}c)}\quad{\Large $12$}\end{tcolorbox}\hfill\begin{tcolorbox}[enhanced,nobeforeafter,width=7.1cm,height=2.1cm,valign=center,colback=cprb!6,colframe=cprb,arc=2.5mm,boxrule=0.9pt,left=3mm]{\Large\bfseries\color{cprb}d)}\quad{\Large $15$}\end{tcolorbox}
\end{tcolorbox}
\newpage
\begin{tcolorbox}[enhanced,width=16cm,height=14cm,colback=white,colframe=cprb,boxrule=2pt,arc=0mm,
left=6mm,right=6mm,top=0mm,bottom=4mm,sharp corners,
overlay={\fill[cprb] (frame.north west) rectangle ([yshift=-1.25cm]frame.north east);
\node[anchor=west,text=white,font=\sffamily\bfseries\LARGE] at ([xshift=5mm,yshift=-0.62cm]frame.north west) {Question 94};
\node[anchor=east,text=white,font=\sffamily\itshape\large] at ([xshift=-5mm,yshift=-0.62cm]frame.north east) {Probabilités et dénombrement};}]
\vspace*{1.75cm}
\begin{minipage}[t][6.9cm][c]{\linewidth}\LARGE\centering
\begin{minipage}{0.95\linewidth}\raggedright\linespread{1.25}\selectfont Combien de codes de 5 caractères peut-on former avec un alphabet de 2 symboles (répétitions autorisées)~?\end{minipage}
\end{minipage}

\vfill
\begin{tcolorbox}[enhanced,nobeforeafter,width=7.1cm,height=2.1cm,valign=center,colback=cprb!6,colframe=cprb,arc=2.5mm,boxrule=0.9pt,left=3mm]{\Large\bfseries\color{cprb}a)}\quad{\Large $6$}\end{tcolorbox}\hfill\begin{tcolorbox}[enhanced,nobeforeafter,width=7.1cm,height=2.1cm,valign=center,colback=cprb!6,colframe=cprb,arc=2.5mm,boxrule=0.9pt,left=3mm]{\Large\bfseries\color{cprb}b)}\quad{\Large $10$}\end{tcolorbox}\par\vspace{2.5mm}
\begin{tcolorbox}[enhanced,nobeforeafter,width=7.1cm,height=2.1cm,valign=center,colback=cprb!6,colframe=cprb,arc=2.5mm,boxrule=0.9pt,left=3mm]{\Large\bfseries\color{cprb}c)}\quad{\Large $32$}\end{tcolorbox}\hfill\begin{tcolorbox}[enhanced,nobeforeafter,width=7.1cm,height=2.1cm,valign=center,colback=cprb!6,colframe=cprb,arc=2.5mm,boxrule=0.9pt,left=3mm]{\Large\bfseries\color{cprb}d)}\quad{\Large $25$}\end{tcolorbox}
\end{tcolorbox}
\newpage
\begin{tcolorbox}[enhanced,width=16cm,height=14cm,colback=white,colframe=cprb,boxrule=2pt,arc=0mm,
left=6mm,right=6mm,top=0mm,bottom=4mm,sharp corners,
overlay={\fill[cprb] (frame.north west) rectangle ([yshift=-1.25cm]frame.north east);
\node[anchor=west,text=white,font=\sffamily\bfseries\LARGE] at ([xshift=5mm,yshift=-0.62cm]frame.north west) {Question 95};
\node[anchor=east,text=white,font=\sffamily\itshape\large] at ([xshift=-5mm,yshift=-0.62cm]frame.north east) {Probabilités et dénombrement};}]
\vspace*{1.75cm}
\begin{minipage}[t][6.9cm][c]{\linewidth}\LARGE\centering
\begin{minipage}{0.95\linewidth}\raggedright\linespread{1.25}\selectfont $X$ suit la loi binomiale $\mathcal{B}\left(3;\frac{1}{2}\right)$. Que vaut $P(X\geq 1)$~?\end{minipage}
\end{minipage}

\vfill
\begin{tcolorbox}[enhanced,nobeforeafter,width=7.1cm,height=2.1cm,valign=center,colback=cprb!6,colframe=cprb,arc=2.5mm,boxrule=0.9pt,left=3mm]{\Large\bfseries\color{cprb}a)}\quad{\Large $\dfrac{7}{8}$}\end{tcolorbox}\hfill\begin{tcolorbox}[enhanced,nobeforeafter,width=7.1cm,height=2.1cm,valign=center,colback=cprb!6,colframe=cprb,arc=2.5mm,boxrule=0.9pt,left=3mm]{\Large\bfseries\color{cprb}b)}\quad{\Large $\dfrac{1}{8}$}\end{tcolorbox}\par\vspace{2.5mm}
\begin{tcolorbox}[enhanced,nobeforeafter,width=7.1cm,height=2.1cm,valign=center,colback=cprb!6,colframe=cprb,arc=2.5mm,boxrule=0.9pt,left=3mm]{\Large\bfseries\color{cprb}c)}\quad{\Large $\dfrac{1}{2}$}\end{tcolorbox}\hfill\begin{tcolorbox}[enhanced,nobeforeafter,width=7.1cm,height=2.1cm,valign=center,colback=cprb!6,colframe=cprb,arc=2.5mm,boxrule=0.9pt,left=3mm]{\Large\bfseries\color{cprb}d)}\quad{\Large $\dfrac{3}{2}$}\end{tcolorbox}
\end{tcolorbox}
\newpage
\begin{tcolorbox}[enhanced,width=16cm,height=14cm,colback=white,colframe=cprb,boxrule=2pt,arc=0mm,
left=6mm,right=6mm,top=0mm,bottom=4mm,sharp corners,
overlay={\fill[cprb] (frame.north west) rectangle ([yshift=-1.25cm]frame.north east);
\node[anchor=west,text=white,font=\sffamily\bfseries\LARGE] at ([xshift=5mm,yshift=-0.62cm]frame.north west) {Question 96};
\node[anchor=east,text=white,font=\sffamily\itshape\large] at ([xshift=-5mm,yshift=-0.62cm]frame.north east) {Probabilités et dénombrement};}]
\vspace*{1.75cm}
\begin{minipage}[t][6.9cm][c]{\linewidth}\LARGE\centering
\begin{minipage}{0.95\linewidth}\raggedright\linespread{1.25}\selectfont $X$ suit la loi binomiale $\mathcal{B}\left(2;\frac{1}{3}\right)$. Que vaut $P(X\geq 1)$~?\end{minipage}
\end{minipage}

\vfill
\begin{tcolorbox}[enhanced,nobeforeafter,width=7.1cm,height=2.1cm,valign=center,colback=cprb!6,colframe=cprb,arc=2.5mm,boxrule=0.9pt,left=3mm]{\Large\bfseries\color{cprb}a)}\quad{\Large $\dfrac{2}{3}$}\end{tcolorbox}\hfill\begin{tcolorbox}[enhanced,nobeforeafter,width=7.1cm,height=2.1cm,valign=center,colback=cprb!6,colframe=cprb,arc=2.5mm,boxrule=0.9pt,left=3mm]{\Large\bfseries\color{cprb}b)}\quad{\Large $\dfrac{5}{9}$}\end{tcolorbox}\par\vspace{2.5mm}
\begin{tcolorbox}[enhanced,nobeforeafter,width=7.1cm,height=2.1cm,valign=center,colback=cprb!6,colframe=cprb,arc=2.5mm,boxrule=0.9pt,left=3mm]{\Large\bfseries\color{cprb}c)}\quad{\Large $\dfrac{4}{9}$}\end{tcolorbox}\hfill\begin{tcolorbox}[enhanced,nobeforeafter,width=7.1cm,height=2.1cm,valign=center,colback=cprb!6,colframe=cprb,arc=2.5mm,boxrule=0.9pt,left=3mm]{\Large\bfseries\color{cprb}d)}\quad{\Large $\dfrac{1}{9}$}\end{tcolorbox}
\end{tcolorbox}
\newpage
\begin{tcolorbox}[enhanced,width=16cm,height=14cm,colback=white,colframe=csui,boxrule=2pt,arc=0mm,
left=6mm,right=6mm,top=0mm,bottom=4mm,sharp corners,
overlay={\fill[csui] (frame.north west) rectangle ([yshift=-1.25cm]frame.north east);
\node[anchor=west,text=white,font=\sffamily\bfseries\LARGE] at ([xshift=5mm,yshift=-0.62cm]frame.north west) {Question 1 B};
\node[anchor=east,text=white,font=\sffamily\itshape\large] at ([xshift=-5mm,yshift=-0.62cm]frame.north east) {Suites};}]
\vspace*{1.75cm}
\begin{minipage}[t][6.9cm][c]{\linewidth}\LARGE\centering
\begin{minipage}{0.95\linewidth}\raggedright\linespread{1.25}\selectfont $(u_n)$ est la suite arithmétique de premier terme $u_0=-4$ et de raison $5$. Que vaut $u_{12}$~?\end{minipage}
\end{minipage}

\vfill
\begin{tcolorbox}[enhanced,nobeforeafter,width=7.1cm,height=2.1cm,valign=center,colback=csui!6,colframe=csui,arc=2.5mm,boxrule=0.9pt,left=3mm]{\Large\bfseries\color{csui}a)}\quad{\Large $61$}\end{tcolorbox}\hfill\begin{tcolorbox}[enhanced,nobeforeafter,width=7.1cm,height=2.1cm,valign=center,colback=csui!6,colframe=csui,arc=2.5mm,boxrule=0.9pt,left=3mm]{\Large\bfseries\color{csui}b)}\quad{\Large $56$}\end{tcolorbox}\par\vspace{2.5mm}
\begin{tcolorbox}[enhanced,nobeforeafter,width=7.1cm,height=2.1cm,valign=center,colback=csui!6,colframe=csui,arc=2.5mm,boxrule=0.9pt,left=3mm]{\Large\bfseries\color{csui}c)}\quad{\Large $51$}\end{tcolorbox}\hfill\begin{tcolorbox}[enhanced,nobeforeafter,width=7.1cm,height=2.1cm,valign=center,colback=csui!6,colframe=csui,arc=2.5mm,boxrule=0.9pt,left=3mm]{\Large\bfseries\color{csui}d)}\quad{\Large $60$}\end{tcolorbox}
\end{tcolorbox}
\newpage
\begin{tcolorbox}[enhanced,width=16cm,height=14cm,colback=white,colframe=csui,boxrule=2pt,arc=0mm,
left=6mm,right=6mm,top=0mm,bottom=4mm,sharp corners,
overlay={\fill[csui] (frame.north west) rectangle ([yshift=-1.25cm]frame.north east);
\node[anchor=west,text=white,font=\sffamily\bfseries\LARGE] at ([xshift=5mm,yshift=-0.62cm]frame.north west) {Question 2 B};
\node[anchor=east,text=white,font=\sffamily\itshape\large] at ([xshift=-5mm,yshift=-0.62cm]frame.north east) {Suites};}]
\vspace*{1.75cm}
\begin{minipage}[t][6.9cm][c]{\linewidth}\LARGE\centering
\begin{minipage}{0.95\linewidth}\raggedright\linespread{1.25}\selectfont $(u_n)$ est la suite arithmétique de premier terme $u_0=7$ et de raison $4$. Que vaut $u_{15}$~?\end{minipage}
\end{minipage}

\vfill
\begin{tcolorbox}[enhanced,nobeforeafter,width=7.1cm,height=2.1cm,valign=center,colback=csui!6,colframe=csui,arc=2.5mm,boxrule=0.9pt,left=3mm]{\Large\bfseries\color{csui}a)}\quad{\Large $67$}\end{tcolorbox}\hfill\begin{tcolorbox}[enhanced,nobeforeafter,width=7.1cm,height=2.1cm,valign=center,colback=csui!6,colframe=csui,arc=2.5mm,boxrule=0.9pt,left=3mm]{\Large\bfseries\color{csui}b)}\quad{\Large $71$}\end{tcolorbox}\par\vspace{2.5mm}
\begin{tcolorbox}[enhanced,nobeforeafter,width=7.1cm,height=2.1cm,valign=center,colback=csui!6,colframe=csui,arc=2.5mm,boxrule=0.9pt,left=3mm]{\Large\bfseries\color{csui}c)}\quad{\Large $60$}\end{tcolorbox}\hfill\begin{tcolorbox}[enhanced,nobeforeafter,width=7.1cm,height=2.1cm,valign=center,colback=csui!6,colframe=csui,arc=2.5mm,boxrule=0.9pt,left=3mm]{\Large\bfseries\color{csui}d)}\quad{\Large $63$}\end{tcolorbox}
\end{tcolorbox}
\newpage
\begin{tcolorbox}[enhanced,width=16cm,height=14cm,colback=white,colframe=csui,boxrule=2pt,arc=0mm,
left=6mm,right=6mm,top=0mm,bottom=4mm,sharp corners,
overlay={\fill[csui] (frame.north west) rectangle ([yshift=-1.25cm]frame.north east);
\node[anchor=west,text=white,font=\sffamily\bfseries\LARGE] at ([xshift=5mm,yshift=-0.62cm]frame.north west) {Question 3 B};
\node[anchor=east,text=white,font=\sffamily\itshape\large] at ([xshift=-5mm,yshift=-0.62cm]frame.north east) {Suites};}]
\vspace*{1.75cm}
\begin{minipage}[t][6.9cm][c]{\linewidth}\LARGE\centering
\begin{minipage}{0.95\linewidth}\raggedright\linespread{1.25}\selectfont Que vaut la somme $1+2+2^2+\dots+2^{7}$~?\end{minipage}
\end{minipage}

\vfill
\begin{tcolorbox}[enhanced,nobeforeafter,width=7.1cm,height=2.1cm,valign=center,colback=csui!6,colframe=csui,arc=2.5mm,boxrule=0.9pt,left=3mm]{\Large\bfseries\color{csui}a)}\quad{\Large $255$}\end{tcolorbox}\hfill\begin{tcolorbox}[enhanced,nobeforeafter,width=7.1cm,height=2.1cm,valign=center,colback=csui!6,colframe=csui,arc=2.5mm,boxrule=0.9pt,left=3mm]{\Large\bfseries\color{csui}b)}\quad{\Large $256$}\end{tcolorbox}\par\vspace{2.5mm}
\begin{tcolorbox}[enhanced,nobeforeafter,width=7.1cm,height=2.1cm,valign=center,colback=csui!6,colframe=csui,arc=2.5mm,boxrule=0.9pt,left=3mm]{\Large\bfseries\color{csui}c)}\quad{\Large $128$}\end{tcolorbox}\hfill\begin{tcolorbox}[enhanced,nobeforeafter,width=7.1cm,height=2.1cm,valign=center,colback=csui!6,colframe=csui,arc=2.5mm,boxrule=0.9pt,left=3mm]{\Large\bfseries\color{csui}d)}\quad{\Large $127$}\end{tcolorbox}
\end{tcolorbox}
\newpage
\begin{tcolorbox}[enhanced,width=16cm,height=14cm,colback=white,colframe=csui,boxrule=2pt,arc=0mm,
left=6mm,right=6mm,top=0mm,bottom=4mm,sharp corners,
overlay={\fill[csui] (frame.north west) rectangle ([yshift=-1.25cm]frame.north east);
\node[anchor=west,text=white,font=\sffamily\bfseries\LARGE] at ([xshift=5mm,yshift=-0.62cm]frame.north west) {Question 4 B};
\node[anchor=east,text=white,font=\sffamily\itshape\large] at ([xshift=-5mm,yshift=-0.62cm]frame.north east) {Suites};}]
\vspace*{1.75cm}
\begin{minipage}[t][6.9cm][c]{\linewidth}\LARGE\centering
\begin{minipage}{0.95\linewidth}\raggedright\linespread{1.25}\selectfont Que vaut la somme $1+3+3^2+\dots+3^{5}$~?\end{minipage}
\end{minipage}

\vfill
\begin{tcolorbox}[enhanced,nobeforeafter,width=7.1cm,height=2.1cm,valign=center,colback=csui!6,colframe=csui,arc=2.5mm,boxrule=0.9pt,left=3mm]{\Large\bfseries\color{csui}a)}\quad{\Large $121$}\end{tcolorbox}\hfill\begin{tcolorbox}[enhanced,nobeforeafter,width=7.1cm,height=2.1cm,valign=center,colback=csui!6,colframe=csui,arc=2.5mm,boxrule=0.9pt,left=3mm]{\Large\bfseries\color{csui}b)}\quad{\Large $729$}\end{tcolorbox}\par\vspace{2.5mm}
\begin{tcolorbox}[enhanced,nobeforeafter,width=7.1cm,height=2.1cm,valign=center,colback=csui!6,colframe=csui,arc=2.5mm,boxrule=0.9pt,left=3mm]{\Large\bfseries\color{csui}c)}\quad{\Large $728$}\end{tcolorbox}\hfill\begin{tcolorbox}[enhanced,nobeforeafter,width=7.1cm,height=2.1cm,valign=center,colback=csui!6,colframe=csui,arc=2.5mm,boxrule=0.9pt,left=3mm]{\Large\bfseries\color{csui}d)}\quad{\Large $364$}\end{tcolorbox}
\end{tcolorbox}
\newpage
\begin{tcolorbox}[enhanced,width=16cm,height=14cm,colback=white,colframe=csui,boxrule=2pt,arc=0mm,
left=6mm,right=6mm,top=0mm,bottom=4mm,sharp corners,
overlay={\fill[csui] (frame.north west) rectangle ([yshift=-1.25cm]frame.north east);
\node[anchor=west,text=white,font=\sffamily\bfseries\LARGE] at ([xshift=5mm,yshift=-0.62cm]frame.north west) {Question 5 B};
\node[anchor=east,text=white,font=\sffamily\itshape\large] at ([xshift=-5mm,yshift=-0.62cm]frame.north east) {Suites};}]
\vspace*{1.75cm}
\begin{minipage}[t][6.9cm][c]{\linewidth}\LARGE\centering
\begin{minipage}{0.95\linewidth}\raggedright\linespread{1.25}\selectfont Que vaut $\displaystyle\lim_{n\to+\infty}\frac{5n^2+3}{2n^2-n}$~?\end{minipage}
\end{minipage}

\vfill
\begin{tcolorbox}[enhanced,nobeforeafter,width=7.1cm,height=2.1cm,valign=center,colback=csui!6,colframe=csui,arc=2.5mm,boxrule=0.9pt,left=3mm]{\Large\bfseries\color{csui}a)}\quad{\Large $0$}\end{tcolorbox}\hfill\begin{tcolorbox}[enhanced,nobeforeafter,width=7.1cm,height=2.1cm,valign=center,colback=csui!6,colframe=csui,arc=2.5mm,boxrule=0.9pt,left=3mm]{\Large\bfseries\color{csui}b)}\quad{\Large $\dfrac{5}{2}$}\end{tcolorbox}\par\vspace{2.5mm}
\begin{tcolorbox}[enhanced,nobeforeafter,width=7.1cm,height=2.1cm,valign=center,colback=csui!6,colframe=csui,arc=2.5mm,boxrule=0.9pt,left=3mm]{\Large\bfseries\color{csui}c)}\quad{\Large $+\infty$}\end{tcolorbox}\hfill\begin{tcolorbox}[enhanced,nobeforeafter,width=7.1cm,height=2.1cm,valign=center,colback=csui!6,colframe=csui,arc=2.5mm,boxrule=0.9pt,left=3mm]{\Large\bfseries\color{csui}d)}\quad{\Large $\dfrac{2}{5}$}\end{tcolorbox}
\end{tcolorbox}
\newpage
\begin{tcolorbox}[enhanced,width=16cm,height=14cm,colback=white,colframe=csui,boxrule=2pt,arc=0mm,
left=6mm,right=6mm,top=0mm,bottom=4mm,sharp corners,
overlay={\fill[csui] (frame.north west) rectangle ([yshift=-1.25cm]frame.north east);
\node[anchor=west,text=white,font=\sffamily\bfseries\LARGE] at ([xshift=5mm,yshift=-0.62cm]frame.north west) {Question 6 B};
\node[anchor=east,text=white,font=\sffamily\itshape\large] at ([xshift=-5mm,yshift=-0.62cm]frame.north east) {Suites};}]
\vspace*{1.75cm}
\begin{minipage}[t][6.9cm][c]{\linewidth}\LARGE\centering
\begin{minipage}{0.95\linewidth}\raggedright\linespread{1.25}\selectfont Que vaut $\displaystyle\lim_{n\to+\infty}\frac{6n^2+1}{4n^2+3n}$~?\end{minipage}
\end{minipage}

\vfill
\begin{tcolorbox}[enhanced,nobeforeafter,width=7.1cm,height=2.1cm,valign=center,colback=csui!6,colframe=csui,arc=2.5mm,boxrule=0.9pt,left=3mm]{\Large\bfseries\color{csui}a)}\quad{\Large $0$}\end{tcolorbox}\hfill\begin{tcolorbox}[enhanced,nobeforeafter,width=7.1cm,height=2.1cm,valign=center,colback=csui!6,colframe=csui,arc=2.5mm,boxrule=0.9pt,left=3mm]{\Large\bfseries\color{csui}b)}\quad{\Large $+\infty$}\end{tcolorbox}\par\vspace{2.5mm}
\begin{tcolorbox}[enhanced,nobeforeafter,width=7.1cm,height=2.1cm,valign=center,colback=csui!6,colframe=csui,arc=2.5mm,boxrule=0.9pt,left=3mm]{\Large\bfseries\color{csui}c)}\quad{\Large $\dfrac{3}{2}$}\end{tcolorbox}\hfill\begin{tcolorbox}[enhanced,nobeforeafter,width=7.1cm,height=2.1cm,valign=center,colback=csui!6,colframe=csui,arc=2.5mm,boxrule=0.9pt,left=3mm]{\Large\bfseries\color{csui}d)}\quad{\Large $\dfrac{2}{3}$}\end{tcolorbox}
\end{tcolorbox}
\newpage
\begin{tcolorbox}[enhanced,width=16cm,height=14cm,colback=white,colframe=csui,boxrule=2pt,arc=0mm,
left=6mm,right=6mm,top=0mm,bottom=4mm,sharp corners,
overlay={\fill[csui] (frame.north west) rectangle ([yshift=-1.25cm]frame.north east);
\node[anchor=west,text=white,font=\sffamily\bfseries\LARGE] at ([xshift=5mm,yshift=-0.62cm]frame.north west) {Question 7 B};
\node[anchor=east,text=white,font=\sffamily\itshape\large] at ([xshift=-5mm,yshift=-0.62cm]frame.north east) {Suites};}]
\vspace*{1.75cm}
\begin{minipage}[t][6.9cm][c]{\linewidth}\LARGE\centering
\begin{minipage}{0.95\linewidth}\raggedright\linespread{1.25}\selectfont La suite définie par $u_0=0$ et $u_{n+1}=\frac{2}{3}\,u_n+2$ converge. Quelle est sa limite~?\end{minipage}
\end{minipage}

\vfill
\begin{tcolorbox}[enhanced,nobeforeafter,width=7.1cm,height=2.1cm,valign=center,colback=csui!6,colframe=csui,arc=2.5mm,boxrule=0.9pt,left=3mm]{\Large\bfseries\color{csui}a)}\quad{\Large $6$}\end{tcolorbox}\hfill\begin{tcolorbox}[enhanced,nobeforeafter,width=7.1cm,height=2.1cm,valign=center,colback=csui!6,colframe=csui,arc=2.5mm,boxrule=0.9pt,left=3mm]{\Large\bfseries\color{csui}b)}\quad{\Large $\dfrac{6}{5}$}\end{tcolorbox}\par\vspace{2.5mm}
\begin{tcolorbox}[enhanced,nobeforeafter,width=7.1cm,height=2.1cm,valign=center,colback=csui!6,colframe=csui,arc=2.5mm,boxrule=0.9pt,left=3mm]{\Large\bfseries\color{csui}c)}\quad{\Large $+\infty$}\end{tcolorbox}\hfill\begin{tcolorbox}[enhanced,nobeforeafter,width=7.1cm,height=2.1cm,valign=center,colback=csui!6,colframe=csui,arc=2.5mm,boxrule=0.9pt,left=3mm]{\Large\bfseries\color{csui}d)}\quad{\Large $2$}\end{tcolorbox}
\end{tcolorbox}
\newpage
\begin{tcolorbox}[enhanced,width=16cm,height=14cm,colback=white,colframe=csui,boxrule=2pt,arc=0mm,
left=6mm,right=6mm,top=0mm,bottom=4mm,sharp corners,
overlay={\fill[csui] (frame.north west) rectangle ([yshift=-1.25cm]frame.north east);
\node[anchor=west,text=white,font=\sffamily\bfseries\LARGE] at ([xshift=5mm,yshift=-0.62cm]frame.north west) {Question 8 B};
\node[anchor=east,text=white,font=\sffamily\itshape\large] at ([xshift=-5mm,yshift=-0.62cm]frame.north east) {Suites};}]
\vspace*{1.75cm}
\begin{minipage}[t][6.9cm][c]{\linewidth}\LARGE\centering
\begin{minipage}{0.95\linewidth}\raggedright\linespread{1.25}\selectfont La suite définie par $u_0=0$ et $u_{n+1}=\frac{3}{4}\,u_n+1$ converge. Quelle est sa limite~?\end{minipage}
\end{minipage}

\vfill
\begin{tcolorbox}[enhanced,nobeforeafter,width=7.1cm,height=2.1cm,valign=center,colback=csui!6,colframe=csui,arc=2.5mm,boxrule=0.9pt,left=3mm]{\Large\bfseries\color{csui}a)}\quad{\Large $4$}\end{tcolorbox}\hfill\begin{tcolorbox}[enhanced,nobeforeafter,width=7.1cm,height=2.1cm,valign=center,colback=csui!6,colframe=csui,arc=2.5mm,boxrule=0.9pt,left=3mm]{\Large\bfseries\color{csui}b)}\quad{\Large $+\infty$}\end{tcolorbox}\par\vspace{2.5mm}
\begin{tcolorbox}[enhanced,nobeforeafter,width=7.1cm,height=2.1cm,valign=center,colback=csui!6,colframe=csui,arc=2.5mm,boxrule=0.9pt,left=3mm]{\Large\bfseries\color{csui}c)}\quad{\Large $1$}\end{tcolorbox}\hfill\begin{tcolorbox}[enhanced,nobeforeafter,width=7.1cm,height=2.1cm,valign=center,colback=csui!6,colframe=csui,arc=2.5mm,boxrule=0.9pt,left=3mm]{\Large\bfseries\color{csui}d)}\quad{\Large $\dfrac{4}{7}$}\end{tcolorbox}
\end{tcolorbox}
\newpage
\begin{tcolorbox}[enhanced,width=16cm,height=14cm,colback=white,colframe=csui,boxrule=2pt,arc=0mm,
left=6mm,right=6mm,top=0mm,bottom=4mm,sharp corners,
overlay={\fill[csui] (frame.north west) rectangle ([yshift=-1.25cm]frame.north east);
\node[anchor=west,text=white,font=\sffamily\bfseries\LARGE] at ([xshift=5mm,yshift=-0.62cm]frame.north west) {Question 9 B};
\node[anchor=east,text=white,font=\sffamily\itshape\large] at ([xshift=-5mm,yshift=-0.62cm]frame.north east) {Suites};}]
\vspace*{1.75cm}
\begin{minipage}[t][6.9cm][c]{\linewidth}\LARGE\centering
\begin{minipage}{0.95\linewidth}\raggedright\linespread{1.25}\selectfont On pose $u_0=2$ et, pour tout entier $n$, $u_{n+1}=u_n+2n+1$. Que vaut $u_{7}$~?\end{minipage}
\end{minipage}

\vfill
\begin{tcolorbox}[enhanced,nobeforeafter,width=7.1cm,height=2.1cm,valign=center,colback=csui!6,colframe=csui,arc=2.5mm,boxrule=0.9pt,left=3mm]{\Large\bfseries\color{csui}a)}\quad{\Large $17$}\end{tcolorbox}\hfill\begin{tcolorbox}[enhanced,nobeforeafter,width=7.1cm,height=2.1cm,valign=center,colback=csui!6,colframe=csui,arc=2.5mm,boxrule=0.9pt,left=3mm]{\Large\bfseries\color{csui}b)}\quad{\Large $66$}\end{tcolorbox}\par\vspace{2.5mm}
\begin{tcolorbox}[enhanced,nobeforeafter,width=7.1cm,height=2.1cm,valign=center,colback=csui!6,colframe=csui,arc=2.5mm,boxrule=0.9pt,left=3mm]{\Large\bfseries\color{csui}c)}\quad{\Large $51$}\end{tcolorbox}\hfill\begin{tcolorbox}[enhanced,nobeforeafter,width=7.1cm,height=2.1cm,valign=center,colback=csui!6,colframe=csui,arc=2.5mm,boxrule=0.9pt,left=3mm]{\Large\bfseries\color{csui}d)}\quad{\Large $50$}\end{tcolorbox}
\end{tcolorbox}
\newpage
\begin{tcolorbox}[enhanced,width=16cm,height=14cm,colback=white,colframe=csui,boxrule=2pt,arc=0mm,
left=6mm,right=6mm,top=0mm,bottom=4mm,sharp corners,
overlay={\fill[csui] (frame.north west) rectangle ([yshift=-1.25cm]frame.north east);
\node[anchor=west,text=white,font=\sffamily\bfseries\LARGE] at ([xshift=5mm,yshift=-0.62cm]frame.north west) {Question 10 B};
\node[anchor=east,text=white,font=\sffamily\itshape\large] at ([xshift=-5mm,yshift=-0.62cm]frame.north east) {Suites};}]
\vspace*{1.75cm}
\begin{minipage}[t][6.9cm][c]{\linewidth}\LARGE\centering
\begin{minipage}{0.95\linewidth}\raggedright\linespread{1.25}\selectfont On pose $u_0=0$ et, pour tout entier $n$, $u_{n+1}=u_n+2n+1$. Que vaut $u_{9}$~?\end{minipage}
\end{minipage}

\vfill
\begin{tcolorbox}[enhanced,nobeforeafter,width=7.1cm,height=2.1cm,valign=center,colback=csui!6,colframe=csui,arc=2.5mm,boxrule=0.9pt,left=3mm]{\Large\bfseries\color{csui}a)}\quad{\Large $19$}\end{tcolorbox}\hfill\begin{tcolorbox}[enhanced,nobeforeafter,width=7.1cm,height=2.1cm,valign=center,colback=csui!6,colframe=csui,arc=2.5mm,boxrule=0.9pt,left=3mm]{\Large\bfseries\color{csui}b)}\quad{\Large $81$}\end{tcolorbox}\par\vspace{2.5mm}
\begin{tcolorbox}[enhanced,nobeforeafter,width=7.1cm,height=2.1cm,valign=center,colback=csui!6,colframe=csui,arc=2.5mm,boxrule=0.9pt,left=3mm]{\Large\bfseries\color{csui}c)}\quad{\Large $100$}\end{tcolorbox}\hfill\begin{tcolorbox}[enhanced,nobeforeafter,width=7.1cm,height=2.1cm,valign=center,colback=csui!6,colframe=csui,arc=2.5mm,boxrule=0.9pt,left=3mm]{\Large\bfseries\color{csui}d)}\quad{\Large $80$}\end{tcolorbox}
\end{tcolorbox}
\newpage
\begin{tcolorbox}[enhanced,width=16cm,height=14cm,colback=white,colframe=csui,boxrule=2pt,arc=0mm,
left=6mm,right=6mm,top=0mm,bottom=4mm,sharp corners,
overlay={\fill[csui] (frame.north west) rectangle ([yshift=-1.25cm]frame.north east);
\node[anchor=west,text=white,font=\sffamily\bfseries\LARGE] at ([xshift=5mm,yshift=-0.62cm]frame.north west) {Question 11 B};
\node[anchor=east,text=white,font=\sffamily\itshape\large] at ([xshift=-5mm,yshift=-0.62cm]frame.north east) {Suites};}]
\vspace*{1.75cm}
\begin{minipage}[t][6.9cm][c]{\linewidth}\LARGE\centering
\begin{minipage}{0.95\linewidth}\raggedright\linespread{1.25}\selectfont Que vaut $\displaystyle\lim_{n\to+\infty}\frac{3n+(-1)^n}{7n+1}$~?\end{minipage}
\end{minipage}

\vfill
\begin{tcolorbox}[enhanced,nobeforeafter,width=7.1cm,height=2.1cm,valign=center,colback=csui!6,colframe=csui,arc=2.5mm,boxrule=0.9pt,left=3mm]{\Large\bfseries\color{csui}a)}\quad{\Large $+\infty$}\end{tcolorbox}\hfill\begin{tcolorbox}[enhanced,nobeforeafter,width=7.1cm,height=2.1cm,valign=center,colback=csui!6,colframe=csui,arc=2.5mm,boxrule=0.9pt,left=3mm]{\Large\bfseries\color{csui}b)}\quad{\Large $\dfrac{3}{7}$}\end{tcolorbox}\par\vspace{2.5mm}
\begin{tcolorbox}[enhanced,nobeforeafter,width=7.1cm,height=2.1cm,valign=center,colback=csui!6,colframe=csui,arc=2.5mm,boxrule=0.9pt,left=3mm]{\Large\bfseries\color{csui}c)}\quad{\Large $0$}\end{tcolorbox}\hfill\begin{tcolorbox}[enhanced,nobeforeafter,width=7.1cm,height=2.1cm,valign=center,colback=csui!6,colframe=csui,arc=2.5mm,boxrule=0.9pt,left=3mm]{\Large\bfseries\color{csui}d)}\quad{\Large $\dfrac{7}{3}$}\end{tcolorbox}
\end{tcolorbox}
\newpage
\begin{tcolorbox}[enhanced,width=16cm,height=14cm,colback=white,colframe=csui,boxrule=2pt,arc=0mm,
left=6mm,right=6mm,top=0mm,bottom=4mm,sharp corners,
overlay={\fill[csui] (frame.north west) rectangle ([yshift=-1.25cm]frame.north east);
\node[anchor=west,text=white,font=\sffamily\bfseries\LARGE] at ([xshift=5mm,yshift=-0.62cm]frame.north west) {Question 12 B};
\node[anchor=east,text=white,font=\sffamily\itshape\large] at ([xshift=-5mm,yshift=-0.62cm]frame.north east) {Suites};}]
\vspace*{1.75cm}
\begin{minipage}[t][6.9cm][c]{\linewidth}\LARGE\centering
\begin{minipage}{0.95\linewidth}\raggedright\linespread{1.25}\selectfont Que vaut $\displaystyle\lim_{n\to+\infty}\frac{4n+(-1)^n}{5n+1}$~?\end{minipage}
\end{minipage}

\vfill
\begin{tcolorbox}[enhanced,nobeforeafter,width=7.1cm,height=2.1cm,valign=center,colback=csui!6,colframe=csui,arc=2.5mm,boxrule=0.9pt,left=3mm]{\Large\bfseries\color{csui}a)}\quad{\Large $\dfrac{4}{5}$}\end{tcolorbox}\hfill\begin{tcolorbox}[enhanced,nobeforeafter,width=7.1cm,height=2.1cm,valign=center,colback=csui!6,colframe=csui,arc=2.5mm,boxrule=0.9pt,left=3mm]{\Large\bfseries\color{csui}b)}\quad{\Large $\dfrac{5}{4}$}\end{tcolorbox}\par\vspace{2.5mm}
\begin{tcolorbox}[enhanced,nobeforeafter,width=7.1cm,height=2.1cm,valign=center,colback=csui!6,colframe=csui,arc=2.5mm,boxrule=0.9pt,left=3mm]{\Large\bfseries\color{csui}c)}\quad{\Large $+\infty$}\end{tcolorbox}\hfill\begin{tcolorbox}[enhanced,nobeforeafter,width=7.1cm,height=2.1cm,valign=center,colback=csui!6,colframe=csui,arc=2.5mm,boxrule=0.9pt,left=3mm]{\Large\bfseries\color{csui}d)}\quad{\Large $0$}\end{tcolorbox}
\end{tcolorbox}
\newpage
\begin{tcolorbox}[enhanced,width=16cm,height=14cm,colback=white,colframe=csui,boxrule=2pt,arc=0mm,
left=6mm,right=6mm,top=0mm,bottom=4mm,sharp corners,
overlay={\fill[csui] (frame.north west) rectangle ([yshift=-1.25cm]frame.north east);
\node[anchor=west,text=white,font=\sffamily\bfseries\LARGE] at ([xshift=5mm,yshift=-0.62cm]frame.north west) {Question 13 B};
\node[anchor=east,text=white,font=\sffamily\itshape\large] at ([xshift=-5mm,yshift=-0.62cm]frame.north east) {Suites};}]
\vspace*{1.75cm}
\begin{minipage}[t][6.9cm][c]{\linewidth}\LARGE\centering
\begin{minipage}{0.95\linewidth}\raggedright\linespread{1.25}\selectfont $(u_n)$ est une suite géométrique de raison $-2$ telle que $u_3=3$. Que vaut $u_5$~?\end{minipage}
\end{minipage}

\vfill
\begin{tcolorbox}[enhanced,nobeforeafter,width=7.1cm,height=2.1cm,valign=center,colback=csui!6,colframe=csui,arc=2.5mm,boxrule=0.9pt,left=3mm]{\Large\bfseries\color{csui}a)}\quad{\Large $12$}\end{tcolorbox}\hfill\begin{tcolorbox}[enhanced,nobeforeafter,width=7.1cm,height=2.1cm,valign=center,colback=csui!6,colframe=csui,arc=2.5mm,boxrule=0.9pt,left=3mm]{\Large\bfseries\color{csui}b)}\quad{\Large $-6$}\end{tcolorbox}\par\vspace{2.5mm}
\begin{tcolorbox}[enhanced,nobeforeafter,width=7.1cm,height=2.1cm,valign=center,colback=csui!6,colframe=csui,arc=2.5mm,boxrule=0.9pt,left=3mm]{\Large\bfseries\color{csui}c)}\quad{\Large $-24$}\end{tcolorbox}\hfill\begin{tcolorbox}[enhanced,nobeforeafter,width=7.1cm,height=2.1cm,valign=center,colback=csui!6,colframe=csui,arc=2.5mm,boxrule=0.9pt,left=3mm]{\Large\bfseries\color{csui}d)}\quad{\Large $-1$}\end{tcolorbox}
\end{tcolorbox}
\newpage
\begin{tcolorbox}[enhanced,width=16cm,height=14cm,colback=white,colframe=csui,boxrule=2pt,arc=0mm,
left=6mm,right=6mm,top=0mm,bottom=4mm,sharp corners,
overlay={\fill[csui] (frame.north west) rectangle ([yshift=-1.25cm]frame.north east);
\node[anchor=west,text=white,font=\sffamily\bfseries\LARGE] at ([xshift=5mm,yshift=-0.62cm]frame.north west) {Question 14 B};
\node[anchor=east,text=white,font=\sffamily\itshape\large] at ([xshift=-5mm,yshift=-0.62cm]frame.north east) {Suites};}]
\vspace*{1.75cm}
\begin{minipage}[t][6.9cm][c]{\linewidth}\LARGE\centering
\begin{minipage}{0.95\linewidth}\raggedright\linespread{1.25}\selectfont $(u_n)$ est une suite géométrique de raison $\frac{1}{2}$ telle que $u_3=8$. Que vaut $u_5$~?\end{minipage}
\end{minipage}

\vfill
\begin{tcolorbox}[enhanced,nobeforeafter,width=7.1cm,height=2.1cm,valign=center,colback=csui!6,colframe=csui,arc=2.5mm,boxrule=0.9pt,left=3mm]{\Large\bfseries\color{csui}a)}\quad{\Large $2$}\end{tcolorbox}\hfill\begin{tcolorbox}[enhanced,nobeforeafter,width=7.1cm,height=2.1cm,valign=center,colback=csui!6,colframe=csui,arc=2.5mm,boxrule=0.9pt,left=3mm]{\Large\bfseries\color{csui}b)}\quad{\Large $4$}\end{tcolorbox}\par\vspace{2.5mm}
\begin{tcolorbox}[enhanced,nobeforeafter,width=7.1cm,height=2.1cm,valign=center,colback=csui!6,colframe=csui,arc=2.5mm,boxrule=0.9pt,left=3mm]{\Large\bfseries\color{csui}c)}\quad{\Large $1$}\end{tcolorbox}\hfill\begin{tcolorbox}[enhanced,nobeforeafter,width=7.1cm,height=2.1cm,valign=center,colback=csui!6,colframe=csui,arc=2.5mm,boxrule=0.9pt,left=3mm]{\Large\bfseries\color{csui}d)}\quad{\Large $9$}\end{tcolorbox}
\end{tcolorbox}
\newpage
\begin{tcolorbox}[enhanced,width=16cm,height=14cm,colback=white,colframe=csui,boxrule=2pt,arc=0mm,
left=6mm,right=6mm,top=0mm,bottom=4mm,sharp corners,
overlay={\fill[csui] (frame.north west) rectangle ([yshift=-1.25cm]frame.north east);
\node[anchor=west,text=white,font=\sffamily\bfseries\LARGE] at ([xshift=5mm,yshift=-0.62cm]frame.north west) {Question 15 B};
\node[anchor=east,text=white,font=\sffamily\itshape\large] at ([xshift=-5mm,yshift=-0.62cm]frame.north east) {Suites};}]
\vspace*{1.75cm}
\begin{minipage}[t][6.9cm][c]{\linewidth}\LARGE\centering
\begin{minipage}{0.95\linewidth}\raggedright\linespread{1.25}\selectfont Quel est le plus petit entier naturel $n$ tel que $2^n>5000$~?\end{minipage}
\end{minipage}

\vfill
\begin{tcolorbox}[enhanced,nobeforeafter,width=7.1cm,height=2.1cm,valign=center,colback=csui!6,colframe=csui,arc=2.5mm,boxrule=0.9pt,left=3mm]{\Large\bfseries\color{csui}a)}\quad{\Large $14$}\end{tcolorbox}\hfill\begin{tcolorbox}[enhanced,nobeforeafter,width=7.1cm,height=2.1cm,valign=center,colback=csui!6,colframe=csui,arc=2.5mm,boxrule=0.9pt,left=3mm]{\Large\bfseries\color{csui}b)}\quad{\Large $12$}\end{tcolorbox}\par\vspace{2.5mm}
\begin{tcolorbox}[enhanced,nobeforeafter,width=7.1cm,height=2.1cm,valign=center,colback=csui!6,colframe=csui,arc=2.5mm,boxrule=0.9pt,left=3mm]{\Large\bfseries\color{csui}c)}\quad{\Large $13$}\end{tcolorbox}\hfill\begin{tcolorbox}[enhanced,nobeforeafter,width=7.1cm,height=2.1cm,valign=center,colback=csui!6,colframe=csui,arc=2.5mm,boxrule=0.9pt,left=3mm]{\Large\bfseries\color{csui}d)}\quad{\Large $15$}\end{tcolorbox}
\end{tcolorbox}
\newpage
\begin{tcolorbox}[enhanced,width=16cm,height=14cm,colback=white,colframe=csui,boxrule=2pt,arc=0mm,
left=6mm,right=6mm,top=0mm,bottom=4mm,sharp corners,
overlay={\fill[csui] (frame.north west) rectangle ([yshift=-1.25cm]frame.north east);
\node[anchor=west,text=white,font=\sffamily\bfseries\LARGE] at ([xshift=5mm,yshift=-0.62cm]frame.north west) {Question 16 B};
\node[anchor=east,text=white,font=\sffamily\itshape\large] at ([xshift=-5mm,yshift=-0.62cm]frame.north east) {Suites};}]
\vspace*{1.75cm}
\begin{minipage}[t][6.9cm][c]{\linewidth}\LARGE\centering
\begin{minipage}{0.95\linewidth}\raggedright\linespread{1.25}\selectfont Quel est le plus petit entier naturel $n$ tel que $3^n>2000$~?\end{minipage}
\end{minipage}

\vfill
\begin{tcolorbox}[enhanced,nobeforeafter,width=7.1cm,height=2.1cm,valign=center,colback=csui!6,colframe=csui,arc=2.5mm,boxrule=0.9pt,left=3mm]{\Large\bfseries\color{csui}a)}\quad{\Large $7$}\end{tcolorbox}\hfill\begin{tcolorbox}[enhanced,nobeforeafter,width=7.1cm,height=2.1cm,valign=center,colback=csui!6,colframe=csui,arc=2.5mm,boxrule=0.9pt,left=3mm]{\Large\bfseries\color{csui}b)}\quad{\Large $8$}\end{tcolorbox}\par\vspace{2.5mm}
\begin{tcolorbox}[enhanced,nobeforeafter,width=7.1cm,height=2.1cm,valign=center,colback=csui!6,colframe=csui,arc=2.5mm,boxrule=0.9pt,left=3mm]{\Large\bfseries\color{csui}c)}\quad{\Large $9$}\end{tcolorbox}\hfill\begin{tcolorbox}[enhanced,nobeforeafter,width=7.1cm,height=2.1cm,valign=center,colback=csui!6,colframe=csui,arc=2.5mm,boxrule=0.9pt,left=3mm]{\Large\bfseries\color{csui}d)}\quad{\Large $6$}\end{tcolorbox}
\end{tcolorbox}
\newpage
\begin{tcolorbox}[enhanced,width=16cm,height=14cm,colback=white,colframe=cfon,boxrule=2pt,arc=0mm,
left=6mm,right=6mm,top=0mm,bottom=4mm,sharp corners,
overlay={\fill[cfon] (frame.north west) rectangle ([yshift=-1.25cm]frame.north east);
\node[anchor=west,text=white,font=\sffamily\bfseries\LARGE] at ([xshift=5mm,yshift=-0.62cm]frame.north west) {Question 17 B};
\node[anchor=east,text=white,font=\sffamily\itshape\large] at ([xshift=-5mm,yshift=-0.62cm]frame.north east) {Fonctions};}]
\vspace*{1.75cm}
\begin{minipage}[t][6.9cm][c]{\linewidth}\LARGE\centering
\begin{minipage}{0.95\linewidth}\raggedright\linespread{1.25}\selectfont Soit $f(x)=(-x+2)^3$. Que vaut $f'(0)$~?\end{minipage}
\end{minipage}

\vfill
\begin{tcolorbox}[enhanced,nobeforeafter,width=7.1cm,height=2.1cm,valign=center,colback=cfon!6,colframe=cfon,arc=2.5mm,boxrule=0.9pt,left=3mm]{\Large\bfseries\color{cfon}a)}\quad{\Large $-12$}\end{tcolorbox}\hfill\begin{tcolorbox}[enhanced,nobeforeafter,width=7.1cm,height=2.1cm,valign=center,colback=cfon!6,colframe=cfon,arc=2.5mm,boxrule=0.9pt,left=3mm]{\Large\bfseries\color{cfon}b)}\quad{\Large $-8$}\end{tcolorbox}\par\vspace{2.5mm}
\begin{tcolorbox}[enhanced,nobeforeafter,width=7.1cm,height=2.1cm,valign=center,colback=cfon!6,colframe=cfon,arc=2.5mm,boxrule=0.9pt,left=3mm]{\Large\bfseries\color{cfon}c)}\quad{\Large $-6$}\end{tcolorbox}\hfill\begin{tcolorbox}[enhanced,nobeforeafter,width=7.1cm,height=2.1cm,valign=center,colback=cfon!6,colframe=cfon,arc=2.5mm,boxrule=0.9pt,left=3mm]{\Large\bfseries\color{cfon}d)}\quad{\Large $12$}\end{tcolorbox}
\end{tcolorbox}
\newpage
\begin{tcolorbox}[enhanced,width=16cm,height=14cm,colback=white,colframe=cfon,boxrule=2pt,arc=0mm,
left=6mm,right=6mm,top=0mm,bottom=4mm,sharp corners,
overlay={\fill[cfon] (frame.north west) rectangle ([yshift=-1.25cm]frame.north east);
\node[anchor=west,text=white,font=\sffamily\bfseries\LARGE] at ([xshift=5mm,yshift=-0.62cm]frame.north west) {Question 18 B};
\node[anchor=east,text=white,font=\sffamily\itshape\large] at ([xshift=-5mm,yshift=-0.62cm]frame.north east) {Fonctions};}]
\vspace*{1.75cm}
\begin{minipage}[t][6.9cm][c]{\linewidth}\LARGE\centering
\begin{minipage}{0.95\linewidth}\raggedright\linespread{1.25}\selectfont Soit $f(x)=(4x+1)^3$. Que vaut $f'(0)$~?\end{minipage}
\end{minipage}

\vfill
\begin{tcolorbox}[enhanced,nobeforeafter,width=7.1cm,height=2.1cm,valign=center,colback=cfon!6,colframe=cfon,arc=2.5mm,boxrule=0.9pt,left=3mm]{\Large\bfseries\color{cfon}a)}\quad{\Large $3$}\end{tcolorbox}\hfill\begin{tcolorbox}[enhanced,nobeforeafter,width=7.1cm,height=2.1cm,valign=center,colback=cfon!6,colframe=cfon,arc=2.5mm,boxrule=0.9pt,left=3mm]{\Large\bfseries\color{cfon}b)}\quad{\Large $12$}\end{tcolorbox}\par\vspace{2.5mm}
\begin{tcolorbox}[enhanced,nobeforeafter,width=7.1cm,height=2.1cm,valign=center,colback=cfon!6,colframe=cfon,arc=2.5mm,boxrule=0.9pt,left=3mm]{\Large\bfseries\color{cfon}c)}\quad{\Large $4$}\end{tcolorbox}\hfill\begin{tcolorbox}[enhanced,nobeforeafter,width=7.1cm,height=2.1cm,valign=center,colback=cfon!6,colframe=cfon,arc=2.5mm,boxrule=0.9pt,left=3mm]{\Large\bfseries\color{cfon}d)}\quad{\Large $1$}\end{tcolorbox}
\end{tcolorbox}
\newpage
\begin{tcolorbox}[enhanced,width=16cm,height=14cm,colback=white,colframe=cfon,boxrule=2pt,arc=0mm,
left=6mm,right=6mm,top=0mm,bottom=4mm,sharp corners,
overlay={\fill[cfon] (frame.north west) rectangle ([yshift=-1.25cm]frame.north east);
\node[anchor=west,text=white,font=\sffamily\bfseries\LARGE] at ([xshift=5mm,yshift=-0.62cm]frame.north west) {Question 19 B};
\node[anchor=east,text=white,font=\sffamily\itshape\large] at ([xshift=-5mm,yshift=-0.62cm]frame.north east) {Fonctions};}]
\vspace*{1.75cm}
\begin{minipage}[t][6.9cm][c]{\linewidth}\LARGE\centering
\begin{minipage}{0.95\linewidth}\raggedright\linespread{1.25}\selectfont Soit $f(x)=\dfrac{3x-1}{2x+1}$. Que vaut $f'(0)$~?\end{minipage}
\end{minipage}

\vfill
\begin{tcolorbox}[enhanced,nobeforeafter,width=7.1cm,height=2.1cm,valign=center,colback=cfon!6,colframe=cfon,arc=2.5mm,boxrule=0.9pt,left=3mm]{\Large\bfseries\color{cfon}a)}\quad{\Large $\dfrac{3}{2}$}\end{tcolorbox}\hfill\begin{tcolorbox}[enhanced,nobeforeafter,width=7.1cm,height=2.1cm,valign=center,colback=cfon!6,colframe=cfon,arc=2.5mm,boxrule=0.9pt,left=3mm]{\Large\bfseries\color{cfon}b)}\quad{\Large $5$}\end{tcolorbox}\par\vspace{2.5mm}
\begin{tcolorbox}[enhanced,nobeforeafter,width=7.1cm,height=2.1cm,valign=center,colback=cfon!6,colframe=cfon,arc=2.5mm,boxrule=0.9pt,left=3mm]{\Large\bfseries\color{cfon}c)}\quad{\Large $1$}\end{tcolorbox}\hfill\begin{tcolorbox}[enhanced,nobeforeafter,width=7.1cm,height=2.1cm,valign=center,colback=cfon!6,colframe=cfon,arc=2.5mm,boxrule=0.9pt,left=3mm]{\Large\bfseries\color{cfon}d)}\quad{\Large $-5$}\end{tcolorbox}
\end{tcolorbox}
\newpage
\begin{tcolorbox}[enhanced,width=16cm,height=14cm,colback=white,colframe=cfon,boxrule=2pt,arc=0mm,
left=6mm,right=6mm,top=0mm,bottom=4mm,sharp corners,
overlay={\fill[cfon] (frame.north west) rectangle ([yshift=-1.25cm]frame.north east);
\node[anchor=west,text=white,font=\sffamily\bfseries\LARGE] at ([xshift=5mm,yshift=-0.62cm]frame.north west) {Question 20 B};
\node[anchor=east,text=white,font=\sffamily\itshape\large] at ([xshift=-5mm,yshift=-0.62cm]frame.north east) {Fonctions};}]
\vspace*{1.75cm}
\begin{minipage}[t][6.9cm][c]{\linewidth}\LARGE\centering
\begin{minipage}{0.95\linewidth}\raggedright\linespread{1.25}\selectfont Soit $f(x)=\dfrac{x+4}{2x+3}$. Que vaut $f'(0)$~?\end{minipage}
\end{minipage}

\vfill
\begin{tcolorbox}[enhanced,nobeforeafter,width=7.1cm,height=2.1cm,valign=center,colback=cfon!6,colframe=cfon,arc=2.5mm,boxrule=0.9pt,left=3mm]{\Large\bfseries\color{cfon}a)}\quad{\Large $\dfrac{1}{2}$}\end{tcolorbox}\hfill\begin{tcolorbox}[enhanced,nobeforeafter,width=7.1cm,height=2.1cm,valign=center,colback=cfon!6,colframe=cfon,arc=2.5mm,boxrule=0.9pt,left=3mm]{\Large\bfseries\color{cfon}b)}\quad{\Large $\dfrac{11}{9}$}\end{tcolorbox}\par\vspace{2.5mm}
\begin{tcolorbox}[enhanced,nobeforeafter,width=7.1cm,height=2.1cm,valign=center,colback=cfon!6,colframe=cfon,arc=2.5mm,boxrule=0.9pt,left=3mm]{\Large\bfseries\color{cfon}c)}\quad{\Large $-\dfrac{5}{9}$}\end{tcolorbox}\hfill\begin{tcolorbox}[enhanced,nobeforeafter,width=7.1cm,height=2.1cm,valign=center,colback=cfon!6,colframe=cfon,arc=2.5mm,boxrule=0.9pt,left=3mm]{\Large\bfseries\color{cfon}d)}\quad{\Large $\dfrac{5}{9}$}\end{tcolorbox}
\end{tcolorbox}
\newpage
\begin{tcolorbox}[enhanced,width=16cm,height=14cm,colback=white,colframe=cfon,boxrule=2pt,arc=0mm,
left=6mm,right=6mm,top=0mm,bottom=4mm,sharp corners,
overlay={\fill[cfon] (frame.north west) rectangle ([yshift=-1.25cm]frame.north east);
\node[anchor=west,text=white,font=\sffamily\bfseries\LARGE] at ([xshift=5mm,yshift=-0.62cm]frame.north west) {Question 21 B};
\node[anchor=east,text=white,font=\sffamily\itshape\large] at ([xshift=-5mm,yshift=-0.62cm]frame.north east) {Fonctions};}]
\vspace*{1.75cm}
\begin{minipage}[t][6.9cm][c]{\linewidth}\LARGE\centering
\begin{minipage}{0.95\linewidth}\raggedright\linespread{1.25}\selectfont Quelle est l'équation de la tangente à la courbe de $f(x)=x^2+3x-1$ au point d'abscisse $-1$~?\end{minipage}
\end{minipage}

\vfill
\begin{tcolorbox}[enhanced,nobeforeafter,width=7.1cm,height=2.1cm,valign=center,colback=cfon!6,colframe=cfon,arc=2.5mm,boxrule=0.9pt,left=3mm]{\Large\bfseries\color{cfon}a)}\quad{\Large $y=x-1$}\end{tcolorbox}\hfill\begin{tcolorbox}[enhanced,nobeforeafter,width=7.1cm,height=2.1cm,valign=center,colback=cfon!6,colframe=cfon,arc=2.5mm,boxrule=0.9pt,left=3mm]{\Large\bfseries\color{cfon}b)}\quad{\Large $y=x-2$}\end{tcolorbox}\par\vspace{2.5mm}
\begin{tcolorbox}[enhanced,nobeforeafter,width=7.1cm,height=2.1cm,valign=center,colback=cfon!6,colframe=cfon,arc=2.5mm,boxrule=0.9pt,left=3mm]{\Large\bfseries\color{cfon}c)}\quad{\Large $y=2x-2$}\end{tcolorbox}\hfill\begin{tcolorbox}[enhanced,nobeforeafter,width=7.1cm,height=2.1cm,valign=center,colback=cfon!6,colframe=cfon,arc=2.5mm,boxrule=0.9pt,left=3mm]{\Large\bfseries\color{cfon}d)}\quad{\Large $y=x+2$}\end{tcolorbox}
\end{tcolorbox}
\newpage
\begin{tcolorbox}[enhanced,width=16cm,height=14cm,colback=white,colframe=cfon,boxrule=2pt,arc=0mm,
left=6mm,right=6mm,top=0mm,bottom=4mm,sharp corners,
overlay={\fill[cfon] (frame.north west) rectangle ([yshift=-1.25cm]frame.north east);
\node[anchor=west,text=white,font=\sffamily\bfseries\LARGE] at ([xshift=5mm,yshift=-0.62cm]frame.north west) {Question 22 B};
\node[anchor=east,text=white,font=\sffamily\itshape\large] at ([xshift=-5mm,yshift=-0.62cm]frame.north east) {Fonctions};}]
\vspace*{1.75cm}
\begin{minipage}[t][6.9cm][c]{\linewidth}\LARGE\centering
\begin{minipage}{0.95\linewidth}\raggedright\linespread{1.25}\selectfont Quelle est l'équation de la tangente à la courbe de $f(x)=x^2+5$ au point d'abscisse $3$~?\end{minipage}
\end{minipage}

\vfill
\begin{tcolorbox}[enhanced,nobeforeafter,width=7.1cm,height=2.1cm,valign=center,colback=cfon!6,colframe=cfon,arc=2.5mm,boxrule=0.9pt,left=3mm]{\Large\bfseries\color{cfon}a)}\quad{\Large $y=6x-4$}\end{tcolorbox}\hfill\begin{tcolorbox}[enhanced,nobeforeafter,width=7.1cm,height=2.1cm,valign=center,colback=cfon!6,colframe=cfon,arc=2.5mm,boxrule=0.9pt,left=3mm]{\Large\bfseries\color{cfon}b)}\quad{\Large $y=6x+4$}\end{tcolorbox}\par\vspace{2.5mm}
\begin{tcolorbox}[enhanced,nobeforeafter,width=7.1cm,height=2.1cm,valign=center,colback=cfon!6,colframe=cfon,arc=2.5mm,boxrule=0.9pt,left=3mm]{\Large\bfseries\color{cfon}c)}\quad{\Large $y=6x+5$}\end{tcolorbox}\hfill\begin{tcolorbox}[enhanced,nobeforeafter,width=7.1cm,height=2.1cm,valign=center,colback=cfon!6,colframe=cfon,arc=2.5mm,boxrule=0.9pt,left=3mm]{\Large\bfseries\color{cfon}d)}\quad{\Large $y=3x-4$}\end{tcolorbox}
\end{tcolorbox}
\newpage
\begin{tcolorbox}[enhanced,width=16cm,height=14cm,colback=white,colframe=cfon,boxrule=2pt,arc=0mm,
left=6mm,right=6mm,top=0mm,bottom=4mm,sharp corners,
overlay={\fill[cfon] (frame.north west) rectangle ([yshift=-1.25cm]frame.north east);
\node[anchor=west,text=white,font=\sffamily\bfseries\LARGE] at ([xshift=5mm,yshift=-0.62cm]frame.north west) {Question 23 B};
\node[anchor=east,text=white,font=\sffamily\itshape\large] at ([xshift=-5mm,yshift=-0.62cm]frame.north east) {Fonctions};}]
\vspace*{1.75cm}
\begin{minipage}[t][6.9cm][c]{\linewidth}\LARGE\centering
\begin{minipage}{0.95\linewidth}\raggedright\linespread{1.25}\selectfont La courbe de $f(x)=\dfrac{x+4}{x-5}$ admet une asymptote verticale. Laquelle~?\end{minipage}
\end{minipage}

\vfill
\begin{tcolorbox}[enhanced,nobeforeafter,width=7.1cm,height=2.1cm,valign=center,colback=cfon!6,colframe=cfon,arc=2.5mm,boxrule=0.9pt,left=3mm]{\Large\bfseries\color{cfon}a)}\quad{\Large $y=1$}\end{tcolorbox}\hfill\begin{tcolorbox}[enhanced,nobeforeafter,width=7.1cm,height=2.1cm,valign=center,colback=cfon!6,colframe=cfon,arc=2.5mm,boxrule=0.9pt,left=3mm]{\Large\bfseries\color{cfon}b)}\quad{\Large $x=-5$}\end{tcolorbox}\par\vspace{2.5mm}
\begin{tcolorbox}[enhanced,nobeforeafter,width=7.1cm,height=2.1cm,valign=center,colback=cfon!6,colframe=cfon,arc=2.5mm,boxrule=0.9pt,left=3mm]{\Large\bfseries\color{cfon}c)}\quad{\Large $x=5$}\end{tcolorbox}\hfill\begin{tcolorbox}[enhanced,nobeforeafter,width=7.1cm,height=2.1cm,valign=center,colback=cfon!6,colframe=cfon,arc=2.5mm,boxrule=0.9pt,left=3mm]{\Large\bfseries\color{cfon}d)}\quad{\Large $x=-4$}\end{tcolorbox}
\end{tcolorbox}
\newpage
\begin{tcolorbox}[enhanced,width=16cm,height=14cm,colback=white,colframe=cfon,boxrule=2pt,arc=0mm,
left=6mm,right=6mm,top=0mm,bottom=4mm,sharp corners,
overlay={\fill[cfon] (frame.north west) rectangle ([yshift=-1.25cm]frame.north east);
\node[anchor=west,text=white,font=\sffamily\bfseries\LARGE] at ([xshift=5mm,yshift=-0.62cm]frame.north west) {Question 24 B};
\node[anchor=east,text=white,font=\sffamily\itshape\large] at ([xshift=-5mm,yshift=-0.62cm]frame.north east) {Fonctions};}]
\vspace*{1.75cm}
\begin{minipage}[t][6.9cm][c]{\linewidth}\LARGE\centering
\begin{minipage}{0.95\linewidth}\raggedright\linespread{1.25}\selectfont La courbe de $f(x)=\dfrac{-2x+3}{x-1}$ admet une asymptote verticale. Laquelle~?\end{minipage}
\end{minipage}

\vfill
\begin{tcolorbox}[enhanced,nobeforeafter,width=7.1cm,height=2.1cm,valign=center,colback=cfon!6,colframe=cfon,arc=2.5mm,boxrule=0.9pt,left=3mm]{\Large\bfseries\color{cfon}a)}\quad{\Large $y=-2$}\end{tcolorbox}\hfill\begin{tcolorbox}[enhanced,nobeforeafter,width=7.1cm,height=2.1cm,valign=center,colback=cfon!6,colframe=cfon,arc=2.5mm,boxrule=0.9pt,left=3mm]{\Large\bfseries\color{cfon}b)}\quad{\Large $x=1$}\end{tcolorbox}\par\vspace{2.5mm}
\begin{tcolorbox}[enhanced,nobeforeafter,width=7.1cm,height=2.1cm,valign=center,colback=cfon!6,colframe=cfon,arc=2.5mm,boxrule=0.9pt,left=3mm]{\Large\bfseries\color{cfon}c)}\quad{\Large $x=\frac{3}{2}$}\end{tcolorbox}\hfill\begin{tcolorbox}[enhanced,nobeforeafter,width=7.1cm,height=2.1cm,valign=center,colback=cfon!6,colframe=cfon,arc=2.5mm,boxrule=0.9pt,left=3mm]{\Large\bfseries\color{cfon}d)}\quad{\Large $x=-1$}\end{tcolorbox}
\end{tcolorbox}
\newpage
\begin{tcolorbox}[enhanced,width=16cm,height=14cm,colback=white,colframe=cfon,boxrule=2pt,arc=0mm,
left=6mm,right=6mm,top=0mm,bottom=4mm,sharp corners,
overlay={\fill[cfon] (frame.north west) rectangle ([yshift=-1.25cm]frame.north east);
\node[anchor=west,text=white,font=\sffamily\bfseries\LARGE] at ([xshift=5mm,yshift=-0.62cm]frame.north west) {Question 25 B};
\node[anchor=east,text=white,font=\sffamily\itshape\large] at ([xshift=-5mm,yshift=-0.62cm]frame.north east) {Fonctions};}]
\vspace*{1.75cm}
\begin{minipage}[t][6.9cm][c]{\linewidth}\LARGE\centering
\begin{minipage}{0.95\linewidth}\raggedright\linespread{1.25}\selectfont La courbe de $f(x)=x^3+2x^2-x$ admet un point d'inflexion. Quelle est son abscisse~?\end{minipage}
\end{minipage}

\vfill
\begin{tcolorbox}[enhanced,nobeforeafter,width=7.1cm,height=2.1cm,valign=center,colback=cfon!6,colframe=cfon,arc=2.5mm,boxrule=0.9pt,left=3mm]{\Large\bfseries\color{cfon}a)}\quad{\Large $-1$}\end{tcolorbox}\hfill\begin{tcolorbox}[enhanced,nobeforeafter,width=7.1cm,height=2.1cm,valign=center,colback=cfon!6,colframe=cfon,arc=2.5mm,boxrule=0.9pt,left=3mm]{\Large\bfseries\color{cfon}b)}\quad{\Large $-\dfrac{1}{3}$}\end{tcolorbox}\par\vspace{2.5mm}
\begin{tcolorbox}[enhanced,nobeforeafter,width=7.1cm,height=2.1cm,valign=center,colback=cfon!6,colframe=cfon,arc=2.5mm,boxrule=0.9pt,left=3mm]{\Large\bfseries\color{cfon}c)}\quad{\Large $\dfrac{2}{3}$}\end{tcolorbox}\hfill\begin{tcolorbox}[enhanced,nobeforeafter,width=7.1cm,height=2.1cm,valign=center,colback=cfon!6,colframe=cfon,arc=2.5mm,boxrule=0.9pt,left=3mm]{\Large\bfseries\color{cfon}d)}\quad{\Large $-\dfrac{2}{3}$}\end{tcolorbox}
\end{tcolorbox}
\newpage
\begin{tcolorbox}[enhanced,width=16cm,height=14cm,colback=white,colframe=cfon,boxrule=2pt,arc=0mm,
left=6mm,right=6mm,top=0mm,bottom=4mm,sharp corners,
overlay={\fill[cfon] (frame.north west) rectangle ([yshift=-1.25cm]frame.north east);
\node[anchor=west,text=white,font=\sffamily\bfseries\LARGE] at ([xshift=5mm,yshift=-0.62cm]frame.north west) {Question 26 B};
\node[anchor=east,text=white,font=\sffamily\itshape\large] at ([xshift=-5mm,yshift=-0.62cm]frame.north east) {Fonctions};}]
\vspace*{1.75cm}
\begin{minipage}[t][6.9cm][c]{\linewidth}\LARGE\centering
\begin{minipage}{0.95\linewidth}\raggedright\linespread{1.25}\selectfont La courbe de $f(x)=x^3-3x^2+5x$ admet un point d'inflexion. Quelle est son abscisse~?\end{minipage}
\end{minipage}

\vfill
\begin{tcolorbox}[enhanced,nobeforeafter,width=7.1cm,height=2.1cm,valign=center,colback=cfon!6,colframe=cfon,arc=2.5mm,boxrule=0.9pt,left=3mm]{\Large\bfseries\color{cfon}a)}\quad{\Large $-1$}\end{tcolorbox}\hfill\begin{tcolorbox}[enhanced,nobeforeafter,width=7.1cm,height=2.1cm,valign=center,colback=cfon!6,colframe=cfon,arc=2.5mm,boxrule=0.9pt,left=3mm]{\Large\bfseries\color{cfon}b)}\quad{\Large $1$}\end{tcolorbox}\par\vspace{2.5mm}
\begin{tcolorbox}[enhanced,nobeforeafter,width=7.1cm,height=2.1cm,valign=center,colback=cfon!6,colframe=cfon,arc=2.5mm,boxrule=0.9pt,left=3mm]{\Large\bfseries\color{cfon}c)}\quad{\Large $\dfrac{3}{2}$}\end{tcolorbox}\hfill\begin{tcolorbox}[enhanced,nobeforeafter,width=7.1cm,height=2.1cm,valign=center,colback=cfon!6,colframe=cfon,arc=2.5mm,boxrule=0.9pt,left=3mm]{\Large\bfseries\color{cfon}d)}\quad{\Large $\dfrac{1}{2}$}\end{tcolorbox}
\end{tcolorbox}
\newpage
\begin{tcolorbox}[enhanced,width=16cm,height=14cm,colback=white,colframe=cfon,boxrule=2pt,arc=0mm,
left=6mm,right=6mm,top=0mm,bottom=4mm,sharp corners,
overlay={\fill[cfon] (frame.north west) rectangle ([yshift=-1.25cm]frame.north east);
\node[anchor=west,text=white,font=\sffamily\bfseries\LARGE] at ([xshift=5mm,yshift=-0.62cm]frame.north west) {Question 27 B};
\node[anchor=east,text=white,font=\sffamily\itshape\large] at ([xshift=-5mm,yshift=-0.62cm]frame.north east) {Fonctions};}]
\vspace*{1.75cm}
\begin{minipage}[t][6.9cm][c]{\linewidth}\LARGE\centering
\begin{minipage}{0.95\linewidth}\raggedright\linespread{1.25}\selectfont Combien de solutions réelles l'équation $x^3-3x=-1$ admet-elle~?\end{minipage}
\end{minipage}

\vfill
\begin{tcolorbox}[enhanced,nobeforeafter,width=7.1cm,height=2.1cm,valign=center,colback=cfon!6,colframe=cfon,arc=2.5mm,boxrule=0.9pt,left=3mm]{\Large\bfseries\color{cfon}a)}\quad{\Large $0$}\end{tcolorbox}\hfill\begin{tcolorbox}[enhanced,nobeforeafter,width=7.1cm,height=2.1cm,valign=center,colback=cfon!6,colframe=cfon,arc=2.5mm,boxrule=0.9pt,left=3mm]{\Large\bfseries\color{cfon}b)}\quad{\Large $2$}\end{tcolorbox}\par\vspace{2.5mm}
\begin{tcolorbox}[enhanced,nobeforeafter,width=7.1cm,height=2.1cm,valign=center,colback=cfon!6,colframe=cfon,arc=2.5mm,boxrule=0.9pt,left=3mm]{\Large\bfseries\color{cfon}c)}\quad{\Large $1$}\end{tcolorbox}\hfill\begin{tcolorbox}[enhanced,nobeforeafter,width=7.1cm,height=2.1cm,valign=center,colback=cfon!6,colframe=cfon,arc=2.5mm,boxrule=0.9pt,left=3mm]{\Large\bfseries\color{cfon}d)}\quad{\Large $3$}\end{tcolorbox}
\end{tcolorbox}
\newpage
\begin{tcolorbox}[enhanced,width=16cm,height=14cm,colback=white,colframe=cfon,boxrule=2pt,arc=0mm,
left=6mm,right=6mm,top=0mm,bottom=4mm,sharp corners,
overlay={\fill[cfon] (frame.north west) rectangle ([yshift=-1.25cm]frame.north east);
\node[anchor=west,text=white,font=\sffamily\bfseries\LARGE] at ([xshift=5mm,yshift=-0.62cm]frame.north west) {Question 28 B};
\node[anchor=east,text=white,font=\sffamily\itshape\large] at ([xshift=-5mm,yshift=-0.62cm]frame.north east) {Fonctions};}]
\vspace*{1.75cm}
\begin{minipage}[t][6.9cm][c]{\linewidth}\LARGE\centering
\begin{minipage}{0.95\linewidth}\raggedright\linespread{1.25}\selectfont Combien de solutions réelles l'équation $x^3-3x=-5$ admet-elle~?\end{minipage}
\end{minipage}

\vfill
\begin{tcolorbox}[enhanced,nobeforeafter,width=7.1cm,height=2.1cm,valign=center,colback=cfon!6,colframe=cfon,arc=2.5mm,boxrule=0.9pt,left=3mm]{\Large\bfseries\color{cfon}a)}\quad{\Large $3$}\end{tcolorbox}\hfill\begin{tcolorbox}[enhanced,nobeforeafter,width=7.1cm,height=2.1cm,valign=center,colback=cfon!6,colframe=cfon,arc=2.5mm,boxrule=0.9pt,left=3mm]{\Large\bfseries\color{cfon}b)}\quad{\Large $1$}\end{tcolorbox}\par\vspace{2.5mm}
\begin{tcolorbox}[enhanced,nobeforeafter,width=7.1cm,height=2.1cm,valign=center,colback=cfon!6,colframe=cfon,arc=2.5mm,boxrule=0.9pt,left=3mm]{\Large\bfseries\color{cfon}c)}\quad{\Large $0$}\end{tcolorbox}\hfill\begin{tcolorbox}[enhanced,nobeforeafter,width=7.1cm,height=2.1cm,valign=center,colback=cfon!6,colframe=cfon,arc=2.5mm,boxrule=0.9pt,left=3mm]{\Large\bfseries\color{cfon}d)}\quad{\Large $2$}\end{tcolorbox}
\end{tcolorbox}
\newpage
\begin{tcolorbox}[enhanced,width=16cm,height=14cm,colback=white,colframe=cfon,boxrule=2pt,arc=0mm,
left=6mm,right=6mm,top=0mm,bottom=4mm,sharp corners,
overlay={\fill[cfon] (frame.north west) rectangle ([yshift=-1.25cm]frame.north east);
\node[anchor=west,text=white,font=\sffamily\bfseries\LARGE] at ([xshift=5mm,yshift=-0.62cm]frame.north west) {Question 29 B};
\node[anchor=east,text=white,font=\sffamily\itshape\large] at ([xshift=-5mm,yshift=-0.62cm]frame.north east) {Fonctions};}]
\vspace*{1.75cm}
\begin{minipage}[t][6.9cm][c]{\linewidth}\LARGE\centering
\begin{minipage}{0.95\linewidth}\raggedright\linespread{1.25}\selectfont Quel est le minimum de la fonction $f(x)=x+\dfrac{25}{x}$ sur $]0;+\infty[$~?\end{minipage}
\end{minipage}

\vfill
\begin{tcolorbox}[enhanced,nobeforeafter,width=7.1cm,height=2.1cm,valign=center,colback=cfon!6,colframe=cfon,arc=2.5mm,boxrule=0.9pt,left=3mm]{\Large\bfseries\color{cfon}a)}\quad{\Large $5$}\end{tcolorbox}\hfill\begin{tcolorbox}[enhanced,nobeforeafter,width=7.1cm,height=2.1cm,valign=center,colback=cfon!6,colframe=cfon,arc=2.5mm,boxrule=0.9pt,left=3mm]{\Large\bfseries\color{cfon}b)}\quad{\Large $25$}\end{tcolorbox}\par\vspace{2.5mm}
\begin{tcolorbox}[enhanced,nobeforeafter,width=7.1cm,height=2.1cm,valign=center,colback=cfon!6,colframe=cfon,arc=2.5mm,boxrule=0.9pt,left=3mm]{\Large\bfseries\color{cfon}c)}\quad{\Large $10$}\end{tcolorbox}\hfill\begin{tcolorbox}[enhanced,nobeforeafter,width=7.1cm,height=2.1cm,valign=center,colback=cfon!6,colframe=cfon,arc=2.5mm,boxrule=0.9pt,left=3mm]{\Large\bfseries\color{cfon}d)}\quad{\Large $26$}\end{tcolorbox}
\end{tcolorbox}
\newpage
\begin{tcolorbox}[enhanced,width=16cm,height=14cm,colback=white,colframe=cfon,boxrule=2pt,arc=0mm,
left=6mm,right=6mm,top=0mm,bottom=4mm,sharp corners,
overlay={\fill[cfon] (frame.north west) rectangle ([yshift=-1.25cm]frame.north east);
\node[anchor=west,text=white,font=\sffamily\bfseries\LARGE] at ([xshift=5mm,yshift=-0.62cm]frame.north west) {Question 30 B};
\node[anchor=east,text=white,font=\sffamily\itshape\large] at ([xshift=-5mm,yshift=-0.62cm]frame.north east) {Fonctions};}]
\vspace*{1.75cm}
\begin{minipage}[t][6.9cm][c]{\linewidth}\LARGE\centering
\begin{minipage}{0.95\linewidth}\raggedright\linespread{1.25}\selectfont Quel est le minimum de la fonction $f(x)=x+\dfrac{16}{x}$ sur $]0;+\infty[$~?\end{minipage}
\end{minipage}

\vfill
\begin{tcolorbox}[enhanced,nobeforeafter,width=7.1cm,height=2.1cm,valign=center,colback=cfon!6,colframe=cfon,arc=2.5mm,boxrule=0.9pt,left=3mm]{\Large\bfseries\color{cfon}a)}\quad{\Large $8$}\end{tcolorbox}\hfill\begin{tcolorbox}[enhanced,nobeforeafter,width=7.1cm,height=2.1cm,valign=center,colback=cfon!6,colframe=cfon,arc=2.5mm,boxrule=0.9pt,left=3mm]{\Large\bfseries\color{cfon}b)}\quad{\Large $17$}\end{tcolorbox}\par\vspace{2.5mm}
\begin{tcolorbox}[enhanced,nobeforeafter,width=7.1cm,height=2.1cm,valign=center,colback=cfon!6,colframe=cfon,arc=2.5mm,boxrule=0.9pt,left=3mm]{\Large\bfseries\color{cfon}c)}\quad{\Large $16$}\end{tcolorbox}\hfill\begin{tcolorbox}[enhanced,nobeforeafter,width=7.1cm,height=2.1cm,valign=center,colback=cfon!6,colframe=cfon,arc=2.5mm,boxrule=0.9pt,left=3mm]{\Large\bfseries\color{cfon}d)}\quad{\Large $4$}\end{tcolorbox}
\end{tcolorbox}
\newpage
\begin{tcolorbox}[enhanced,width=16cm,height=14cm,colback=white,colframe=cfon,boxrule=2pt,arc=0mm,
left=6mm,right=6mm,top=0mm,bottom=4mm,sharp corners,
overlay={\fill[cfon] (frame.north west) rectangle ([yshift=-1.25cm]frame.north east);
\node[anchor=west,text=white,font=\sffamily\bfseries\LARGE] at ([xshift=5mm,yshift=-0.62cm]frame.north west) {Question 31 B};
\node[anchor=east,text=white,font=\sffamily\itshape\large] at ([xshift=-5mm,yshift=-0.62cm]frame.north east) {Fonctions};}]
\vspace*{1.75cm}
\begin{minipage}[t][6.9cm][c]{\linewidth}\LARGE\centering
\begin{minipage}{0.95\linewidth}\raggedright\linespread{1.25}\selectfont Soit $f(x)=\sqrt{x+3}$. Que vaut $f'(1)$~?\end{minipage}
\end{minipage}

\vfill
\begin{tcolorbox}[enhanced,nobeforeafter,width=7.1cm,height=2.1cm,valign=center,colback=cfon!6,colframe=cfon,arc=2.5mm,boxrule=0.9pt,left=3mm]{\Large\bfseries\color{cfon}a)}\quad{\Large $\dfrac{1}{2}$}\end{tcolorbox}\hfill\begin{tcolorbox}[enhanced,nobeforeafter,width=7.1cm,height=2.1cm,valign=center,colback=cfon!6,colframe=cfon,arc=2.5mm,boxrule=0.9pt,left=3mm]{\Large\bfseries\color{cfon}b)}\quad{\Large $4$}\end{tcolorbox}\par\vspace{2.5mm}
\begin{tcolorbox}[enhanced,nobeforeafter,width=7.1cm,height=2.1cm,valign=center,colback=cfon!6,colframe=cfon,arc=2.5mm,boxrule=0.9pt,left=3mm]{\Large\bfseries\color{cfon}c)}\quad{\Large $\dfrac{1}{8}$}\end{tcolorbox}\hfill\begin{tcolorbox}[enhanced,nobeforeafter,width=7.1cm,height=2.1cm,valign=center,colback=cfon!6,colframe=cfon,arc=2.5mm,boxrule=0.9pt,left=3mm]{\Large\bfseries\color{cfon}d)}\quad{\Large $\dfrac{1}{4}$}\end{tcolorbox}
\end{tcolorbox}
\newpage
\begin{tcolorbox}[enhanced,width=16cm,height=14cm,colback=white,colframe=cfon,boxrule=2pt,arc=0mm,
left=6mm,right=6mm,top=0mm,bottom=4mm,sharp corners,
overlay={\fill[cfon] (frame.north west) rectangle ([yshift=-1.25cm]frame.north east);
\node[anchor=west,text=white,font=\sffamily\bfseries\LARGE] at ([xshift=5mm,yshift=-0.62cm]frame.north west) {Question 32 B};
\node[anchor=east,text=white,font=\sffamily\itshape\large] at ([xshift=-5mm,yshift=-0.62cm]frame.north east) {Fonctions};}]
\vspace*{1.75cm}
\begin{minipage}[t][6.9cm][c]{\linewidth}\LARGE\centering
\begin{minipage}{0.95\linewidth}\raggedright\linespread{1.25}\selectfont Soit $f(x)=\sqrt{4x}$. Que vaut $f'(1)$~?\end{minipage}
\end{minipage}

\vfill
\begin{tcolorbox}[enhanced,nobeforeafter,width=7.1cm,height=2.1cm,valign=center,colback=cfon!6,colframe=cfon,arc=2.5mm,boxrule=0.9pt,left=3mm]{\Large\bfseries\color{cfon}a)}\quad{\Large $2$}\end{tcolorbox}\hfill\begin{tcolorbox}[enhanced,nobeforeafter,width=7.1cm,height=2.1cm,valign=center,colback=cfon!6,colframe=cfon,arc=2.5mm,boxrule=0.9pt,left=3mm]{\Large\bfseries\color{cfon}b)}\quad{\Large $1$}\end{tcolorbox}\par\vspace{2.5mm}
\begin{tcolorbox}[enhanced,nobeforeafter,width=7.1cm,height=2.1cm,valign=center,colback=cfon!6,colframe=cfon,arc=2.5mm,boxrule=0.9pt,left=3mm]{\Large\bfseries\color{cfon}c)}\quad{\Large $\dfrac{1}{4}$}\end{tcolorbox}\hfill\begin{tcolorbox}[enhanced,nobeforeafter,width=7.1cm,height=2.1cm,valign=center,colback=cfon!6,colframe=cfon,arc=2.5mm,boxrule=0.9pt,left=3mm]{\Large\bfseries\color{cfon}d)}\quad{\Large $\dfrac{1}{2}$}\end{tcolorbox}
\end{tcolorbox}
\newpage
\begin{tcolorbox}[enhanced,width=16cm,height=14cm,colback=white,colframe=cexl,boxrule=2pt,arc=0mm,
left=6mm,right=6mm,top=0mm,bottom=4mm,sharp corners,
overlay={\fill[cexl] (frame.north west) rectangle ([yshift=-1.25cm]frame.north east);
\node[anchor=west,text=white,font=\sffamily\bfseries\LARGE] at ([xshift=5mm,yshift=-0.62cm]frame.north west) {Question 33 B};
\node[anchor=east,text=white,font=\sffamily\itshape\large] at ([xshift=-5mm,yshift=-0.62cm]frame.north east) {Exponentielle et logarithme};}]
\vspace*{1.75cm}
\begin{minipage}[t][6.9cm][c]{\linewidth}\LARGE\centering
\begin{minipage}{0.95\linewidth}\raggedright\linespread{1.25}\selectfont Simplifier $\dfrac{e\times e^{4}}{e^{3}}$.\end{minipage}
\end{minipage}

\vfill
\begin{tcolorbox}[enhanced,nobeforeafter,width=7.1cm,height=2.1cm,valign=center,colback=cexl!6,colframe=cexl,arc=2.5mm,boxrule=0.9pt,left=3mm]{\Large\bfseries\color{cexl}a)}\quad{\Large $1$}\end{tcolorbox}\hfill\begin{tcolorbox}[enhanced,nobeforeafter,width=7.1cm,height=2.1cm,valign=center,colback=cexl!6,colframe=cexl,arc=2.5mm,boxrule=0.9pt,left=3mm]{\Large\bfseries\color{cexl}b)}\quad{\Large $e^{8}$}\end{tcolorbox}\par\vspace{2.5mm}
\begin{tcolorbox}[enhanced,nobeforeafter,width=7.1cm,height=2.1cm,valign=center,colback=cexl!6,colframe=cexl,arc=2.5mm,boxrule=0.9pt,left=3mm]{\Large\bfseries\color{cexl}c)}\quad{\Large $e$}\end{tcolorbox}\hfill\begin{tcolorbox}[enhanced,nobeforeafter,width=7.1cm,height=2.1cm,valign=center,colback=cexl!6,colframe=cexl,arc=2.5mm,boxrule=0.9pt,left=3mm]{\Large\bfseries\color{cexl}d)}\quad{\Large $e^{2}$}\end{tcolorbox}
\end{tcolorbox}
\newpage
\begin{tcolorbox}[enhanced,width=16cm,height=14cm,colback=white,colframe=cexl,boxrule=2pt,arc=0mm,
left=6mm,right=6mm,top=0mm,bottom=4mm,sharp corners,
overlay={\fill[cexl] (frame.north west) rectangle ([yshift=-1.25cm]frame.north east);
\node[anchor=west,text=white,font=\sffamily\bfseries\LARGE] at ([xshift=5mm,yshift=-0.62cm]frame.north west) {Question 34 B};
\node[anchor=east,text=white,font=\sffamily\itshape\large] at ([xshift=-5mm,yshift=-0.62cm]frame.north east) {Exponentielle et logarithme};}]
\vspace*{1.75cm}
\begin{minipage}[t][6.9cm][c]{\linewidth}\LARGE\centering
\begin{minipage}{0.95\linewidth}\raggedright\linespread{1.25}\selectfont Simplifier $\dfrac{e^{3}\times e^{3}}{e^{2}}$.\end{minipage}
\end{minipage}

\vfill
\begin{tcolorbox}[enhanced,nobeforeafter,width=7.1cm,height=2.1cm,valign=center,colback=cexl!6,colframe=cexl,arc=2.5mm,boxrule=0.9pt,left=3mm]{\Large\bfseries\color{cexl}a)}\quad{\Large $e^{4}$}\end{tcolorbox}\hfill\begin{tcolorbox}[enhanced,nobeforeafter,width=7.1cm,height=2.1cm,valign=center,colback=cexl!6,colframe=cexl,arc=2.5mm,boxrule=0.9pt,left=3mm]{\Large\bfseries\color{cexl}b)}\quad{\Large $e^{7}$}\end{tcolorbox}\par\vspace{2.5mm}
\begin{tcolorbox}[enhanced,nobeforeafter,width=7.1cm,height=2.1cm,valign=center,colback=cexl!6,colframe=cexl,arc=2.5mm,boxrule=0.9pt,left=3mm]{\Large\bfseries\color{cexl}c)}\quad{\Large $e^{2}$}\end{tcolorbox}\hfill\begin{tcolorbox}[enhanced,nobeforeafter,width=7.1cm,height=2.1cm,valign=center,colback=cexl!6,colframe=cexl,arc=2.5mm,boxrule=0.9pt,left=3mm]{\Large\bfseries\color{cexl}d)}\quad{\Large $e^{8}$}\end{tcolorbox}
\end{tcolorbox}
\newpage
\begin{tcolorbox}[enhanced,width=16cm,height=14cm,colback=white,colframe=cexl,boxrule=2pt,arc=0mm,
left=6mm,right=6mm,top=0mm,bottom=4mm,sharp corners,
overlay={\fill[cexl] (frame.north west) rectangle ([yshift=-1.25cm]frame.north east);
\node[anchor=west,text=white,font=\sffamily\bfseries\LARGE] at ([xshift=5mm,yshift=-0.62cm]frame.north west) {Question 35 B};
\node[anchor=east,text=white,font=\sffamily\itshape\large] at ([xshift=-5mm,yshift=-0.62cm]frame.north east) {Exponentielle et logarithme};}]
\vspace*{1.75cm}
\begin{minipage}[t][6.9cm][c]{\linewidth}\LARGE\centering
\begin{minipage}{0.95\linewidth}\raggedright\linespread{1.25}\selectfont Quelle est la solution de l'équation $\ln(5x+1)=0$~?\end{minipage}
\end{minipage}

\vfill
\begin{tcolorbox}[enhanced,nobeforeafter,width=7.1cm,height=2.1cm,valign=center,colback=cexl!6,colframe=cexl,arc=2.5mm,boxrule=0.9pt,left=3mm]{\Large\bfseries\color{cexl}a)}\quad{\Large $-\dfrac{1}{5}$}\end{tcolorbox}\hfill\begin{tcolorbox}[enhanced,nobeforeafter,width=7.1cm,height=2.1cm,valign=center,colback=cexl!6,colframe=cexl,arc=2.5mm,boxrule=0.9pt,left=3mm]{\Large\bfseries\color{cexl}b)}\quad{\Large $\dfrac{e-1}{5}$}\end{tcolorbox}\par\vspace{2.5mm}
\begin{tcolorbox}[enhanced,nobeforeafter,width=7.1cm,height=2.1cm,valign=center,colback=cexl!6,colframe=cexl,arc=2.5mm,boxrule=0.9pt,left=3mm]{\Large\bfseries\color{cexl}c)}\quad{\Large $0$}\end{tcolorbox}\hfill\begin{tcolorbox}[enhanced,nobeforeafter,width=7.1cm,height=2.1cm,valign=center,colback=cexl!6,colframe=cexl,arc=2.5mm,boxrule=0.9pt,left=3mm]{\Large\bfseries\color{cexl}d)}\quad{\Large $\dfrac{2}{5}$}\end{tcolorbox}
\end{tcolorbox}
\newpage
\begin{tcolorbox}[enhanced,width=16cm,height=14cm,colback=white,colframe=cexl,boxrule=2pt,arc=0mm,
left=6mm,right=6mm,top=0mm,bottom=4mm,sharp corners,
overlay={\fill[cexl] (frame.north west) rectangle ([yshift=-1.25cm]frame.north east);
\node[anchor=west,text=white,font=\sffamily\bfseries\LARGE] at ([xshift=5mm,yshift=-0.62cm]frame.north west) {Question 36 B};
\node[anchor=east,text=white,font=\sffamily\itshape\large] at ([xshift=-5mm,yshift=-0.62cm]frame.north east) {Exponentielle et logarithme};}]
\vspace*{1.75cm}
\begin{minipage}[t][6.9cm][c]{\linewidth}\LARGE\centering
\begin{minipage}{0.95\linewidth}\raggedright\linespread{1.25}\selectfont Quelle est la solution de l'équation $\ln(4x-7)=0$~?\end{minipage}
\end{minipage}

\vfill
\begin{tcolorbox}[enhanced,nobeforeafter,width=7.1cm,height=2.1cm,valign=center,colback=cexl!6,colframe=cexl,arc=2.5mm,boxrule=0.9pt,left=3mm]{\Large\bfseries\color{cexl}a)}\quad{\Large $\dfrac{7}{4}$}\end{tcolorbox}\hfill\begin{tcolorbox}[enhanced,nobeforeafter,width=7.1cm,height=2.1cm,valign=center,colback=cexl!6,colframe=cexl,arc=2.5mm,boxrule=0.9pt,left=3mm]{\Large\bfseries\color{cexl}b)}\quad{\Large $\dfrac{e+7}{4}$}\end{tcolorbox}\par\vspace{2.5mm}
\begin{tcolorbox}[enhanced,nobeforeafter,width=7.1cm,height=2.1cm,valign=center,colback=cexl!6,colframe=cexl,arc=2.5mm,boxrule=0.9pt,left=3mm]{\Large\bfseries\color{cexl}c)}\quad{\Large $2$}\end{tcolorbox}\hfill\begin{tcolorbox}[enhanced,nobeforeafter,width=7.1cm,height=2.1cm,valign=center,colback=cexl!6,colframe=cexl,arc=2.5mm,boxrule=0.9pt,left=3mm]{\Large\bfseries\color{cexl}d)}\quad{\Large $-\dfrac{3}{2}$}\end{tcolorbox}
\end{tcolorbox}
\newpage
\begin{tcolorbox}[enhanced,width=16cm,height=14cm,colback=white,colframe=cexl,boxrule=2pt,arc=0mm,
left=6mm,right=6mm,top=0mm,bottom=4mm,sharp corners,
overlay={\fill[cexl] (frame.north west) rectangle ([yshift=-1.25cm]frame.north east);
\node[anchor=west,text=white,font=\sffamily\bfseries\LARGE] at ([xshift=5mm,yshift=-0.62cm]frame.north west) {Question 37 B};
\node[anchor=east,text=white,font=\sffamily\itshape\large] at ([xshift=-5mm,yshift=-0.62cm]frame.north east) {Exponentielle et logarithme};}]
\vspace*{1.75cm}
\begin{minipage}[t][6.9cm][c]{\linewidth}\LARGE\centering
\begin{minipage}{0.95\linewidth}\raggedright\linespread{1.25}\selectfont Simplifier $\ln 32-\ln 2$.\end{minipage}
\end{minipage}

\vfill
\begin{tcolorbox}[enhanced,nobeforeafter,width=7.1cm,height=2.1cm,valign=center,colback=cexl!6,colframe=cexl,arc=2.5mm,boxrule=0.9pt,left=3mm]{\Large\bfseries\color{cexl}a)}\quad{\Large $5\ln 2$}\end{tcolorbox}\hfill\begin{tcolorbox}[enhanced,nobeforeafter,width=7.1cm,height=2.1cm,valign=center,colback=cexl!6,colframe=cexl,arc=2.5mm,boxrule=0.9pt,left=3mm]{\Large\bfseries\color{cexl}b)}\quad{\Large $4\ln 2$}\end{tcolorbox}\par\vspace{2.5mm}
\begin{tcolorbox}[enhanced,nobeforeafter,width=7.1cm,height=2.1cm,valign=center,colback=cexl!6,colframe=cexl,arc=2.5mm,boxrule=0.9pt,left=3mm]{\Large\bfseries\color{cexl}c)}\quad{\Large $\ln 30$}\end{tcolorbox}\hfill\begin{tcolorbox}[enhanced,nobeforeafter,width=7.1cm,height=2.1cm,valign=center,colback=cexl!6,colframe=cexl,arc=2.5mm,boxrule=0.9pt,left=3mm]{\Large\bfseries\color{cexl}d)}\quad{\Large $6\ln 2$}\end{tcolorbox}
\end{tcolorbox}
\newpage
\begin{tcolorbox}[enhanced,width=16cm,height=14cm,colback=white,colframe=cexl,boxrule=2pt,arc=0mm,
left=6mm,right=6mm,top=0mm,bottom=4mm,sharp corners,
overlay={\fill[cexl] (frame.north west) rectangle ([yshift=-1.25cm]frame.north east);
\node[anchor=west,text=white,font=\sffamily\bfseries\LARGE] at ([xshift=5mm,yshift=-0.62cm]frame.north west) {Question 38 B};
\node[anchor=east,text=white,font=\sffamily\itshape\large] at ([xshift=-5mm,yshift=-0.62cm]frame.north east) {Exponentielle et logarithme};}]
\vspace*{1.75cm}
\begin{minipage}[t][6.9cm][c]{\linewidth}\LARGE\centering
\begin{minipage}{0.95\linewidth}\raggedright\linespread{1.25}\selectfont Simplifier $\ln 125-\ln 5$.\end{minipage}
\end{minipage}

\vfill
\begin{tcolorbox}[enhanced,nobeforeafter,width=7.1cm,height=2.1cm,valign=center,colback=cexl!6,colframe=cexl,arc=2.5mm,boxrule=0.9pt,left=3mm]{\Large\bfseries\color{cexl}a)}\quad{\Large $\ln 120$}\end{tcolorbox}\hfill\begin{tcolorbox}[enhanced,nobeforeafter,width=7.1cm,height=2.1cm,valign=center,colback=cexl!6,colframe=cexl,arc=2.5mm,boxrule=0.9pt,left=3mm]{\Large\bfseries\color{cexl}b)}\quad{\Large $3\ln 5$}\end{tcolorbox}\par\vspace{2.5mm}
\begin{tcolorbox}[enhanced,nobeforeafter,width=7.1cm,height=2.1cm,valign=center,colback=cexl!6,colframe=cexl,arc=2.5mm,boxrule=0.9pt,left=3mm]{\Large\bfseries\color{cexl}c)}\quad{\Large $4\ln 5$}\end{tcolorbox}\hfill\begin{tcolorbox}[enhanced,nobeforeafter,width=7.1cm,height=2.1cm,valign=center,colback=cexl!6,colframe=cexl,arc=2.5mm,boxrule=0.9pt,left=3mm]{\Large\bfseries\color{cexl}d)}\quad{\Large $2\ln 5$}\end{tcolorbox}
\end{tcolorbox}
\newpage
\begin{tcolorbox}[enhanced,width=16cm,height=14cm,colback=white,colframe=cexl,boxrule=2pt,arc=0mm,
left=6mm,right=6mm,top=0mm,bottom=4mm,sharp corners,
overlay={\fill[cexl] (frame.north west) rectangle ([yshift=-1.25cm]frame.north east);
\node[anchor=west,text=white,font=\sffamily\bfseries\LARGE] at ([xshift=5mm,yshift=-0.62cm]frame.north west) {Question 39 B};
\node[anchor=east,text=white,font=\sffamily\itshape\large] at ([xshift=-5mm,yshift=-0.62cm]frame.north east) {Exponentielle et logarithme};}]
\vspace*{1.75cm}
\begin{minipage}[t][6.9cm][c]{\linewidth}\LARGE\centering
\begin{minipage}{0.95\linewidth}\raggedright\linespread{1.25}\selectfont Soit $f(x)=x\,e^{-x}$. Que vaut $f'(1)$~?\end{minipage}
\end{minipage}

\vfill
\begin{tcolorbox}[enhanced,nobeforeafter,width=7.1cm,height=2.1cm,valign=center,colback=cexl!6,colframe=cexl,arc=2.5mm,boxrule=0.9pt,left=3mm]{\Large\bfseries\color{cexl}a)}\quad{\Large $e^{-1}$}\end{tcolorbox}\hfill\begin{tcolorbox}[enhanced,nobeforeafter,width=7.1cm,height=2.1cm,valign=center,colback=cexl!6,colframe=cexl,arc=2.5mm,boxrule=0.9pt,left=3mm]{\Large\bfseries\color{cexl}b)}\quad{\Large $-e^{-1}$}\end{tcolorbox}\par\vspace{2.5mm}
\begin{tcolorbox}[enhanced,nobeforeafter,width=7.1cm,height=2.1cm,valign=center,colback=cexl!6,colframe=cexl,arc=2.5mm,boxrule=0.9pt,left=3mm]{\Large\bfseries\color{cexl}c)}\quad{\Large $0$}\end{tcolorbox}\hfill\begin{tcolorbox}[enhanced,nobeforeafter,width=7.1cm,height=2.1cm,valign=center,colback=cexl!6,colframe=cexl,arc=2.5mm,boxrule=0.9pt,left=3mm]{\Large\bfseries\color{cexl}d)}\quad{\Large $2e^{-1}$}\end{tcolorbox}
\end{tcolorbox}
\newpage
\begin{tcolorbox}[enhanced,width=16cm,height=14cm,colback=white,colframe=cexl,boxrule=2pt,arc=0mm,
left=6mm,right=6mm,top=0mm,bottom=4mm,sharp corners,
overlay={\fill[cexl] (frame.north west) rectangle ([yshift=-1.25cm]frame.north east);
\node[anchor=west,text=white,font=\sffamily\bfseries\LARGE] at ([xshift=5mm,yshift=-0.62cm]frame.north west) {Question 40 B};
\node[anchor=east,text=white,font=\sffamily\itshape\large] at ([xshift=-5mm,yshift=-0.62cm]frame.north east) {Exponentielle et logarithme};}]
\vspace*{1.75cm}
\begin{minipage}[t][6.9cm][c]{\linewidth}\LARGE\centering
\begin{minipage}{0.95\linewidth}\raggedright\linespread{1.25}\selectfont Soit $f(x)=x\,e^{3x}$. Que vaut $f'(1)$~?\end{minipage}
\end{minipage}

\vfill
\begin{tcolorbox}[enhanced,nobeforeafter,width=7.1cm,height=2.1cm,valign=center,colback=cexl!6,colframe=cexl,arc=2.5mm,boxrule=0.9pt,left=3mm]{\Large\bfseries\color{cexl}a)}\quad{\Large $3e^{3}$}\end{tcolorbox}\hfill\begin{tcolorbox}[enhanced,nobeforeafter,width=7.1cm,height=2.1cm,valign=center,colback=cexl!6,colframe=cexl,arc=2.5mm,boxrule=0.9pt,left=3mm]{\Large\bfseries\color{cexl}b)}\quad{\Large $4e^{3}$}\end{tcolorbox}\par\vspace{2.5mm}
\begin{tcolorbox}[enhanced,nobeforeafter,width=7.1cm,height=2.1cm,valign=center,colback=cexl!6,colframe=cexl,arc=2.5mm,boxrule=0.9pt,left=3mm]{\Large\bfseries\color{cexl}c)}\quad{\Large $-2e^{3}$}\end{tcolorbox}\hfill\begin{tcolorbox}[enhanced,nobeforeafter,width=7.1cm,height=2.1cm,valign=center,colback=cexl!6,colframe=cexl,arc=2.5mm,boxrule=0.9pt,left=3mm]{\Large\bfseries\color{cexl}d)}\quad{\Large $e^{3}$}\end{tcolorbox}
\end{tcolorbox}
\newpage
\begin{tcolorbox}[enhanced,width=16cm,height=14cm,colback=white,colframe=cexl,boxrule=2pt,arc=0mm,
left=6mm,right=6mm,top=0mm,bottom=4mm,sharp corners,
overlay={\fill[cexl] (frame.north west) rectangle ([yshift=-1.25cm]frame.north east);
\node[anchor=west,text=white,font=\sffamily\bfseries\LARGE] at ([xshift=5mm,yshift=-0.62cm]frame.north west) {Question 41 B};
\node[anchor=east,text=white,font=\sffamily\itshape\large] at ([xshift=-5mm,yshift=-0.62cm]frame.north east) {Exponentielle et logarithme};}]
\vspace*{1.75cm}
\begin{minipage}[t][6.9cm][c]{\linewidth}\LARGE\centering
\begin{minipage}{0.95\linewidth}\raggedright\linespread{1.25}\selectfont Que vaut $\displaystyle\lim_{x\to+\infty}\frac{3e^x+x^2}{4e^x+1}$~?\end{minipage}
\end{minipage}

\vfill
\begin{tcolorbox}[enhanced,nobeforeafter,width=7.1cm,height=2.1cm,valign=center,colback=cexl!6,colframe=cexl,arc=2.5mm,boxrule=0.9pt,left=3mm]{\Large\bfseries\color{cexl}a)}\quad{\Large $0$}\end{tcolorbox}\hfill\begin{tcolorbox}[enhanced,nobeforeafter,width=7.1cm,height=2.1cm,valign=center,colback=cexl!6,colframe=cexl,arc=2.5mm,boxrule=0.9pt,left=3mm]{\Large\bfseries\color{cexl}b)}\quad{\Large $\dfrac{3}{4}$}\end{tcolorbox}\par\vspace{2.5mm}
\begin{tcolorbox}[enhanced,nobeforeafter,width=7.1cm,height=2.1cm,valign=center,colback=cexl!6,colframe=cexl,arc=2.5mm,boxrule=0.9pt,left=3mm]{\Large\bfseries\color{cexl}c)}\quad{\Large $+\infty$}\end{tcolorbox}\hfill\begin{tcolorbox}[enhanced,nobeforeafter,width=7.1cm,height=2.1cm,valign=center,colback=cexl!6,colframe=cexl,arc=2.5mm,boxrule=0.9pt,left=3mm]{\Large\bfseries\color{cexl}d)}\quad{\Large $\dfrac{4}{3}$}\end{tcolorbox}
\end{tcolorbox}
\newpage
\begin{tcolorbox}[enhanced,width=16cm,height=14cm,colback=white,colframe=cexl,boxrule=2pt,arc=0mm,
left=6mm,right=6mm,top=0mm,bottom=4mm,sharp corners,
overlay={\fill[cexl] (frame.north west) rectangle ([yshift=-1.25cm]frame.north east);
\node[anchor=west,text=white,font=\sffamily\bfseries\LARGE] at ([xshift=5mm,yshift=-0.62cm]frame.north west) {Question 42 B};
\node[anchor=east,text=white,font=\sffamily\itshape\large] at ([xshift=-5mm,yshift=-0.62cm]frame.north east) {Exponentielle et logarithme};}]
\vspace*{1.75cm}
\begin{minipage}[t][6.9cm][c]{\linewidth}\LARGE\centering
\begin{minipage}{0.95\linewidth}\raggedright\linespread{1.25}\selectfont Que vaut $\displaystyle\lim_{x\to+\infty}\frac{7e^x+x^2}{5e^x+1}$~?\end{minipage}
\end{minipage}

\vfill
\begin{tcolorbox}[enhanced,nobeforeafter,width=7.1cm,height=2.1cm,valign=center,colback=cexl!6,colframe=cexl,arc=2.5mm,boxrule=0.9pt,left=3mm]{\Large\bfseries\color{cexl}a)}\quad{\Large $+\infty$}\end{tcolorbox}\hfill\begin{tcolorbox}[enhanced,nobeforeafter,width=7.1cm,height=2.1cm,valign=center,colback=cexl!6,colframe=cexl,arc=2.5mm,boxrule=0.9pt,left=3mm]{\Large\bfseries\color{cexl}b)}\quad{\Large $\dfrac{7}{5}$}\end{tcolorbox}\par\vspace{2.5mm}
\begin{tcolorbox}[enhanced,nobeforeafter,width=7.1cm,height=2.1cm,valign=center,colback=cexl!6,colframe=cexl,arc=2.5mm,boxrule=0.9pt,left=3mm]{\Large\bfseries\color{cexl}c)}\quad{\Large $\dfrac{5}{7}$}\end{tcolorbox}\hfill\begin{tcolorbox}[enhanced,nobeforeafter,width=7.1cm,height=2.1cm,valign=center,colback=cexl!6,colframe=cexl,arc=2.5mm,boxrule=0.9pt,left=3mm]{\Large\bfseries\color{cexl}d)}\quad{\Large $0$}\end{tcolorbox}
\end{tcolorbox}
\newpage
\begin{tcolorbox}[enhanced,width=16cm,height=14cm,colback=white,colframe=cexl,boxrule=2pt,arc=0mm,
left=6mm,right=6mm,top=0mm,bottom=4mm,sharp corners,
overlay={\fill[cexl] (frame.north west) rectangle ([yshift=-1.25cm]frame.north east);
\node[anchor=west,text=white,font=\sffamily\bfseries\LARGE] at ([xshift=5mm,yshift=-0.62cm]frame.north west) {Question 43 B};
\node[anchor=east,text=white,font=\sffamily\itshape\large] at ([xshift=-5mm,yshift=-0.62cm]frame.north east) {Exponentielle et logarithme};}]
\vspace*{1.75cm}
\begin{minipage}[t][6.9cm][c]{\linewidth}\LARGE\centering
\begin{minipage}{0.95\linewidth}\raggedright\linespread{1.25}\selectfont Soit $f(x)=\ln(3x^2+1)$. Que vaut $f'(1)$~?\end{minipage}
\end{minipage}

\vfill
\begin{tcolorbox}[enhanced,nobeforeafter,width=7.1cm,height=2.1cm,valign=center,colback=cexl!6,colframe=cexl,arc=2.5mm,boxrule=0.9pt,left=3mm]{\Large\bfseries\color{cexl}a)}\quad{\Large $\dfrac{1}{4}$}\end{tcolorbox}\hfill\begin{tcolorbox}[enhanced,nobeforeafter,width=7.1cm,height=2.1cm,valign=center,colback=cexl!6,colframe=cexl,arc=2.5mm,boxrule=0.9pt,left=3mm]{\Large\bfseries\color{cexl}b)}\quad{\Large $\dfrac{3}{2}$}\end{tcolorbox}\par\vspace{2.5mm}
\begin{tcolorbox}[enhanced,nobeforeafter,width=7.1cm,height=2.1cm,valign=center,colback=cexl!6,colframe=cexl,arc=2.5mm,boxrule=0.9pt,left=3mm]{\Large\bfseries\color{cexl}c)}\quad{\Large $\dfrac{3}{4}$}\end{tcolorbox}\hfill\begin{tcolorbox}[enhanced,nobeforeafter,width=7.1cm,height=2.1cm,valign=center,colback=cexl!6,colframe=cexl,arc=2.5mm,boxrule=0.9pt,left=3mm]{\Large\bfseries\color{cexl}d)}\quad{\Large $6$}\end{tcolorbox}
\end{tcolorbox}
\newpage
\begin{tcolorbox}[enhanced,width=16cm,height=14cm,colback=white,colframe=cexl,boxrule=2pt,arc=0mm,
left=6mm,right=6mm,top=0mm,bottom=4mm,sharp corners,
overlay={\fill[cexl] (frame.north west) rectangle ([yshift=-1.25cm]frame.north east);
\node[anchor=west,text=white,font=\sffamily\bfseries\LARGE] at ([xshift=5mm,yshift=-0.62cm]frame.north west) {Question 44 B};
\node[anchor=east,text=white,font=\sffamily\itshape\large] at ([xshift=-5mm,yshift=-0.62cm]frame.north east) {Exponentielle et logarithme};}]
\vspace*{1.75cm}
\begin{minipage}[t][6.9cm][c]{\linewidth}\LARGE\centering
\begin{minipage}{0.95\linewidth}\raggedright\linespread{1.25}\selectfont Soit $f(x)=\ln(x^2+4)$. Que vaut $f'(1)$~?\end{minipage}
\end{minipage}

\vfill
\begin{tcolorbox}[enhanced,nobeforeafter,width=7.1cm,height=2.1cm,valign=center,colback=cexl!6,colframe=cexl,arc=2.5mm,boxrule=0.9pt,left=3mm]{\Large\bfseries\color{cexl}a)}\quad{\Large $2$}\end{tcolorbox}\hfill\begin{tcolorbox}[enhanced,nobeforeafter,width=7.1cm,height=2.1cm,valign=center,colback=cexl!6,colframe=cexl,arc=2.5mm,boxrule=0.9pt,left=3mm]{\Large\bfseries\color{cexl}b)}\quad{\Large $\dfrac{2}{5}$}\end{tcolorbox}\par\vspace{2.5mm}
\begin{tcolorbox}[enhanced,nobeforeafter,width=7.1cm,height=2.1cm,valign=center,colback=cexl!6,colframe=cexl,arc=2.5mm,boxrule=0.9pt,left=3mm]{\Large\bfseries\color{cexl}c)}\quad{\Large $\dfrac{3}{5}$}\end{tcolorbox}\hfill\begin{tcolorbox}[enhanced,nobeforeafter,width=7.1cm,height=2.1cm,valign=center,colback=cexl!6,colframe=cexl,arc=2.5mm,boxrule=0.9pt,left=3mm]{\Large\bfseries\color{cexl}d)}\quad{\Large $\dfrac{1}{5}$}\end{tcolorbox}
\end{tcolorbox}
\newpage
\begin{tcolorbox}[enhanced,width=16cm,height=14cm,colback=white,colframe=cexl,boxrule=2pt,arc=0mm,
left=6mm,right=6mm,top=0mm,bottom=4mm,sharp corners,
overlay={\fill[cexl] (frame.north west) rectangle ([yshift=-1.25cm]frame.north east);
\node[anchor=west,text=white,font=\sffamily\bfseries\LARGE] at ([xshift=5mm,yshift=-0.62cm]frame.north west) {Question 45 B};
\node[anchor=east,text=white,font=\sffamily\itshape\large] at ([xshift=-5mm,yshift=-0.62cm]frame.north east) {Exponentielle et logarithme};}]
\vspace*{1.75cm}
\begin{minipage}[t][6.9cm][c]{\linewidth}\LARGE\centering
\begin{minipage}{0.95\linewidth}\raggedright\linespread{1.25}\selectfont Quel est le plus petit entier naturel $n$ tel que $0{,}95^n<0{,}5$~?\end{minipage}
\end{minipage}

\vfill
\begin{tcolorbox}[enhanced,nobeforeafter,width=7.1cm,height=2.1cm,valign=center,colback=cexl!6,colframe=cexl,arc=2.5mm,boxrule=0.9pt,left=3mm]{\Large\bfseries\color{cexl}a)}\quad{\Large $14$}\end{tcolorbox}\hfill\begin{tcolorbox}[enhanced,nobeforeafter,width=7.1cm,height=2.1cm,valign=center,colback=cexl!6,colframe=cexl,arc=2.5mm,boxrule=0.9pt,left=3mm]{\Large\bfseries\color{cexl}b)}\quad{\Large $13$}\end{tcolorbox}\par\vspace{2.5mm}
\begin{tcolorbox}[enhanced,nobeforeafter,width=7.1cm,height=2.1cm,valign=center,colback=cexl!6,colframe=cexl,arc=2.5mm,boxrule=0.9pt,left=3mm]{\Large\bfseries\color{cexl}c)}\quad{\Large $16$}\end{tcolorbox}\hfill\begin{tcolorbox}[enhanced,nobeforeafter,width=7.1cm,height=2.1cm,valign=center,colback=cexl!6,colframe=cexl,arc=2.5mm,boxrule=0.9pt,left=3mm]{\Large\bfseries\color{cexl}d)}\quad{\Large $15$}\end{tcolorbox}
\end{tcolorbox}
\newpage
\begin{tcolorbox}[enhanced,width=16cm,height=14cm,colback=white,colframe=cexl,boxrule=2pt,arc=0mm,
left=6mm,right=6mm,top=0mm,bottom=4mm,sharp corners,
overlay={\fill[cexl] (frame.north west) rectangle ([yshift=-1.25cm]frame.north east);
\node[anchor=west,text=white,font=\sffamily\bfseries\LARGE] at ([xshift=5mm,yshift=-0.62cm]frame.north west) {Question 46 B};
\node[anchor=east,text=white,font=\sffamily\itshape\large] at ([xshift=-5mm,yshift=-0.62cm]frame.north east) {Exponentielle et logarithme};}]
\vspace*{1.75cm}
\begin{minipage}[t][6.9cm][c]{\linewidth}\LARGE\centering
\begin{minipage}{0.95\linewidth}\raggedright\linespread{1.25}\selectfont Quel est le plus petit entier naturel $n$ tel que $0{,}7^n<0{,}01$~?\end{minipage}
\end{minipage}

\vfill
\begin{tcolorbox}[enhanced,nobeforeafter,width=7.1cm,height=2.1cm,valign=center,colback=cexl!6,colframe=cexl,arc=2.5mm,boxrule=0.9pt,left=3mm]{\Large\bfseries\color{cexl}a)}\quad{\Large $12$}\end{tcolorbox}\hfill\begin{tcolorbox}[enhanced,nobeforeafter,width=7.1cm,height=2.1cm,valign=center,colback=cexl!6,colframe=cexl,arc=2.5mm,boxrule=0.9pt,left=3mm]{\Large\bfseries\color{cexl}b)}\quad{\Large $14$}\end{tcolorbox}\par\vspace{2.5mm}
\begin{tcolorbox}[enhanced,nobeforeafter,width=7.1cm,height=2.1cm,valign=center,colback=cexl!6,colframe=cexl,arc=2.5mm,boxrule=0.9pt,left=3mm]{\Large\bfseries\color{cexl}c)}\quad{\Large $15$}\end{tcolorbox}\hfill\begin{tcolorbox}[enhanced,nobeforeafter,width=7.1cm,height=2.1cm,valign=center,colback=cexl!6,colframe=cexl,arc=2.5mm,boxrule=0.9pt,left=3mm]{\Large\bfseries\color{cexl}d)}\quad{\Large $13$}\end{tcolorbox}
\end{tcolorbox}
\newpage
\begin{tcolorbox}[enhanced,width=16cm,height=14cm,colback=white,colframe=cexl,boxrule=2pt,arc=0mm,
left=6mm,right=6mm,top=0mm,bottom=4mm,sharp corners,
overlay={\fill[cexl] (frame.north west) rectangle ([yshift=-1.25cm]frame.north east);
\node[anchor=west,text=white,font=\sffamily\bfseries\LARGE] at ([xshift=5mm,yshift=-0.62cm]frame.north west) {Question 47 B};
\node[anchor=east,text=white,font=\sffamily\itshape\large] at ([xshift=-5mm,yshift=-0.62cm]frame.north east) {Exponentielle et logarithme};}]
\vspace*{1.75cm}
\begin{minipage}[t][6.9cm][c]{\linewidth}\LARGE\centering
\begin{minipage}{0.95\linewidth}\raggedright\linespread{1.25}\selectfont Quelle est la solution de l'équation $e^{x+4}=2$~?\end{minipage}
\end{minipage}

\vfill
\begin{tcolorbox}[enhanced,nobeforeafter,width=7.1cm,height=2.1cm,valign=center,colback=cexl!6,colframe=cexl,arc=2.5mm,boxrule=0.9pt,left=3mm]{\Large\bfseries\color{cexl}a)}\quad{\Large $\ln 2-4$}\end{tcolorbox}\hfill\begin{tcolorbox}[enhanced,nobeforeafter,width=7.1cm,height=2.1cm,valign=center,colback=cexl!6,colframe=cexl,arc=2.5mm,boxrule=0.9pt,left=3mm]{\Large\bfseries\color{cexl}b)}\quad{\Large $e^{2}$}\end{tcolorbox}\par\vspace{2.5mm}
\begin{tcolorbox}[enhanced,nobeforeafter,width=7.1cm,height=2.1cm,valign=center,colback=cexl!6,colframe=cexl,arc=2.5mm,boxrule=0.9pt,left=3mm]{\Large\bfseries\color{cexl}c)}\quad{\Large $\ln 2+4$}\end{tcolorbox}\hfill\begin{tcolorbox}[enhanced,nobeforeafter,width=7.1cm,height=2.1cm,valign=center,colback=cexl!6,colframe=cexl,arc=2.5mm,boxrule=0.9pt,left=3mm]{\Large\bfseries\color{cexl}d)}\quad{\Large $\ln 2$}\end{tcolorbox}
\end{tcolorbox}
\newpage
\begin{tcolorbox}[enhanced,width=16cm,height=14cm,colback=white,colframe=cexl,boxrule=2pt,arc=0mm,
left=6mm,right=6mm,top=0mm,bottom=4mm,sharp corners,
overlay={\fill[cexl] (frame.north west) rectangle ([yshift=-1.25cm]frame.north east);
\node[anchor=west,text=white,font=\sffamily\bfseries\LARGE] at ([xshift=5mm,yshift=-0.62cm]frame.north west) {Question 48 B};
\node[anchor=east,text=white,font=\sffamily\itshape\large] at ([xshift=-5mm,yshift=-0.62cm]frame.north east) {Exponentielle et logarithme};}]
\vspace*{1.75cm}
\begin{minipage}[t][6.9cm][c]{\linewidth}\LARGE\centering
\begin{minipage}{0.95\linewidth}\raggedright\linespread{1.25}\selectfont Quelle est la solution de l'équation $e^{2x-1}=7$~?\end{minipage}
\end{minipage}

\vfill
\begin{tcolorbox}[enhanced,nobeforeafter,width=7.1cm,height=2.1cm,valign=center,colback=cexl!6,colframe=cexl,arc=2.5mm,boxrule=0.9pt,left=3mm]{\Large\bfseries\color{cexl}a)}\quad{\Large $\dfrac{\ln 7-1}{2}$}\end{tcolorbox}\hfill\begin{tcolorbox}[enhanced,nobeforeafter,width=7.1cm,height=2.1cm,valign=center,colback=cexl!6,colframe=cexl,arc=2.5mm,boxrule=0.9pt,left=3mm]{\Large\bfseries\color{cexl}b)}\quad{\Large $\dfrac{e^{7}}{2}$}\end{tcolorbox}\par\vspace{2.5mm}
\begin{tcolorbox}[enhanced,nobeforeafter,width=7.1cm,height=2.1cm,valign=center,colback=cexl!6,colframe=cexl,arc=2.5mm,boxrule=0.9pt,left=3mm]{\Large\bfseries\color{cexl}c)}\quad{\Large $\dfrac{\ln 7}{2}+1$}\end{tcolorbox}\hfill\begin{tcolorbox}[enhanced,nobeforeafter,width=7.1cm,height=2.1cm,valign=center,colback=cexl!6,colframe=cexl,arc=2.5mm,boxrule=0.9pt,left=3mm]{\Large\bfseries\color{cexl}d)}\quad{\Large $\dfrac{\ln 7+1}{2}$}\end{tcolorbox}
\end{tcolorbox}
\newpage
\begin{tcolorbox}[enhanced,width=16cm,height=14cm,colback=white,colframe=cint,boxrule=2pt,arc=0mm,
left=6mm,right=6mm,top=0mm,bottom=4mm,sharp corners,
overlay={\fill[cint] (frame.north west) rectangle ([yshift=-1.25cm]frame.north east);
\node[anchor=west,text=white,font=\sffamily\bfseries\LARGE] at ([xshift=5mm,yshift=-0.62cm]frame.north west) {Question 49 B};
\node[anchor=east,text=white,font=\sffamily\itshape\large] at ([xshift=-5mm,yshift=-0.62cm]frame.north east) {Intégrales et équations différentielles};}]
\vspace*{1.75cm}
\begin{minipage}[t][6.9cm][c]{\linewidth}\LARGE\centering
\begin{minipage}{0.95\linewidth}\raggedright\linespread{1.25}\selectfont Que vaut $\displaystyle\int_0^{1} x^{3}\,\mathrm{d}x$~?\end{minipage}
\end{minipage}

\vfill
\begin{tcolorbox}[enhanced,nobeforeafter,width=7.1cm,height=2.1cm,valign=center,colback=cint!6,colframe=cint,arc=2.5mm,boxrule=0.9pt,left=3mm]{\Large\bfseries\color{cint}a)}\quad{\Large $\dfrac{1}{3}$}\end{tcolorbox}\hfill\begin{tcolorbox}[enhanced,nobeforeafter,width=7.1cm,height=2.1cm,valign=center,colback=cint!6,colframe=cint,arc=2.5mm,boxrule=0.9pt,left=3mm]{\Large\bfseries\color{cint}b)}\quad{\Large $3$}\end{tcolorbox}\par\vspace{2.5mm}
\begin{tcolorbox}[enhanced,nobeforeafter,width=7.1cm,height=2.1cm,valign=center,colback=cint!6,colframe=cint,arc=2.5mm,boxrule=0.9pt,left=3mm]{\Large\bfseries\color{cint}c)}\quad{\Large $1$}\end{tcolorbox}\hfill\begin{tcolorbox}[enhanced,nobeforeafter,width=7.1cm,height=2.1cm,valign=center,colback=cint!6,colframe=cint,arc=2.5mm,boxrule=0.9pt,left=3mm]{\Large\bfseries\color{cint}d)}\quad{\Large $\dfrac{1}{4}$}\end{tcolorbox}
\end{tcolorbox}
\newpage
\begin{tcolorbox}[enhanced,width=16cm,height=14cm,colback=white,colframe=cint,boxrule=2pt,arc=0mm,
left=6mm,right=6mm,top=0mm,bottom=4mm,sharp corners,
overlay={\fill[cint] (frame.north west) rectangle ([yshift=-1.25cm]frame.north east);
\node[anchor=west,text=white,font=\sffamily\bfseries\LARGE] at ([xshift=5mm,yshift=-0.62cm]frame.north west) {Question 50 B};
\node[anchor=east,text=white,font=\sffamily\itshape\large] at ([xshift=-5mm,yshift=-0.62cm]frame.north east) {Intégrales et équations différentielles};}]
\vspace*{1.75cm}
\begin{minipage}[t][6.9cm][c]{\linewidth}\LARGE\centering
\begin{minipage}{0.95\linewidth}\raggedright\linespread{1.25}\selectfont Que vaut $\displaystyle\int_0^{2} x^{3}\,\mathrm{d}x$~?\end{minipage}
\end{minipage}

\vfill
\begin{tcolorbox}[enhanced,nobeforeafter,width=7.1cm,height=2.1cm,valign=center,colback=cint!6,colframe=cint,arc=2.5mm,boxrule=0.9pt,left=3mm]{\Large\bfseries\color{cint}a)}\quad{\Large $4$}\end{tcolorbox}\hfill\begin{tcolorbox}[enhanced,nobeforeafter,width=7.1cm,height=2.1cm,valign=center,colback=cint!6,colframe=cint,arc=2.5mm,boxrule=0.9pt,left=3mm]{\Large\bfseries\color{cint}b)}\quad{\Large $16$}\end{tcolorbox}\par\vspace{2.5mm}
\begin{tcolorbox}[enhanced,nobeforeafter,width=7.1cm,height=2.1cm,valign=center,colback=cint!6,colframe=cint,arc=2.5mm,boxrule=0.9pt,left=3mm]{\Large\bfseries\color{cint}c)}\quad{\Large $12$}\end{tcolorbox}\hfill\begin{tcolorbox}[enhanced,nobeforeafter,width=7.1cm,height=2.1cm,valign=center,colback=cint!6,colframe=cint,arc=2.5mm,boxrule=0.9pt,left=3mm]{\Large\bfseries\color{cint}d)}\quad{\Large $\dfrac{8}{3}$}\end{tcolorbox}
\end{tcolorbox}
\newpage
\begin{tcolorbox}[enhanced,width=16cm,height=14cm,colback=white,colframe=cint,boxrule=2pt,arc=0mm,
left=6mm,right=6mm,top=0mm,bottom=4mm,sharp corners,
overlay={\fill[cint] (frame.north west) rectangle ([yshift=-1.25cm]frame.north east);
\node[anchor=west,text=white,font=\sffamily\bfseries\LARGE] at ([xshift=5mm,yshift=-0.62cm]frame.north west) {Question 51 B};
\node[anchor=east,text=white,font=\sffamily\itshape\large] at ([xshift=-5mm,yshift=-0.62cm]frame.north east) {Intégrales et équations différentielles};}]
\vspace*{1.75cm}
\begin{minipage}[t][6.9cm][c]{\linewidth}\LARGE\centering
\begin{minipage}{0.95\linewidth}\raggedright\linespread{1.25}\selectfont Que vaut $\displaystyle\int_0^{1} e^{x}\,\mathrm{d}x$~?\end{minipage}
\end{minipage}

\vfill
\begin{tcolorbox}[enhanced,nobeforeafter,width=7.1cm,height=2.1cm,valign=center,colback=cint!6,colframe=cint,arc=2.5mm,boxrule=0.9pt,left=3mm]{\Large\bfseries\color{cint}a)}\quad{\Large $e-1$}\end{tcolorbox}\hfill\begin{tcolorbox}[enhanced,nobeforeafter,width=7.1cm,height=2.1cm,valign=center,colback=cint!6,colframe=cint,arc=2.5mm,boxrule=0.9pt,left=3mm]{\Large\bfseries\color{cint}b)}\quad{\Large $e$}\end{tcolorbox}\par\vspace{2.5mm}
\begin{tcolorbox}[enhanced,nobeforeafter,width=7.1cm,height=2.1cm,valign=center,colback=cint!6,colframe=cint,arc=2.5mm,boxrule=0.9pt,left=3mm]{\Large\bfseries\color{cint}c)}\quad{\Large $e+1$}\end{tcolorbox}\hfill\begin{tcolorbox}[enhanced,nobeforeafter,width=7.1cm,height=2.1cm,valign=center,colback=cint!6,colframe=cint,arc=2.5mm,boxrule=0.9pt,left=3mm]{\Large\bfseries\color{cint}d)}\quad{\Large $\dfrac{e-1}{2}$}\end{tcolorbox}
\end{tcolorbox}
\newpage
\begin{tcolorbox}[enhanced,width=16cm,height=14cm,colback=white,colframe=cint,boxrule=2pt,arc=0mm,
left=6mm,right=6mm,top=0mm,bottom=4mm,sharp corners,
overlay={\fill[cint] (frame.north west) rectangle ([yshift=-1.25cm]frame.north east);
\node[anchor=west,text=white,font=\sffamily\bfseries\LARGE] at ([xshift=5mm,yshift=-0.62cm]frame.north west) {Question 52 B};
\node[anchor=east,text=white,font=\sffamily\itshape\large] at ([xshift=-5mm,yshift=-0.62cm]frame.north east) {Intégrales et équations différentielles};}]
\vspace*{1.75cm}
\begin{minipage}[t][6.9cm][c]{\linewidth}\LARGE\centering
\begin{minipage}{0.95\linewidth}\raggedright\linespread{1.25}\selectfont Que vaut $\displaystyle\int_0^{1} e^{4x}\,\mathrm{d}x$~?\end{minipage}
\end{minipage}

\vfill
\begin{tcolorbox}[enhanced,nobeforeafter,width=7.1cm,height=2.1cm,valign=center,colback=cint!6,colframe=cint,arc=2.5mm,boxrule=0.9pt,left=3mm]{\Large\bfseries\color{cint}a)}\quad{\Large $\dfrac{e^{4}-1}{4}$}\end{tcolorbox}\hfill\begin{tcolorbox}[enhanced,nobeforeafter,width=7.1cm,height=2.1cm,valign=center,colback=cint!6,colframe=cint,arc=2.5mm,boxrule=0.9pt,left=3mm]{\Large\bfseries\color{cint}b)}\quad{\Large $e^{4}-1$}\end{tcolorbox}\par\vspace{2.5mm}
\begin{tcolorbox}[enhanced,nobeforeafter,width=7.1cm,height=2.1cm,valign=center,colback=cint!6,colframe=cint,arc=2.5mm,boxrule=0.9pt,left=3mm]{\Large\bfseries\color{cint}c)}\quad{\Large $4(e^{4}-1)$}\end{tcolorbox}\hfill\begin{tcolorbox}[enhanced,nobeforeafter,width=7.1cm,height=2.1cm,valign=center,colback=cint!6,colframe=cint,arc=2.5mm,boxrule=0.9pt,left=3mm]{\Large\bfseries\color{cint}d)}\quad{\Large $\dfrac{e^{4}}{4}$}\end{tcolorbox}
\end{tcolorbox}
\newpage
\begin{tcolorbox}[enhanced,width=16cm,height=14cm,colback=white,colframe=cint,boxrule=2pt,arc=0mm,
left=6mm,right=6mm,top=0mm,bottom=4mm,sharp corners,
overlay={\fill[cint] (frame.north west) rectangle ([yshift=-1.25cm]frame.north east);
\node[anchor=west,text=white,font=\sffamily\bfseries\LARGE] at ([xshift=5mm,yshift=-0.62cm]frame.north west) {Question 53 B};
\node[anchor=east,text=white,font=\sffamily\itshape\large] at ([xshift=-5mm,yshift=-0.62cm]frame.north east) {Intégrales et équations différentielles};}]
\vspace*{1.75cm}
\begin{minipage}[t][6.9cm][c]{\linewidth}\LARGE\centering
\begin{minipage}{0.95\linewidth}\raggedright\linespread{1.25}\selectfont Que vaut $\displaystyle\int_1^{e} \left(\frac{1}{x}+3\right)\mathrm{d}x$~?\end{minipage}
\end{minipage}

\vfill
\begin{tcolorbox}[enhanced,nobeforeafter,width=7.1cm,height=2.1cm,valign=center,colback=cint!6,colframe=cint,arc=2.5mm,boxrule=0.9pt,left=3mm]{\Large\bfseries\color{cint}a)}\quad{\Large $3e+1$}\end{tcolorbox}\hfill\begin{tcolorbox}[enhanced,nobeforeafter,width=7.1cm,height=2.1cm,valign=center,colback=cint!6,colframe=cint,arc=2.5mm,boxrule=0.9pt,left=3mm]{\Large\bfseries\color{cint}b)}\quad{\Large $3e-2$}\end{tcolorbox}\par\vspace{2.5mm}
\begin{tcolorbox}[enhanced,nobeforeafter,width=7.1cm,height=2.1cm,valign=center,colback=cint!6,colframe=cint,arc=2.5mm,boxrule=0.9pt,left=3mm]{\Large\bfseries\color{cint}c)}\quad{\Large $e+3$}\end{tcolorbox}\hfill\begin{tcolorbox}[enhanced,nobeforeafter,width=7.1cm,height=2.1cm,valign=center,colback=cint!6,colframe=cint,arc=2.5mm,boxrule=0.9pt,left=3mm]{\Large\bfseries\color{cint}d)}\quad{\Large $4e$}\end{tcolorbox}
\end{tcolorbox}
\newpage
\begin{tcolorbox}[enhanced,width=16cm,height=14cm,colback=white,colframe=cint,boxrule=2pt,arc=0mm,
left=6mm,right=6mm,top=0mm,bottom=4mm,sharp corners,
overlay={\fill[cint] (frame.north west) rectangle ([yshift=-1.25cm]frame.north east);
\node[anchor=west,text=white,font=\sffamily\bfseries\LARGE] at ([xshift=5mm,yshift=-0.62cm]frame.north west) {Question 54 B};
\node[anchor=east,text=white,font=\sffamily\itshape\large] at ([xshift=-5mm,yshift=-0.62cm]frame.north east) {Intégrales et équations différentielles};}]
\vspace*{1.75cm}
\begin{minipage}[t][6.9cm][c]{\linewidth}\LARGE\centering
\begin{minipage}{0.95\linewidth}\raggedright\linespread{1.25}\selectfont Que vaut $\displaystyle\int_1^{e} \left(\frac{4}{x}+1\right)\mathrm{d}x$~?\end{minipage}
\end{minipage}

\vfill
\begin{tcolorbox}[enhanced,nobeforeafter,width=7.1cm,height=2.1cm,valign=center,colback=cint!6,colframe=cint,arc=2.5mm,boxrule=0.9pt,left=3mm]{\Large\bfseries\color{cint}a)}\quad{\Large $4e+1$}\end{tcolorbox}\hfill\begin{tcolorbox}[enhanced,nobeforeafter,width=7.1cm,height=2.1cm,valign=center,colback=cint!6,colframe=cint,arc=2.5mm,boxrule=0.9pt,left=3mm]{\Large\bfseries\color{cint}b)}\quad{\Large $e+4$}\end{tcolorbox}\par\vspace{2.5mm}
\begin{tcolorbox}[enhanced,nobeforeafter,width=7.1cm,height=2.1cm,valign=center,colback=cint!6,colframe=cint,arc=2.5mm,boxrule=0.9pt,left=3mm]{\Large\bfseries\color{cint}c)}\quad{\Large $5e$}\end{tcolorbox}\hfill\begin{tcolorbox}[enhanced,nobeforeafter,width=7.1cm,height=2.1cm,valign=center,colback=cint!6,colframe=cint,arc=2.5mm,boxrule=0.9pt,left=3mm]{\Large\bfseries\color{cint}d)}\quad{\Large $e+3$}\end{tcolorbox}
\end{tcolorbox}
\newpage
\begin{tcolorbox}[enhanced,width=16cm,height=14cm,colback=white,colframe=cint,boxrule=2pt,arc=0mm,
left=6mm,right=6mm,top=0mm,bottom=4mm,sharp corners,
overlay={\fill[cint] (frame.north west) rectangle ([yshift=-1.25cm]frame.north east);
\node[anchor=west,text=white,font=\sffamily\bfseries\LARGE] at ([xshift=5mm,yshift=-0.62cm]frame.north west) {Question 55 B};
\node[anchor=east,text=white,font=\sffamily\itshape\large] at ([xshift=-5mm,yshift=-0.62cm]frame.north east) {Intégrales et équations différentielles};}]
\vspace*{1.75cm}
\begin{minipage}[t][6.9cm][c]{\linewidth}\LARGE\centering
\begin{minipage}{0.95\linewidth}\raggedright\linespread{1.25}\selectfont Quelle est la primitive $F$ de $f(x)=e^{-x}$ sur $\mathbb{R}$ telle que $F(0)=2$~?\end{minipage}
\end{minipage}

\vfill
\begin{tcolorbox}[enhanced,nobeforeafter,width=7.1cm,height=2.1cm,valign=center,colback=cint!6,colframe=cint,arc=2.5mm,boxrule=0.9pt,left=3mm]{\Large\bfseries\color{cint}a)}\quad{\Large $F(x)=-e^{-x}+1$}\end{tcolorbox}\hfill\begin{tcolorbox}[enhanced,nobeforeafter,width=7.1cm,height=2.1cm,valign=center,colback=cint!6,colframe=cint,arc=2.5mm,boxrule=0.9pt,left=3mm]{\Large\bfseries\color{cint}b)}\quad{\Large $F(x)=e^{-x}+1$}\end{tcolorbox}\par\vspace{2.5mm}
\begin{tcolorbox}[enhanced,nobeforeafter,width=7.1cm,height=2.1cm,valign=center,colback=cint!6,colframe=cint,arc=2.5mm,boxrule=0.9pt,left=3mm]{\Large\bfseries\color{cint}c)}\quad{\Large $F(x)=-e^{-x}+3$}\end{tcolorbox}\hfill\begin{tcolorbox}[enhanced,nobeforeafter,width=7.1cm,height=2.1cm,valign=center,colback=cint!6,colframe=cint,arc=2.5mm,boxrule=0.9pt,left=3mm]{\Large\bfseries\color{cint}d)}\quad{\Large $F(x)=-e^{-x}+2$}\end{tcolorbox}
\end{tcolorbox}
\newpage
\begin{tcolorbox}[enhanced,width=16cm,height=14cm,colback=white,colframe=cint,boxrule=2pt,arc=0mm,
left=6mm,right=6mm,top=0mm,bottom=4mm,sharp corners,
overlay={\fill[cint] (frame.north west) rectangle ([yshift=-1.25cm]frame.north east);
\node[anchor=west,text=white,font=\sffamily\bfseries\LARGE] at ([xshift=5mm,yshift=-0.62cm]frame.north west) {Question 56 B};
\node[anchor=east,text=white,font=\sffamily\itshape\large] at ([xshift=-5mm,yshift=-0.62cm]frame.north east) {Intégrales et équations différentielles};}]
\vspace*{1.75cm}
\begin{minipage}[t][6.9cm][c]{\linewidth}\LARGE\centering
\begin{minipage}{0.95\linewidth}\raggedright\linespread{1.25}\selectfont Quelle est la primitive $F$ de $f(x)=e^{4x}$ sur $\mathbb{R}$ telle que $F(0)=1$~?\end{minipage}
\end{minipage}

\vfill
\begin{tcolorbox}[enhanced,nobeforeafter,width=7.1cm,height=2.1cm,valign=center,colback=cint!6,colframe=cint,arc=2.5mm,boxrule=0.9pt,left=3mm]{\Large\bfseries\color{cint}a)}\quad{\Large $F(x)=\frac{1}{4}\,e^{4x}+\frac{3}{4}$}\end{tcolorbox}\hfill\begin{tcolorbox}[enhanced,nobeforeafter,width=7.1cm,height=2.1cm,valign=center,colback=cint!6,colframe=cint,arc=2.5mm,boxrule=0.9pt,left=3mm]{\Large\bfseries\color{cint}b)}\quad{\Large $F(x)=\frac{1}{4}\,e^{4x}+1$}\end{tcolorbox}\par\vspace{2.5mm}
\begin{tcolorbox}[enhanced,nobeforeafter,width=7.1cm,height=2.1cm,valign=center,colback=cint!6,colframe=cint,arc=2.5mm,boxrule=0.9pt,left=3mm]{\Large\bfseries\color{cint}c)}\quad{\Large $F(x)=\frac{1}{4}\,e^{4x}+\frac{5}{4}$}\end{tcolorbox}\hfill\begin{tcolorbox}[enhanced,nobeforeafter,width=7.1cm,height=2.1cm,valign=center,colback=cint!6,colframe=cint,arc=2.5mm,boxrule=0.9pt,left=3mm]{\Large\bfseries\color{cint}d)}\quad{\Large $F(x)=4\,e^{4x}-3$}\end{tcolorbox}
\end{tcolorbox}
\newpage
\begin{tcolorbox}[enhanced,width=16cm,height=14cm,colback=white,colframe=cint,boxrule=2pt,arc=0mm,
left=6mm,right=6mm,top=0mm,bottom=4mm,sharp corners,
overlay={\fill[cint] (frame.north west) rectangle ([yshift=-1.25cm]frame.north east);
\node[anchor=west,text=white,font=\sffamily\bfseries\LARGE] at ([xshift=5mm,yshift=-0.62cm]frame.north west) {Question 57 B};
\node[anchor=east,text=white,font=\sffamily\itshape\large] at ([xshift=-5mm,yshift=-0.62cm]frame.north east) {Intégrales et équations différentielles};}]
\vspace*{1.75cm}
\begin{minipage}[t][6.9cm][c]{\linewidth}\LARGE\centering
\begin{minipage}{0.95\linewidth}\raggedright\linespread{1.25}\selectfont Quelle est la valeur moyenne de la fonction $x\mapsto x^2$ sur l'intervalle $[0;2]$~?\end{minipage}
\end{minipage}

\vfill
\begin{tcolorbox}[enhanced,nobeforeafter,width=7.1cm,height=2.1cm,valign=center,colback=cint!6,colframe=cint,arc=2.5mm,boxrule=0.9pt,left=3mm]{\Large\bfseries\color{cint}a)}\quad{\Large $4$}\end{tcolorbox}\hfill\begin{tcolorbox}[enhanced,nobeforeafter,width=7.1cm,height=2.1cm,valign=center,colback=cint!6,colframe=cint,arc=2.5mm,boxrule=0.9pt,left=3mm]{\Large\bfseries\color{cint}b)}\quad{\Large $2$}\end{tcolorbox}\par\vspace{2.5mm}
\begin{tcolorbox}[enhanced,nobeforeafter,width=7.1cm,height=2.1cm,valign=center,colback=cint!6,colframe=cint,arc=2.5mm,boxrule=0.9pt,left=3mm]{\Large\bfseries\color{cint}c)}\quad{\Large $\dfrac{4}{3}$}\end{tcolorbox}\hfill\begin{tcolorbox}[enhanced,nobeforeafter,width=7.1cm,height=2.1cm,valign=center,colback=cint!6,colframe=cint,arc=2.5mm,boxrule=0.9pt,left=3mm]{\Large\bfseries\color{cint}d)}\quad{\Large $\dfrac{8}{3}$}\end{tcolorbox}
\end{tcolorbox}
\newpage
\begin{tcolorbox}[enhanced,width=16cm,height=14cm,colback=white,colframe=cint,boxrule=2pt,arc=0mm,
left=6mm,right=6mm,top=0mm,bottom=4mm,sharp corners,
overlay={\fill[cint] (frame.north west) rectangle ([yshift=-1.25cm]frame.north east);
\node[anchor=west,text=white,font=\sffamily\bfseries\LARGE] at ([xshift=5mm,yshift=-0.62cm]frame.north west) {Question 58 B};
\node[anchor=east,text=white,font=\sffamily\itshape\large] at ([xshift=-5mm,yshift=-0.62cm]frame.north east) {Intégrales et équations différentielles};}]
\vspace*{1.75cm}
\begin{minipage}[t][6.9cm][c]{\linewidth}\LARGE\centering
\begin{minipage}{0.95\linewidth}\raggedright\linespread{1.25}\selectfont Quelle est la valeur moyenne de la fonction $x\mapsto x^2$ sur l'intervalle $[0;1]$~?\end{minipage}
\end{minipage}

\vfill
\begin{tcolorbox}[enhanced,nobeforeafter,width=7.1cm,height=2.1cm,valign=center,colback=cint!6,colframe=cint,arc=2.5mm,boxrule=0.9pt,left=3mm]{\Large\bfseries\color{cint}a)}\quad{\Large $\dfrac{1}{6}$}\end{tcolorbox}\hfill\begin{tcolorbox}[enhanced,nobeforeafter,width=7.1cm,height=2.1cm,valign=center,colback=cint!6,colframe=cint,arc=2.5mm,boxrule=0.9pt,left=3mm]{\Large\bfseries\color{cint}b)}\quad{\Large $\dfrac{1}{3}$}\end{tcolorbox}\par\vspace{2.5mm}
\begin{tcolorbox}[enhanced,nobeforeafter,width=7.1cm,height=2.1cm,valign=center,colback=cint!6,colframe=cint,arc=2.5mm,boxrule=0.9pt,left=3mm]{\Large\bfseries\color{cint}c)}\quad{\Large $\dfrac{1}{2}$}\end{tcolorbox}\hfill\begin{tcolorbox}[enhanced,nobeforeafter,width=7.1cm,height=2.1cm,valign=center,colback=cint!6,colframe=cint,arc=2.5mm,boxrule=0.9pt,left=3mm]{\Large\bfseries\color{cint}d)}\quad{\Large $1$}\end{tcolorbox}
\end{tcolorbox}
\newpage
\begin{tcolorbox}[enhanced,width=16cm,height=14cm,colback=white,colframe=cint,boxrule=2pt,arc=0mm,
left=6mm,right=6mm,top=0mm,bottom=4mm,sharp corners,
overlay={\fill[cint] (frame.north west) rectangle ([yshift=-1.25cm]frame.north east);
\node[anchor=west,text=white,font=\sffamily\bfseries\LARGE] at ([xshift=5mm,yshift=-0.62cm]frame.north west) {Question 59 B};
\node[anchor=east,text=white,font=\sffamily\itshape\large] at ([xshift=-5mm,yshift=-0.62cm]frame.north east) {Intégrales et équations différentielles};}]
\vspace*{1.75cm}
\begin{minipage}[t][6.9cm][c]{\linewidth}\LARGE\centering
\begin{minipage}{0.95\linewidth}\raggedright\linespread{1.25}\selectfont Quelle est la solution constante de l'équation différentielle $y'=-y+4$~?\end{minipage}
\end{minipage}

\vfill
\begin{tcolorbox}[enhanced,nobeforeafter,width=7.1cm,height=2.1cm,valign=center,colback=cint!6,colframe=cint,arc=2.5mm,boxrule=0.9pt,left=3mm]{\Large\bfseries\color{cint}a)}\quad{\Large $y=-4$}\end{tcolorbox}\hfill\begin{tcolorbox}[enhanced,nobeforeafter,width=7.1cm,height=2.1cm,valign=center,colback=cint!6,colframe=cint,arc=2.5mm,boxrule=0.9pt,left=3mm]{\Large\bfseries\color{cint}b)}\quad{\Large $y=4$}\end{tcolorbox}\par\vspace{2.5mm}
\begin{tcolorbox}[enhanced,nobeforeafter,width=7.1cm,height=2.1cm,valign=center,colback=cint!6,colframe=cint,arc=2.5mm,boxrule=0.9pt,left=3mm]{\Large\bfseries\color{cint}c)}\quad{\Large $y=\dfrac{1}{4}$}\end{tcolorbox}\hfill\begin{tcolorbox}[enhanced,nobeforeafter,width=7.1cm,height=2.1cm,valign=center,colback=cint!6,colframe=cint,arc=2.5mm,boxrule=0.9pt,left=3mm]{\Large\bfseries\color{cint}d)}\quad{\Large $y=0$}\end{tcolorbox}
\end{tcolorbox}
\newpage
\begin{tcolorbox}[enhanced,width=16cm,height=14cm,colback=white,colframe=cint,boxrule=2pt,arc=0mm,
left=6mm,right=6mm,top=0mm,bottom=4mm,sharp corners,
overlay={\fill[cint] (frame.north west) rectangle ([yshift=-1.25cm]frame.north east);
\node[anchor=west,text=white,font=\sffamily\bfseries\LARGE] at ([xshift=5mm,yshift=-0.62cm]frame.north west) {Question 60 B};
\node[anchor=east,text=white,font=\sffamily\itshape\large] at ([xshift=-5mm,yshift=-0.62cm]frame.north east) {Intégrales et équations différentielles};}]
\vspace*{1.75cm}
\begin{minipage}[t][6.9cm][c]{\linewidth}\LARGE\centering
\begin{minipage}{0.95\linewidth}\raggedright\linespread{1.25}\selectfont Quelle est la solution constante de l'équation différentielle $y'=5y+10$~?\end{minipage}
\end{minipage}

\vfill
\begin{tcolorbox}[enhanced,nobeforeafter,width=7.1cm,height=2.1cm,valign=center,colback=cint!6,colframe=cint,arc=2.5mm,boxrule=0.9pt,left=3mm]{\Large\bfseries\color{cint}a)}\quad{\Large $y=-2$}\end{tcolorbox}\hfill\begin{tcolorbox}[enhanced,nobeforeafter,width=7.1cm,height=2.1cm,valign=center,colback=cint!6,colframe=cint,arc=2.5mm,boxrule=0.9pt,left=3mm]{\Large\bfseries\color{cint}b)}\quad{\Large $y=-\dfrac{1}{2}$}\end{tcolorbox}\par\vspace{2.5mm}
\begin{tcolorbox}[enhanced,nobeforeafter,width=7.1cm,height=2.1cm,valign=center,colback=cint!6,colframe=cint,arc=2.5mm,boxrule=0.9pt,left=3mm]{\Large\bfseries\color{cint}c)}\quad{\Large $y=10$}\end{tcolorbox}\hfill\begin{tcolorbox}[enhanced,nobeforeafter,width=7.1cm,height=2.1cm,valign=center,colback=cint!6,colframe=cint,arc=2.5mm,boxrule=0.9pt,left=3mm]{\Large\bfseries\color{cint}d)}\quad{\Large $y=2$}\end{tcolorbox}
\end{tcolorbox}
\newpage
\begin{tcolorbox}[enhanced,width=16cm,height=14cm,colback=white,colframe=cint,boxrule=2pt,arc=0mm,
left=6mm,right=6mm,top=0mm,bottom=4mm,sharp corners,
overlay={\fill[cint] (frame.north west) rectangle ([yshift=-1.25cm]frame.north east);
\node[anchor=west,text=white,font=\sffamily\bfseries\LARGE] at ([xshift=5mm,yshift=-0.62cm]frame.north west) {Question 61 B};
\node[anchor=east,text=white,font=\sffamily\itshape\large] at ([xshift=-5mm,yshift=-0.62cm]frame.north east) {Intégrales et équations différentielles};}]
\vspace*{1.75cm}
\begin{minipage}[t][6.9cm][c]{\linewidth}\LARGE\centering
\begin{minipage}{0.95\linewidth}\raggedright\linespread{1.25}\selectfont $f$ est la solution de $y'=3y$ telle que $f(0)=1$. Que vaut $f(1)$~?\end{minipage}
\end{minipage}

\vfill
\begin{tcolorbox}[enhanced,nobeforeafter,width=7.1cm,height=2.1cm,valign=center,colback=cint!6,colframe=cint,arc=2.5mm,boxrule=0.9pt,left=3mm]{\Large\bfseries\color{cint}a)}\quad{\Large $3e$}\end{tcolorbox}\hfill\begin{tcolorbox}[enhanced,nobeforeafter,width=7.1cm,height=2.1cm,valign=center,colback=cint!6,colframe=cint,arc=2.5mm,boxrule=0.9pt,left=3mm]{\Large\bfseries\color{cint}b)}\quad{\Large $1+e^{3}$}\end{tcolorbox}\par\vspace{2.5mm}
\begin{tcolorbox}[enhanced,nobeforeafter,width=7.1cm,height=2.1cm,valign=center,colback=cint!6,colframe=cint,arc=2.5mm,boxrule=0.9pt,left=3mm]{\Large\bfseries\color{cint}c)}\quad{\Large $e^{3}$}\end{tcolorbox}\hfill\begin{tcolorbox}[enhanced,nobeforeafter,width=7.1cm,height=2.1cm,valign=center,colback=cint!6,colframe=cint,arc=2.5mm,boxrule=0.9pt,left=3mm]{\Large\bfseries\color{cint}d)}\quad{\Large $e^{-3}$}\end{tcolorbox}
\end{tcolorbox}
\newpage
\begin{tcolorbox}[enhanced,width=16cm,height=14cm,colback=white,colframe=cint,boxrule=2pt,arc=0mm,
left=6mm,right=6mm,top=0mm,bottom=4mm,sharp corners,
overlay={\fill[cint] (frame.north west) rectangle ([yshift=-1.25cm]frame.north east);
\node[anchor=west,text=white,font=\sffamily\bfseries\LARGE] at ([xshift=5mm,yshift=-0.62cm]frame.north west) {Question 62 B};
\node[anchor=east,text=white,font=\sffamily\itshape\large] at ([xshift=-5mm,yshift=-0.62cm]frame.north east) {Intégrales et équations différentielles};}]
\vspace*{1.75cm}
\begin{minipage}[t][6.9cm][c]{\linewidth}\LARGE\centering
\begin{minipage}{0.95\linewidth}\raggedright\linespread{1.25}\selectfont $f$ est la solution de $y'=y$ telle que $f(0)=4$. Que vaut $f(1)$~?\end{minipage}
\end{minipage}

\vfill
\begin{tcolorbox}[enhanced,nobeforeafter,width=7.1cm,height=2.1cm,valign=center,colback=cint!6,colframe=cint,arc=2.5mm,boxrule=0.9pt,left=3mm]{\Large\bfseries\color{cint}a)}\quad{\Large $4+e$}\end{tcolorbox}\hfill\begin{tcolorbox}[enhanced,nobeforeafter,width=7.1cm,height=2.1cm,valign=center,colback=cint!6,colframe=cint,arc=2.5mm,boxrule=0.9pt,left=3mm]{\Large\bfseries\color{cint}b)}\quad{\Large $e^{4}$}\end{tcolorbox}\par\vspace{2.5mm}
\begin{tcolorbox}[enhanced,nobeforeafter,width=7.1cm,height=2.1cm,valign=center,colback=cint!6,colframe=cint,arc=2.5mm,boxrule=0.9pt,left=3mm]{\Large\bfseries\color{cint}c)}\quad{\Large $4e$}\end{tcolorbox}\hfill\begin{tcolorbox}[enhanced,nobeforeafter,width=7.1cm,height=2.1cm,valign=center,colback=cint!6,colframe=cint,arc=2.5mm,boxrule=0.9pt,left=3mm]{\Large\bfseries\color{cint}d)}\quad{\Large $4e^{-1}$}\end{tcolorbox}
\end{tcolorbox}
\newpage
\begin{tcolorbox}[enhanced,width=16cm,height=14cm,colback=white,colframe=cint,boxrule=2pt,arc=0mm,
left=6mm,right=6mm,top=0mm,bottom=4mm,sharp corners,
overlay={\fill[cint] (frame.north west) rectangle ([yshift=-1.25cm]frame.north east);
\node[anchor=west,text=white,font=\sffamily\bfseries\LARGE] at ([xshift=5mm,yshift=-0.62cm]frame.north west) {Question 63 B};
\node[anchor=east,text=white,font=\sffamily\itshape\large] at ([xshift=-5mm,yshift=-0.62cm]frame.north east) {Intégrales et équations différentielles};}]
\vspace*{1.75cm}
\begin{minipage}[t][6.9cm][c]{\linewidth}\LARGE\centering
\begin{minipage}{0.95\linewidth}\raggedright\linespread{1.25}\selectfont Que vaut $\displaystyle\int_0^{3} x\,e^{x}\,\mathrm{d}x$~?\end{minipage}
\end{minipage}

\vfill
\begin{tcolorbox}[enhanced,nobeforeafter,width=7.1cm,height=2.1cm,valign=center,colback=cint!6,colframe=cint,arc=2.5mm,boxrule=0.9pt,left=3mm]{\Large\bfseries\color{cint}a)}\quad{\Large $4e^{3}-1$}\end{tcolorbox}\hfill\begin{tcolorbox}[enhanced,nobeforeafter,width=7.1cm,height=2.1cm,valign=center,colback=cint!6,colframe=cint,arc=2.5mm,boxrule=0.9pt,left=3mm]{\Large\bfseries\color{cint}b)}\quad{\Large $2e^{3}+1$}\end{tcolorbox}\par\vspace{2.5mm}
\begin{tcolorbox}[enhanced,nobeforeafter,width=7.1cm,height=2.1cm,valign=center,colback=cint!6,colframe=cint,arc=2.5mm,boxrule=0.9pt,left=3mm]{\Large\bfseries\color{cint}c)}\quad{\Large $2e^{3}$}\end{tcolorbox}\hfill\begin{tcolorbox}[enhanced,nobeforeafter,width=7.1cm,height=2.1cm,valign=center,colback=cint!6,colframe=cint,arc=2.5mm,boxrule=0.9pt,left=3mm]{\Large\bfseries\color{cint}d)}\quad{\Large $3e^{3}$}\end{tcolorbox}
\end{tcolorbox}
\newpage
\begin{tcolorbox}[enhanced,width=16cm,height=14cm,colback=white,colframe=cint,boxrule=2pt,arc=0mm,
left=6mm,right=6mm,top=0mm,bottom=4mm,sharp corners,
overlay={\fill[cint] (frame.north west) rectangle ([yshift=-1.25cm]frame.north east);
\node[anchor=west,text=white,font=\sffamily\bfseries\LARGE] at ([xshift=5mm,yshift=-0.62cm]frame.north west) {Question 64 B};
\node[anchor=east,text=white,font=\sffamily\itshape\large] at ([xshift=-5mm,yshift=-0.62cm]frame.north east) {Intégrales et équations différentielles};}]
\vspace*{1.75cm}
\begin{minipage}[t][6.9cm][c]{\linewidth}\LARGE\centering
\begin{minipage}{0.95\linewidth}\raggedright\linespread{1.25}\selectfont Que vaut $\displaystyle\int_0^{4} x\,e^{x}\,\mathrm{d}x$~?\end{minipage}
\end{minipage}

\vfill
\begin{tcolorbox}[enhanced,nobeforeafter,width=7.1cm,height=2.1cm,valign=center,colback=cint!6,colframe=cint,arc=2.5mm,boxrule=0.9pt,left=3mm]{\Large\bfseries\color{cint}a)}\quad{\Large $5e^{4}-1$}\end{tcolorbox}\hfill\begin{tcolorbox}[enhanced,nobeforeafter,width=7.1cm,height=2.1cm,valign=center,colback=cint!6,colframe=cint,arc=2.5mm,boxrule=0.9pt,left=3mm]{\Large\bfseries\color{cint}b)}\quad{\Large $3e^{4}+1$}\end{tcolorbox}\par\vspace{2.5mm}
\begin{tcolorbox}[enhanced,nobeforeafter,width=7.1cm,height=2.1cm,valign=center,colback=cint!6,colframe=cint,arc=2.5mm,boxrule=0.9pt,left=3mm]{\Large\bfseries\color{cint}c)}\quad{\Large $3e^{4}$}\end{tcolorbox}\hfill\begin{tcolorbox}[enhanced,nobeforeafter,width=7.1cm,height=2.1cm,valign=center,colback=cint!6,colframe=cint,arc=2.5mm,boxrule=0.9pt,left=3mm]{\Large\bfseries\color{cint}d)}\quad{\Large $4e^{4}$}\end{tcolorbox}
\end{tcolorbox}
\newpage
\begin{tcolorbox}[enhanced,width=16cm,height=14cm,colback=white,colframe=cesp,boxrule=2pt,arc=0mm,
left=6mm,right=6mm,top=0mm,bottom=4mm,sharp corners,
overlay={\fill[cesp] (frame.north west) rectangle ([yshift=-1.25cm]frame.north east);
\node[anchor=west,text=white,font=\sffamily\bfseries\LARGE] at ([xshift=5mm,yshift=-0.62cm]frame.north west) {Question 65 B};
\node[anchor=east,text=white,font=\sffamily\itshape\large] at ([xshift=-5mm,yshift=-0.62cm]frame.north east) {Géométrie dans l'espace};}]
\vspace*{1.75cm}
\begin{minipage}[t][6.9cm][c]{\linewidth}\Large\centering
\begin{minipage}{0.95\linewidth}\raggedright\linespread{1.25}\selectfont Dans un repère orthonormé, que vaut le produit scalaire $\vec u\cdot\vec v$ pour $\vec u\begin{pmatrix}-2\\1\\3\end{pmatrix}$ et $\vec v\begin{pmatrix}4\\2\\-1\end{pmatrix}$~?\end{minipage}
\end{minipage}

\vfill
\begin{tcolorbox}[enhanced,nobeforeafter,width=7.1cm,height=2.1cm,valign=center,colback=cesp!6,colframe=cesp,arc=2.5mm,boxrule=0.9pt,left=3mm]{\Large\bfseries\color{cesp}a)}\quad{\Large $9$}\end{tcolorbox}\hfill\begin{tcolorbox}[enhanced,nobeforeafter,width=7.1cm,height=2.1cm,valign=center,colback=cesp!6,colframe=cesp,arc=2.5mm,boxrule=0.9pt,left=3mm]{\Large\bfseries\color{cesp}b)}\quad{\Large $10$}\end{tcolorbox}\par\vspace{2.5mm}
\begin{tcolorbox}[enhanced,nobeforeafter,width=7.1cm,height=2.1cm,valign=center,colback=cesp!6,colframe=cesp,arc=2.5mm,boxrule=0.9pt,left=3mm]{\Large\bfseries\color{cesp}c)}\quad{\Large $-7$}\end{tcolorbox}\hfill\begin{tcolorbox}[enhanced,nobeforeafter,width=7.1cm,height=2.1cm,valign=center,colback=cesp!6,colframe=cesp,arc=2.5mm,boxrule=0.9pt,left=3mm]{\Large\bfseries\color{cesp}d)}\quad{\Large $-9$}\end{tcolorbox}
\end{tcolorbox}
\newpage
\begin{tcolorbox}[enhanced,width=16cm,height=14cm,colback=white,colframe=cesp,boxrule=2pt,arc=0mm,
left=6mm,right=6mm,top=0mm,bottom=4mm,sharp corners,
overlay={\fill[cesp] (frame.north west) rectangle ([yshift=-1.25cm]frame.north east);
\node[anchor=west,text=white,font=\sffamily\bfseries\LARGE] at ([xshift=5mm,yshift=-0.62cm]frame.north west) {Question 66 B};
\node[anchor=east,text=white,font=\sffamily\itshape\large] at ([xshift=-5mm,yshift=-0.62cm]frame.north east) {Géométrie dans l'espace};}]
\vspace*{1.75cm}
\begin{minipage}[t][6.9cm][c]{\linewidth}\Large\centering
\begin{minipage}{0.95\linewidth}\raggedright\linespread{1.25}\selectfont Dans un repère orthonormé, que vaut le produit scalaire $\vec u\cdot\vec v$ pour $\vec u\begin{pmatrix}3\\1\\-2\end{pmatrix}$ et $\vec v\begin{pmatrix}2\\-5\\1\end{pmatrix}$~?\end{minipage}
\end{minipage}

\vfill
\begin{tcolorbox}[enhanced,nobeforeafter,width=7.1cm,height=2.1cm,valign=center,colback=cesp!6,colframe=cesp,arc=2.5mm,boxrule=0.9pt,left=3mm]{\Large\bfseries\color{cesp}a)}\quad{\Large $-4$}\end{tcolorbox}\hfill\begin{tcolorbox}[enhanced,nobeforeafter,width=7.1cm,height=2.1cm,valign=center,colback=cesp!6,colframe=cesp,arc=2.5mm,boxrule=0.9pt,left=3mm]{\Large\bfseries\color{cesp}b)}\quad{\Large $-1$}\end{tcolorbox}\par\vspace{2.5mm}
\begin{tcolorbox}[enhanced,nobeforeafter,width=7.1cm,height=2.1cm,valign=center,colback=cesp!6,colframe=cesp,arc=2.5mm,boxrule=0.9pt,left=3mm]{\Large\bfseries\color{cesp}c)}\quad{\Large $0$}\end{tcolorbox}\hfill\begin{tcolorbox}[enhanced,nobeforeafter,width=7.1cm,height=2.1cm,valign=center,colback=cesp!6,colframe=cesp,arc=2.5mm,boxrule=0.9pt,left=3mm]{\Large\bfseries\color{cesp}d)}\quad{\Large $1$}\end{tcolorbox}
\end{tcolorbox}
\newpage
\begin{tcolorbox}[enhanced,width=16cm,height=14cm,colback=white,colframe=cesp,boxrule=2pt,arc=0mm,
left=6mm,right=6mm,top=0mm,bottom=4mm,sharp corners,
overlay={\fill[cesp] (frame.north west) rectangle ([yshift=-1.25cm]frame.north east);
\node[anchor=west,text=white,font=\sffamily\bfseries\LARGE] at ([xshift=5mm,yshift=-0.62cm]frame.north west) {Question 67 B};
\node[anchor=east,text=white,font=\sffamily\itshape\large] at ([xshift=-5mm,yshift=-0.62cm]frame.north east) {Géométrie dans l'espace};}]
\vspace*{1.75cm}
\begin{minipage}[t][6.9cm][c]{\linewidth}\LARGE\centering
\begin{minipage}{0.95\linewidth}\raggedright\linespread{1.25}\selectfont Dans un repère orthonormé, quelle est la norme du vecteur $\vec u\begin{pmatrix}1\\1\\1\end{pmatrix}$~?\end{minipage}
\end{minipage}

\vfill
\begin{tcolorbox}[enhanced,nobeforeafter,width=7.1cm,height=2.1cm,valign=center,colback=cesp!6,colframe=cesp,arc=2.5mm,boxrule=0.9pt,left=3mm]{\Large\bfseries\color{cesp}a)}\quad{\Large $3$}\end{tcolorbox}\hfill\begin{tcolorbox}[enhanced,nobeforeafter,width=7.1cm,height=2.1cm,valign=center,colback=cesp!6,colframe=cesp,arc=2.5mm,boxrule=0.9pt,left=3mm]{\Large\bfseries\color{cesp}b)}\quad{\Large $\sqrt{3}$}\end{tcolorbox}\par\vspace{2.5mm}
\begin{tcolorbox}[enhanced,nobeforeafter,width=7.1cm,height=2.1cm,valign=center,colback=cesp!6,colframe=cesp,arc=2.5mm,boxrule=0.9pt,left=3mm]{\Large\bfseries\color{cesp}c)}\quad{\Large $2$}\end{tcolorbox}\hfill\begin{tcolorbox}[enhanced,nobeforeafter,width=7.1cm,height=2.1cm,valign=center,colback=cesp!6,colframe=cesp,arc=2.5mm,boxrule=0.9pt,left=3mm]{\Large\bfseries\color{cesp}d)}\quad{\Large $2\sqrt{3}$}\end{tcolorbox}
\end{tcolorbox}
\newpage
\begin{tcolorbox}[enhanced,width=16cm,height=14cm,colback=white,colframe=cesp,boxrule=2pt,arc=0mm,
left=6mm,right=6mm,top=0mm,bottom=4mm,sharp corners,
overlay={\fill[cesp] (frame.north west) rectangle ([yshift=-1.25cm]frame.north east);
\node[anchor=west,text=white,font=\sffamily\bfseries\LARGE] at ([xshift=5mm,yshift=-0.62cm]frame.north west) {Question 68 B};
\node[anchor=east,text=white,font=\sffamily\itshape\large] at ([xshift=-5mm,yshift=-0.62cm]frame.north east) {Géométrie dans l'espace};}]
\vspace*{1.75cm}
\begin{minipage}[t][6.9cm][c]{\linewidth}\LARGE\centering
\begin{minipage}{0.95\linewidth}\raggedright\linespread{1.25}\selectfont Dans un repère orthonormé, quelle est la norme du vecteur $\vec u\begin{pmatrix}2\\-1\\2\end{pmatrix}$~?\end{minipage}
\end{minipage}

\vfill
\begin{tcolorbox}[enhanced,nobeforeafter,width=7.1cm,height=2.1cm,valign=center,colback=cesp!6,colframe=cesp,arc=2.5mm,boxrule=0.9pt,left=3mm]{\Large\bfseries\color{cesp}a)}\quad{\Large $9$}\end{tcolorbox}\hfill\begin{tcolorbox}[enhanced,nobeforeafter,width=7.1cm,height=2.1cm,valign=center,colback=cesp!6,colframe=cesp,arc=2.5mm,boxrule=0.9pt,left=3mm]{\Large\bfseries\color{cesp}b)}\quad{\Large $3$}\end{tcolorbox}\par\vspace{2.5mm}
\begin{tcolorbox}[enhanced,nobeforeafter,width=7.1cm,height=2.1cm,valign=center,colback=cesp!6,colframe=cesp,arc=2.5mm,boxrule=0.9pt,left=3mm]{\Large\bfseries\color{cesp}c)}\quad{\Large $\sqrt{10}$}\end{tcolorbox}\hfill\begin{tcolorbox}[enhanced,nobeforeafter,width=7.1cm,height=2.1cm,valign=center,colback=cesp!6,colframe=cesp,arc=2.5mm,boxrule=0.9pt,left=3mm]{\Large\bfseries\color{cesp}d)}\quad{\Large $5$}\end{tcolorbox}
\end{tcolorbox}
\newpage
\begin{tcolorbox}[enhanced,width=16cm,height=14cm,colback=white,colframe=cesp,boxrule=2pt,arc=0mm,
left=6mm,right=6mm,top=0mm,bottom=4mm,sharp corners,
overlay={\fill[cesp] (frame.north west) rectangle ([yshift=-1.25cm]frame.north east);
\node[anchor=west,text=white,font=\sffamily\bfseries\LARGE] at ([xshift=5mm,yshift=-0.62cm]frame.north west) {Question 69 B};
\node[anchor=east,text=white,font=\sffamily\itshape\large] at ([xshift=-5mm,yshift=-0.62cm]frame.north east) {Géométrie dans l'espace};}]
\vspace*{1.75cm}
\begin{minipage}[t][6.9cm][c]{\linewidth}\Large\centering
\begin{minipage}{0.95\linewidth}\raggedright\linespread{1.25}\selectfont Pour quelle valeur de $m$ les vecteurs $\vec u\begin{pmatrix}1\\m\\2\end{pmatrix}$ et $\vec v\begin{pmatrix}4\\-1\\1\end{pmatrix}$ sont-ils orthogonaux~?\end{minipage}
\end{minipage}

\vfill
\begin{tcolorbox}[enhanced,nobeforeafter,width=7.1cm,height=2.1cm,valign=center,colback=cesp!6,colframe=cesp,arc=2.5mm,boxrule=0.9pt,left=3mm]{\Large\bfseries\color{cesp}a)}\quad{\Large $m=5$}\end{tcolorbox}\hfill\begin{tcolorbox}[enhanced,nobeforeafter,width=7.1cm,height=2.1cm,valign=center,colback=cesp!6,colframe=cesp,arc=2.5mm,boxrule=0.9pt,left=3mm]{\Large\bfseries\color{cesp}b)}\quad{\Large $m=7$}\end{tcolorbox}\par\vspace{2.5mm}
\begin{tcolorbox}[enhanced,nobeforeafter,width=7.1cm,height=2.1cm,valign=center,colback=cesp!6,colframe=cesp,arc=2.5mm,boxrule=0.9pt,left=3mm]{\Large\bfseries\color{cesp}c)}\quad{\Large $m=6$}\end{tcolorbox}\hfill\begin{tcolorbox}[enhanced,nobeforeafter,width=7.1cm,height=2.1cm,valign=center,colback=cesp!6,colframe=cesp,arc=2.5mm,boxrule=0.9pt,left=3mm]{\Large\bfseries\color{cesp}d)}\quad{\Large $m=-6$}\end{tcolorbox}
\end{tcolorbox}
\newpage
\begin{tcolorbox}[enhanced,width=16cm,height=14cm,colback=white,colframe=cesp,boxrule=2pt,arc=0mm,
left=6mm,right=6mm,top=0mm,bottom=4mm,sharp corners,
overlay={\fill[cesp] (frame.north west) rectangle ([yshift=-1.25cm]frame.north east);
\node[anchor=west,text=white,font=\sffamily\bfseries\LARGE] at ([xshift=5mm,yshift=-0.62cm]frame.north west) {Question 70 B};
\node[anchor=east,text=white,font=\sffamily\itshape\large] at ([xshift=-5mm,yshift=-0.62cm]frame.north east) {Géométrie dans l'espace};}]
\vspace*{1.75cm}
\begin{minipage}[t][6.9cm][c]{\linewidth}\Large\centering
\begin{minipage}{0.95\linewidth}\raggedright\linespread{1.25}\selectfont Pour quelle valeur de $m$ les vecteurs $\vec u\begin{pmatrix}1\\m\\2\end{pmatrix}$ et $\vec v\begin{pmatrix}3\\3\\-6\end{pmatrix}$ sont-ils orthogonaux~?\end{minipage}
\end{minipage}

\vfill
\begin{tcolorbox}[enhanced,nobeforeafter,width=7.1cm,height=2.1cm,valign=center,colback=cesp!6,colframe=cesp,arc=2.5mm,boxrule=0.9pt,left=3mm]{\Large\bfseries\color{cesp}a)}\quad{\Large $m=1$}\end{tcolorbox}\hfill\begin{tcolorbox}[enhanced,nobeforeafter,width=7.1cm,height=2.1cm,valign=center,colback=cesp!6,colframe=cesp,arc=2.5mm,boxrule=0.9pt,left=3mm]{\Large\bfseries\color{cesp}b)}\quad{\Large $m=-3$}\end{tcolorbox}\par\vspace{2.5mm}
\begin{tcolorbox}[enhanced,nobeforeafter,width=7.1cm,height=2.1cm,valign=center,colback=cesp!6,colframe=cesp,arc=2.5mm,boxrule=0.9pt,left=3mm]{\Large\bfseries\color{cesp}c)}\quad{\Large $m=4$}\end{tcolorbox}\hfill\begin{tcolorbox}[enhanced,nobeforeafter,width=7.1cm,height=2.1cm,valign=center,colback=cesp!6,colframe=cesp,arc=2.5mm,boxrule=0.9pt,left=3mm]{\Large\bfseries\color{cesp}d)}\quad{\Large $m=3$}\end{tcolorbox}
\end{tcolorbox}
\newpage
\begin{tcolorbox}[enhanced,width=16cm,height=14cm,colback=white,colframe=cesp,boxrule=2pt,arc=0mm,
left=6mm,right=6mm,top=0mm,bottom=4mm,sharp corners,
overlay={\fill[cesp] (frame.north west) rectangle ([yshift=-1.25cm]frame.north east);
\node[anchor=west,text=white,font=\sffamily\bfseries\LARGE] at ([xshift=5mm,yshift=-0.62cm]frame.north west) {Question 71 B};
\node[anchor=east,text=white,font=\sffamily\itshape\large] at ([xshift=-5mm,yshift=-0.62cm]frame.north east) {Géométrie dans l'espace};}]
\vspace*{1.75cm}
\begin{minipage}[t][6.9cm][c]{\linewidth}\Large\centering
\begin{minipage}{0.95\linewidth}\raggedright\linespread{1.25}\selectfont Quelle est une équation du plan passant par $A(2;-1;3)$ et de vecteur normal $\vec n\begin{pmatrix}1\\1\\1\end{pmatrix}$~?\end{minipage}
\end{minipage}

\vfill
\begin{tcolorbox}[enhanced,nobeforeafter,width=7.1cm,height=2.1cm,valign=center,colback=cesp!6,colframe=cesp,arc=2.5mm,boxrule=0.9pt,left=3mm]{\Large\bfseries\color{cesp}a)}\quad{\Large $2x-y+3z-4=0$}\end{tcolorbox}\hfill\begin{tcolorbox}[enhanced,nobeforeafter,width=7.1cm,height=2.1cm,valign=center,colback=cesp!6,colframe=cesp,arc=2.5mm,boxrule=0.9pt,left=3mm]{\Large\bfseries\color{cesp}b)}\quad{\Large $x+y+z-3=0$}\end{tcolorbox}\par\vspace{2.5mm}
\begin{tcolorbox}[enhanced,nobeforeafter,width=7.1cm,height=2.1cm,valign=center,colback=cesp!6,colframe=cesp,arc=2.5mm,boxrule=0.9pt,left=3mm]{\Large\bfseries\color{cesp}c)}\quad{\Large $x+y+z-4=0$}\end{tcolorbox}\hfill\begin{tcolorbox}[enhanced,nobeforeafter,width=7.1cm,height=2.1cm,valign=center,colback=cesp!6,colframe=cesp,arc=2.5mm,boxrule=0.9pt,left=3mm]{\Large\bfseries\color{cesp}d)}\quad{\Large $x+y+z+4=0$}\end{tcolorbox}
\end{tcolorbox}
\newpage
\begin{tcolorbox}[enhanced,width=16cm,height=14cm,colback=white,colframe=cesp,boxrule=2pt,arc=0mm,
left=6mm,right=6mm,top=0mm,bottom=4mm,sharp corners,
overlay={\fill[cesp] (frame.north west) rectangle ([yshift=-1.25cm]frame.north east);
\node[anchor=west,text=white,font=\sffamily\bfseries\LARGE] at ([xshift=5mm,yshift=-0.62cm]frame.north west) {Question 72 B};
\node[anchor=east,text=white,font=\sffamily\itshape\large] at ([xshift=-5mm,yshift=-0.62cm]frame.north east) {Géométrie dans l'espace};}]
\vspace*{1.75cm}
\begin{minipage}[t][6.9cm][c]{\linewidth}\Large\centering
\begin{minipage}{0.95\linewidth}\raggedright\linespread{1.25}\selectfont Quelle est une équation du plan passant par $A(1;4;1)$ et de vecteur normal $\vec n\begin{pmatrix}3\\0\\-2\end{pmatrix}$~?\end{minipage}
\end{minipage}

\vfill
\begin{tcolorbox}[enhanced,nobeforeafter,width=7.1cm,height=2.1cm,valign=center,colback=cesp!6,colframe=cesp,arc=2.5mm,boxrule=0.9pt,left=3mm]{\Large\bfseries\color{cesp}a)}\quad{\Large $3x-2z-1=0$}\end{tcolorbox}\hfill\begin{tcolorbox}[enhanced,nobeforeafter,width=7.1cm,height=2.1cm,valign=center,colback=cesp!6,colframe=cesp,arc=2.5mm,boxrule=0.9pt,left=3mm]{\Large\bfseries\color{cesp}b)}\quad{\Large $x+4y+z-1=0$}\end{tcolorbox}\par\vspace{2.5mm}
\begin{tcolorbox}[enhanced,nobeforeafter,width=7.1cm,height=2.1cm,valign=center,colback=cesp!6,colframe=cesp,arc=2.5mm,boxrule=0.9pt,left=3mm]{\Large\bfseries\color{cesp}c)}\quad{\Large $3x-2z=0$}\end{tcolorbox}\hfill\begin{tcolorbox}[enhanced,nobeforeafter,width=7.1cm,height=2.1cm,valign=center,colback=cesp!6,colframe=cesp,arc=2.5mm,boxrule=0.9pt,left=3mm]{\Large\bfseries\color{cesp}d)}\quad{\Large $3x-2z+1=0$}\end{tcolorbox}
\end{tcolorbox}
\newpage
\begin{tcolorbox}[enhanced,width=16cm,height=14cm,colback=white,colframe=cesp,boxrule=2pt,arc=0mm,
left=6mm,right=6mm,top=0mm,bottom=4mm,sharp corners,
overlay={\fill[cesp] (frame.north west) rectangle ([yshift=-1.25cm]frame.north east);
\node[anchor=west,text=white,font=\sffamily\bfseries\LARGE] at ([xshift=5mm,yshift=-0.62cm]frame.north west) {Question 73 B};
\node[anchor=east,text=white,font=\sffamily\itshape\large] at ([xshift=-5mm,yshift=-0.62cm]frame.north east) {Géométrie dans l'espace};}]
\vspace*{1.75cm}
\begin{minipage}[t][6.9cm][c]{\linewidth}\Large\centering
\begin{minipage}{0.95\linewidth}\raggedright\linespread{1.25}\selectfont Lequel de ces points appartient à la droite de représentation paramétrique $\left\{\begin{array}{l}x=2-t\\y=2t\\z=-1+t\end{array}\right.$ ($t\in\mathbb{R}$)~?\end{minipage}
\end{minipage}

\vfill
\begin{tcolorbox}[enhanced,nobeforeafter,width=7.1cm,height=2.1cm,valign=center,colback=cesp!6,colframe=cesp,arc=2.5mm,boxrule=0.9pt,left=3mm]{\Large\bfseries\color{cesp}a)}\quad{\Large $(-1;6;3)$}\end{tcolorbox}\hfill\begin{tcolorbox}[enhanced,nobeforeafter,width=7.1cm,height=2.1cm,valign=center,colback=cesp!6,colframe=cesp,arc=2.5mm,boxrule=0.9pt,left=3mm]{\Large\bfseries\color{cesp}b)}\quad{\Large $(0;6;2)$}\end{tcolorbox}\par\vspace{2.5mm}
\begin{tcolorbox}[enhanced,nobeforeafter,width=7.1cm,height=2.1cm,valign=center,colback=cesp!6,colframe=cesp,arc=2.5mm,boxrule=0.9pt,left=3mm]{\Large\bfseries\color{cesp}c)}\quad{\Large $(-1;7;2)$}\end{tcolorbox}\hfill\begin{tcolorbox}[enhanced,nobeforeafter,width=7.1cm,height=2.1cm,valign=center,colback=cesp!6,colframe=cesp,arc=2.5mm,boxrule=0.9pt,left=3mm]{\Large\bfseries\color{cesp}d)}\quad{\Large $(-1;6;2)$}\end{tcolorbox}
\end{tcolorbox}
\newpage
\begin{tcolorbox}[enhanced,width=16cm,height=14cm,colback=white,colframe=cesp,boxrule=2pt,arc=0mm,
left=6mm,right=6mm,top=0mm,bottom=4mm,sharp corners,
overlay={\fill[cesp] (frame.north west) rectangle ([yshift=-1.25cm]frame.north east);
\node[anchor=west,text=white,font=\sffamily\bfseries\LARGE] at ([xshift=5mm,yshift=-0.62cm]frame.north west) {Question 74 B};
\node[anchor=east,text=white,font=\sffamily\itshape\large] at ([xshift=-5mm,yshift=-0.62cm]frame.north east) {Géométrie dans l'espace};}]
\vspace*{1.75cm}
\begin{minipage}[t][6.9cm][c]{\linewidth}\Large\centering
\begin{minipage}{0.95\linewidth}\raggedright\linespread{1.25}\selectfont Lequel de ces points appartient à la droite de représentation paramétrique $\left\{\begin{array}{l}x=-1+3t\\y=1+t\\z=2-2t\end{array}\right.$ ($t\in\mathbb{R}$)~?\end{minipage}
\end{minipage}

\vfill
\begin{tcolorbox}[enhanced,nobeforeafter,width=7.1cm,height=2.1cm,valign=center,colback=cesp!6,colframe=cesp,arc=2.5mm,boxrule=0.9pt,left=3mm]{\Large\bfseries\color{cesp}a)}\quad{\Large $(2;2;0)$}\end{tcolorbox}\hfill\begin{tcolorbox}[enhanced,nobeforeafter,width=7.1cm,height=2.1cm,valign=center,colback=cesp!6,colframe=cesp,arc=2.5mm,boxrule=0.9pt,left=3mm]{\Large\bfseries\color{cesp}b)}\quad{\Large $(2;3;0)$}\end{tcolorbox}\par\vspace{2.5mm}
\begin{tcolorbox}[enhanced,nobeforeafter,width=7.1cm,height=2.1cm,valign=center,colback=cesp!6,colframe=cesp,arc=2.5mm,boxrule=0.9pt,left=3mm]{\Large\bfseries\color{cesp}c)}\quad{\Large $(2;2;1)$}\end{tcolorbox}\hfill\begin{tcolorbox}[enhanced,nobeforeafter,width=7.1cm,height=2.1cm,valign=center,colback=cesp!6,colframe=cesp,arc=2.5mm,boxrule=0.9pt,left=3mm]{\Large\bfseries\color{cesp}d)}\quad{\Large $(3;2;0)$}\end{tcolorbox}
\end{tcolorbox}
\newpage
\begin{tcolorbox}[enhanced,width=16cm,height=14cm,colback=white,colframe=cesp,boxrule=2pt,arc=0mm,
left=6mm,right=6mm,top=0mm,bottom=4mm,sharp corners,
overlay={\fill[cesp] (frame.north west) rectangle ([yshift=-1.25cm]frame.north east);
\node[anchor=west,text=white,font=\sffamily\bfseries\LARGE] at ([xshift=5mm,yshift=-0.62cm]frame.north west) {Question 75 B};
\node[anchor=east,text=white,font=\sffamily\itshape\large] at ([xshift=-5mm,yshift=-0.62cm]frame.north east) {Géométrie dans l'espace};}]
\vspace*{1.75cm}
\begin{minipage}[t][6.9cm][c]{\linewidth}\Large\centering
\begin{minipage}{0.95\linewidth}\raggedright\linespread{1.25}\selectfont Pour quelle valeur de $m$ les vecteurs $\vec u\begin{pmatrix}-2\\m\\-4\end{pmatrix}$ et $\vec v\begin{pmatrix}1\\4\\2\end{pmatrix}$ sont-ils colinéaires~?\end{minipage}
\end{minipage}

\vfill
\begin{tcolorbox}[enhanced,nobeforeafter,width=7.1cm,height=2.1cm,valign=center,colback=cesp!6,colframe=cesp,arc=2.5mm,boxrule=0.9pt,left=3mm]{\Large\bfseries\color{cesp}a)}\quad{\Large $m=8$}\end{tcolorbox}\hfill\begin{tcolorbox}[enhanced,nobeforeafter,width=7.1cm,height=2.1cm,valign=center,colback=cesp!6,colframe=cesp,arc=2.5mm,boxrule=0.9pt,left=3mm]{\Large\bfseries\color{cesp}b)}\quad{\Large $m=-2$}\end{tcolorbox}\par\vspace{2.5mm}
\begin{tcolorbox}[enhanced,nobeforeafter,width=7.1cm,height=2.1cm,valign=center,colback=cesp!6,colframe=cesp,arc=2.5mm,boxrule=0.9pt,left=3mm]{\Large\bfseries\color{cesp}c)}\quad{\Large $m=2$}\end{tcolorbox}\hfill\begin{tcolorbox}[enhanced,nobeforeafter,width=7.1cm,height=2.1cm,valign=center,colback=cesp!6,colframe=cesp,arc=2.5mm,boxrule=0.9pt,left=3mm]{\Large\bfseries\color{cesp}d)}\quad{\Large $m=-8$}\end{tcolorbox}
\end{tcolorbox}
\newpage
\begin{tcolorbox}[enhanced,width=16cm,height=14cm,colback=white,colframe=cesp,boxrule=2pt,arc=0mm,
left=6mm,right=6mm,top=0mm,bottom=4mm,sharp corners,
overlay={\fill[cesp] (frame.north west) rectangle ([yshift=-1.25cm]frame.north east);
\node[anchor=west,text=white,font=\sffamily\bfseries\LARGE] at ([xshift=5mm,yshift=-0.62cm]frame.north west) {Question 76 B};
\node[anchor=east,text=white,font=\sffamily\itshape\large] at ([xshift=-5mm,yshift=-0.62cm]frame.north east) {Géométrie dans l'espace};}]
\vspace*{1.75cm}
\begin{minipage}[t][6.9cm][c]{\linewidth}\Large\centering
\begin{minipage}{0.95\linewidth}\raggedright\linespread{1.25}\selectfont Pour quelle valeur de $m$ les vecteurs $\vec u\begin{pmatrix}4\\m\\-4\end{pmatrix}$ et $\vec v\begin{pmatrix}1\\2\\-1\end{pmatrix}$ sont-ils colinéaires~?\end{minipage}
\end{minipage}

\vfill
\begin{tcolorbox}[enhanced,nobeforeafter,width=7.1cm,height=2.1cm,valign=center,colback=cesp!6,colframe=cesp,arc=2.5mm,boxrule=0.9pt,left=3mm]{\Large\bfseries\color{cesp}a)}\quad{\Large $m=6$}\end{tcolorbox}\hfill\begin{tcolorbox}[enhanced,nobeforeafter,width=7.1cm,height=2.1cm,valign=center,colback=cesp!6,colframe=cesp,arc=2.5mm,boxrule=0.9pt,left=3mm]{\Large\bfseries\color{cesp}b)}\quad{\Large $m=8$}\end{tcolorbox}\par\vspace{2.5mm}
\begin{tcolorbox}[enhanced,nobeforeafter,width=7.1cm,height=2.1cm,valign=center,colback=cesp!6,colframe=cesp,arc=2.5mm,boxrule=0.9pt,left=3mm]{\Large\bfseries\color{cesp}c)}\quad{\Large $m=\dfrac{1}{2}$}\end{tcolorbox}\hfill\begin{tcolorbox}[enhanced,nobeforeafter,width=7.1cm,height=2.1cm,valign=center,colback=cesp!6,colframe=cesp,arc=2.5mm,boxrule=0.9pt,left=3mm]{\Large\bfseries\color{cesp}d)}\quad{\Large $m=-8$}\end{tcolorbox}
\end{tcolorbox}
\newpage
\begin{tcolorbox}[enhanced,width=16cm,height=14cm,colback=white,colframe=cesp,boxrule=2pt,arc=0mm,
left=6mm,right=6mm,top=0mm,bottom=4mm,sharp corners,
overlay={\fill[cesp] (frame.north west) rectangle ([yshift=-1.25cm]frame.north east);
\node[anchor=west,text=white,font=\sffamily\bfseries\LARGE] at ([xshift=5mm,yshift=-0.62cm]frame.north west) {Question 77 B};
\node[anchor=east,text=white,font=\sffamily\itshape\large] at ([xshift=-5mm,yshift=-0.62cm]frame.north east) {Géométrie dans l'espace};}]
\vspace*{1.75cm}
\begin{minipage}[t][6.9cm][c]{\linewidth}\LARGE\centering
\begin{minipage}{0.95\linewidth}\raggedright\linespread{1.25}\selectfont Dans un repère orthonormé, quelle est la distance entre $A(2;0;1)$ et $B(5;4;1)$~?\end{minipage}
\end{minipage}

\vfill
\begin{tcolorbox}[enhanced,nobeforeafter,width=7.1cm,height=2.1cm,valign=center,colback=cesp!6,colframe=cesp,arc=2.5mm,boxrule=0.9pt,left=3mm]{\Large\bfseries\color{cesp}a)}\quad{\Large $5$}\end{tcolorbox}\hfill\begin{tcolorbox}[enhanced,nobeforeafter,width=7.1cm,height=2.1cm,valign=center,colback=cesp!6,colframe=cesp,arc=2.5mm,boxrule=0.9pt,left=3mm]{\Large\bfseries\color{cesp}b)}\quad{\Large $3\sqrt{3}$}\end{tcolorbox}\par\vspace{2.5mm}
\begin{tcolorbox}[enhanced,nobeforeafter,width=7.1cm,height=2.1cm,valign=center,colback=cesp!6,colframe=cesp,arc=2.5mm,boxrule=0.9pt,left=3mm]{\Large\bfseries\color{cesp}c)}\quad{\Large $7$}\end{tcolorbox}\hfill\begin{tcolorbox}[enhanced,nobeforeafter,width=7.1cm,height=2.1cm,valign=center,colback=cesp!6,colframe=cesp,arc=2.5mm,boxrule=0.9pt,left=3mm]{\Large\bfseries\color{cesp}d)}\quad{\Large $25$}\end{tcolorbox}
\end{tcolorbox}
\newpage
\begin{tcolorbox}[enhanced,width=16cm,height=14cm,colback=white,colframe=cesp,boxrule=2pt,arc=0mm,
left=6mm,right=6mm,top=0mm,bottom=4mm,sharp corners,
overlay={\fill[cesp] (frame.north west) rectangle ([yshift=-1.25cm]frame.north east);
\node[anchor=west,text=white,font=\sffamily\bfseries\LARGE] at ([xshift=5mm,yshift=-0.62cm]frame.north west) {Question 78 B};
\node[anchor=east,text=white,font=\sffamily\itshape\large] at ([xshift=-5mm,yshift=-0.62cm]frame.north east) {Géométrie dans l'espace};}]
\vspace*{1.75cm}
\begin{minipage}[t][6.9cm][c]{\linewidth}\LARGE\centering
\begin{minipage}{0.95\linewidth}\raggedright\linespread{1.25}\selectfont Dans un repère orthonormé, quelle est la distance entre $A(-1;1;0)$ et $B(1;2;2)$~?\end{minipage}
\end{minipage}

\vfill
\begin{tcolorbox}[enhanced,nobeforeafter,width=7.1cm,height=2.1cm,valign=center,colback=cesp!6,colframe=cesp,arc=2.5mm,boxrule=0.9pt,left=3mm]{\Large\bfseries\color{cesp}a)}\quad{\Large $\sqrt{11}$}\end{tcolorbox}\hfill\begin{tcolorbox}[enhanced,nobeforeafter,width=7.1cm,height=2.1cm,valign=center,colback=cesp!6,colframe=cesp,arc=2.5mm,boxrule=0.9pt,left=3mm]{\Large\bfseries\color{cesp}b)}\quad{\Large $9$}\end{tcolorbox}\par\vspace{2.5mm}
\begin{tcolorbox}[enhanced,nobeforeafter,width=7.1cm,height=2.1cm,valign=center,colback=cesp!6,colframe=cesp,arc=2.5mm,boxrule=0.9pt,left=3mm]{\Large\bfseries\color{cesp}c)}\quad{\Large $3$}\end{tcolorbox}\hfill\begin{tcolorbox}[enhanced,nobeforeafter,width=7.1cm,height=2.1cm,valign=center,colback=cesp!6,colframe=cesp,arc=2.5mm,boxrule=0.9pt,left=3mm]{\Large\bfseries\color{cesp}d)}\quad{\Large $5$}\end{tcolorbox}
\end{tcolorbox}
\newpage
\begin{tcolorbox}[enhanced,width=16cm,height=14cm,colback=white,colframe=cesp,boxrule=2pt,arc=0mm,
left=6mm,right=6mm,top=0mm,bottom=4mm,sharp corners,
overlay={\fill[cesp] (frame.north west) rectangle ([yshift=-1.25cm]frame.north east);
\node[anchor=west,text=white,font=\sffamily\bfseries\LARGE] at ([xshift=5mm,yshift=-0.62cm]frame.north west) {Question 79 B};
\node[anchor=east,text=white,font=\sffamily\itshape\large] at ([xshift=-5mm,yshift=-0.62cm]frame.north east) {Géométrie dans l'espace};}]
\vspace*{1.75cm}
\begin{minipage}[t][6.9cm][c]{\linewidth}\Large\centering
\begin{minipage}{0.95\linewidth}\raggedright\linespread{1.25}\selectfont Quel est le point d'intersection de la droite $\left\{\begin{array}{l}x=2+t\\y=-1+t\\z=-t\end{array}\right.$ ($t\in\mathbb{R}$) et du plan $x+y+z=3$~?\end{minipage}
\end{minipage}

\vfill
\begin{tcolorbox}[enhanced,nobeforeafter,width=7.1cm,height=2.1cm,valign=center,colback=cesp!6,colframe=cesp,arc=2.5mm,boxrule=0.9pt,left=3mm]{\Large\bfseries\color{cesp}a)}\quad{\Large $(3;0;-1)$}\end{tcolorbox}\hfill\begin{tcolorbox}[enhanced,nobeforeafter,width=7.1cm,height=2.1cm,valign=center,colback=cesp!6,colframe=cesp,arc=2.5mm,boxrule=0.9pt,left=3mm]{\Large\bfseries\color{cesp}b)}\quad{\Large $(4;1;-2)$}\end{tcolorbox}\par\vspace{2.5mm}
\begin{tcolorbox}[enhanced,nobeforeafter,width=7.1cm,height=2.1cm,valign=center,colback=cesp!6,colframe=cesp,arc=2.5mm,boxrule=0.9pt,left=3mm]{\Large\bfseries\color{cesp}c)}\quad{\Large $(2;-1;0)$}\end{tcolorbox}\hfill\begin{tcolorbox}[enhanced,nobeforeafter,width=7.1cm,height=2.1cm,valign=center,colback=cesp!6,colframe=cesp,arc=2.5mm,boxrule=0.9pt,left=3mm]{\Large\bfseries\color{cesp}d)}\quad{\Large $(5;2;-3)$}\end{tcolorbox}
\end{tcolorbox}
\newpage
\begin{tcolorbox}[enhanced,width=16cm,height=14cm,colback=white,colframe=cesp,boxrule=2pt,arc=0mm,
left=6mm,right=6mm,top=0mm,bottom=4mm,sharp corners,
overlay={\fill[cesp] (frame.north west) rectangle ([yshift=-1.25cm]frame.north east);
\node[anchor=west,text=white,font=\sffamily\bfseries\LARGE] at ([xshift=5mm,yshift=-0.62cm]frame.north west) {Question 80 B};
\node[anchor=east,text=white,font=\sffamily\itshape\large] at ([xshift=-5mm,yshift=-0.62cm]frame.north east) {Géométrie dans l'espace};}]
\vspace*{1.75cm}
\begin{minipage}[t][6.9cm][c]{\linewidth}\Large\centering
\begin{minipage}{0.95\linewidth}\raggedright\linespread{1.25}\selectfont Quel est le point d'intersection de la droite $\left\{\begin{array}{l}x=1+2t\\y=2+t\\z=3+t\end{array}\right.$ ($t\in\mathbb{R}$) et du plan $x+y+z=10$~?\end{minipage}
\end{minipage}

\vfill
\begin{tcolorbox}[enhanced,nobeforeafter,width=7.1cm,height=2.1cm,valign=center,colback=cesp!6,colframe=cesp,arc=2.5mm,boxrule=0.9pt,left=3mm]{\Large\bfseries\color{cesp}a)}\quad{\Large $(1;2;3)$}\end{tcolorbox}\hfill\begin{tcolorbox}[enhanced,nobeforeafter,width=7.1cm,height=2.1cm,valign=center,colback=cesp!6,colframe=cesp,arc=2.5mm,boxrule=0.9pt,left=3mm]{\Large\bfseries\color{cesp}b)}\quad{\Large $(3;3;4)$}\end{tcolorbox}\par\vspace{2.5mm}
\begin{tcolorbox}[enhanced,nobeforeafter,width=7.1cm,height=2.1cm,valign=center,colback=cesp!6,colframe=cesp,arc=2.5mm,boxrule=0.9pt,left=3mm]{\Large\bfseries\color{cesp}c)}\quad{\Large $(5;4;5)$}\end{tcolorbox}\hfill\begin{tcolorbox}[enhanced,nobeforeafter,width=7.1cm,height=2.1cm,valign=center,colback=cesp!6,colframe=cesp,arc=2.5mm,boxrule=0.9pt,left=3mm]{\Large\bfseries\color{cesp}d)}\quad{\Large $(-1;1;2)$}\end{tcolorbox}
\end{tcolorbox}
\newpage
\begin{tcolorbox}[enhanced,width=16cm,height=14cm,colback=white,colframe=cprb,boxrule=2pt,arc=0mm,
left=6mm,right=6mm,top=0mm,bottom=4mm,sharp corners,
overlay={\fill[cprb] (frame.north west) rectangle ([yshift=-1.25cm]frame.north east);
\node[anchor=west,text=white,font=\sffamily\bfseries\LARGE] at ([xshift=5mm,yshift=-0.62cm]frame.north west) {Question 81 B};
\node[anchor=east,text=white,font=\sffamily\itshape\large] at ([xshift=-5mm,yshift=-0.62cm]frame.north east) {Probabilités et dénombrement};}]
\vspace*{1.75cm}
\begin{minipage}[t][6.9cm][c]{\linewidth}\LARGE\centering
\begin{minipage}{0.95\linewidth}\raggedright\linespread{1.25}\selectfont $X$ suit la loi binomiale $\mathcal{B}\left(6;\frac12\right)$. Que vaut $P(X=3)$~?\end{minipage}
\end{minipage}

\vfill
\begin{tcolorbox}[enhanced,nobeforeafter,width=7.1cm,height=2.1cm,valign=center,colback=cprb!6,colframe=cprb,arc=2.5mm,boxrule=0.9pt,left=3mm]{\Large\bfseries\color{cprb}a)}\quad{\Large $\dfrac{5}{16}$}\end{tcolorbox}\hfill\begin{tcolorbox}[enhanced,nobeforeafter,width=7.1cm,height=2.1cm,valign=center,colback=cprb!6,colframe=cprb,arc=2.5mm,boxrule=0.9pt,left=3mm]{\Large\bfseries\color{cprb}b)}\quad{\Large $\dfrac{5}{2}$}\end{tcolorbox}\par\vspace{2.5mm}
\begin{tcolorbox}[enhanced,nobeforeafter,width=7.1cm,height=2.1cm,valign=center,colback=cprb!6,colframe=cprb,arc=2.5mm,boxrule=0.9pt,left=3mm]{\Large\bfseries\color{cprb}c)}\quad{\Large $\dfrac{1}{64}$}\end{tcolorbox}\hfill\begin{tcolorbox}[enhanced,nobeforeafter,width=7.1cm,height=2.1cm,valign=center,colback=cprb!6,colframe=cprb,arc=2.5mm,boxrule=0.9pt,left=3mm]{\Large\bfseries\color{cprb}d)}\quad{\Large $\dfrac{1}{2}$}\end{tcolorbox}
\end{tcolorbox}
\newpage
\begin{tcolorbox}[enhanced,width=16cm,height=14cm,colback=white,colframe=cprb,boxrule=2pt,arc=0mm,
left=6mm,right=6mm,top=0mm,bottom=4mm,sharp corners,
overlay={\fill[cprb] (frame.north west) rectangle ([yshift=-1.25cm]frame.north east);
\node[anchor=west,text=white,font=\sffamily\bfseries\LARGE] at ([xshift=5mm,yshift=-0.62cm]frame.north west) {Question 82 B};
\node[anchor=east,text=white,font=\sffamily\itshape\large] at ([xshift=-5mm,yshift=-0.62cm]frame.north east) {Probabilités et dénombrement};}]
\vspace*{1.75cm}
\begin{minipage}[t][6.9cm][c]{\linewidth}\LARGE\centering
\begin{minipage}{0.95\linewidth}\raggedright\linespread{1.25}\selectfont $X$ suit la loi binomiale $\mathcal{B}\left(5;\frac12\right)$. Que vaut $P(X=2)$~?\end{minipage}
\end{minipage}

\vfill
\begin{tcolorbox}[enhanced,nobeforeafter,width=7.1cm,height=2.1cm,valign=center,colback=cprb!6,colframe=cprb,arc=2.5mm,boxrule=0.9pt,left=3mm]{\Large\bfseries\color{cprb}a)}\quad{\Large $\dfrac{2}{5}$}\end{tcolorbox}\hfill\begin{tcolorbox}[enhanced,nobeforeafter,width=7.1cm,height=2.1cm,valign=center,colback=cprb!6,colframe=cprb,arc=2.5mm,boxrule=0.9pt,left=3mm]{\Large\bfseries\color{cprb}b)}\quad{\Large $\dfrac{5}{16}$}\end{tcolorbox}\par\vspace{2.5mm}
\begin{tcolorbox}[enhanced,nobeforeafter,width=7.1cm,height=2.1cm,valign=center,colback=cprb!6,colframe=cprb,arc=2.5mm,boxrule=0.9pt,left=3mm]{\Large\bfseries\color{cprb}c)}\quad{\Large $\dfrac{5}{2}$}\end{tcolorbox}\hfill\begin{tcolorbox}[enhanced,nobeforeafter,width=7.1cm,height=2.1cm,valign=center,colback=cprb!6,colframe=cprb,arc=2.5mm,boxrule=0.9pt,left=3mm]{\Large\bfseries\color{cprb}d)}\quad{\Large $\dfrac{1}{32}$}\end{tcolorbox}
\end{tcolorbox}
\newpage
\begin{tcolorbox}[enhanced,width=16cm,height=14cm,colback=white,colframe=cprb,boxrule=2pt,arc=0mm,
left=6mm,right=6mm,top=0mm,bottom=4mm,sharp corners,
overlay={\fill[cprb] (frame.north west) rectangle ([yshift=-1.25cm]frame.north east);
\node[anchor=west,text=white,font=\sffamily\bfseries\LARGE] at ([xshift=5mm,yshift=-0.62cm]frame.north west) {Question 83 B};
\node[anchor=east,text=white,font=\sffamily\itshape\large] at ([xshift=-5mm,yshift=-0.62cm]frame.north east) {Probabilités et dénombrement};}]
\vspace*{1.75cm}
\begin{minipage}[t][6.9cm][c]{\linewidth}\LARGE\centering
\begin{minipage}{0.95\linewidth}\raggedright\linespread{1.25}\selectfont $X$ suit la loi binomiale $\mathcal{B}\left(12;\frac{1}{6}\right)$. Quelle est son espérance $E(X)$~?\end{minipage}
\end{minipage}

\vfill
\begin{tcolorbox}[enhanced,nobeforeafter,width=7.1cm,height=2.1cm,valign=center,colback=cprb!6,colframe=cprb,arc=2.5mm,boxrule=0.9pt,left=3mm]{\Large\bfseries\color{cprb}a)}\quad{\Large $\dfrac{5}{3}$}\end{tcolorbox}\hfill\begin{tcolorbox}[enhanced,nobeforeafter,width=7.1cm,height=2.1cm,valign=center,colback=cprb!6,colframe=cprb,arc=2.5mm,boxrule=0.9pt,left=3mm]{\Large\bfseries\color{cprb}b)}\quad{\Large $\dfrac{1}{6}$}\end{tcolorbox}\par\vspace{2.5mm}
\begin{tcolorbox}[enhanced,nobeforeafter,width=7.1cm,height=2.1cm,valign=center,colback=cprb!6,colframe=cprb,arc=2.5mm,boxrule=0.9pt,left=3mm]{\Large\bfseries\color{cprb}c)}\quad{\Large $10$}\end{tcolorbox}\hfill\begin{tcolorbox}[enhanced,nobeforeafter,width=7.1cm,height=2.1cm,valign=center,colback=cprb!6,colframe=cprb,arc=2.5mm,boxrule=0.9pt,left=3mm]{\Large\bfseries\color{cprb}d)}\quad{\Large $2$}\end{tcolorbox}
\end{tcolorbox}
\newpage
\begin{tcolorbox}[enhanced,width=16cm,height=14cm,colback=white,colframe=cprb,boxrule=2pt,arc=0mm,
left=6mm,right=6mm,top=0mm,bottom=4mm,sharp corners,
overlay={\fill[cprb] (frame.north west) rectangle ([yshift=-1.25cm]frame.north east);
\node[anchor=west,text=white,font=\sffamily\bfseries\LARGE] at ([xshift=5mm,yshift=-0.62cm]frame.north west) {Question 84 B};
\node[anchor=east,text=white,font=\sffamily\itshape\large] at ([xshift=-5mm,yshift=-0.62cm]frame.north east) {Probabilités et dénombrement};}]
\vspace*{1.75cm}
\begin{minipage}[t][6.9cm][c]{\linewidth}\LARGE\centering
\begin{minipage}{0.95\linewidth}\raggedright\linespread{1.25}\selectfont $X$ suit la loi binomiale $\mathcal{B}\left(40;\frac{3}{8}\right)$. Quelle est son espérance $E(X)$~?\end{minipage}
\end{minipage}

\vfill
\begin{tcolorbox}[enhanced,nobeforeafter,width=7.1cm,height=2.1cm,valign=center,colback=cprb!6,colframe=cprb,arc=2.5mm,boxrule=0.9pt,left=3mm]{\Large\bfseries\color{cprb}a)}\quad{\Large $15$}\end{tcolorbox}\hfill\begin{tcolorbox}[enhanced,nobeforeafter,width=7.1cm,height=2.1cm,valign=center,colback=cprb!6,colframe=cprb,arc=2.5mm,boxrule=0.9pt,left=3mm]{\Large\bfseries\color{cprb}b)}\quad{\Large $25$}\end{tcolorbox}\par\vspace{2.5mm}
\begin{tcolorbox}[enhanced,nobeforeafter,width=7.1cm,height=2.1cm,valign=center,colback=cprb!6,colframe=cprb,arc=2.5mm,boxrule=0.9pt,left=3mm]{\Large\bfseries\color{cprb}c)}\quad{\Large $\dfrac{75}{8}$}\end{tcolorbox}\hfill\begin{tcolorbox}[enhanced,nobeforeafter,width=7.1cm,height=2.1cm,valign=center,colback=cprb!6,colframe=cprb,arc=2.5mm,boxrule=0.9pt,left=3mm]{\Large\bfseries\color{cprb}d)}\quad{\Large $\dfrac{3}{8}$}\end{tcolorbox}
\end{tcolorbox}
\newpage
\begin{tcolorbox}[enhanced,width=16cm,height=14cm,colback=white,colframe=cprb,boxrule=2pt,arc=0mm,
left=6mm,right=6mm,top=0mm,bottom=4mm,sharp corners,
overlay={\fill[cprb] (frame.north west) rectangle ([yshift=-1.25cm]frame.north east);
\node[anchor=west,text=white,font=\sffamily\bfseries\LARGE] at ([xshift=5mm,yshift=-0.62cm]frame.north west) {Question 85 B};
\node[anchor=east,text=white,font=\sffamily\itshape\large] at ([xshift=-5mm,yshift=-0.62cm]frame.north east) {Probabilités et dénombrement};}]
\vspace*{1.75cm}
\begin{minipage}[t][6.9cm][c]{\linewidth}\LARGE\centering
\begin{minipage}{0.95\linewidth}\raggedright\linespread{1.25}\selectfont $X$ suit la loi binomiale $\mathcal{B}\left(18;\frac{1}{3}\right)$. Quelle est sa variance $V(X)$~?\end{minipage}
\end{minipage}

\vfill
\begin{tcolorbox}[enhanced,nobeforeafter,width=7.1cm,height=2.1cm,valign=center,colback=cprb!6,colframe=cprb,arc=2.5mm,boxrule=0.9pt,left=3mm]{\Large\bfseries\color{cprb}a)}\quad{\Large $4$}\end{tcolorbox}\hfill\begin{tcolorbox}[enhanced,nobeforeafter,width=7.1cm,height=2.1cm,valign=center,colback=cprb!6,colframe=cprb,arc=2.5mm,boxrule=0.9pt,left=3mm]{\Large\bfseries\color{cprb}b)}\quad{\Large $2$}\end{tcolorbox}\par\vspace{2.5mm}
\begin{tcolorbox}[enhanced,nobeforeafter,width=7.1cm,height=2.1cm,valign=center,colback=cprb!6,colframe=cprb,arc=2.5mm,boxrule=0.9pt,left=3mm]{\Large\bfseries\color{cprb}c)}\quad{\Large $6$}\end{tcolorbox}\hfill\begin{tcolorbox}[enhanced,nobeforeafter,width=7.1cm,height=2.1cm,valign=center,colback=cprb!6,colframe=cprb,arc=2.5mm,boxrule=0.9pt,left=3mm]{\Large\bfseries\color{cprb}d)}\quad{\Large $8$}\end{tcolorbox}
\end{tcolorbox}
\newpage
\begin{tcolorbox}[enhanced,width=16cm,height=14cm,colback=white,colframe=cprb,boxrule=2pt,arc=0mm,
left=6mm,right=6mm,top=0mm,bottom=4mm,sharp corners,
overlay={\fill[cprb] (frame.north west) rectangle ([yshift=-1.25cm]frame.north east);
\node[anchor=west,text=white,font=\sffamily\bfseries\LARGE] at ([xshift=5mm,yshift=-0.62cm]frame.north west) {Question 86 B};
\node[anchor=east,text=white,font=\sffamily\itshape\large] at ([xshift=-5mm,yshift=-0.62cm]frame.north east) {Probabilités et dénombrement};}]
\vspace*{1.75cm}
\begin{minipage}[t][6.9cm][c]{\linewidth}\LARGE\centering
\begin{minipage}{0.95\linewidth}\raggedright\linespread{1.25}\selectfont $X$ suit la loi binomiale $\mathcal{B}\left(16;\frac{3}{4}\right)$. Quelle est sa variance $V(X)$~?\end{minipage}
\end{minipage}

\vfill
\begin{tcolorbox}[enhanced,nobeforeafter,width=7.1cm,height=2.1cm,valign=center,colback=cprb!6,colframe=cprb,arc=2.5mm,boxrule=0.9pt,left=3mm]{\Large\bfseries\color{cprb}a)}\quad{\Large $1$}\end{tcolorbox}\hfill\begin{tcolorbox}[enhanced,nobeforeafter,width=7.1cm,height=2.1cm,valign=center,colback=cprb!6,colframe=cprb,arc=2.5mm,boxrule=0.9pt,left=3mm]{\Large\bfseries\color{cprb}b)}\quad{\Large $3$}\end{tcolorbox}\par\vspace{2.5mm}
\begin{tcolorbox}[enhanced,nobeforeafter,width=7.1cm,height=2.1cm,valign=center,colback=cprb!6,colframe=cprb,arc=2.5mm,boxrule=0.9pt,left=3mm]{\Large\bfseries\color{cprb}c)}\quad{\Large $9$}\end{tcolorbox}\hfill\begin{tcolorbox}[enhanced,nobeforeafter,width=7.1cm,height=2.1cm,valign=center,colback=cprb!6,colframe=cprb,arc=2.5mm,boxrule=0.9pt,left=3mm]{\Large\bfseries\color{cprb}d)}\quad{\Large $12$}\end{tcolorbox}
\end{tcolorbox}
\newpage
\begin{tcolorbox}[enhanced,width=16cm,height=14cm,colback=white,colframe=cprb,boxrule=2pt,arc=0mm,
left=6mm,right=6mm,top=0mm,bottom=4mm,sharp corners,
overlay={\fill[cprb] (frame.north west) rectangle ([yshift=-1.25cm]frame.north east);
\node[anchor=west,text=white,font=\sffamily\bfseries\LARGE] at ([xshift=5mm,yshift=-0.62cm]frame.north west) {Question 87 B};
\node[anchor=east,text=white,font=\sffamily\itshape\large] at ([xshift=-5mm,yshift=-0.62cm]frame.north east) {Probabilités et dénombrement};}]
\vspace*{1.75cm}
\begin{minipage}[t][6.9cm][c]{\linewidth}\LARGE\centering
\begin{minipage}{0.95\linewidth}\raggedright\linespread{1.25}\selectfont On a $P(A)=0{,}7$ et $P_A(B)=0{,}2$. Que vaut $P(A\cap B)$~?\end{minipage}
\end{minipage}

\vfill
\begin{tcolorbox}[enhanced,nobeforeafter,width=7.1cm,height=2.1cm,valign=center,colback=cprb!6,colframe=cprb,arc=2.5mm,boxrule=0.9pt,left=3mm]{\Large\bfseries\color{cprb}a)}\quad{\Large $0{,}56$}\end{tcolorbox}\hfill\begin{tcolorbox}[enhanced,nobeforeafter,width=7.1cm,height=2.1cm,valign=center,colback=cprb!6,colframe=cprb,arc=2.5mm,boxrule=0.9pt,left=3mm]{\Large\bfseries\color{cprb}b)}\quad{\Large $0{,}5$}\end{tcolorbox}\par\vspace{2.5mm}
\begin{tcolorbox}[enhanced,nobeforeafter,width=7.1cm,height=2.1cm,valign=center,colback=cprb!6,colframe=cprb,arc=2.5mm,boxrule=0.9pt,left=3mm]{\Large\bfseries\color{cprb}c)}\quad{\Large $0{,}14$}\end{tcolorbox}\hfill\begin{tcolorbox}[enhanced,nobeforeafter,width=7.1cm,height=2.1cm,valign=center,colback=cprb!6,colframe=cprb,arc=2.5mm,boxrule=0.9pt,left=3mm]{\Large\bfseries\color{cprb}d)}\quad{\Large $0{,}9$}\end{tcolorbox}
\end{tcolorbox}
\newpage
\begin{tcolorbox}[enhanced,width=16cm,height=14cm,colback=white,colframe=cprb,boxrule=2pt,arc=0mm,
left=6mm,right=6mm,top=0mm,bottom=4mm,sharp corners,
overlay={\fill[cprb] (frame.north west) rectangle ([yshift=-1.25cm]frame.north east);
\node[anchor=west,text=white,font=\sffamily\bfseries\LARGE] at ([xshift=5mm,yshift=-0.62cm]frame.north west) {Question 88 B};
\node[anchor=east,text=white,font=\sffamily\itshape\large] at ([xshift=-5mm,yshift=-0.62cm]frame.north east) {Probabilités et dénombrement};}]
\vspace*{1.75cm}
\begin{minipage}[t][6.9cm][c]{\linewidth}\LARGE\centering
\begin{minipage}{0.95\linewidth}\raggedright\linespread{1.25}\selectfont On a $P(A)=0{,}25$ et $P_A(B)=0{,}8$. Que vaut $P(A\cap B)$~?\end{minipage}
\end{minipage}

\vfill
\begin{tcolorbox}[enhanced,nobeforeafter,width=7.1cm,height=2.1cm,valign=center,colback=cprb!6,colframe=cprb,arc=2.5mm,boxrule=0.9pt,left=3mm]{\Large\bfseries\color{cprb}a)}\quad{\Large $0{,}55$}\end{tcolorbox}\hfill\begin{tcolorbox}[enhanced,nobeforeafter,width=7.1cm,height=2.1cm,valign=center,colback=cprb!6,colframe=cprb,arc=2.5mm,boxrule=0.9pt,left=3mm]{\Large\bfseries\color{cprb}b)}\quad{\Large $0{,}2$}\end{tcolorbox}\par\vspace{2.5mm}
\begin{tcolorbox}[enhanced,nobeforeafter,width=7.1cm,height=2.1cm,valign=center,colback=cprb!6,colframe=cprb,arc=2.5mm,boxrule=0.9pt,left=3mm]{\Large\bfseries\color{cprb}c)}\quad{\Large $0{,}05$}\end{tcolorbox}\hfill\begin{tcolorbox}[enhanced,nobeforeafter,width=7.1cm,height=2.1cm,valign=center,colback=cprb!6,colframe=cprb,arc=2.5mm,boxrule=0.9pt,left=3mm]{\Large\bfseries\color{cprb}d)}\quad{\Large $0{,}8$}\end{tcolorbox}
\end{tcolorbox}
\newpage
\begin{tcolorbox}[enhanced,width=16cm,height=14cm,colback=white,colframe=cprb,boxrule=2pt,arc=0mm,
left=6mm,right=6mm,top=0mm,bottom=4mm,sharp corners,
overlay={\fill[cprb] (frame.north west) rectangle ([yshift=-1.25cm]frame.north east);
\node[anchor=west,text=white,font=\sffamily\bfseries\LARGE] at ([xshift=5mm,yshift=-0.62cm]frame.north west) {Question 89 B};
\node[anchor=east,text=white,font=\sffamily\itshape\large] at ([xshift=-5mm,yshift=-0.62cm]frame.north east) {Probabilités et dénombrement};}]
\vspace*{1.75cm}
\begin{minipage}[t][6.9cm][c]{\linewidth}\LARGE\centering
\begin{minipage}{0.95\linewidth}\raggedright\linespread{1.25}\selectfont On a $P(A)=0{,}5$, $P_A(B)=0{,}8$ et $P_{\overline{A}}(B)=0{,}4$. Que vaut $P(B)$~?\end{minipage}
\end{minipage}

\vfill
\begin{tcolorbox}[enhanced,nobeforeafter,width=7.1cm,height=2.1cm,valign=center,colback=cprb!6,colframe=cprb,arc=2.5mm,boxrule=0.9pt,left=3mm]{\Large\bfseries\color{cprb}a)}\quad{\Large $0{,}6$}\end{tcolorbox}\hfill\begin{tcolorbox}[enhanced,nobeforeafter,width=7.1cm,height=2.1cm,valign=center,colback=cprb!6,colframe=cprb,arc=2.5mm,boxrule=0.9pt,left=3mm]{\Large\bfseries\color{cprb}b)}\quad{\Large $0{,}4$}\end{tcolorbox}\par\vspace{2.5mm}
\begin{tcolorbox}[enhanced,nobeforeafter,width=7.1cm,height=2.1cm,valign=center,colback=cprb!6,colframe=cprb,arc=2.5mm,boxrule=0.9pt,left=3mm]{\Large\bfseries\color{cprb}c)}\quad{\Large $0{,}32$}\end{tcolorbox}\hfill\begin{tcolorbox}[enhanced,nobeforeafter,width=7.1cm,height=2.1cm,valign=center,colback=cprb!6,colframe=cprb,arc=2.5mm,boxrule=0.9pt,left=3mm]{\Large\bfseries\color{cprb}d)}\quad{\Large $0{,}2$}\end{tcolorbox}
\end{tcolorbox}
\newpage
\begin{tcolorbox}[enhanced,width=16cm,height=14cm,colback=white,colframe=cprb,boxrule=2pt,arc=0mm,
left=6mm,right=6mm,top=0mm,bottom=4mm,sharp corners,
overlay={\fill[cprb] (frame.north west) rectangle ([yshift=-1.25cm]frame.north east);
\node[anchor=west,text=white,font=\sffamily\bfseries\LARGE] at ([xshift=5mm,yshift=-0.62cm]frame.north west) {Question 90 B};
\node[anchor=east,text=white,font=\sffamily\itshape\large] at ([xshift=-5mm,yshift=-0.62cm]frame.north east) {Probabilités et dénombrement};}]
\vspace*{1.75cm}
\begin{minipage}[t][6.9cm][c]{\linewidth}\LARGE\centering
\begin{minipage}{0.95\linewidth}\raggedright\linespread{1.25}\selectfont On a $P(A)=0{,}2$, $P_A(B)=0{,}9$ et $P_{\overline{A}}(B)=0{,}3$. Que vaut $P(B)$~?\end{minipage}
\end{minipage}

\vfill
\begin{tcolorbox}[enhanced,nobeforeafter,width=7.1cm,height=2.1cm,valign=center,colback=cprb!6,colframe=cprb,arc=2.5mm,boxrule=0.9pt,left=3mm]{\Large\bfseries\color{cprb}a)}\quad{\Large $0{,}42$}\end{tcolorbox}\hfill\begin{tcolorbox}[enhanced,nobeforeafter,width=7.1cm,height=2.1cm,valign=center,colback=cprb!6,colframe=cprb,arc=2.5mm,boxrule=0.9pt,left=3mm]{\Large\bfseries\color{cprb}b)}\quad{\Large $0{,}27$}\end{tcolorbox}\par\vspace{2.5mm}
\begin{tcolorbox}[enhanced,nobeforeafter,width=7.1cm,height=2.1cm,valign=center,colback=cprb!6,colframe=cprb,arc=2.5mm,boxrule=0.9pt,left=3mm]{\Large\bfseries\color{cprb}c)}\quad{\Large $0{,}24$}\end{tcolorbox}\hfill\begin{tcolorbox}[enhanced,nobeforeafter,width=7.1cm,height=2.1cm,valign=center,colback=cprb!6,colframe=cprb,arc=2.5mm,boxrule=0.9pt,left=3mm]{\Large\bfseries\color{cprb}d)}\quad{\Large $0{,}18$}\end{tcolorbox}
\end{tcolorbox}
\newpage
\begin{tcolorbox}[enhanced,width=16cm,height=14cm,colback=white,colframe=cprb,boxrule=2pt,arc=0mm,
left=6mm,right=6mm,top=0mm,bottom=4mm,sharp corners,
overlay={\fill[cprb] (frame.north west) rectangle ([yshift=-1.25cm]frame.north east);
\node[anchor=west,text=white,font=\sffamily\bfseries\LARGE] at ([xshift=5mm,yshift=-0.62cm]frame.north west) {Question 91 B};
\node[anchor=east,text=white,font=\sffamily\itshape\large] at ([xshift=-5mm,yshift=-0.62cm]frame.north east) {Probabilités et dénombrement};}]
\vspace*{1.75cm}
\begin{minipage}[t][6.9cm][c]{\linewidth}\LARGE\centering
\begin{minipage}{0.95\linewidth}\raggedright\linespread{1.25}\selectfont Combien y a-t-il de façons de choisir 2 élèves parmi 8 (sans tenir compte de l'ordre)~?\end{minipage}
\end{minipage}

\vfill
\begin{tcolorbox}[enhanced,nobeforeafter,width=7.1cm,height=2.1cm,valign=center,colback=cprb!6,colframe=cprb,arc=2.5mm,boxrule=0.9pt,left=3mm]{\Large\bfseries\color{cprb}a)}\quad{\Large $28$}\end{tcolorbox}\hfill\begin{tcolorbox}[enhanced,nobeforeafter,width=7.1cm,height=2.1cm,valign=center,colback=cprb!6,colframe=cprb,arc=2.5mm,boxrule=0.9pt,left=3mm]{\Large\bfseries\color{cprb}b)}\quad{\Large $56$}\end{tcolorbox}\par\vspace{2.5mm}
\begin{tcolorbox}[enhanced,nobeforeafter,width=7.1cm,height=2.1cm,valign=center,colback=cprb!6,colframe=cprb,arc=2.5mm,boxrule=0.9pt,left=3mm]{\Large\bfseries\color{cprb}c)}\quad{\Large $16$}\end{tcolorbox}\hfill\begin{tcolorbox}[enhanced,nobeforeafter,width=7.1cm,height=2.1cm,valign=center,colback=cprb!6,colframe=cprb,arc=2.5mm,boxrule=0.9pt,left=3mm]{\Large\bfseries\color{cprb}d)}\quad{\Large $64$}\end{tcolorbox}
\end{tcolorbox}
\newpage
\begin{tcolorbox}[enhanced,width=16cm,height=14cm,colback=white,colframe=cprb,boxrule=2pt,arc=0mm,
left=6mm,right=6mm,top=0mm,bottom=4mm,sharp corners,
overlay={\fill[cprb] (frame.north west) rectangle ([yshift=-1.25cm]frame.north east);
\node[anchor=west,text=white,font=\sffamily\bfseries\LARGE] at ([xshift=5mm,yshift=-0.62cm]frame.north west) {Question 92 B};
\node[anchor=east,text=white,font=\sffamily\itshape\large] at ([xshift=-5mm,yshift=-0.62cm]frame.north east) {Probabilités et dénombrement};}]
\vspace*{1.75cm}
\begin{minipage}[t][6.9cm][c]{\linewidth}\LARGE\centering
\begin{minipage}{0.95\linewidth}\raggedright\linespread{1.25}\selectfont Combien y a-t-il de façons de choisir 3 élèves parmi 7 (sans tenir compte de l'ordre)~?\end{minipage}
\end{minipage}

\vfill
\begin{tcolorbox}[enhanced,nobeforeafter,width=7.1cm,height=2.1cm,valign=center,colback=cprb!6,colframe=cprb,arc=2.5mm,boxrule=0.9pt,left=3mm]{\Large\bfseries\color{cprb}a)}\quad{\Large $21$}\end{tcolorbox}\hfill\begin{tcolorbox}[enhanced,nobeforeafter,width=7.1cm,height=2.1cm,valign=center,colback=cprb!6,colframe=cprb,arc=2.5mm,boxrule=0.9pt,left=3mm]{\Large\bfseries\color{cprb}b)}\quad{\Large $35$}\end{tcolorbox}\par\vspace{2.5mm}
\begin{tcolorbox}[enhanced,nobeforeafter,width=7.1cm,height=2.1cm,valign=center,colback=cprb!6,colframe=cprb,arc=2.5mm,boxrule=0.9pt,left=3mm]{\Large\bfseries\color{cprb}c)}\quad{\Large $343$}\end{tcolorbox}\hfill\begin{tcolorbox}[enhanced,nobeforeafter,width=7.1cm,height=2.1cm,valign=center,colback=cprb!6,colframe=cprb,arc=2.5mm,boxrule=0.9pt,left=3mm]{\Large\bfseries\color{cprb}d)}\quad{\Large $210$}\end{tcolorbox}
\end{tcolorbox}
\newpage
\begin{tcolorbox}[enhanced,width=16cm,height=14cm,colback=white,colframe=cprb,boxrule=2pt,arc=0mm,
left=6mm,right=6mm,top=0mm,bottom=4mm,sharp corners,
overlay={\fill[cprb] (frame.north west) rectangle ([yshift=-1.25cm]frame.north east);
\node[anchor=west,text=white,font=\sffamily\bfseries\LARGE] at ([xshift=5mm,yshift=-0.62cm]frame.north west) {Question 93 B};
\node[anchor=east,text=white,font=\sffamily\itshape\large] at ([xshift=-5mm,yshift=-0.62cm]frame.north east) {Probabilités et dénombrement};}]
\vspace*{1.75cm}
\begin{minipage}[t][6.9cm][c]{\linewidth}\LARGE\centering
\begin{minipage}{0.95\linewidth}\raggedright\linespread{1.25}\selectfont Combien de codes de 3 caractères peut-on former avec un alphabet de 4 symboles (répétitions autorisées)~?\end{minipage}
\end{minipage}

\vfill
\begin{tcolorbox}[enhanced,nobeforeafter,width=7.1cm,height=2.1cm,valign=center,colback=cprb!6,colframe=cprb,arc=2.5mm,boxrule=0.9pt,left=3mm]{\Large\bfseries\color{cprb}a)}\quad{\Large $64$}\end{tcolorbox}\hfill\begin{tcolorbox}[enhanced,nobeforeafter,width=7.1cm,height=2.1cm,valign=center,colback=cprb!6,colframe=cprb,arc=2.5mm,boxrule=0.9pt,left=3mm]{\Large\bfseries\color{cprb}b)}\quad{\Large $12$}\end{tcolorbox}\par\vspace{2.5mm}
\begin{tcolorbox}[enhanced,nobeforeafter,width=7.1cm,height=2.1cm,valign=center,colback=cprb!6,colframe=cprb,arc=2.5mm,boxrule=0.9pt,left=3mm]{\Large\bfseries\color{cprb}c)}\quad{\Large $81$}\end{tcolorbox}\hfill\begin{tcolorbox}[enhanced,nobeforeafter,width=7.1cm,height=2.1cm,valign=center,colback=cprb!6,colframe=cprb,arc=2.5mm,boxrule=0.9pt,left=3mm]{\Large\bfseries\color{cprb}d)}\quad{\Large $20$}\end{tcolorbox}
\end{tcolorbox}
\newpage
\begin{tcolorbox}[enhanced,width=16cm,height=14cm,colback=white,colframe=cprb,boxrule=2pt,arc=0mm,
left=6mm,right=6mm,top=0mm,bottom=4mm,sharp corners,
overlay={\fill[cprb] (frame.north west) rectangle ([yshift=-1.25cm]frame.north east);
\node[anchor=west,text=white,font=\sffamily\bfseries\LARGE] at ([xshift=5mm,yshift=-0.62cm]frame.north west) {Question 94 B};
\node[anchor=east,text=white,font=\sffamily\itshape\large] at ([xshift=-5mm,yshift=-0.62cm]frame.north east) {Probabilités et dénombrement};}]
\vspace*{1.75cm}
\begin{minipage}[t][6.9cm][c]{\linewidth}\LARGE\centering
\begin{minipage}{0.95\linewidth}\raggedright\linespread{1.25}\selectfont Combien de codes de 2 caractères peut-on former avec un alphabet de 5 symboles (répétitions autorisées)~?\end{minipage}
\end{minipage}

\vfill
\begin{tcolorbox}[enhanced,nobeforeafter,width=7.1cm,height=2.1cm,valign=center,colback=cprb!6,colframe=cprb,arc=2.5mm,boxrule=0.9pt,left=3mm]{\Large\bfseries\color{cprb}a)}\quad{\Large $15$}\end{tcolorbox}\hfill\begin{tcolorbox}[enhanced,nobeforeafter,width=7.1cm,height=2.1cm,valign=center,colback=cprb!6,colframe=cprb,arc=2.5mm,boxrule=0.9pt,left=3mm]{\Large\bfseries\color{cprb}b)}\quad{\Large $25$}\end{tcolorbox}\par\vspace{2.5mm}
\begin{tcolorbox}[enhanced,nobeforeafter,width=7.1cm,height=2.1cm,valign=center,colback=cprb!6,colframe=cprb,arc=2.5mm,boxrule=0.9pt,left=3mm]{\Large\bfseries\color{cprb}c)}\quad{\Large $32$}\end{tcolorbox}\hfill\begin{tcolorbox}[enhanced,nobeforeafter,width=7.1cm,height=2.1cm,valign=center,colback=cprb!6,colframe=cprb,arc=2.5mm,boxrule=0.9pt,left=3mm]{\Large\bfseries\color{cprb}d)}\quad{\Large $10$}\end{tcolorbox}
\end{tcolorbox}
\newpage
\begin{tcolorbox}[enhanced,width=16cm,height=14cm,colback=white,colframe=cprb,boxrule=2pt,arc=0mm,
left=6mm,right=6mm,top=0mm,bottom=4mm,sharp corners,
overlay={\fill[cprb] (frame.north west) rectangle ([yshift=-1.25cm]frame.north east);
\node[anchor=west,text=white,font=\sffamily\bfseries\LARGE] at ([xshift=5mm,yshift=-0.62cm]frame.north west) {Question 95 B};
\node[anchor=east,text=white,font=\sffamily\itshape\large] at ([xshift=-5mm,yshift=-0.62cm]frame.north east) {Probabilités et dénombrement};}]
\vspace*{1.75cm}
\begin{minipage}[t][6.9cm][c]{\linewidth}\LARGE\centering
\begin{minipage}{0.95\linewidth}\raggedright\linespread{1.25}\selectfont $X$ suit la loi binomiale $\mathcal{B}\left(4;\frac{1}{2}\right)$. Que vaut $P(X\geq 1)$~?\end{minipage}
\end{minipage}

\vfill
\begin{tcolorbox}[enhanced,nobeforeafter,width=7.1cm,height=2.1cm,valign=center,colback=cprb!6,colframe=cprb,arc=2.5mm,boxrule=0.9pt,left=3mm]{\Large\bfseries\color{cprb}a)}\quad{\Large $\dfrac{1}{16}$}\end{tcolorbox}\hfill\begin{tcolorbox}[enhanced,nobeforeafter,width=7.1cm,height=2.1cm,valign=center,colback=cprb!6,colframe=cprb,arc=2.5mm,boxrule=0.9pt,left=3mm]{\Large\bfseries\color{cprb}b)}\quad{\Large $\dfrac{15}{16}$}\end{tcolorbox}\par\vspace{2.5mm}
\begin{tcolorbox}[enhanced,nobeforeafter,width=7.1cm,height=2.1cm,valign=center,colback=cprb!6,colframe=cprb,arc=2.5mm,boxrule=0.9pt,left=3mm]{\Large\bfseries\color{cprb}c)}\quad{\Large $2$}\end{tcolorbox}\hfill\begin{tcolorbox}[enhanced,nobeforeafter,width=7.1cm,height=2.1cm,valign=center,colback=cprb!6,colframe=cprb,arc=2.5mm,boxrule=0.9pt,left=3mm]{\Large\bfseries\color{cprb}d)}\quad{\Large $\dfrac{1}{2}$}\end{tcolorbox}
\end{tcolorbox}
\newpage
\begin{tcolorbox}[enhanced,width=16cm,height=14cm,colback=white,colframe=cprb,boxrule=2pt,arc=0mm,
left=6mm,right=6mm,top=0mm,bottom=4mm,sharp corners,
overlay={\fill[cprb] (frame.north west) rectangle ([yshift=-1.25cm]frame.north east);
\node[anchor=west,text=white,font=\sffamily\bfseries\LARGE] at ([xshift=5mm,yshift=-0.62cm]frame.north west) {Question 96 B};
\node[anchor=east,text=white,font=\sffamily\itshape\large] at ([xshift=-5mm,yshift=-0.62cm]frame.north east) {Probabilités et dénombrement};}]
\vspace*{1.75cm}
\begin{minipage}[t][6.9cm][c]{\linewidth}\LARGE\centering
\begin{minipage}{0.95\linewidth}\raggedright\linespread{1.25}\selectfont $X$ suit la loi binomiale $\mathcal{B}\left(3;\frac{1}{4}\right)$. Que vaut $P(X\geq 1)$~?\end{minipage}
\end{minipage}

\vfill
\begin{tcolorbox}[enhanced,nobeforeafter,width=7.1cm,height=2.1cm,valign=center,colback=cprb!6,colframe=cprb,arc=2.5mm,boxrule=0.9pt,left=3mm]{\Large\bfseries\color{cprb}a)}\quad{\Large $\dfrac{37}{64}$}\end{tcolorbox}\hfill\begin{tcolorbox}[enhanced,nobeforeafter,width=7.1cm,height=2.1cm,valign=center,colback=cprb!6,colframe=cprb,arc=2.5mm,boxrule=0.9pt,left=3mm]{\Large\bfseries\color{cprb}b)}\quad{\Large $\dfrac{1}{64}$}\end{tcolorbox}\par\vspace{2.5mm}
\begin{tcolorbox}[enhanced,nobeforeafter,width=7.1cm,height=2.1cm,valign=center,colback=cprb!6,colframe=cprb,arc=2.5mm,boxrule=0.9pt,left=3mm]{\Large\bfseries\color{cprb}c)}\quad{\Large $\dfrac{3}{4}$}\end{tcolorbox}\hfill\begin{tcolorbox}[enhanced,nobeforeafter,width=7.1cm,height=2.1cm,valign=center,colback=cprb!6,colframe=cprb,arc=2.5mm,boxrule=0.9pt,left=3mm]{\Large\bfseries\color{cprb}d)}\quad{\Large $\dfrac{27}{64}$}\end{tcolorbox}
\end{tcolorbox}
\newpage
\end{document}